\documentclass[11pt]{article}

    \usepackage[breakable]{tcolorbox}
    \usepackage{parskip} % Stop auto-indenting (to mimic markdown behaviour)
    

    % Basic figure setup, for now with no caption control since it's done
    % automatically by Pandoc (which extracts ![](path) syntax from Markdown).
    \usepackage{graphicx}
    % Keep aspect ratio if custom image width or height is specified
    \setkeys{Gin}{keepaspectratio}
    % Maintain compatibility with old templates. Remove in nbconvert 6.0
    \let\Oldincludegraphics\includegraphics
    % Ensure that by default, figures have no caption (until we provide a
    % proper Figure object with a Caption API and a way to capture that
    % in the conversion process - todo).
    \usepackage{caption}
    \DeclareCaptionFormat{nocaption}{}
    \captionsetup{format=nocaption,aboveskip=0pt,belowskip=0pt}

    \usepackage{float}
    \floatplacement{figure}{H} % forces figures to be placed at the correct location
    \usepackage{xcolor} % Allow colors to be defined
    \usepackage{enumerate} % Needed for markdown enumerations to work
    \usepackage{geometry} % Used to adjust the document margins
    \usepackage{amsmath} % Equations
    \usepackage{amssymb} % Equations
    \usepackage{textcomp} % defines textquotesingle
    % Hack from http://tex.stackexchange.com/a/47451/13684:
    \AtBeginDocument{%
        \def\PYZsq{\textquotesingle}% Upright quotes in Pygmentized code
    }
    \usepackage{upquote} % Upright quotes for verbatim code
    \usepackage{eurosym} % defines \euro

    \usepackage{iftex}
    \ifPDFTeX
        \usepackage[T2A]{fontenc}
        \IfFileExists{alphabeta.sty}{
              \usepackage{alphabeta}
          }{
              \usepackage[mathletters]{ucs}
              \usepackage[utf8x]{inputenc}
          }
    \else
        \usepackage{fontspec}
        \usepackage{unicode-math}
    \fi

    \usepackage{fancyvrb} % verbatim replacement that allows latex
    \usepackage{grffile} % extends the file name processing of package graphics
                         % to support a larger range
    \makeatletter % fix for old versions of grffile with XeLaTeX
    \@ifpackagelater{grffile}{2019/11/01}
    {
      % Do nothing on new versions
    }
    {
      \def\Gread@@xetex#1{%
        \IfFileExists{"\Gin@base".bb}%
        {\Gread@eps{\Gin@base.bb}}%
        {\Gread@@xetex@aux#1}%
      }
    }
    \makeatother
    \usepackage[Export]{adjustbox} % Used to constrain images to a maximum size
    \adjustboxset{max size={0.9\linewidth}{0.9\paperheight}}

    % The hyperref package gives us a pdf with properly built
    % internal navigation ('pdf bookmarks' for the table of contents,
    % internal cross-reference links, web links for URLs, etc.)
    \usepackage{hyperref}
    % The default LaTeX title has an obnoxious amount of whitespace. By default,
    % titling removes some of it. It also provides customization options.
    \usepackage{titling}
    \usepackage{longtable} % longtable support required by pandoc >1.10
    \usepackage{booktabs}  % table support for pandoc > 1.12.2
    \usepackage{array}     % table support for pandoc >= 2.11.3
    \usepackage{calc}      % table minipage width calculation for pandoc >= 2.11.1
    \usepackage[inline]{enumitem} % IRkernel/repr support (it uses the enumerate* environment)
    \usepackage[normalem]{ulem} % ulem is needed to support strikethroughs (\sout)
                                % normalem makes italics be italics, not underlines
    \usepackage{soul}      % strikethrough (\st) support for pandoc >= 3.0.0
    \usepackage{mathrsfs}
    
    \usepackage[english,russian]{babel}

    
    % Colors for the hyperref package
    \definecolor{urlcolor}{rgb}{0,.145,.698}
    \definecolor{linkcolor}{rgb}{.71,0.21,0.01}
    \definecolor{citecolor}{rgb}{.12,.54,.11}

    % ANSI colors
    \definecolor{ansi-black}{HTML}{3E424D}
    \definecolor{ansi-black-intense}{HTML}{282C36}
    \definecolor{ansi-red}{HTML}{E75C58}
    \definecolor{ansi-red-intense}{HTML}{B22B31}
    \definecolor{ansi-green}{HTML}{00A250}
    \definecolor{ansi-green-intense}{HTML}{007427}
    \definecolor{ansi-yellow}{HTML}{DDB62B}
    \definecolor{ansi-yellow-intense}{HTML}{B27D12}
    \definecolor{ansi-blue}{HTML}{208FFB}
    \definecolor{ansi-blue-intense}{HTML}{0065CA}
    \definecolor{ansi-magenta}{HTML}{D160C4}
    \definecolor{ansi-magenta-intense}{HTML}{A03196}
    \definecolor{ansi-cyan}{HTML}{60C6C8}
    \definecolor{ansi-cyan-intense}{HTML}{258F8F}
    \definecolor{ansi-white}{HTML}{C5C1B4}
    \definecolor{ansi-white-intense}{HTML}{A1A6B2}
    \definecolor{ansi-default-inverse-fg}{HTML}{FFFFFF}
    \definecolor{ansi-default-inverse-bg}{HTML}{000000}

    % common color for the border for error outputs.
    \definecolor{outerrorbackground}{HTML}{FFDFDF}

    % commands and environments needed by pandoc snippets
    % extracted from the output of `pandoc -s`
    \providecommand{\tightlist}{%
      \setlength{\itemsep}{0pt}\setlength{\parskip}{0pt}}
    \DefineVerbatimEnvironment{Highlighting}{Verbatim}{commandchars=\\\{\}}
    % Add ',fontsize=\small' for more characters per line
    \newenvironment{Shaded}{}{}
    \newcommand{\KeywordTok}[1]{\textcolor[rgb]{0.00,0.44,0.13}{\textbf{{#1}}}}
    \newcommand{\DataTypeTok}[1]{\textcolor[rgb]{0.56,0.13,0.00}{{#1}}}
    \newcommand{\DecValTok}[1]{\textcolor[rgb]{0.25,0.63,0.44}{{#1}}}
    \newcommand{\BaseNTok}[1]{\textcolor[rgb]{0.25,0.63,0.44}{{#1}}}
    \newcommand{\FloatTok}[1]{\textcolor[rgb]{0.25,0.63,0.44}{{#1}}}
    \newcommand{\CharTok}[1]{\textcolor[rgb]{0.25,0.44,0.63}{{#1}}}
    \newcommand{\StringTok}[1]{\textcolor[rgb]{0.25,0.44,0.63}{{#1}}}
    \newcommand{\CommentTok}[1]{\textcolor[rgb]{0.38,0.63,0.69}{\textit{{#1}}}}
    \newcommand{\OtherTok}[1]{\textcolor[rgb]{0.00,0.44,0.13}{{#1}}}
    \newcommand{\AlertTok}[1]{\textcolor[rgb]{1.00,0.00,0.00}{\textbf{{#1}}}}
    \newcommand{\FunctionTok}[1]{\textcolor[rgb]{0.02,0.16,0.49}{{#1}}}
    \newcommand{\RegionMarkerTok}[1]{{#1}}
    \newcommand{\ErrorTok}[1]{\textcolor[rgb]{1.00,0.00,0.00}{\textbf{{#1}}}}
    \newcommand{\NormalTok}[1]{{#1}}

    % Additional commands for more recent versions of Pandoc
    \newcommand{\ConstantTok}[1]{\textcolor[rgb]{0.53,0.00,0.00}{{#1}}}
    \newcommand{\SpecialCharTok}[1]{\textcolor[rgb]{0.25,0.44,0.63}{{#1}}}
    \newcommand{\VerbatimStringTok}[1]{\textcolor[rgb]{0.25,0.44,0.63}{{#1}}}
    \newcommand{\SpecialStringTok}[1]{\textcolor[rgb]{0.73,0.40,0.53}{{#1}}}
    \newcommand{\ImportTok}[1]{{#1}}
    \newcommand{\DocumentationTok}[1]{\textcolor[rgb]{0.73,0.13,0.13}{\textit{{#1}}}}
    \newcommand{\AnnotationTok}[1]{\textcolor[rgb]{0.38,0.63,0.69}{\textbf{\textit{{#1}}}}}
    \newcommand{\CommentVarTok}[1]{\textcolor[rgb]{0.38,0.63,0.69}{\textbf{\textit{{#1}}}}}
    \newcommand{\VariableTok}[1]{\textcolor[rgb]{0.10,0.09,0.49}{{#1}}}
    \newcommand{\ControlFlowTok}[1]{\textcolor[rgb]{0.00,0.44,0.13}{\textbf{{#1}}}}
    \newcommand{\OperatorTok}[1]{\textcolor[rgb]{0.40,0.40,0.40}{{#1}}}
    \newcommand{\BuiltInTok}[1]{{#1}}
    \newcommand{\ExtensionTok}[1]{{#1}}
    \newcommand{\PreprocessorTok}[1]{\textcolor[rgb]{0.74,0.48,0.00}{{#1}}}
    \newcommand{\AttributeTok}[1]{\textcolor[rgb]{0.49,0.56,0.16}{{#1}}}
    \newcommand{\InformationTok}[1]{\textcolor[rgb]{0.38,0.63,0.69}{\textbf{\textit{{#1}}}}}
    \newcommand{\WarningTok}[1]{\textcolor[rgb]{0.38,0.63,0.69}{\textbf{\textit{{#1}}}}}
    \makeatletter
    \newsavebox\pandoc@box
    \newcommand*\pandocbounded[1]{%
      \sbox\pandoc@box{#1}%
      % scaling factors for width and height
      \Gscale@div\@tempa\textheight{\dimexpr\ht\pandoc@box+\dp\pandoc@box\relax}%
      \Gscale@div\@tempb\linewidth{\wd\pandoc@box}%
      % select the smaller of both
      \ifdim\@tempb\p@<\@tempa\p@
        \let\@tempa\@tempb
      \fi
      % scaling accordingly (\@tempa < 1)
      \ifdim\@tempa\p@<\p@
        \scalebox{\@tempa}{\usebox\pandoc@box}%
      % scaling not needed, use as it is
      \else
        \usebox{\pandoc@box}%
      \fi
    }
    \makeatother

    % Define a nice break command that doesn't care if a line doesn't already
    % exist.
    \def\br{\hspace*{\fill} \\* }
    % Math Jax compatibility definitions
    \def\gt{>}
    \def\lt{<}
    \let\Oldtex\TeX
    \let\Oldlatex\LaTeX
    \renewcommand{\TeX}{\textrm{\Oldtex}}
    \renewcommand{\LaTeX}{\textrm{\Oldlatex}}
    % Document parameters
    % Document title
    \title{Notebook\_BVP\_k\_M\_structure}
    
    
    
    
    
    
    
% Pygments definitions
\makeatletter
\def\PY@reset{\let\PY@it=\relax \let\PY@bf=\relax%
    \let\PY@ul=\relax \let\PY@tc=\relax%
    \let\PY@bc=\relax \let\PY@ff=\relax}
\def\PY@tok#1{\csname PY@tok@#1\endcsname}
\def\PY@toks#1+{\ifx\relax#1\empty\else%
    \PY@tok{#1}\expandafter\PY@toks\fi}
\def\PY@do#1{\PY@bc{\PY@tc{\PY@ul{%
    \PY@it{\PY@bf{\PY@ff{#1}}}}}}}
\def\PY#1#2{\PY@reset\PY@toks#1+\relax+\PY@do{#2}}

\@namedef{PY@tok@w}{\def\PY@tc##1{\textcolor[rgb]{0.73,0.73,0.73}{##1}}}
\@namedef{PY@tok@c}{\let\PY@it=\textit\def\PY@tc##1{\textcolor[rgb]{0.24,0.48,0.48}{##1}}}
\@namedef{PY@tok@cp}{\def\PY@tc##1{\textcolor[rgb]{0.61,0.40,0.00}{##1}}}
\@namedef{PY@tok@k}{\let\PY@bf=\textbf\def\PY@tc##1{\textcolor[rgb]{0.00,0.50,0.00}{##1}}}
\@namedef{PY@tok@kp}{\def\PY@tc##1{\textcolor[rgb]{0.00,0.50,0.00}{##1}}}
\@namedef{PY@tok@kt}{\def\PY@tc##1{\textcolor[rgb]{0.69,0.00,0.25}{##1}}}
\@namedef{PY@tok@o}{\def\PY@tc##1{\textcolor[rgb]{0.40,0.40,0.40}{##1}}}
\@namedef{PY@tok@ow}{\let\PY@bf=\textbf\def\PY@tc##1{\textcolor[rgb]{0.67,0.13,1.00}{##1}}}
\@namedef{PY@tok@nb}{\def\PY@tc##1{\textcolor[rgb]{0.00,0.50,0.00}{##1}}}
\@namedef{PY@tok@nf}{\def\PY@tc##1{\textcolor[rgb]{0.00,0.00,1.00}{##1}}}
\@namedef{PY@tok@nc}{\let\PY@bf=\textbf\def\PY@tc##1{\textcolor[rgb]{0.00,0.00,1.00}{##1}}}
\@namedef{PY@tok@nn}{\let\PY@bf=\textbf\def\PY@tc##1{\textcolor[rgb]{0.00,0.00,1.00}{##1}}}
\@namedef{PY@tok@ne}{\let\PY@bf=\textbf\def\PY@tc##1{\textcolor[rgb]{0.80,0.25,0.22}{##1}}}
\@namedef{PY@tok@nv}{\def\PY@tc##1{\textcolor[rgb]{0.10,0.09,0.49}{##1}}}
\@namedef{PY@tok@no}{\def\PY@tc##1{\textcolor[rgb]{0.53,0.00,0.00}{##1}}}
\@namedef{PY@tok@nl}{\def\PY@tc##1{\textcolor[rgb]{0.46,0.46,0.00}{##1}}}
\@namedef{PY@tok@ni}{\let\PY@bf=\textbf\def\PY@tc##1{\textcolor[rgb]{0.44,0.44,0.44}{##1}}}
\@namedef{PY@tok@na}{\def\PY@tc##1{\textcolor[rgb]{0.41,0.47,0.13}{##1}}}
\@namedef{PY@tok@nt}{\let\PY@bf=\textbf\def\PY@tc##1{\textcolor[rgb]{0.00,0.50,0.00}{##1}}}
\@namedef{PY@tok@nd}{\def\PY@tc##1{\textcolor[rgb]{0.67,0.13,1.00}{##1}}}
\@namedef{PY@tok@s}{\def\PY@tc##1{\textcolor[rgb]{0.73,0.13,0.13}{##1}}}
\@namedef{PY@tok@sd}{\let\PY@it=\textit\def\PY@tc##1{\textcolor[rgb]{0.73,0.13,0.13}{##1}}}
\@namedef{PY@tok@si}{\let\PY@bf=\textbf\def\PY@tc##1{\textcolor[rgb]{0.64,0.35,0.47}{##1}}}
\@namedef{PY@tok@se}{\let\PY@bf=\textbf\def\PY@tc##1{\textcolor[rgb]{0.67,0.36,0.12}{##1}}}
\@namedef{PY@tok@sr}{\def\PY@tc##1{\textcolor[rgb]{0.64,0.35,0.47}{##1}}}
\@namedef{PY@tok@ss}{\def\PY@tc##1{\textcolor[rgb]{0.10,0.09,0.49}{##1}}}
\@namedef{PY@tok@sx}{\def\PY@tc##1{\textcolor[rgb]{0.00,0.50,0.00}{##1}}}
\@namedef{PY@tok@m}{\def\PY@tc##1{\textcolor[rgb]{0.40,0.40,0.40}{##1}}}
\@namedef{PY@tok@gh}{\let\PY@bf=\textbf\def\PY@tc##1{\textcolor[rgb]{0.00,0.00,0.50}{##1}}}
\@namedef{PY@tok@gu}{\let\PY@bf=\textbf\def\PY@tc##1{\textcolor[rgb]{0.50,0.00,0.50}{##1}}}
\@namedef{PY@tok@gd}{\def\PY@tc##1{\textcolor[rgb]{0.63,0.00,0.00}{##1}}}
\@namedef{PY@tok@gi}{\def\PY@tc##1{\textcolor[rgb]{0.00,0.52,0.00}{##1}}}
\@namedef{PY@tok@gr}{\def\PY@tc##1{\textcolor[rgb]{0.89,0.00,0.00}{##1}}}
\@namedef{PY@tok@ge}{\let\PY@it=\textit}
\@namedef{PY@tok@gs}{\let\PY@bf=\textbf}
\@namedef{PY@tok@gp}{\let\PY@bf=\textbf\def\PY@tc##1{\textcolor[rgb]{0.00,0.00,0.50}{##1}}}
\@namedef{PY@tok@go}{\def\PY@tc##1{\textcolor[rgb]{0.44,0.44,0.44}{##1}}}
\@namedef{PY@tok@gt}{\def\PY@tc##1{\textcolor[rgb]{0.00,0.27,0.87}{##1}}}
\@namedef{PY@tok@err}{\def\PY@bc##1{{\setlength{\fboxsep}{\string -\fboxrule}\fcolorbox[rgb]{1.00,0.00,0.00}{1,1,1}{\strut ##1}}}}
\@namedef{PY@tok@kc}{\let\PY@bf=\textbf\def\PY@tc##1{\textcolor[rgb]{0.00,0.50,0.00}{##1}}}
\@namedef{PY@tok@kd}{\let\PY@bf=\textbf\def\PY@tc##1{\textcolor[rgb]{0.00,0.50,0.00}{##1}}}
\@namedef{PY@tok@kn}{\let\PY@bf=\textbf\def\PY@tc##1{\textcolor[rgb]{0.00,0.50,0.00}{##1}}}
\@namedef{PY@tok@kr}{\let\PY@bf=\textbf\def\PY@tc##1{\textcolor[rgb]{0.00,0.50,0.00}{##1}}}
\@namedef{PY@tok@bp}{\def\PY@tc##1{\textcolor[rgb]{0.00,0.50,0.00}{##1}}}
\@namedef{PY@tok@fm}{\def\PY@tc##1{\textcolor[rgb]{0.00,0.00,1.00}{##1}}}
\@namedef{PY@tok@vc}{\def\PY@tc##1{\textcolor[rgb]{0.10,0.09,0.49}{##1}}}
\@namedef{PY@tok@vg}{\def\PY@tc##1{\textcolor[rgb]{0.10,0.09,0.49}{##1}}}
\@namedef{PY@tok@vi}{\def\PY@tc##1{\textcolor[rgb]{0.10,0.09,0.49}{##1}}}
\@namedef{PY@tok@vm}{\def\PY@tc##1{\textcolor[rgb]{0.10,0.09,0.49}{##1}}}
\@namedef{PY@tok@sa}{\def\PY@tc##1{\textcolor[rgb]{0.73,0.13,0.13}{##1}}}
\@namedef{PY@tok@sb}{\def\PY@tc##1{\textcolor[rgb]{0.73,0.13,0.13}{##1}}}
\@namedef{PY@tok@sc}{\def\PY@tc##1{\textcolor[rgb]{0.73,0.13,0.13}{##1}}}
\@namedef{PY@tok@dl}{\def\PY@tc##1{\textcolor[rgb]{0.73,0.13,0.13}{##1}}}
\@namedef{PY@tok@s2}{\def\PY@tc##1{\textcolor[rgb]{0.73,0.13,0.13}{##1}}}
\@namedef{PY@tok@sh}{\def\PY@tc##1{\textcolor[rgb]{0.73,0.13,0.13}{##1}}}
\@namedef{PY@tok@s1}{\def\PY@tc##1{\textcolor[rgb]{0.73,0.13,0.13}{##1}}}
\@namedef{PY@tok@mb}{\def\PY@tc##1{\textcolor[rgb]{0.40,0.40,0.40}{##1}}}
\@namedef{PY@tok@mf}{\def\PY@tc##1{\textcolor[rgb]{0.40,0.40,0.40}{##1}}}
\@namedef{PY@tok@mh}{\def\PY@tc##1{\textcolor[rgb]{0.40,0.40,0.40}{##1}}}
\@namedef{PY@tok@mi}{\def\PY@tc##1{\textcolor[rgb]{0.40,0.40,0.40}{##1}}}
\@namedef{PY@tok@il}{\def\PY@tc##1{\textcolor[rgb]{0.40,0.40,0.40}{##1}}}
\@namedef{PY@tok@mo}{\def\PY@tc##1{\textcolor[rgb]{0.40,0.40,0.40}{##1}}}
\@namedef{PY@tok@ch}{\let\PY@it=\textit\def\PY@tc##1{\textcolor[rgb]{0.24,0.48,0.48}{##1}}}
\@namedef{PY@tok@cm}{\let\PY@it=\textit\def\PY@tc##1{\textcolor[rgb]{0.24,0.48,0.48}{##1}}}
\@namedef{PY@tok@cpf}{\let\PY@it=\textit\def\PY@tc##1{\textcolor[rgb]{0.24,0.48,0.48}{##1}}}
\@namedef{PY@tok@c1}{\let\PY@it=\textit\def\PY@tc##1{\textcolor[rgb]{0.24,0.48,0.48}{##1}}}
\@namedef{PY@tok@cs}{\let\PY@it=\textit\def\PY@tc##1{\textcolor[rgb]{0.24,0.48,0.48}{##1}}}

\def\PYZbs{\char`\\}
\def\PYZus{\char`\_}
\def\PYZob{\char`\{}
\def\PYZcb{\char`\}}
\def\PYZca{\char`\^}
\def\PYZam{\char`\&}
\def\PYZlt{\char`\<}
\def\PYZgt{\char`\>}
\def\PYZsh{\char`\#}
\def\PYZpc{\char`\%}
\def\PYZdl{\char`\$}
\def\PYZhy{\char`\-}
\def\PYZsq{\char`\'}
\def\PYZdq{\char`\"}
\def\PYZti{\char`\~}
% for compatibility with earlier versions
\def\PYZat{@}
\def\PYZlb{[}
\def\PYZrb{]}
\makeatother


    % For linebreaks inside Verbatim environment from package fancyvrb.
    \makeatletter
        \newbox\Wrappedcontinuationbox
        \newbox\Wrappedvisiblespacebox
        \newcommand*\Wrappedvisiblespace {\textcolor{red}{\textvisiblespace}}
        \newcommand*\Wrappedcontinuationsymbol {\textcolor{red}{\llap{\tiny$\m@th\hookrightarrow$}}}
        \newcommand*\Wrappedcontinuationindent {3ex }
        \newcommand*\Wrappedafterbreak {\kern\Wrappedcontinuationindent\copy\Wrappedcontinuationbox}
        % Take advantage of the already applied Pygments mark-up to insert
        % potential linebreaks for TeX processing.
        %        {, <, #, %, $, ' and ": go to next line.
        %        _, }, ^, &, >, - and ~: stay at end of broken line.
        % Use of \textquotesingle for straight quote.
        \newcommand*\Wrappedbreaksatspecials {%
            \def\PYGZus{\discretionary{\char`\_}{\Wrappedafterbreak}{\char`\_}}%
            \def\PYGZob{\discretionary{}{\Wrappedafterbreak\char`\{}{\char`\{}}%
            \def\PYGZcb{\discretionary{\char`\}}{\Wrappedafterbreak}{\char`\}}}%
            \def\PYGZca{\discretionary{\char`\^}{\Wrappedafterbreak}{\char`\^}}%
            \def\PYGZam{\discretionary{\char`\&}{\Wrappedafterbreak}{\char`\&}}%
            \def\PYGZlt{\discretionary{}{\Wrappedafterbreak\char`\<}{\char`\<}}%
            \def\PYGZgt{\discretionary{\char`\>}{\Wrappedafterbreak}{\char`\>}}%
            \def\PYGZsh{\discretionary{}{\Wrappedafterbreak\char`\#}{\char`\#}}%
            \def\PYGZpc{\discretionary{}{\Wrappedafterbreak\char`\%}{\char`\%}}%
            \def\PYGZdl{\discretionary{}{\Wrappedafterbreak\char`\$}{\char`\$}}%
            \def\PYGZhy{\discretionary{\char`\-}{\Wrappedafterbreak}{\char`\-}}%
            \def\PYGZsq{\discretionary{}{\Wrappedafterbreak\textquotesingle}{\textquotesingle}}%
            \def\PYGZdq{\discretionary{}{\Wrappedafterbreak\char`\"}{\char`\"}}%
            \def\PYGZti{\discretionary{\char`\~}{\Wrappedafterbreak}{\char`\~}}%
        }
        % Some characters . , ; ? ! / are not pygmentized.
        % This macro makes them "active" and they will insert potential linebreaks
        \newcommand*\Wrappedbreaksatpunct {%
            \lccode`\~`\.\lowercase{\def~}{\discretionary{\hbox{\char`\.}}{\Wrappedafterbreak}{\hbox{\char`\.}}}%
            \lccode`\~`\,\lowercase{\def~}{\discretionary{\hbox{\char`\,}}{\Wrappedafterbreak}{\hbox{\char`\,}}}%
            \lccode`\~`\;\lowercase{\def~}{\discretionary{\hbox{\char`\;}}{\Wrappedafterbreak}{\hbox{\char`\;}}}%
            \lccode`\~`\:\lowercase{\def~}{\discretionary{\hbox{\char`\:}}{\Wrappedafterbreak}{\hbox{\char`\:}}}%
            \lccode`\~`\?\lowercase{\def~}{\discretionary{\hbox{\char`\?}}{\Wrappedafterbreak}{\hbox{\char`\?}}}%
            \lccode`\~`\!\lowercase{\def~}{\discretionary{\hbox{\char`\!}}{\Wrappedafterbreak}{\hbox{\char`\!}}}%
            \lccode`\~`\/\lowercase{\def~}{\discretionary{\hbox{\char`\/}}{\Wrappedafterbreak}{\hbox{\char`\/}}}%
            \catcode`\.\active
            \catcode`\,\active
            \catcode`\;\active
            \catcode`\:\active
            \catcode`\?\active
            \catcode`\!\active
            \catcode`\/\active
            \lccode`\~`\~
        }
    \makeatother

    \let\OriginalVerbatim=\Verbatim
    \makeatletter
    \renewcommand{\Verbatim}[1][1]{%
        %\parskip\z@skip
        \sbox\Wrappedcontinuationbox {\Wrappedcontinuationsymbol}%
        \sbox\Wrappedvisiblespacebox {\FV@SetupFont\Wrappedvisiblespace}%
        \def\FancyVerbFormatLine ##1{\hsize\linewidth
            \vtop{\raggedright\hyphenpenalty\z@\exhyphenpenalty\z@
                \doublehyphendemerits\z@\finalhyphendemerits\z@
                \strut ##1\strut}%
        }%
        % If the linebreak is at a space, the latter will be displayed as visible
        % space at end of first line, and a continuation symbol starts next line.
        % Stretch/shrink are however usually zero for typewriter font.
        \def\FV@Space {%
            \nobreak\hskip\z@ plus\fontdimen3\font minus\fontdimen4\font
            \discretionary{\copy\Wrappedvisiblespacebox}{\Wrappedafterbreak}
            {\kern\fontdimen2\font}%
        }%

        % Allow breaks at special characters using \PYG... macros.
        \Wrappedbreaksatspecials
        % Breaks at punctuation characters . , ; ? ! and / need catcode=\active
        \OriginalVerbatim[#1,codes*=\Wrappedbreaksatpunct]%
    }
    \makeatother

    % Exact colors from NB
    \definecolor{incolor}{HTML}{303F9F}
    \definecolor{outcolor}{HTML}{D84315}
    \definecolor{cellborder}{HTML}{CFCFCF}
    \definecolor{cellbackground}{HTML}{F7F7F7}

    % prompt
    \makeatletter
    \newcommand{\boxspacing}{\kern\kvtcb@left@rule\kern\kvtcb@boxsep}
    \makeatother
    \newcommand{\prompt}[4]{
        {\ttfamily\llap{{\color{#2}[#3]:\hspace{3pt}#4}}\vspace{-\baselineskip}}
    }
    

    
    % Prevent overflowing lines due to hard-to-break entities
    \sloppy
    % Setup hyperref package
    \hypersetup{
      breaklinks=true,  % so long urls are correctly broken across lines
      colorlinks=true,
      urlcolor=urlcolor,
      linkcolor=linkcolor,
      citecolor=citecolor,
      }
    % Slightly bigger margins than the latex defaults
    
    \geometry{verbose,tmargin=1in,bmargin=1in,lmargin=1in,rmargin=1in}
    
    

\begin{document}
    
    \maketitle
    
    

    
    \hypertarget{introduction}{%
\section{Introduction}\label{introduction}}

    \textbf{File:} Notebook\_BVP\_k\_M\_structure.ipynb\\
\textbf{Author:} Elizabeth Gould\\
\textbf{Date:} 15.03.2026\\
\textbf{Problem:} Find \(k_{\mu}\left(B, R, dk, \mu\right)\) and
\(M_{\mu}\left(B, R, dk, \mu\right)\)

    This Jupyter notebook is used for modeling the changes in our
wavenumbers \(k_j\) and \(M_j\) of the two oscillations of the solution
of the nonlinear Schrödinger equation for persistant current in
micrometer scale gold loops which solve the BVP requiring continuity of
the wavefunction at the beginning and end of the loop.

For our observed approximate solution to the nonlinear Schrödinger
equation \begin{equation}
\psi_j(x) = e^{i M_{j} x} \left( a_j e^{i k_{j} x} + b_j e^{- i k_{j} x} \right),
\end{equation} we know that \(\psi_j(0) = \psi_j(2\pi R_j)\) when
\begin{equation}
\psi_j^\prime(0) = i M_{j} + k_j \left( \frac{e^{- 2 i M_{j}\pi R} - \cos\left( 2\pi k_{j} R_j \right)}{\sin\left(2\pi k_{j} R_j\right)}\right),
\end{equation}

where \(k_{j}\left(\mu, R_j, B, k_0, \psi_j^\prime(0) \right)\) and
\(M_{j}\left(\mu, R_j, B, k_0, \psi_j^\prime(0) \right)\) are given by
our previously solved model. Since the equations for \(k_j\) and \(M_j\)
of our model are dependent on \(\psi_j^\prime(0)\), we have obtained the
values of \(\psi_j^\prime(0)\) which solve this BVP equation with a
numerical root finding algorithm. We now wish to reexamine our modeling
for \(k_j\) and \(M_j\), and figure out the structure of the changes of
\(k_{j}\left(\mu, R_j, B, k_0 \right)\) and
\(M_{j}\left(\mu, R_j, B, k_0\right)\) from their linear solutions for
the cases where \(k_j\) and \(M_j\) solve the BVP.

In our calculations, we vary \(R_j\), \(B\), \(dk\), and \(\mu\), with
\(k_0\) given by \(k_0 = k_{Fermi, Au} + \frac{dk}{2 R_{max}}\).

    \hypertarget{loading-grid}{%
\section{Loading Grid}\label{loading-grid}}

    To start with, we need to load our data into memory.

As typical, we start by loading our libraries.

    \begin{tcolorbox}[breakable, size=fbox, boxrule=1pt, pad at break*=1mm,colback=cellbackground, colframe=cellborder]
\prompt{In}{incolor}{1}{\boxspacing}
\begin{Verbatim}[commandchars=\\\{\}]
\PY{c+c1}{\PYZsh{}\PYZhy{}\PYZhy{}LIBRARIES\PYZhy{}\PYZhy{}\PYZhy{}\PYZhy{}\PYZhy{}\PYZhy{}\PYZhy{}\PYZhy{}}

\PY{k+kn}{import} \PY{n+nn}{eelib}
\PY{k+kn}{from} \PY{n+nn}{eelib} \PY{k+kn}{import} \PY{n}{pi}\PY{p}{,} \PY{n}{B\PYZus{}max}\PY{p}{,} \PY{n}{phi0inv}\PY{p}{,} \PY{n}{R\PYZus{}max}
\PY{k+kn}{from} \PY{n+nn}{eelib} \PY{k+kn}{import} \PY{n}{find\PYZus{}ind\PYZus{}var}\PY{p}{,} \PY{n}{clean\PYZus{}table}

\PY{k+kn}{import} \PY{n+nn}{numpy} \PY{k}{as} \PY{n+nn}{np}
\PY{k+kn}{import} \PY{n+nn}{pickle}
\PY{k+kn}{import} \PY{n+nn}{pandas} \PY{k}{as} \PY{n+nn}{pd}
\PY{k+kn}{from} \PY{n+nn}{sklearn}\PY{n+nn}{.}\PY{n+nn}{linear\PYZus{}model} \PY{k+kn}{import} \PY{n}{LinearRegression}
\PY{k+kn}{import} \PY{n+nn}{warnings}

\PY{c+c1}{\PYZsh{}plotting}
\PY{k+kn}{import} \PY{n+nn}{matplotlib}\PY{n+nn}{.}\PY{n+nn}{pyplot} \PY{k}{as} \PY{n+nn}{plt}
\PY{k+kn}{import} \PY{n+nn}{seaborn} \PY{k}{as} \PY{n+nn}{sns}
\end{Verbatim}
\end{tcolorbox}

    For this notebook, we will make a list of our grid objects.

    \begin{tcolorbox}[breakable, size=fbox, boxrule=1pt, pad at break*=1mm,colback=cellbackground, colframe=cellborder]
\prompt{In}{incolor}{2}{\boxspacing}
\begin{Verbatim}[commandchars=\\\{\}]
\PY{n}{gridl}\PY{o}{=}\PY{p}{[}\PY{p}{]}
\end{Verbatim}
\end{tcolorbox}

    And load each of our grids into the object, in order.

    \begin{tcolorbox}[breakable, size=fbox, boxrule=1pt, pad at break*=1mm,colback=cellbackground, colframe=cellborder]
\prompt{In}{incolor}{3}{\boxspacing}
\begin{Verbatim}[commandchars=\\\{\}]
\PY{c+c1}{\PYZsh{} loading script; be careful of python version}
\PY{n}{file} \PY{o}{=} \PY{n+nb}{open}\PY{p}{(}\PY{l+s+s1}{\PYZsq{}}\PY{l+s+s1}{grid030}\PY{l+s+s1}{\PYZsq{}}\PY{p}{,} \PY{l+s+s1}{\PYZsq{}}\PY{l+s+s1}{rb}\PY{l+s+s1}{\PYZsq{}}\PY{p}{)}    
\PY{n}{gridl}\PY{o}{.}\PY{n}{append}\PY{p}{(}\PY{n}{pickle}\PY{o}{.}\PY{n}{load}\PY{p}{(}\PY{n}{file}\PY{p}{)}\PY{p}{)}
\PY{n}{file}\PY{o}{.}\PY{n}{close}\PY{p}{(}\PY{p}{)}

\PY{c+c1}{\PYZsh{} loading script; be careful of python version}
\PY{n}{file} \PY{o}{=} \PY{n+nb}{open}\PY{p}{(}\PY{l+s+s1}{\PYZsq{}}\PY{l+s+s1}{grid033}\PY{l+s+s1}{\PYZsq{}}\PY{p}{,} \PY{l+s+s1}{\PYZsq{}}\PY{l+s+s1}{rb}\PY{l+s+s1}{\PYZsq{}}\PY{p}{)}    
\PY{n}{gridl}\PY{o}{.}\PY{n}{append}\PY{p}{(}\PY{n}{pickle}\PY{o}{.}\PY{n}{load}\PY{p}{(}\PY{n}{file}\PY{p}{)}\PY{p}{)}
\PY{n}{file}\PY{o}{.}\PY{n}{close}\PY{p}{(}\PY{p}{)}

\PY{c+c1}{\PYZsh{} loading script; be careful of python version}
\PY{n}{file} \PY{o}{=} \PY{n+nb}{open}\PY{p}{(}\PY{l+s+s1}{\PYZsq{}}\PY{l+s+s1}{grid040}\PY{l+s+s1}{\PYZsq{}}\PY{p}{,} \PY{l+s+s1}{\PYZsq{}}\PY{l+s+s1}{rb}\PY{l+s+s1}{\PYZsq{}}\PY{p}{)}    
\PY{n}{gridl}\PY{o}{.}\PY{n}{append}\PY{p}{(}\PY{n}{pickle}\PY{o}{.}\PY{n}{load}\PY{p}{(}\PY{n}{file}\PY{p}{)}\PY{p}{)}
\PY{n}{file}\PY{o}{.}\PY{n}{close}\PY{p}{(}\PY{p}{)}

\PY{c+c1}{\PYZsh{} loading script; be careful of python version}
\PY{n}{file} \PY{o}{=} \PY{n+nb}{open}\PY{p}{(}\PY{l+s+s1}{\PYZsq{}}\PY{l+s+s1}{grid041}\PY{l+s+s1}{\PYZsq{}}\PY{p}{,} \PY{l+s+s1}{\PYZsq{}}\PY{l+s+s1}{rb}\PY{l+s+s1}{\PYZsq{}}\PY{p}{)}    
\PY{n}{gridl}\PY{o}{.}\PY{n}{append}\PY{p}{(}\PY{n}{pickle}\PY{o}{.}\PY{n}{load}\PY{p}{(}\PY{n}{file}\PY{p}{)}\PY{p}{)}
\PY{n}{file}\PY{o}{.}\PY{n}{close}\PY{p}{(}\PY{p}{)}

\PY{c+c1}{\PYZsh{} loading script; be careful of python version}
\PY{n}{file} \PY{o}{=} \PY{n+nb}{open}\PY{p}{(}\PY{l+s+s1}{\PYZsq{}}\PY{l+s+s1}{grid026}\PY{l+s+s1}{\PYZsq{}}\PY{p}{,} \PY{l+s+s1}{\PYZsq{}}\PY{l+s+s1}{rb}\PY{l+s+s1}{\PYZsq{}}\PY{p}{)}    
\PY{n}{gridl}\PY{o}{.}\PY{n}{append}\PY{p}{(}\PY{n}{pickle}\PY{o}{.}\PY{n}{load}\PY{p}{(}\PY{n}{file}\PY{p}{)}\PY{p}{)}
\PY{n}{file}\PY{o}{.}\PY{n}{close}\PY{p}{(}\PY{p}{)}
\end{Verbatim}
\end{tcolorbox}

    \begin{tcolorbox}[breakable, size=fbox, boxrule=1pt, pad at break*=1mm,colback=cellbackground, colframe=cellborder]
\prompt{In}{incolor}{4}{\boxspacing}
\begin{Verbatim}[commandchars=\\\{\}]
\PY{k}{for} \PY{n}{ob} \PY{o+ow}{in} \PY{n}{gridl}\PY{p}{:} \PY{n+nb}{print}\PY{p}{(}\PY{n}{ob}\PY{p}{)}
\end{Verbatim}
\end{tcolorbox}

    \begin{Verbatim}[commandchars=\\\{\}]
Grid object to measure current:
mu is: 1.5e-08
dk is: 0.45
B has 100 points from 0.5 to 0.6.
R has 100 points from 0.9 to 0.9002.
A is: [1.]
k0 is: 12000000000.0
Total number of parameters: 10000
Grid object to measure current:
mu has 100 points from 3.1622776601683795e-09 to 3.162277660168379e-08.
dk has 100 points from 0.0 to 3.0.
B is: 0.5
R is: 0.9
A is: [1.]
k0 is: 12000000000.0
Total number of parameters: 10000
Grid object to measure current:
mu is: 1.5e-08
dk is: 0.45
B has 20 points from 0.5 to 0.6.
R has 511 points from 0.5709233333333332 to 0.9959233333333333.
A is: [1.]
k0 is: 12000000000.0
Total number of parameters: 10220
Grid object to measure current:
mu is: 1.5e-08
dk is: 0.45
B has 300 points from 0.11 to 1.0.
R has 30 points from 0.9 to 0.9002.
A is: [1.]
k0 is: 12000000000.0
Total number of parameters: 9000
Grid object to measure current:
mu has 11 points from 3.1622776601683795e-10 to 3.162277660168379e-08.
dk has 11 points from 0.1 to 0.9.
B has 11 points from 0.5 to 0.506.
R has 11 points from 0.9 to 0.9002.
A is: [1.]
k0 is: 12000000000.0
Total number of parameters: 14641
    \end{Verbatim}

    Here we are working with the following 5 objects:

\begin{enumerate}
\def\labelenumi{\arabic{enumi})}
\setcounter{enumi}{-1}
\tightlist
\item
  grid 30 - \(B\) and small \(R\) variations to find oscillation
  structure

  \begin{itemize}
  \tightlist
  \item
    mu is: 1.5e-08
  \item
    dk is: 0.45
  \item
    B has 100 points from 0.5 to 0.6.
  \item
    R has 100 points from 0.9 to 0.9002.
  \end{itemize}
\item
  grid 33 - \(dk\) and \(\mu\) variations, to examine effect of \(\mu\)
  and \(dk\)

  \begin{itemize}
  \tightlist
  \item
    mu has 100 points from 3.1622776601683795e-09 to
    3.162277660168379e-08.
  \item
    dk has 100 points from 0.0 to 3.0.
  \item
    B is: 0.5
  \item
    R is: 0.9
  \end{itemize}
\item
  grid 40 - full range of \(R\)

  \begin{itemize}
  \tightlist
  \item
    mu is: 1.5e-08
  \item
    dk is: 0.45
  \item
    B has 20 points from 0.5 to 0.6.
  \item
    R has 511 points from 0.5709233333333332 to 0.9959233333333333.
  \end{itemize}
\item
  grid 41 - full range of \(B\)

  \begin{itemize}
  \tightlist
  \item
    mu is: 1.5e-08
  \item
    dk is: 0.45
  \item
    B has 300 points from 0.11 to 1.0.
  \item
    R has 30 points from 0.9 to 0.9002.
  \end{itemize}
\item
  grid 26 - small grid varied over all 4 parameters

  \begin{itemize}
  \tightlist
  \item
    mu has 11 points from 3.1622776601683795e-10 to
    3.162277660168379e-08.
  \item
    dk has 11 points from 0.1 to 0.9.
  \item
    B has 11 points from 0.5 to 0.506.
  \item
    R has 11 points from 0.9 to 0.9002.
  \end{itemize}
\end{enumerate}

    I want to transfer all of my data to pandas dataframes for analysis.

    \begin{tcolorbox}[breakable, size=fbox, boxrule=1pt, pad at break*=1mm,colback=cellbackground, colframe=cellborder]
\prompt{In}{incolor}{5}{\boxspacing}
\begin{Verbatim}[commandchars=\\\{\}]
\PY{n}{tbl} \PY{o}{=} \PY{p}{[}\PY{p}{]}
\end{Verbatim}
\end{tcolorbox}

    \begin{tcolorbox}[breakable, size=fbox, boxrule=1pt, pad at break*=1mm,colback=cellbackground, colframe=cellborder]
\prompt{In}{incolor}{6}{\boxspacing}
\begin{Verbatim}[commandchars=\\\{\}]
\PY{k}{for} \PY{n}{grid} \PY{o+ow}{in} \PY{n}{gridl}\PY{p}{:}
    \PY{n}{tbl}\PY{o}{.}\PY{n}{append}\PY{p}{(}\PY{n}{pd}\PY{o}{.}\PY{n}{DataFrame}\PY{p}{(}\PY{n}{grid}\PY{o}{.}\PY{n}{derivs}\PY{p}{)}\PY{p}{)}
\end{Verbatim}
\end{tcolorbox}

    \begin{tcolorbox}[breakable, size=fbox, boxrule=1pt, pad at break*=1mm,colback=cellbackground, colframe=cellborder]
\prompt{In}{incolor}{7}{\boxspacing}
\begin{Verbatim}[commandchars=\\\{\}]
\PY{k}{del} \PY{n}{gridl} \PY{c+c1}{\PYZsh{} To clear up some memory. We have all the data already stored in the tables.}
\end{Verbatim}
\end{tcolorbox}

    \begin{tcolorbox}[breakable, size=fbox, boxrule=1pt, pad at break*=1mm,colback=cellbackground, colframe=cellborder]
\prompt{In}{incolor}{8}{\boxspacing}
\begin{Verbatim}[commandchars=\\\{\}]
\PY{n}{tbl}\PY{p}{[}\PY{l+m+mi}{1}\PY{p}{]}\PY{o}{.}\PY{n}{head}\PY{p}{(}\PY{p}{)} \PY{c+c1}{\PYZsh{} And check a grid.}
\end{Verbatim}
\end{tcolorbox}

            \begin{tcolorbox}[breakable, size=fbox, boxrule=.5pt, pad at break*=1mm, opacityfill=0]
\prompt{Out}{outcolor}{8}{\boxspacing}
\begin{Verbatim}[commandchars=\\\{\}]
     R    B        dk            mu             k    A  \textbackslash{}
0  0.9  0.5  0.000000  3.162278e-09  1.200000e+10  1.0
1  0.9  0.5  0.030303  3.162278e-09  1.200000e+10  1.0
2  0.9  0.5  0.060606  3.162278e-09  1.200000e+10  1.0
3  0.9  0.5  0.090909  3.162278e-09  1.200000e+10  1.0
4  0.9  0.5  0.121212  3.162278e-09  1.200000e+10  1.0

                                            dpsi0  Exact End  \textbackslash{}
0  8.863435e+09+4.613110e+                    07j   1.0+0.0j
1  5.136212e+08+4.700433e+                    08j   1.0-0.0j
2  1.030296e+09+2.381010e+                    08j   1.0+0.0j
3  1.550530e+09+1.611744e+                    08j   1.0+0.0j
4  2.076495e+09+1.230092e+                    08j   1.0+0.0j

                   a0                  b0  {\ldots}        I0  \textbackslash{}
0  0.501737-0.369303j  0.498263+0.369303j  {\ldots}  0.003474
1  0.519400-0.021401j  0.480600+0.021401j  {\ldots}  0.038800
2  0.509736-0.042929j  0.490264+0.042929j  {\ldots}  0.019471
3  0.506530-0.064605j  0.493470+0.064605j  {\ldots}  0.013061
4  0.504940-0.086520j  0.495060+0.086520j  {\ldots}  0.009880

                                             dpsi   New End  \textbackslash{}
0  8.452251e+09+4.579073e+                    07j  1.0+0.0j
1  3.322678e+08+5.318948e+                    08j  1.0-0.0j
2  8.474647e+08+2.520906e+                    08j  1.0-0.0j
3  1.364701e+09+1.670318e+                    08j  1.0-0.0j
4  1.886200e+09+1.261143e+                    08j  1.0-0.0j

                    a                   b  A max new      I v1      I v2  \textbackslash{}
0  0.501723-0.352171j  0.498277+0.352171j   1.223152  0.003446  0.003446
1  0.521977-0.013844j  0.478023+0.013844j   1.000384  0.043954  0.043954
2  0.510319-0.035311j  0.489681+0.035311j   1.002492  0.020637  0.020637
3  0.506774-0.056862j  0.493226+0.056862j   1.006447  0.013549  0.013549
4  0.505070-0.078591j  0.494930+0.078591j   1.012279  0.010139  0.010139

                 I v3  effective mu
0  0.003446+0.000000j  4.731084e-09
1  0.043954+0.000000j  3.164707e-09
2  0.020637+0.000000j  3.178056e-09
3  0.013549+0.000000j  3.203184e-09
4  0.010139+0.000000j  3.240414e-09

[5 rows x 21 columns]
\end{Verbatim}
\end{tcolorbox}
        
    \hypertarget{defining-variables-for-analysis}{%
\section{Defining variables for
analysis}\label{defining-variables-for-analysis}}

    I have all my data in table form, but I know that it was all produced
from grids of values. I therefore want to load my potential parameter
values into arrays, so that I can more easily select all data of a given
value of a given parameter.

    \begin{tcolorbox}[breakable, size=fbox, boxrule=1pt, pad at break*=1mm,colback=cellbackground, colframe=cellborder]
\prompt{In}{incolor}{9}{\boxspacing}
\begin{Verbatim}[commandchars=\\\{\}]
\PY{n}{R\PYZus{}g} \PY{o}{=} \PY{p}{[}\PY{p}{]}
\PY{n}{B\PYZus{}g} \PY{o}{=} \PY{p}{[}\PY{p}{]}
\PY{n}{dk\PYZus{}g} \PY{o}{=} \PY{p}{[}\PY{p}{]}
\PY{n}{mu\PYZus{}g} \PY{o}{=} \PY{p}{[}\PY{p}{]}
\end{Verbatim}
\end{tcolorbox}

    \begin{tcolorbox}[breakable, size=fbox, boxrule=1pt, pad at break*=1mm,colback=cellbackground, colframe=cellborder]
\prompt{In}{incolor}{10}{\boxspacing}
\begin{Verbatim}[commandchars=\\\{\}]
\PY{k}{for} \PY{n}{df} \PY{o+ow}{in} \PY{n}{tbl}\PY{p}{:}
    \PY{n}{R\PYZus{}g}\PY{o}{.}\PY{n}{append}\PY{p}{(}\PY{n}{df}\PY{p}{[}\PY{l+s+s2}{\PYZdq{}}\PY{l+s+s2}{R}\PY{l+s+s2}{\PYZdq{}}\PY{p}{]}\PY{o}{.}\PY{n}{unique}\PY{p}{(}\PY{p}{)}\PY{p}{)}
    \PY{n}{B\PYZus{}g}\PY{o}{.}\PY{n}{append}\PY{p}{(}\PY{n}{df}\PY{p}{[}\PY{l+s+s2}{\PYZdq{}}\PY{l+s+s2}{B}\PY{l+s+s2}{\PYZdq{}}\PY{p}{]}\PY{o}{.}\PY{n}{unique}\PY{p}{(}\PY{p}{)}\PY{p}{)}
    \PY{n}{dk\PYZus{}g}\PY{o}{.}\PY{n}{append}\PY{p}{(}\PY{n}{df}\PY{p}{[}\PY{l+s+s2}{\PYZdq{}}\PY{l+s+s2}{dk}\PY{l+s+s2}{\PYZdq{}}\PY{p}{]}\PY{o}{.}\PY{n}{unique}\PY{p}{(}\PY{p}{)}\PY{p}{)}
    \PY{n}{mu\PYZus{}g}\PY{o}{.}\PY{n}{append}\PY{p}{(}\PY{n}{df}\PY{p}{[}\PY{l+s+s2}{\PYZdq{}}\PY{l+s+s2}{mu}\PY{l+s+s2}{\PYZdq{}}\PY{p}{]}\PY{o}{.}\PY{n}{unique}\PY{p}{(}\PY{p}{)}\PY{p}{)}
\end{Verbatim}
\end{tcolorbox}

    \begin{tcolorbox}[breakable, size=fbox, boxrule=1pt, pad at break*=1mm,colback=cellbackground, colframe=cellborder]
\prompt{In}{incolor}{11}{\boxspacing}
\begin{Verbatim}[commandchars=\\\{\}]
\PY{n}{k0}   \PY{o}{=} \PY{n}{tbl}\PY{p}{[}\PY{l+m+mi}{0}\PY{p}{]}\PY{p}{[}\PY{l+s+s2}{\PYZdq{}}\PY{l+s+s2}{k}\PY{l+s+s2}{\PYZdq{}}\PY{p}{]}\PY{o}{.}\PY{n}{unique}\PY{p}{(}\PY{p}{)}\PY{p}{[}\PY{l+m+mi}{0}\PY{p}{]}
\PY{n}{A}    \PY{o}{=} \PY{n}{tbl}\PY{p}{[}\PY{l+m+mi}{0}\PY{p}{]}\PY{p}{[}\PY{l+s+s2}{\PYZdq{}}\PY{l+s+s2}{A}\PY{l+s+s2}{\PYZdq{}}\PY{p}{]}\PY{o}{.}\PY{n}{unique}\PY{p}{(}\PY{p}{)}\PY{p}{[}\PY{l+m+mi}{0}\PY{p}{]}
\end{Verbatim}
\end{tcolorbox}

    \begin{tcolorbox}[breakable, size=fbox, boxrule=1pt, pad at break*=1mm,colback=cellbackground, colframe=cellborder]
\prompt{In}{incolor}{12}{\boxspacing}
\begin{Verbatim}[commandchars=\\\{\}]
\PY{c+c1}{\PYZsh{} The first 2 grids have two varied parameters of length 100.}
\PY{n}{grid\PYZus{}size} \PY{o}{=} \PY{n}{R\PYZus{}g}\PY{p}{[}\PY{l+m+mi}{0}\PY{p}{]}\PY{o}{.}\PY{n}{shape}\PY{p}{[}\PY{l+m+mi}{0}\PY{p}{]}
\PY{c+c1}{\PYZsh{} For the 3rd grid, my grid sizes are 20 for B and 511 for R.}
\PY{c+c1}{\PYZsh{} My fourth grid has 300 points for B and 30 for R. }
\PY{c+c1}{\PYZsh{} I need to save these lengths separately.}
\PY{n}{gs\PYZus{}r1} \PY{o}{=} \PY{n}{R\PYZus{}g}\PY{p}{[}\PY{l+m+mi}{2}\PY{p}{]}\PY{o}{.}\PY{n}{shape}\PY{p}{[}\PY{l+m+mi}{0}\PY{p}{]}
\PY{n}{gs\PYZus{}r2} \PY{o}{=} \PY{n}{R\PYZus{}g}\PY{p}{[}\PY{l+m+mi}{3}\PY{p}{]}\PY{o}{.}\PY{n}{shape}\PY{p}{[}\PY{l+m+mi}{0}\PY{p}{]}
\PY{n}{gs\PYZus{}b1} \PY{o}{=} \PY{n}{B\PYZus{}g}\PY{p}{[}\PY{l+m+mi}{2}\PY{p}{]}\PY{o}{.}\PY{n}{shape}\PY{p}{[}\PY{l+m+mi}{0}\PY{p}{]}
\PY{n}{gs\PYZus{}b2} \PY{o}{=} \PY{n}{B\PYZus{}g}\PY{p}{[}\PY{l+m+mi}{3}\PY{p}{]}\PY{o}{.}\PY{n}{shape}\PY{p}{[}\PY{l+m+mi}{0}\PY{p}{]}
\PY{c+c1}{\PYZsh{} The last grid has a grid size of 11 for all 4 parameters.}
\PY{n}{gs2} \PY{o}{=} \PY{n}{R\PYZus{}g}\PY{p}{[}\PY{l+m+mi}{4}\PY{p}{]}\PY{o}{.}\PY{n}{shape}\PY{p}{[}\PY{l+m+mi}{0}\PY{p}{]}
\end{Verbatim}
\end{tcolorbox}

    \begin{tcolorbox}[breakable, size=fbox, boxrule=1pt, pad at break*=1mm,colback=cellbackground, colframe=cellborder]
\prompt{In}{incolor}{13}{\boxspacing}
\begin{Verbatim}[commandchars=\\\{\}]
\PY{c+c1}{\PYZsh{} display grid sizes, to check that they are correct}
\PY{n}{grid\PYZus{}size}\PY{p}{,} \PY{n}{gs\PYZus{}r1}\PY{p}{,} \PY{n}{gs\PYZus{}r2}\PY{p}{,} \PY{n}{gs\PYZus{}b1}\PY{p}{,} \PY{n}{gs\PYZus{}b2}\PY{p}{,} \PY{n}{gs2}
\end{Verbatim}
\end{tcolorbox}

            \begin{tcolorbox}[breakable, size=fbox, boxrule=.5pt, pad at break*=1mm, opacityfill=0]
\prompt{Out}{outcolor}{13}{\boxspacing}
\begin{Verbatim}[commandchars=\\\{\}]
(100, 511, 30, 20, 300, 11)
\end{Verbatim}
\end{tcolorbox}
        
    Ok, I saved my array sizes. Now I want to check how well I have found a
solution to the BVP. My tolerance here is 1 \%, so all of the nonlinear
solutions should give an error of around 1 \%.

    \begin{tcolorbox}[breakable, size=fbox, boxrule=1pt, pad at break*=1mm,colback=cellbackground, colframe=cellborder]
\prompt{In}{incolor}{14}{\boxspacing}
\begin{Verbatim}[commandchars=\\\{\}]
\PY{k}{for} \PY{n}{df} \PY{o+ow}{in} \PY{n}{tbl}\PY{p}{:}
    \PY{n+nb}{print}\PY{p}{(}\PY{l+s+s2}{\PYZdq{}}\PY{l+s+s2}{Max error in new endpoint:}\PY{l+s+s2}{\PYZdq{}}\PY{p}{,} \PY{n}{np}\PY{o}{.}\PY{n}{max}\PY{p}{(}\PY{n}{np}\PY{o}{.}\PY{n}{abs}\PY{p}{(}\PY{n}{df}\PY{p}{[}\PY{l+s+s2}{\PYZdq{}}\PY{l+s+s2}{New End}\PY{l+s+s2}{\PYZdq{}}\PY{p}{]}\PY{o}{.}\PY{n}{to\PYZus{}numpy}\PY{p}{(}\PY{p}{)}\PY{o}{\PYZhy{}}\PY{p}{(}\PY{l+m+mf}{1.0}\PY{o}{+}\PY{l+m+mf}{0.0}\PY{n}{j}\PY{p}{)}\PY{p}{)}\PY{p}{)}\PY{p}{)}
    \PY{n+nb}{print}\PY{p}{(}\PY{l+s+s2}{\PYZdq{}}\PY{l+s+s2}{Max error in old endpoint:}\PY{l+s+s2}{\PYZdq{}}\PY{p}{,} \PY{n}{np}\PY{o}{.}\PY{n}{max}\PY{p}{(}\PY{n}{np}\PY{o}{.}\PY{n}{abs}\PY{p}{(}\PY{n}{df}\PY{p}{[}\PY{l+s+s2}{\PYZdq{}}\PY{l+s+s2}{Exact End}\PY{l+s+s2}{\PYZdq{}}\PY{p}{]}\PY{o}{.}\PY{n}{to\PYZus{}numpy}\PY{p}{(}\PY{p}{)}\PY{o}{\PYZhy{}}\PY{p}{(}\PY{l+m+mf}{1.0}\PY{o}{+}\PY{l+m+mf}{0.0}\PY{n}{j}\PY{p}{)}\PY{p}{)}\PY{p}{)}\PY{p}{)}
\end{Verbatim}
\end{tcolorbox}

    \begin{Verbatim}[commandchars=\\\{\}]
Max error in new endpoint: 0.009995294828413707
Max error in old endpoint: 9.668703732403756e-10
Max error in new endpoint: 0.00999759070608627
Max error in old endpoint: 1.1897410141962145e-10
Max error in new endpoint: 0.009998719201844994
Max error in old endpoint: 2.6063175770843754e-11
Max error in new endpoint: 0.009998713933494732
Max error in old endpoint: 2.27591748523702e-10
Max error in new endpoint: 0.009999028548925254
Max error in old endpoint: 2.1603538026081975e-09
    \end{Verbatim}

    Now lets clean the tables of excess data.

    \begin{tcolorbox}[breakable, size=fbox, boxrule=1pt, pad at break*=1mm,colback=cellbackground, colframe=cellborder]
\prompt{In}{incolor}{15}{\boxspacing}
\begin{Verbatim}[commandchars=\\\{\}]
\PY{n}{tbl\PYZus{}red} \PY{o}{=} \PY{p}{[}\PY{p}{]}
\PY{k}{for} \PY{n}{df} \PY{o+ow}{in} \PY{n}{tbl}\PY{p}{:}
    \PY{n}{tbl\PYZus{}red}\PY{o}{.}\PY{n}{append}\PY{p}{(}\PY{n}{df}\PY{p}{[}\PY{p}{[}\PY{l+s+s2}{\PYZdq{}}\PY{l+s+s2}{R}\PY{l+s+s2}{\PYZdq{}}\PY{p}{,} \PY{l+s+s2}{\PYZdq{}}\PY{l+s+s2}{B}\PY{l+s+s2}{\PYZdq{}}\PY{p}{,} \PY{l+s+s2}{\PYZdq{}}\PY{l+s+s2}{dk}\PY{l+s+s2}{\PYZdq{}}\PY{p}{,} \PY{l+s+s2}{\PYZdq{}}\PY{l+s+s2}{mu}\PY{l+s+s2}{\PYZdq{}}\PY{p}{,} \PY{l+s+s1}{\PYZsq{}}\PY{l+s+s1}{dpsi0}\PY{l+s+s1}{\PYZsq{}}\PY{p}{,} \PY{l+s+s1}{\PYZsq{}}\PY{l+s+s1}{a0}\PY{l+s+s1}{\PYZsq{}}\PY{p}{,} \PY{l+s+s1}{\PYZsq{}}\PY{l+s+s1}{b0}\PY{l+s+s1}{\PYZsq{}}\PY{p}{,} \PY{l+s+s1}{\PYZsq{}}\PY{l+s+s1}{A max 0}\PY{l+s+s1}{\PYZsq{}}\PY{p}{,} \PY{l+s+s1}{\PYZsq{}}\PY{l+s+s1}{I0}\PY{l+s+s1}{\PYZsq{}}\PY{p}{,} \PY{l+s+s1}{\PYZsq{}}\PY{l+s+s1}{dpsi}\PY{l+s+s1}{\PYZsq{}}\PY{p}{,} \PY{l+s+s1}{\PYZsq{}}\PY{l+s+s1}{a}\PY{l+s+s1}{\PYZsq{}}\PY{p}{,} \PY{l+s+s1}{\PYZsq{}}\PY{l+s+s1}{b}\PY{l+s+s1}{\PYZsq{}}\PY{p}{,} \PY{l+s+s1}{\PYZsq{}}\PY{l+s+s1}{A max new}\PY{l+s+s1}{\PYZsq{}}\PY{p}{,} \PY{l+s+s1}{\PYZsq{}}\PY{l+s+s1}{I v2}\PY{l+s+s1}{\PYZsq{}}\PY{p}{,} \PY{l+s+s1}{\PYZsq{}}\PY{l+s+s1}{effective mu}\PY{l+s+s1}{\PYZsq{}}\PY{p}{]}\PY{p}{]}\PY{p}{)}
\end{Verbatim}
\end{tcolorbox}

    \begin{tcolorbox}[breakable, size=fbox, boxrule=1pt, pad at break*=1mm,colback=cellbackground, colframe=cellborder]
\prompt{In}{incolor}{16}{\boxspacing}
\begin{Verbatim}[commandchars=\\\{\}]
\PY{n}{tbl\PYZus{}cl} \PY{o}{=} \PY{p}{[}\PY{p}{]}

\PY{k}{for} \PY{n}{df} \PY{o+ow}{in} \PY{n}{tbl\PYZus{}red}\PY{p}{:}
    \PY{n}{tbl\PYZus{}cl}\PY{o}{.}\PY{n}{append}\PY{p}{(}\PY{n}{clean\PYZus{}table}\PY{p}{(}\PY{n}{df}\PY{p}{)}\PY{p}{)}
\end{Verbatim}
\end{tcolorbox}

    I need to check the length of my tables, as I have probably lost some
solutions, and this could cause problems with the rest of my code, which
may not know how to handle this.

    \begin{tcolorbox}[breakable, size=fbox, boxrule=1pt, pad at break*=1mm,colback=cellbackground, colframe=cellborder]
\prompt{In}{incolor}{17}{\boxspacing}
\begin{Verbatim}[commandchars=\\\{\}]
\PY{k}{for} \PY{n}{df} \PY{o+ow}{in} \PY{n}{tbl\PYZus{}cl}\PY{p}{:}
    \PY{n+nb}{print}\PY{p}{(}\PY{n}{df}\PY{o}{.}\PY{n}{shape}\PY{p}{)}
\end{Verbatim}
\end{tcolorbox}

    \begin{Verbatim}[commandchars=\\\{\}]
(10000, 15)
(10000, 15)
(10220, 15)
(9000, 15)
(14641, 15)
    \end{Verbatim}

    Nope. All good. Fitting with multiple algorithms and starting points
works well. Lets check for erroneous points on the graphs however.

    Now I want to calculate my \(k_j\) and \(M_j\) values, using my model
and my known initial derivatives of \(\psi\). The difference between the
linear and nonlinear solutions will be added to all of the tables.

    \begin{tcolorbox}[breakable, size=fbox, boxrule=1pt, pad at break*=1mm,colback=cellbackground, colframe=cellborder]
\prompt{In}{incolor}{18}{\boxspacing}
\begin{Verbatim}[commandchars=\\\{\}]
\PY{k}{del} \PY{n}{tbl\PYZus{}red} \PY{c+c1}{\PYZsh{} Clear old data.}
\PY{n}{tbl\PYZus{}red} \PY{o}{=} \PY{p}{[}\PY{p}{]} \PY{c+c1}{\PYZsh{} Reset our variable to an empty list.}
\end{Verbatim}
\end{tcolorbox}

    \begin{tcolorbox}[breakable, size=fbox, boxrule=1pt, pad at break*=1mm,colback=cellbackground, colframe=cellborder]
\prompt{In}{incolor}{ }{\boxspacing}
\begin{Verbatim}[commandchars=\\\{\}]
\PY{c+c1}{\PYZsh{} We are creating arrays in which to store our predicted nonlinear wave numbers.}
\PY{n}{kn} \PY{o}{=} \PY{p}{[}\PY{p}{]}
\PY{n}{Mn} \PY{o}{=} \PY{p}{[}\PY{p}{]}
\PY{k}{for} \PY{n}{i}\PY{p}{,} \PY{n}{df} \PY{o+ow}{in} \PY{n+nb}{enumerate}\PY{p}{(}\PY{n}{tbl\PYZus{}cl}\PY{p}{)}\PY{p}{:}
    \PY{n}{kn}\PY{o}{.}\PY{n}{append}\PY{p}{(}\PY{n}{np}\PY{o}{.}\PY{n}{zeros}\PY{p}{(}\PY{n}{df}\PY{o}{.}\PY{n}{shape}\PY{p}{[}\PY{l+m+mi}{0}\PY{p}{]}\PY{p}{)}\PY{p}{)}
    \PY{n}{Mn}\PY{o}{.}\PY{n}{append}\PY{p}{(}\PY{n}{np}\PY{o}{.}\PY{n}{zeros}\PY{p}{(}\PY{n}{df}\PY{o}{.}\PY{n}{shape}\PY{p}{[}\PY{l+m+mi}{0}\PY{p}{]}\PY{p}{)}\PY{p}{)}
\end{Verbatim}
\end{tcolorbox}

    \begin{tcolorbox}[breakable, size=fbox, boxrule=1pt, pad at break*=1mm,colback=cellbackground, colframe=cellborder]
\prompt{In}{incolor}{ }{\boxspacing}
\begin{Verbatim}[commandchars=\\\{\}]
\PY{c+c1}{\PYZsh{} For each row in each of our tables, I am adding a model prediction for }
\PY{c+c1}{\PYZsh{} both of our wave numbers. I then want to subtract the linear }
\PY{c+c1}{\PYZsh{} expectation for that wave number, and add a column to our table with the }
\PY{c+c1}{\PYZsh{} difference between linear and nonlinear wave number.}
\PY{k}{for} \PY{n}{i}\PY{p}{,} \PY{n}{df} \PY{o+ow}{in} \PY{n+nb}{enumerate}\PY{p}{(}\PY{n}{tbl\PYZus{}cl}\PY{p}{)}\PY{p}{:}
    \PY{k}{for} \PY{n}{iel} \PY{o+ow}{in} \PY{n+nb}{range}\PY{p}{(}\PY{n}{df}\PY{o}{.}\PY{n}{shape}\PY{p}{[}\PY{l+m+mi}{0}\PY{p}{]}\PY{p}{)}\PY{p}{:}
        \PY{n}{rr} \PY{o}{=} \PY{n}{df}\PY{p}{[}\PY{l+s+s2}{\PYZdq{}}\PY{l+s+s2}{R}\PY{l+s+s2}{\PYZdq{}}\PY{p}{]}\PY{o}{.}\PY{n}{iloc}\PY{p}{[}\PY{n}{iel}\PY{p}{]}
        \PY{n}{bb} \PY{o}{=} \PY{n}{df}\PY{p}{[}\PY{l+s+s2}{\PYZdq{}}\PY{l+s+s2}{B}\PY{l+s+s2}{\PYZdq{}}\PY{p}{]}\PY{o}{.}\PY{n}{iloc}\PY{p}{[}\PY{n}{iel}\PY{p}{]}
        \PY{n}{kk} \PY{o}{=} \PY{n}{df}\PY{p}{[}\PY{l+s+s2}{\PYZdq{}}\PY{l+s+s2}{dk}\PY{l+s+s2}{\PYZdq{}}\PY{p}{]}\PY{o}{.}\PY{n}{iloc}\PY{p}{[}\PY{n}{iel}\PY{p}{]}
        \PY{n}{mm} \PY{o}{=} \PY{n}{df}\PY{p}{[}\PY{l+s+s2}{\PYZdq{}}\PY{l+s+s2}{mu}\PY{l+s+s2}{\PYZdq{}}\PY{p}{]}\PY{o}{.}\PY{n}{iloc}\PY{p}{[}\PY{n}{iel}\PY{p}{]}
        
        \PY{n}{RR} \PY{o}{=} \PY{n}{rr} \PY{o}{*} \PY{n}{R\PYZus{}max}
        \PY{n}{BB} \PY{o}{=} \PY{n}{bb} \PY{o}{*} \PY{n}{B\PYZus{}max}
        
        \PY{n}{dpsi} \PY{o}{=} \PY{n}{df}\PY{p}{[}\PY{l+s+s1}{\PYZsq{}}\PY{l+s+s1}{dpsi}\PY{l+s+s1}{\PYZsq{}}\PY{p}{]}\PY{o}{.}\PY{n}{iloc}\PY{p}{[}\PY{n}{iel}\PY{p}{]}

        \PY{n}{k\PYZus{}new} \PY{o}{=} \PY{n}{eelib}\PY{o}{.}\PY{n}{pred\PYZus{}fast\PYZus{}k\PYZus{}true}\PY{p}{(}\PY{n}{dpsi}\PY{p}{,} \PY{n}{mm}\PY{p}{,} \PY{n}{kk}\PY{p}{,} \PY{n}{BB}\PY{p}{,} \PY{n}{RR}\PY{p}{,} \PY{n}{A}\PY{o}{=}\PY{n}{A}\PY{p}{,} \PY{n}{k0}\PY{o}{=}\PY{n}{k0}\PY{p}{)}        
        \PY{n}{M\PYZus{}new} \PY{o}{=} \PY{n}{eelib}\PY{o}{.}\PY{n}{pred\PYZus{}slow\PYZus{}k\PYZus{}v3}\PY{p}{(}\PY{n}{dpsi}\PY{p}{,} \PY{n}{mm}\PY{p}{,} \PY{n}{kk}\PY{p}{,} \PY{n}{BB}\PY{p}{,} \PY{n}{RR}\PY{p}{,} \PY{n}{A}\PY{o}{=}\PY{n}{A}\PY{p}{,} \PY{n}{k0}\PY{o}{=}\PY{n}{k0}\PY{p}{)}

        \PY{n}{dk\PYZus{}new} \PY{o}{=} \PY{n}{k\PYZus{}new} \PY{o}{\PYZhy{}} \PY{p}{(}\PY{n}{k0}\PY{o}{+}\PY{n}{kk}\PY{o}{/}\PY{l+m+mf}{1e\PYZhy{}6}\PY{o}{/}\PY{l+m+mf}{2.0}\PY{p}{)}
        \PY{n}{dM\PYZus{}new} \PY{o}{=} \PY{n}{M\PYZus{}new} \PY{o}{\PYZhy{}} \PY{n}{rr}\PY{o}{*}\PY{n}{R\PYZus{}max}\PY{o}{*}\PY{n}{bb}\PY{o}{*}\PY{n}{B\PYZus{}max}\PY{o}{*}\PY{n}{phi0inv}

        \PY{n}{kn}\PY{p}{[}\PY{n}{i}\PY{p}{]}\PY{p}{[}\PY{n}{iel}\PY{p}{]} \PY{o}{=} \PY{n}{dk\PYZus{}new}
        \PY{n}{Mn}\PY{p}{[}\PY{n}{i}\PY{p}{]}\PY{p}{[}\PY{n}{iel}\PY{p}{]} \PY{o}{=} \PY{n}{dM\PYZus{}new}

    \PY{n}{tbl\PYZus{}cl}\PY{p}{[}\PY{n}{i}\PY{p}{]}\PY{p}{[}\PY{l+s+s2}{\PYZdq{}}\PY{l+s+s2}{dkn}\PY{l+s+s2}{\PYZdq{}}\PY{p}{]} \PY{o}{=} \PY{n}{kn}\PY{p}{[}\PY{n}{i}\PY{p}{]}
    \PY{n}{tbl\PYZus{}cl}\PY{p}{[}\PY{n}{i}\PY{p}{]}\PY{p}{[}\PY{l+s+s2}{\PYZdq{}}\PY{l+s+s2}{dMn}\PY{l+s+s2}{\PYZdq{}}\PY{p}{]} \PY{o}{=} \PY{n}{Mn}\PY{p}{[}\PY{n}{i}\PY{p}{]}
\end{Verbatim}
\end{tcolorbox}

    \begin{tcolorbox}[breakable, size=fbox, boxrule=1pt, pad at break*=1mm,colback=cellbackground, colframe=cellborder]
\prompt{In}{incolor}{ }{\boxspacing}
\begin{Verbatim}[commandchars=\\\{\}]
\PY{c+c1}{\PYZsh{} This determines which of our table columns we wish to save.}
\PY{n}{cols} \PY{o}{=} \PY{p}{[}\PY{p}{]}
\PY{k}{for} \PY{n}{i}\PY{p}{,} \PY{n}{df} \PY{o+ow}{in} \PY{n+nb}{enumerate}\PY{p}{(}\PY{n}{tbl\PYZus{}cl}\PY{p}{)}\PY{p}{:}
    \PY{n}{col\PYZus{}names} \PY{o}{=} \PY{n}{find\PYZus{}ind\PYZus{}var}\PY{p}{(}\PY{n}{i}\PY{p}{,} \PY{n}{R\PYZus{}g}\PY{p}{,} \PY{n}{B\PYZus{}g}\PY{p}{,} \PY{n}{dk\PYZus{}g}\PY{p}{,} \PY{n}{mu\PYZus{}g}\PY{p}{)} \PY{c+c1}{\PYZsh{} To determine which parameters to plot.}
    \PY{n}{col\PYZus{}names}\PY{o}{.}\PY{n}{append}\PY{p}{(}\PY{l+s+s2}{\PYZdq{}}\PY{l+s+s2}{dkn}\PY{l+s+s2}{\PYZdq{}}\PY{p}{)} \PY{c+c1}{\PYZsh{} Also plot the change in fast oscillation wave number.}
    \PY{n}{col\PYZus{}names}\PY{o}{.}\PY{n}{append}\PY{p}{(}\PY{l+s+s2}{\PYZdq{}}\PY{l+s+s2}{dMn}\PY{l+s+s2}{\PYZdq{}}\PY{p}{)} \PY{c+c1}{\PYZsh{} Also plot the change in slow oscillation wave number.}
    \PY{n}{cols}\PY{o}{.}\PY{n}{append}\PY{p}{(}\PY{n}{col\PYZus{}names}\PY{p}{)}
\end{Verbatim}
\end{tcolorbox}

    \begin{tcolorbox}[breakable, size=fbox, boxrule=1pt, pad at break*=1mm,colback=cellbackground, colframe=cellborder]
\prompt{In}{incolor}{22}{\boxspacing}
\begin{Verbatim}[commandchars=\\\{\}]
\PY{c+c1}{\PYZsh{} Our reduced tables will have only the columns which are useful for our analysis.}
\PY{k}{for} \PY{n}{i}\PY{p}{,}\PY{n}{df} \PY{o+ow}{in} \PY{n+nb}{enumerate}\PY{p}{(}\PY{n}{tbl\PYZus{}cl}\PY{p}{)}\PY{p}{:}
    \PY{n}{tbl\PYZus{}red}\PY{o}{.}\PY{n}{append}\PY{p}{(}\PY{n}{df}\PY{p}{[}\PY{n}{cols}\PY{p}{[}\PY{n}{i}\PY{p}{]}\PY{p}{]}\PY{p}{)}
\end{Verbatim}
\end{tcolorbox}

    Ok, now that we have cleaned up datasets, lets move on to visualization,
and try to figure out with what we are dealing.

    \hypertarget{visualization}{%
\section{Visualization}\label{visualization}}

    Now lets look at the triangle plots for all of our datasets.

    \begin{tcolorbox}[breakable, size=fbox, boxrule=1pt, pad at break*=1mm,colback=cellbackground, colframe=cellborder]
\prompt{In}{incolor}{23}{\boxspacing}
\begin{Verbatim}[commandchars=\\\{\}]
\PY{n}{g} \PY{o}{=} \PY{n}{sns}\PY{o}{.}\PY{n}{PairGrid}\PY{p}{(}\PY{n}{tbl\PYZus{}red}\PY{p}{[}\PY{l+m+mi}{0}\PY{p}{]}\PY{p}{,} \PY{n}{diag\PYZus{}sharey}\PY{o}{=}\PY{k+kc}{False}\PY{p}{,} \PY{n}{corner}\PY{o}{=}\PY{k+kc}{True}\PY{p}{)}

\PY{n}{g}\PY{o}{.}\PY{n}{map\PYZus{}diag}\PY{p}{(}\PY{n}{sns}\PY{o}{.}\PY{n}{kdeplot}\PY{p}{)}
\PY{n}{g}\PY{o}{.}\PY{n}{map\PYZus{}lower}\PY{p}{(}\PY{n}{sns}\PY{o}{.}\PY{n}{scatterplot}\PY{p}{)}
\end{Verbatim}
\end{tcolorbox}

            \begin{tcolorbox}[breakable, size=fbox, boxrule=.5pt, pad at break*=1mm, opacityfill=0]
\prompt{Out}{outcolor}{23}{\boxspacing}
\begin{Verbatim}[commandchars=\\\{\}]
<seaborn.axisgrid.PairGrid at 0x1d19b929070>
\end{Verbatim}
\end{tcolorbox}
        
    \begin{center}
    \adjustimage{max size={0.9\linewidth}{0.9\paperheight}}{figs/bvp_kM_en/output_41_1.png}
    \end{center}
    { \hspace*{\fill} \\}
    
    \begin{tcolorbox}[breakable, size=fbox, boxrule=1pt, pad at break*=1mm,colback=cellbackground, colframe=cellborder]
\prompt{In}{incolor}{24}{\boxspacing}
\begin{Verbatim}[commandchars=\\\{\}]
\PY{c+c1}{\PYZsh{} I want to delete the outliers from my dataset to make the pattern more visible.}
\PY{n}{tbl\PYZus{}red}\PY{p}{[}\PY{l+m+mi}{0}\PY{p}{]}\PY{o}{.}\PY{n}{loc}\PY{p}{[}\PY{n}{tbl\PYZus{}red}\PY{p}{[}\PY{l+m+mi}{0}\PY{p}{]}\PY{p}{[}\PY{l+s+s2}{\PYZdq{}}\PY{l+s+s2}{dkn}\PY{l+s+s2}{\PYZdq{}}\PY{p}{]}\PY{o}{\PYZlt{}} \PY{o}{\PYZhy{}}\PY{l+m+mf}{0.75e6}\PY{p}{,}\PY{l+s+s2}{\PYZdq{}}\PY{l+s+s2}{dkn}\PY{l+s+s2}{\PYZdq{}}\PY{p}{]}\PY{o}{=} \PY{n}{np}\PY{o}{.}\PY{n}{nan}
\PY{n}{tbl\PYZus{}red}\PY{p}{[}\PY{l+m+mi}{0}\PY{p}{]}\PY{o}{.}\PY{n}{loc}\PY{p}{[}\PY{n}{tbl\PYZus{}red}\PY{p}{[}\PY{l+m+mi}{0}\PY{p}{]}\PY{p}{[}\PY{l+s+s2}{\PYZdq{}}\PY{l+s+s2}{dMn}\PY{l+s+s2}{\PYZdq{}}\PY{p}{]}\PY{o}{\PYZgt{}} \PY{l+m+mi}{50000}\PY{p}{,}\PY{l+s+s2}{\PYZdq{}}\PY{l+s+s2}{dMn}\PY{l+s+s2}{\PYZdq{}}\PY{p}{]}\PY{o}{=} \PY{n}{np}\PY{o}{.}\PY{n}{nan}
\end{Verbatim}
\end{tcolorbox}

    \begin{tcolorbox}[breakable, size=fbox, boxrule=1pt, pad at break*=1mm,colback=cellbackground, colframe=cellborder]
\prompt{In}{incolor}{25}{\boxspacing}
\begin{Verbatim}[commandchars=\\\{\}]
\PY{n}{g} \PY{o}{=} \PY{n}{sns}\PY{o}{.}\PY{n}{PairGrid}\PY{p}{(}\PY{n}{tbl\PYZus{}red}\PY{p}{[}\PY{l+m+mi}{0}\PY{p}{]}\PY{p}{,} \PY{n}{diag\PYZus{}sharey}\PY{o}{=}\PY{k+kc}{False}\PY{p}{,} \PY{n}{corner}\PY{o}{=}\PY{k+kc}{True}\PY{p}{)}

\PY{n}{g}\PY{o}{.}\PY{n}{map\PYZus{}diag}\PY{p}{(}\PY{n}{sns}\PY{o}{.}\PY{n}{kdeplot}\PY{p}{)}
\PY{n}{g}\PY{o}{.}\PY{n}{map\PYZus{}lower}\PY{p}{(}\PY{n}{sns}\PY{o}{.}\PY{n}{scatterplot}\PY{p}{)}
\end{Verbatim}
\end{tcolorbox}

            \begin{tcolorbox}[breakable, size=fbox, boxrule=.5pt, pad at break*=1mm, opacityfill=0]
\prompt{Out}{outcolor}{25}{\boxspacing}
\begin{Verbatim}[commandchars=\\\{\}]
<seaborn.axisgrid.PairGrid at 0x1d1a892cd10>
\end{Verbatim}
\end{tcolorbox}
        
    \begin{center}
    \adjustimage{max size={0.9\linewidth}{0.9\paperheight}}{figs/bvp_kM_en/output_43_1.png}
    \end{center}
    { \hspace*{\fill} \\}
    
    Oscillatory behavior clearly visible. It looks like our \(B\) range has
not completed an oscillation. I see for \(R\) vs \(dkn\) some sort of
tilted secant squared pattern, while for \(B\) vs dMn, I see a tilted
sine wave.

    These are the equations I found for the skewed sine wave, which has the
sawtooth pattern as its limit. They may be useful later on for fitting
the data to a functional form.

\begin{equation}
\frac{2^{2n-1}}{\binom{2n}{n}}\cos^{2n}\left(\frac{x}2\right)-\frac12
\end{equation}

\begin{equation}
\sum_{k=1}^n\frac{\binom{2n}{n-k}}{\binom{2n}{n}}\frac{\sin(kx)}k
\end{equation}

    \begin{tcolorbox}[breakable, size=fbox, boxrule=1pt, pad at break*=1mm,colback=cellbackground, colframe=cellborder]
\prompt{In}{incolor}{26}{\boxspacing}
\begin{Verbatim}[commandchars=\\\{\}]
\PY{n}{g} \PY{o}{=} \PY{n}{sns}\PY{o}{.}\PY{n}{PairGrid}\PY{p}{(}\PY{n}{tbl\PYZus{}red}\PY{p}{[}\PY{l+m+mi}{1}\PY{p}{]}\PY{p}{,} \PY{n}{diag\PYZus{}sharey}\PY{o}{=}\PY{k+kc}{False}\PY{p}{,} \PY{n}{corner}\PY{o}{=}\PY{k+kc}{True}\PY{p}{)}

\PY{n}{g}\PY{o}{.}\PY{n}{map\PYZus{}diag}\PY{p}{(}\PY{n}{sns}\PY{o}{.}\PY{n}{kdeplot}\PY{p}{)}
\PY{n}{g}\PY{o}{.}\PY{n}{map\PYZus{}lower}\PY{p}{(}\PY{n}{sns}\PY{o}{.}\PY{n}{scatterplot}\PY{p}{)}
\end{Verbatim}
\end{tcolorbox}

            \begin{tcolorbox}[breakable, size=fbox, boxrule=.5pt, pad at break*=1mm, opacityfill=0]
\prompt{Out}{outcolor}{26}{\boxspacing}
\begin{Verbatim}[commandchars=\\\{\}]
<seaborn.axisgrid.PairGrid at 0x1d1a95107d0>
\end{Verbatim}
\end{tcolorbox}
        
    \begin{center}
    \adjustimage{max size={0.9\linewidth}{0.9\paperheight}}{figs/bvp_kM_en/output_46_1.png}
    \end{center}
    { \hspace*{\fill} \\}
    
    dk variations cause the same sort of pattern in the data as my small
scale \(R\) variations of the previous dataset. This is expected on our
scale as they are equivalent. I want to extract the period and sample on
their minima to calculate the other oscillation.

The dMn behavior here looks like a flat curve, with spikes on the period
boundaries, as if it were the limit to some sort of function.

\(\mu\) causes linear growth of the minima, but slower than linear for
the maxima. Most of my analyses have shown a linear dependence on
\(\mu\) for the amplitude of my oscillations, but I know \(\mu\) also
changes my parameters, which appears to shorten the spikes for my tilted
secant squared functions. From this graph, this effect is not so
visible, but I will come back to \(\mu\) analysis later.

    \begin{tcolorbox}[breakable, size=fbox, boxrule=1pt, pad at break*=1mm,colback=cellbackground, colframe=cellborder]
\prompt{In}{incolor}{27}{\boxspacing}
\begin{Verbatim}[commandchars=\\\{\}]
\PY{n}{g} \PY{o}{=} \PY{n}{sns}\PY{o}{.}\PY{n}{PairGrid}\PY{p}{(}\PY{n}{tbl\PYZus{}red}\PY{p}{[}\PY{l+m+mi}{2}\PY{p}{]}\PY{p}{,} \PY{n}{diag\PYZus{}sharey}\PY{o}{=}\PY{k+kc}{False}\PY{p}{,} \PY{n}{corner}\PY{o}{=}\PY{k+kc}{True}\PY{p}{)}

\PY{n}{g}\PY{o}{.}\PY{n}{map\PYZus{}diag}\PY{p}{(}\PY{n}{sns}\PY{o}{.}\PY{n}{kdeplot}\PY{p}{)}
\PY{n}{g}\PY{o}{.}\PY{n}{map\PYZus{}lower}\PY{p}{(}\PY{n}{sns}\PY{o}{.}\PY{n}{scatterplot}\PY{p}{)}
\end{Verbatim}
\end{tcolorbox}

            \begin{tcolorbox}[breakable, size=fbox, boxrule=.5pt, pad at break*=1mm, opacityfill=0]
\prompt{Out}{outcolor}{27}{\boxspacing}
\begin{Verbatim}[commandchars=\\\{\}]
<seaborn.axisgrid.PairGrid at 0x1d1ac25e120>
\end{Verbatim}
\end{tcolorbox}
        
    \begin{center}
    \adjustimage{max size={0.9\linewidth}{0.9\paperheight}}{figs/bvp_kM_en/output_48_1.png}
    \end{center}
    { \hspace*{\fill} \\}
    
    \begin{tcolorbox}[breakable, size=fbox, boxrule=1pt, pad at break*=1mm,colback=cellbackground, colframe=cellborder]
\prompt{In}{incolor}{28}{\boxspacing}
\begin{Verbatim}[commandchars=\\\{\}]
\PY{c+c1}{\PYZsh{} I need to delete outliers again. I was hoping mu was small enough to avoid this.}
\PY{n}{tbl\PYZus{}red}\PY{p}{[}\PY{l+m+mi}{2}\PY{p}{]}\PY{o}{.}\PY{n}{loc}\PY{p}{[}\PY{n}{tbl\PYZus{}red}\PY{p}{[}\PY{l+m+mi}{2}\PY{p}{]}\PY{p}{[}\PY{l+s+s2}{\PYZdq{}}\PY{l+s+s2}{dMn}\PY{l+s+s2}{\PYZdq{}}\PY{p}{]} \PY{o}{\PYZlt{}} \PY{o}{\PYZhy{}}\PY{l+m+mi}{20000}\PY{p}{,}\PY{l+s+s2}{\PYZdq{}}\PY{l+s+s2}{dMn}\PY{l+s+s2}{\PYZdq{}}\PY{p}{]} \PY{o}{=} \PY{k+kc}{None}
\PY{n}{tbl\PYZus{}red}\PY{p}{[}\PY{l+m+mi}{2}\PY{p}{]}\PY{o}{.}\PY{n}{loc}\PY{p}{[}\PY{n}{tbl\PYZus{}red}\PY{p}{[}\PY{l+m+mi}{2}\PY{p}{]}\PY{p}{[}\PY{l+s+s2}{\PYZdq{}}\PY{l+s+s2}{dkn}\PY{l+s+s2}{\PYZdq{}}\PY{p}{]} \PY{o}{\PYZlt{}} \PY{o}{\PYZhy{}}\PY{l+m+mf}{4.0e5}\PY{p}{,}\PY{l+s+s2}{\PYZdq{}}\PY{l+s+s2}{dkn}\PY{l+s+s2}{\PYZdq{}}\PY{p}{]} \PY{o}{=} \PY{k+kc}{None}
\end{Verbatim}
\end{tcolorbox}

    \begin{tcolorbox}[breakable, size=fbox, boxrule=1pt, pad at break*=1mm,colback=cellbackground, colframe=cellborder]
\prompt{In}{incolor}{29}{\boxspacing}
\begin{Verbatim}[commandchars=\\\{\}]
\PY{n}{g} \PY{o}{=} \PY{n}{sns}\PY{o}{.}\PY{n}{PairGrid}\PY{p}{(}\PY{n}{tbl\PYZus{}red}\PY{p}{[}\PY{l+m+mi}{2}\PY{p}{]}\PY{p}{,} \PY{n}{diag\PYZus{}sharey}\PY{o}{=}\PY{k+kc}{False}\PY{p}{,} \PY{n}{corner}\PY{o}{=}\PY{k+kc}{True}\PY{p}{)}

\PY{n}{g}\PY{o}{.}\PY{n}{map\PYZus{}diag}\PY{p}{(}\PY{n}{sns}\PY{o}{.}\PY{n}{kdeplot}\PY{p}{)}
\PY{n}{g}\PY{o}{.}\PY{n}{map\PYZus{}lower}\PY{p}{(}\PY{n}{sns}\PY{o}{.}\PY{n}{scatterplot}\PY{p}{)}
\end{Verbatim}
\end{tcolorbox}

            \begin{tcolorbox}[breakable, size=fbox, boxrule=.5pt, pad at break*=1mm, opacityfill=0]
\prompt{Out}{outcolor}{29}{\boxspacing}
\begin{Verbatim}[commandchars=\\\{\}]
<seaborn.axisgrid.PairGrid at 0x1d1ac462420>
\end{Verbatim}
\end{tcolorbox}
        
    \begin{center}
    \adjustimage{max size={0.9\linewidth}{0.9\paperheight}}{figs/bvp_kM_en/output_50_1.png}
    \end{center}
    { \hspace*{\fill} \\}
    
    The patterns are not quite obvious here. I know that I have my
small-scale (fast) oscillations fixed for this dataset, while the slow
\(B\) dependent oscillations should be visible. By the time I had chosen
this dataset, I already know that I have one pattern based on the
decimal part of \(kR\) where \(k\) and \(R\) are in natural units, and a
second pattern based on the decimal part of \(\pi^2BR^2\) where \(B\)
and \(R\) are a percent of their maximum value. The second pattern
should be responsible for the oscillatory behavior visible here, while
the remaining behavior is what I wish to extract from this dataset. This
pattern should be found by following the envelope of these curves. There
appears to be amplitude dependence on \(R\).

    \begin{tcolorbox}[breakable, size=fbox, boxrule=1pt, pad at break*=1mm,colback=cellbackground, colframe=cellborder]
\prompt{In}{incolor}{30}{\boxspacing}
\begin{Verbatim}[commandchars=\\\{\}]
\PY{n}{g} \PY{o}{=} \PY{n}{sns}\PY{o}{.}\PY{n}{PairGrid}\PY{p}{(}\PY{n}{tbl\PYZus{}red}\PY{p}{[}\PY{l+m+mi}{3}\PY{p}{]}\PY{p}{,} \PY{n}{diag\PYZus{}sharey}\PY{o}{=}\PY{k+kc}{False}\PY{p}{,} \PY{n}{corner}\PY{o}{=}\PY{k+kc}{True}\PY{p}{)}

\PY{n}{g}\PY{o}{.}\PY{n}{map\PYZus{}diag}\PY{p}{(}\PY{n}{sns}\PY{o}{.}\PY{n}{kdeplot}\PY{p}{)}
\PY{n}{g}\PY{o}{.}\PY{n}{map\PYZus{}lower}\PY{p}{(}\PY{n}{sns}\PY{o}{.}\PY{n}{scatterplot}\PY{p}{)}
\end{Verbatim}
\end{tcolorbox}

            \begin{tcolorbox}[breakable, size=fbox, boxrule=.5pt, pad at break*=1mm, opacityfill=0]
\prompt{Out}{outcolor}{30}{\boxspacing}
\begin{Verbatim}[commandchars=\\\{\}]
<seaborn.axisgrid.PairGrid at 0x1d19cdcf590>
\end{Verbatim}
\end{tcolorbox}
        
    \begin{center}
    \adjustimage{max size={0.9\linewidth}{0.9\paperheight}}{figs/bvp_kM_en/output_52_1.png}
    \end{center}
    { \hspace*{\fill} \\}
    
    Full \(B\) range for the small \(R\) range. Periodic behavior is clearly
visible in this image. Unlike the full \(R\) range graph, there appears
to be no \(B\) dependence on the parameters which give the shape of the
oscillations. \(\mu\) and \(R\) determine the shape and amplitude of
oscillations. \(B\) and \(dk\) only determine where in the period this
point is.

    \begin{tcolorbox}[breakable, size=fbox, boxrule=1pt, pad at break*=1mm,colback=cellbackground, colframe=cellborder]
\prompt{In}{incolor}{31}{\boxspacing}
\begin{Verbatim}[commandchars=\\\{\}]
\PY{n}{g} \PY{o}{=} \PY{n}{sns}\PY{o}{.}\PY{n}{PairGrid}\PY{p}{(}\PY{n}{tbl\PYZus{}red}\PY{p}{[}\PY{l+m+mi}{4}\PY{p}{]}\PY{p}{,} \PY{n}{diag\PYZus{}sharey}\PY{o}{=}\PY{k+kc}{False}\PY{p}{,} \PY{n}{corner}\PY{o}{=}\PY{k+kc}{True}\PY{p}{)}

\PY{n}{g}\PY{o}{.}\PY{n}{map\PYZus{}diag}\PY{p}{(}\PY{n}{sns}\PY{o}{.}\PY{n}{kdeplot}\PY{p}{)}
\PY{n}{g}\PY{o}{.}\PY{n}{map\PYZus{}lower}\PY{p}{(}\PY{n}{sns}\PY{o}{.}\PY{n}{scatterplot}\PY{p}{)}
\end{Verbatim}
\end{tcolorbox}

            \begin{tcolorbox}[breakable, size=fbox, boxrule=.5pt, pad at break*=1mm, opacityfill=0]
\prompt{Out}{outcolor}{31}{\boxspacing}
\begin{Verbatim}[commandchars=\\\{\}]
<seaborn.axisgrid.PairGrid at 0x1d1a8883bf0>
\end{Verbatim}
\end{tcolorbox}
        
    \begin{center}
    \adjustimage{max size={0.9\linewidth}{0.9\paperheight}}{figs/bvp_kM_en/output_54_1.png}
    \end{center}
    { \hspace*{\fill} \\}
    
    This large grid is difficult to decipher here, but contains \(\mu\)
variation for all of my parameters, which will be useful later. However,
there are too few datapoints to use it for any other analysis.

    The analysis will proceed as follows: 

\begin{itemize}
\item Analysis of the \(dk\)-dependent structure can give us our period of the small-scale oscillations, as well as their functional form.
\item Analysis of the \(B\)-dependent structure can give us our period of the large-scale oscillations, as well as the shape and behavior of all 4 of our oscillations: 
\begin{itemize}
\item \(dkn\) vs \(dk\) or \(dR\) 
\item \(dMn\) vs \(dk\) or \(dR\) 
\item  \(dkn\) vs \(B\) 
\item \(dMn\) vs \(B\)
\end{itemize}
\item Analysis of the \(\mu\)-dependent structure from our
\(dk\) and \(\mu\) varying dataset can show how \(\mu\) affects the
structure. 
\item Analysis of the large-scale \(R\)-dependent structure can
complete our picture, verifying our large-scale oscillations, and
showing how \(R\) affects the structure.

\end{itemize}

Actual numerical analysis of \(\mu\) and \(R\) dependence will need to
be saved for later.

    \hypertarget{small-scale-oscillations}{%
\section{Small-Scale Oscillations}\label{small-scale-oscillations}}

    We are going to start by working with our \(k\) and \(\mu\) varying
dataset.

    \begin{tcolorbox}[breakable, size=fbox, boxrule=1pt, pad at break*=1mm,colback=cellbackground, colframe=cellborder]
\prompt{In}{incolor}{33}{\boxspacing}
\begin{Verbatim}[commandchars=\\\{\}]
\PY{n}{tbl\PYZus{}red}\PY{p}{[}\PY{l+m+mi}{1}\PY{p}{]}
\end{Verbatim}
\end{tcolorbox}

            \begin{tcolorbox}[breakable, size=fbox, boxrule=.5pt, pad at break*=1mm, opacityfill=0]
\prompt{Out}{outcolor}{33}{\boxspacing}
\begin{Verbatim}[commandchars=\\\{\}]
       dk            mu            dkn         dMn
0     0.0  3.162278e-09   -7969.899509   13.457522
1     0.0  3.236689e-09  -21241.450718   14.759672
2     0.0  3.312852e-09  -21655.040102   15.081138
3     0.0  3.390807e-09  -22075.471048   15.409241
4     0.0  3.470596e-09  -22502.823629   15.744108
{\ldots}   {\ldots}           {\ldots}            {\ldots}         {\ldots}
9995  3.0  2.881350e-08 -102227.695108  100.520093
9996  3.0  2.949151e-08 -103715.570360  102.482218
9997  3.0  3.018547e-08 -105221.666897  104.480650
9998  3.0  3.089577e-08 -106746.223455  106.516039
9999  3.0  3.162278e-08 -108289.484694  108.589044

[10000 rows x 4 columns]
\end{Verbatim}
\end{tcolorbox}
        
    \begin{tcolorbox}[breakable, size=fbox, boxrule=1pt, pad at break*=1mm,colback=cellbackground, colframe=cellborder]
\prompt{In}{incolor}{32}{\boxspacing}
\begin{Verbatim}[commandchars=\\\{\}]
\PY{n}{g} \PY{o}{=} \PY{n}{sns}\PY{o}{.}\PY{n}{PairGrid}\PY{p}{(}\PY{n}{tbl\PYZus{}red}\PY{p}{[}\PY{l+m+mi}{1}\PY{p}{]}\PY{p}{,} \PY{n}{diag\PYZus{}sharey}\PY{o}{=}\PY{k+kc}{False}\PY{p}{,} \PY{n}{corner}\PY{o}{=}\PY{k+kc}{True}\PY{p}{)}

\PY{n}{g}\PY{o}{.}\PY{n}{map\PYZus{}diag}\PY{p}{(}\PY{n}{sns}\PY{o}{.}\PY{n}{kdeplot}\PY{p}{)}
\PY{n}{g}\PY{o}{.}\PY{n}{map\PYZus{}lower}\PY{p}{(}\PY{n}{sns}\PY{o}{.}\PY{n}{scatterplot}\PY{p}{)}
\end{Verbatim}
\end{tcolorbox}

            \begin{tcolorbox}[breakable, size=fbox, boxrule=.5pt, pad at break*=1mm, opacityfill=0]
\prompt{Out}{outcolor}{32}{\boxspacing}
\begin{Verbatim}[commandchars=\\\{\}]
<seaborn.axisgrid.PairGrid at 0x1d19bff6120>
\end{Verbatim}
\end{tcolorbox}
        
    \begin{center}
    \adjustimage{max size={0.9\linewidth}{0.9\paperheight}}{figs/bvp_kM_en/output_60_1.png}
    \end{center}
    { \hspace*{\fill} \\}
    
    I am renaming my variables here, so that I don't need to remember which
table number I am working with. All of the data is extracted from my
table and put into arrays.

    \begin{tcolorbox}[breakable, size=fbox, boxrule=1pt, pad at break*=1mm,colback=cellbackground, colframe=cellborder]
\prompt{In}{incolor}{55}{\boxspacing}
\begin{Verbatim}[commandchars=\\\{\}]
\PY{n}{gind} \PY{o}{=} \PY{l+m+mi}{1}
\end{Verbatim}
\end{tcolorbox}

    \begin{tcolorbox}[breakable, size=fbox, boxrule=1pt, pad at break*=1mm,colback=cellbackground, colframe=cellborder]
\prompt{In}{incolor}{56}{\boxspacing}
\begin{Verbatim}[commandchars=\\\{\}]
\PY{n}{dk} \PY{o}{=} \PY{n}{dk\PYZus{}g}\PY{p}{[}\PY{n}{gind}\PY{p}{]}
\PY{n}{mu} \PY{o}{=} \PY{n}{mu\PYZus{}g}\PY{p}{[}\PY{n}{gind}\PY{p}{]}

\PY{n}{dkn\PYZus{}arr} \PY{o}{=} \PY{n}{np}\PY{o}{.}\PY{n}{zeros}\PY{p}{(}\PY{p}{(}\PY{n}{grid\PYZus{}size}\PY{p}{,} \PY{n}{grid\PYZus{}size}\PY{p}{)}\PY{p}{)}
\PY{n}{dMn\PYZus{}arr} \PY{o}{=} \PY{n}{np}\PY{o}{.}\PY{n}{zeros}\PY{p}{(}\PY{p}{(}\PY{n}{grid\PYZus{}size}\PY{p}{,} \PY{n}{grid\PYZus{}size}\PY{p}{)}\PY{p}{)}

\PY{k}{for} \PY{n}{r}\PY{p}{,} \PY{n}{rr} \PY{o+ow}{in} \PY{n+nb}{enumerate}\PY{p}{(}\PY{n}{dk\PYZus{}g}\PY{p}{[}\PY{n}{gind}\PY{p}{]}\PY{p}{)}\PY{p}{:}
    \PY{k}{for} \PY{n}{b}\PY{p}{,} \PY{n}{bb} \PY{o+ow}{in} \PY{n+nb}{enumerate}\PY{p}{(}\PY{n}{mu\PYZus{}g}\PY{p}{[}\PY{n}{gind}\PY{p}{]}\PY{p}{)}\PY{p}{:}
        \PY{n}{tbl\PYZus{}sel} \PY{o}{=} \PY{n}{tbl\PYZus{}red}\PY{p}{[}\PY{n}{gind}\PY{p}{]}\PY{p}{[}\PY{p}{(}\PY{n}{tbl\PYZus{}red}\PY{p}{[}\PY{n}{gind}\PY{p}{]}\PY{p}{[}\PY{l+s+s2}{\PYZdq{}}\PY{l+s+s2}{dk}\PY{l+s+s2}{\PYZdq{}}\PY{p}{]}\PY{o}{==}\PY{n}{rr}\PY{p}{)} \PY{o}{\PYZam{}}\PY{p}{(}\PY{n}{tbl\PYZus{}red}\PY{p}{[}\PY{n}{gind}\PY{p}{]}\PY{p}{[}\PY{l+s+s2}{\PYZdq{}}\PY{l+s+s2}{mu}\PY{l+s+s2}{\PYZdq{}}\PY{p}{]}\PY{o}{==}\PY{n}{bb}\PY{p}{)}\PY{p}{]}
        \PY{n}{dkn\PYZus{}arr}\PY{p}{[}\PY{n}{r}\PY{p}{,}\PY{n}{b}\PY{p}{]} \PY{o}{=} \PY{n}{tbl\PYZus{}sel}\PY{p}{[}\PY{l+s+s2}{\PYZdq{}}\PY{l+s+s2}{dkn}\PY{l+s+s2}{\PYZdq{}}\PY{p}{]}\PY{o}{.}\PY{n}{to\PYZus{}numpy}\PY{p}{(}\PY{p}{)}\PY{p}{[}\PY{l+m+mi}{0}\PY{p}{]}
        \PY{n}{dMn\PYZus{}arr}\PY{p}{[}\PY{n}{r}\PY{p}{,}\PY{n}{b}\PY{p}{]} \PY{o}{=} \PY{n}{tbl\PYZus{}sel}\PY{p}{[}\PY{l+s+s2}{\PYZdq{}}\PY{l+s+s2}{dMn}\PY{l+s+s2}{\PYZdq{}}\PY{p}{]}\PY{o}{.}\PY{n}{to\PYZus{}numpy}\PY{p}{(}\PY{p}{)}\PY{p}{[}\PY{l+m+mi}{0}\PY{p}{]}
\end{Verbatim}
\end{tcolorbox}

    \begin{tcolorbox}[breakable, size=fbox, boxrule=1pt, pad at break*=1mm,colback=cellbackground, colframe=cellborder]
\prompt{In}{incolor}{ }{\boxspacing}
\begin{Verbatim}[commandchars=\\\{\}]
\PY{c+c1}{\PYZsh{} The indicies to plot are in an array, so that one can easily change which plots are displayed.}

\PY{n}{jj} \PY{o}{=} \PY{p}{[}\PY{l+m+mi}{10}\PY{p}{,} \PY{l+m+mi}{20}\PY{p}{,} \PY{l+m+mi}{30}\PY{p}{,} \PY{l+m+mi}{40}\PY{p}{,} \PY{l+m+mi}{50}\PY{p}{,} \PY{l+m+mi}{60}\PY{p}{,} \PY{l+m+mi}{70}\PY{p}{,} \PY{l+m+mi}{80}\PY{p}{,} \PY{l+m+mi}{90}\PY{p}{]}

\PY{n}{ilim} \PY{o}{=} \PY{n+nb}{len}\PY{p}{(}\PY{n}{jj}\PY{p}{)}

\PY{n}{fig}\PY{p}{,} \PY{n}{ax} \PY{o}{=} \PY{n}{plt}\PY{o}{.}\PY{n}{subplots}\PY{p}{(}\PY{p}{)}
\PY{n}{plt}\PY{o}{.}\PY{n}{title}\PY{p}{(}\PY{l+s+sa}{f}\PY{l+s+s2}{\PYZdq{}}\PY{l+s+s2}{Plot of the Change in B\PYZhy{}Oscillation Wave number Due}\PY{l+s+se}{\PYZbs{}n}\PY{l+s+s2}{to Nonlinearity vs Linear k Variation}\PY{l+s+s2}{\PYZdq{}}\PY{p}{)}
\PY{n}{ax}\PY{o}{.}\PY{n}{set\PYZus{}ylabel}\PY{p}{(}\PY{l+s+s1}{\PYZsq{}}\PY{l+s+s1}{Change in M due to nonlinearity (dMn) (m}\PY{l+s+se}{\PYZbs{}u207b}\PY{l+s+se}{\PYZbs{}u00b9}\PY{l+s+s1}{)}\PY{l+s+s1}{\PYZsq{}}\PY{p}{)}
\PY{n}{ax}\PY{o}{.}\PY{n}{set\PYZus{}xlabel}\PY{p}{(}\PY{l+s+s2}{\PYZdq{}}\PY{l+s+s2}{Small\PYZhy{}scale variation of linear wave number k (dk (}\PY{l+s+s2}{\PYZpc{}}\PY{l+s+s2}{))}\PY{l+s+s2}{\PYZdq{}}\PY{p}{)}

\PY{k}{for} \PY{n}{i} \PY{o+ow}{in} \PY{n+nb}{range}\PY{p}{(}\PY{n}{ilim}\PY{p}{)}\PY{p}{:}
    \PY{n}{ax}\PY{o}{.}\PY{n}{plot}\PY{p}{(}\PY{n}{dk}\PY{p}{,} \PY{n}{dMn\PYZus{}arr}\PY{p}{[}\PY{p}{:}\PY{p}{,} \PY{n}{jj}\PY{p}{[}\PY{n}{i}\PY{p}{]}\PY{p}{]}\PY{p}{)}
\PY{n}{plt}\PY{o}{.}\PY{n}{axline}\PY{p}{(}\PY{p}{(}\PY{l+m+mi}{0}\PY{p}{,}\PY{l+m+mi}{0}\PY{p}{)}\PY{p}{,} \PY{p}{(}\PY{l+m+mf}{1e\PYZhy{}15}\PY{p}{,}\PY{l+m+mi}{0}\PY{p}{)}\PY{p}{,} \PY{n}{color}\PY{o}{=}\PY{l+s+s1}{\PYZsq{}}\PY{l+s+s1}{k}\PY{l+s+s1}{\PYZsq{}}\PY{p}{)}
\PY{n}{plt}\PY{o}{.}\PY{n}{axline}\PY{p}{(}\PY{p}{(}\PY{l+m+mi}{0}\PY{p}{,}\PY{l+m+mi}{0}\PY{p}{)}\PY{p}{,} \PY{p}{(}\PY{l+m+mi}{0}\PY{p}{,}\PY{l+m+mf}{1e\PYZhy{}15}\PY{p}{)}\PY{p}{,} \PY{n}{color}\PY{o}{=}\PY{l+s+s1}{\PYZsq{}}\PY{l+s+s1}{k}\PY{l+s+s1}{\PYZsq{}}\PY{p}{)}
\PY{n}{plt}\PY{o}{.}\PY{n}{axline}\PY{p}{(}\PY{p}{(}\PY{l+m+mi}{2}\PY{o}{/}\PY{n}{R\PYZus{}g}\PY{p}{[}\PY{n}{gind}\PY{p}{]}\PY{p}{[}\PY{l+m+mi}{0}\PY{p}{]}\PY{p}{,}\PY{l+m+mi}{0}\PY{p}{)}\PY{p}{,} \PY{p}{(}\PY{l+m+mi}{2}\PY{o}{/}\PY{n}{R\PYZus{}g}\PY{p}{[}\PY{n}{gind}\PY{p}{]}\PY{p}{[}\PY{l+m+mi}{0}\PY{p}{]}\PY{p}{,}\PY{l+m+mf}{1e\PYZhy{}15}\PY{p}{)}\PY{p}{,} \PY{n}{color}\PY{o}{=}\PY{l+s+s1}{\PYZsq{}}\PY{l+s+s1}{k}\PY{l+s+s1}{\PYZsq{}}\PY{p}{)}
\PY{n}{plt}\PY{o}{.}\PY{n}{show}\PY{p}{(}\PY{p}{)}
\end{Verbatim}
\end{tcolorbox}

    \begin{center}
    \adjustimage{max size={0.9\linewidth}{0.9\paperheight}}{figs/bvp_kM_en/output_64_0.png}
    \end{center}
    { \hspace*{\fill} \\}
    
    \begin{tcolorbox}[breakable, size=fbox, boxrule=1pt, pad at break*=1mm,colback=cellbackground, colframe=cellborder]
\prompt{In}{incolor}{ }{\boxspacing}
\begin{Verbatim}[commandchars=\\\{\}]
\PY{c+c1}{\PYZsh{} The indicies to plot are in an array, so that one can easily change which plots are displayed.}

\PY{n}{jj} \PY{o}{=} \PY{p}{[}\PY{l+m+mi}{10}\PY{p}{,} \PY{l+m+mi}{20}\PY{p}{,} \PY{l+m+mi}{30}\PY{p}{,} \PY{l+m+mi}{40}\PY{p}{,} \PY{l+m+mi}{50}\PY{p}{,} \PY{l+m+mi}{60}\PY{p}{,} \PY{l+m+mi}{70}\PY{p}{,} \PY{l+m+mi}{80}\PY{p}{,} \PY{l+m+mi}{90}\PY{p}{]}

\PY{n}{ilim} \PY{o}{=} \PY{n+nb}{len}\PY{p}{(}\PY{n}{jj}\PY{p}{)}

\PY{c+c1}{\PYZsh{} Define our plot object}
\PY{n}{fig}\PY{p}{,} \PY{n}{ax} \PY{o}{=} \PY{n}{plt}\PY{o}{.}\PY{n}{subplots}\PY{p}{(}\PY{p}{)}
\PY{n}{plt}\PY{o}{.}\PY{n}{title}\PY{p}{(}\PY{l+s+sa}{f}\PY{l+s+s2}{\PYZdq{}}\PY{l+s+s2}{Plot of the Change in k Due to Nonlinearity vs Linear k}\PY{l+s+se}{\PYZbs{}n}\PY{l+s+s2}{Variation For Various Values of }\PY{l+s+se}{\PYZbs{}u03BC}\PY{l+s+s2}{\PYZdq{}}\PY{p}{)}
\PY{n}{ax}\PY{o}{.}\PY{n}{set\PYZus{}ylabel}\PY{p}{(}\PY{l+s+s1}{\PYZsq{}}\PY{l+s+s1}{Change in k due to nonlinearity (dkn) (m}\PY{l+s+se}{\PYZbs{}u207b}\PY{l+s+se}{\PYZbs{}u00b9}\PY{l+s+s1}{)}\PY{l+s+s1}{\PYZsq{}}\PY{p}{)}
\PY{n}{ax}\PY{o}{.}\PY{n}{set\PYZus{}xlabel}\PY{p}{(}\PY{l+s+s2}{\PYZdq{}}\PY{l+s+s2}{Small\PYZhy{}scale variation of linear wave number k (dk (}\PY{l+s+s2}{\PYZpc{}}\PY{l+s+s2}{))}\PY{l+s+s2}{\PYZdq{}}\PY{p}{)}

\PY{n}{h\PYZus{}list} \PY{o}{=} \PY{p}{[}\PY{p}{]} \PY{c+c1}{\PYZsh{} For making the legend.}

\PY{c+c1}{\PYZsh{} Our functions to plot}
\PY{k}{for} \PY{n}{i} \PY{o+ow}{in} \PY{n+nb}{range}\PY{p}{(}\PY{n}{ilim}\PY{p}{)}\PY{p}{:}
    \PY{n}{line}\PY{p}{,} \PY{o}{=} \PY{n}{ax}\PY{o}{.}\PY{n}{plot}\PY{p}{(}\PY{n}{dk}\PY{p}{,} \PY{n}{dkn\PYZus{}arr}\PY{p}{[}\PY{p}{:}\PY{p}{,} \PY{n}{jj}\PY{p}{[}\PY{n}{i}\PY{p}{]}\PY{p}{]}\PY{p}{,} \PY{n}{label} \PY{o}{=} \PY{l+s+sa}{f}\PY{l+s+s1}{\PYZsq{}}\PY{l+s+se}{\PYZbs{}u03BC}\PY{l+s+s1}{ = }\PY{l+s+si}{\PYZob{}}\PY{n}{mu}\PY{p}{[}\PY{n}{jj}\PY{p}{[}\PY{n}{i}\PY{p}{]}\PY{p}{]}\PY{l+s+si}{:}\PY{l+s+s1}{.2e}\PY{l+s+si}{\PYZcb{}}\PY{l+s+s1}{\PYZsq{}}\PY{p}{)}
    \PY{n}{h\PYZus{}list}\PY{o}{.}\PY{n}{append}\PY{p}{(}\PY{n}{line}\PY{p}{)}

\PY{c+c1}{\PYZsh{} Lines for the axes and oscillation boundary.}
\PY{n}{plt}\PY{o}{.}\PY{n}{axline}\PY{p}{(}\PY{p}{(}\PY{l+m+mi}{0}\PY{p}{,}\PY{l+m+mi}{0}\PY{p}{)}\PY{p}{,} \PY{p}{(}\PY{l+m+mf}{1e\PYZhy{}15}\PY{p}{,}\PY{l+m+mi}{0}\PY{p}{)}\PY{p}{,} \PY{n}{color}\PY{o}{=}\PY{l+s+s1}{\PYZsq{}}\PY{l+s+s1}{k}\PY{l+s+s1}{\PYZsq{}}\PY{p}{)}
\PY{n}{plt}\PY{o}{.}\PY{n}{axline}\PY{p}{(}\PY{p}{(}\PY{l+m+mi}{0}\PY{p}{,}\PY{l+m+mi}{0}\PY{p}{)}\PY{p}{,} \PY{p}{(}\PY{l+m+mi}{0}\PY{p}{,}\PY{l+m+mf}{1e\PYZhy{}15}\PY{p}{)}\PY{p}{,} \PY{n}{color}\PY{o}{=}\PY{l+s+s1}{\PYZsq{}}\PY{l+s+s1}{k}\PY{l+s+s1}{\PYZsq{}}\PY{p}{)}
\PY{n}{plt}\PY{o}{.}\PY{n}{axline}\PY{p}{(}\PY{p}{(}\PY{l+m+mi}{2}\PY{o}{/}\PY{n}{R\PYZus{}g}\PY{p}{[}\PY{n}{gind}\PY{p}{]}\PY{p}{[}\PY{l+m+mi}{0}\PY{p}{]}\PY{p}{,}\PY{l+m+mi}{0}\PY{p}{)}\PY{p}{,} \PY{p}{(}\PY{l+m+mi}{2}\PY{o}{/}\PY{n}{R\PYZus{}g}\PY{p}{[}\PY{n}{gind}\PY{p}{]}\PY{p}{[}\PY{l+m+mi}{0}\PY{p}{]}\PY{p}{,}\PY{l+m+mf}{1e\PYZhy{}15}\PY{p}{)}\PY{p}{,} \PY{n}{color}\PY{o}{=}\PY{l+s+s1}{\PYZsq{}}\PY{l+s+s1}{k}\PY{l+s+s1}{\PYZsq{}}\PY{p}{)}

\PY{c+c1}{\PYZsh{} Legend}
\PY{n}{ax}\PY{o}{.}\PY{n}{legend}\PY{p}{(}\PY{n}{handles}\PY{o}{=}\PY{n}{h\PYZus{}list}\PY{p}{)}
\PY{n}{plt}\PY{o}{.}\PY{n}{show}\PY{p}{(}\PY{p}{)}
\end{Verbatim}
\end{tcolorbox}

    \begin{center}
    \adjustimage{max size={0.9\linewidth}{0.9\paperheight}}{figs/bvp_kM_en/output_65_0.png}
    \end{center}
    { \hspace*{\fill} \\}
    
    My structure here appears to have a repetition on the scale of the
decimal part of \(\frac{kR}{2}\) as my period of \(dk\) oscillations is
\(\frac{2}{R}\). However, \(k= k_{F} + \frac{dk}{2 R_{max}}\), so my
repetition is actually based on
\(k R R_{max} + \frac{dk R R_{max}}{2 R_{max}}\). The first term is
always a whole number for \(k_F = 1.2 \times 10^{10}\) and
\(R R_{max} = 9 \times 10^{-7}\), leaving us with the second term,
\(\frac{dk R}{2}\), and my repeating pattern is a function of the
decimal part of \(kR\), in real units.

This is not quite true, as the shape of the oscillations most likely has
a dependence on large-scale variations \(R\) and \(k\), but for a fixed
large scale variation of these parameters, our small-scale oscillations
have their value based on the decimal part of \(kR\), which produces
very fast oscillations, given that $ 6.84 \times 10^{3} \le kR
\le 1.2 \times 10^4$ for gold and our range of values for \(R\).

    Before moving into trying to fit this shape to a tilted squared secant
function, I will show this same oscillatory effect on the \(R\)
dependent graphs. I will extract this data from my \(R\) and \(B\)
dependent dataset.

    \begin{tcolorbox}[breakable, size=fbox, boxrule=1pt, pad at break*=1mm,colback=cellbackground, colframe=cellborder]
\prompt{In}{incolor}{59}{\boxspacing}
\begin{Verbatim}[commandchars=\\\{\}]
\PY{n}{gind} \PY{o}{=} \PY{l+m+mi}{0}

\PY{n}{dR} \PY{o}{=} \PY{p}{(}\PY{n}{R\PYZus{}g}\PY{p}{[}\PY{n}{gind}\PY{p}{]} \PY{o}{\PYZhy{}} \PY{l+m+mf}{0.9}\PY{p}{)}\PY{o}{*}\PY{l+m+mf}{1e4} \PY{c+c1}{\PYZsh{} To remove the constant part of R and analyze only the change}
\PY{n}{B}  \PY{o}{=} \PY{n}{B\PYZus{}g}\PY{p}{[}\PY{n}{gind}\PY{p}{]}

\PY{k}{del} \PY{n}{dkn\PYZus{}arr}\PY{p}{,} \PY{n}{dMn\PYZus{}arr}

\PY{n}{dkn\PYZus{}arr} \PY{o}{=} \PY{n}{np}\PY{o}{.}\PY{n}{zeros}\PY{p}{(}\PY{p}{(}\PY{n}{grid\PYZus{}size}\PY{p}{,} \PY{n}{grid\PYZus{}size}\PY{p}{)}\PY{p}{)}
\PY{n}{dMn\PYZus{}arr} \PY{o}{=} \PY{n}{np}\PY{o}{.}\PY{n}{zeros}\PY{p}{(}\PY{p}{(}\PY{n}{grid\PYZus{}size}\PY{p}{,} \PY{n}{grid\PYZus{}size}\PY{p}{)}\PY{p}{)}

\PY{k}{for} \PY{n}{r}\PY{p}{,} \PY{n}{rr} \PY{o+ow}{in} \PY{n+nb}{enumerate}\PY{p}{(}\PY{n}{R\PYZus{}g}\PY{p}{[}\PY{n}{gind}\PY{p}{]}\PY{p}{)}\PY{p}{:}
    \PY{k}{for} \PY{n}{b}\PY{p}{,} \PY{n}{bb} \PY{o+ow}{in} \PY{n+nb}{enumerate}\PY{p}{(}\PY{n}{B\PYZus{}g}\PY{p}{[}\PY{n}{gind}\PY{p}{]}\PY{p}{)}\PY{p}{:}
        \PY{n}{tbl\PYZus{}sel} \PY{o}{=} \PY{n}{tbl\PYZus{}red}\PY{p}{[}\PY{n}{gind}\PY{p}{]}\PY{p}{[}\PY{p}{(}\PY{n}{tbl\PYZus{}red}\PY{p}{[}\PY{n}{gind}\PY{p}{]}\PY{p}{[}\PY{l+s+s2}{\PYZdq{}}\PY{l+s+s2}{R}\PY{l+s+s2}{\PYZdq{}}\PY{p}{]}\PY{o}{==}\PY{n}{rr}\PY{p}{)} \PY{o}{\PYZam{}}\PY{p}{(}\PY{n}{tbl\PYZus{}red}\PY{p}{[}\PY{n}{gind}\PY{p}{]}\PY{p}{[}\PY{l+s+s2}{\PYZdq{}}\PY{l+s+s2}{B}\PY{l+s+s2}{\PYZdq{}}\PY{p}{]}\PY{o}{==}\PY{n}{bb}\PY{p}{)}\PY{p}{]}
        \PY{n}{dkn\PYZus{}arr}\PY{p}{[}\PY{n}{r}\PY{p}{,}\PY{n}{b}\PY{p}{]} \PY{o}{=} \PY{n}{tbl\PYZus{}sel}\PY{p}{[}\PY{l+s+s2}{\PYZdq{}}\PY{l+s+s2}{dkn}\PY{l+s+s2}{\PYZdq{}}\PY{p}{]}\PY{o}{.}\PY{n}{to\PYZus{}numpy}\PY{p}{(}\PY{p}{)}\PY{p}{[}\PY{l+m+mi}{0}\PY{p}{]} \PY{o}{/} \PY{l+m+mi}{100000} \PY{c+c1}{\PYZsh{} rescaling to make plots easier to read}
        \PY{n}{dMn\PYZus{}arr}\PY{p}{[}\PY{n}{r}\PY{p}{,}\PY{n}{b}\PY{p}{]} \PY{o}{=} \PY{n}{tbl\PYZus{}sel}\PY{p}{[}\PY{l+s+s2}{\PYZdq{}}\PY{l+s+s2}{dMn}\PY{l+s+s2}{\PYZdq{}}\PY{p}{]}\PY{o}{.}\PY{n}{to\PYZus{}numpy}\PY{p}{(}\PY{p}{)}\PY{p}{[}\PY{l+m+mi}{0}\PY{p}{]} \PY{o}{/} \PY{l+m+mi}{10000} \PY{c+c1}{\PYZsh{} rescaling to make plots easier to read}
\end{Verbatim}
\end{tcolorbox}

    Period calculations. We know that \(dk = 0.45\) or
\(0.225 \times 10^6\), \(k_0 = 1.2 \times 10^{10}\), \(R_0 = 0.9\) or
\(0.9 \times 10^{-6}\) and \(dR = dR_{shown} \times 10^{-10}\).
\begin{equation}
kR = k_0 R_0 + k_0 dR + dk R_0 + dk dR
\end{equation} The \(k_0 R_0\) and \(dk dR\) terms are irrelevant, as
the first always gives a whole number and the second is too small to be
seen. So we have \begin{equation}
kR = k_0 dR + dk R_0 = 1.2 dR_{shown} + 0.225 * 0.9 =  0.2025 + 1.2 R_{shown}.
\end{equation} This is 0 when \(R_{shown} = - 0.2025 / 1.2 = -0.16875\),
and 1 when \(R_{shown} = -0.16875 + 1/1.2 = 1.0208\).

    \begin{tcolorbox}[breakable, size=fbox, boxrule=1pt, pad at break*=1mm,colback=cellbackground, colframe=cellborder]
\prompt{In}{incolor}{ }{\boxspacing}
\begin{Verbatim}[commandchars=\\\{\}]
\PY{c+c1}{\PYZsh{} The indexes to plot are in an array, so that one can easily change which plots are displayed.}
\PY{c+c1}{\PYZsh{} I have chosen the B value which sits at our node, like for the dk graph case. }
\PY{c+c1}{\PYZsh{} I will analyze the rest of this data later. }

\PY{n}{jj} \PY{o}{=} \PY{p}{[}\PY{l+m+mi}{0}\PY{p}{]}

\PY{n}{ilim} \PY{o}{=} \PY{n+nb}{len}\PY{p}{(}\PY{n}{jj}\PY{p}{)}

\PY{n}{fig}\PY{p}{,} \PY{n}{ax} \PY{o}{=} \PY{n}{plt}\PY{o}{.}\PY{n}{subplots}\PY{p}{(}\PY{p}{)}
\PY{n}{plt}\PY{o}{.}\PY{n}{title}\PY{p}{(}\PY{l+s+sa}{f}\PY{l+s+s2}{\PYZdq{}}\PY{l+s+s2}{Plot of the Change in k and M Due to Nonlinearity}\PY{l+s+se}{\PYZbs{}n}\PY{l+s+s2}{vs R for dk = 0.45 / 2 * 10}\PY{l+s+se}{\PYZbs{}u2076}\PY{l+s+s2}{ m}\PY{l+s+se}{\PYZbs{}u207b}\PY{l+s+se}{\PYZbs{}u00b9}\PY{l+s+s2}{\PYZdq{}}\PY{p}{)}
\PY{n}{ax}\PY{o}{.}\PY{n}{set\PYZus{}ylabel}\PY{p}{(}\PY{l+s+s1}{\PYZsq{}}\PY{l+s+s1}{Change in k (dkn) or M (dMn) due to nonlinearity}\PY{l+s+s1}{\PYZsq{}}\PY{p}{)}
\PY{n}{ax}\PY{o}{.}\PY{n}{set\PYZus{}xlabel}\PY{p}{(}\PY{l+s+s2}{\PYZdq{}}\PY{l+s+s2}{Change dr in radius R of ring (R = 9 * 10}\PY{l+s+se}{\PYZbs{}u207b}\PY{l+s+se}{\PYZbs{}u2077}\PY{l+s+s2}{ m + dr * 10}\PY{l+s+se}{\PYZbs{}u207b}\PY{l+s+se}{\PYZbs{}u00b9}\PY{l+s+se}{\PYZbs{}u2070}\PY{l+s+s2}{ m)}\PY{l+s+s2}{\PYZdq{}}\PY{p}{)}

\PY{c+c1}{\PYZsh{} Graphs}
\PY{k}{for} \PY{n}{i} \PY{o+ow}{in} \PY{n+nb}{range}\PY{p}{(}\PY{n}{ilim}\PY{p}{)}\PY{p}{:}
    \PY{n}{ax}\PY{o}{.}\PY{n}{plot}\PY{p}{(}\PY{n}{dR}\PY{p}{,} \PY{n}{dkn\PYZus{}arr}\PY{p}{[}\PY{p}{:}\PY{p}{,}\PY{n}{jj}\PY{p}{[}\PY{n}{i}\PY{p}{]}\PY{p}{]}\PY{p}{,} \PY{n}{label}\PY{o}{=}\PY{l+s+s2}{\PYZdq{}}\PY{l+s+s2}{dkn (10}\PY{l+s+se}{\PYZbs{}u2075}\PY{l+s+s2}{ m}\PY{l+s+se}{\PYZbs{}u207b}\PY{l+s+se}{\PYZbs{}u00b9}\PY{l+s+s2}{)}\PY{l+s+s2}{\PYZdq{}}\PY{p}{)}
    \PY{n}{ax}\PY{o}{.}\PY{n}{plot}\PY{p}{(}\PY{n}{dR}\PY{p}{,} \PY{n}{dMn\PYZus{}arr}\PY{p}{[}\PY{p}{:}\PY{p}{,}\PY{n}{jj}\PY{p}{[}\PY{n}{i}\PY{p}{]}\PY{p}{]} \PY{o}{\PYZhy{}} \PY{l+m+mi}{2}\PY{p}{,} \PY{n}{label}\PY{o}{=}\PY{l+s+s2}{\PYZdq{}}\PY{l+s+s2}{dMn (10}\PY{l+s+se}{\PYZbs{}u2074}\PY{l+s+s2}{ m}\PY{l+s+se}{\PYZbs{}u207b}\PY{l+s+se}{\PYZbs{}u00b9}\PY{l+s+s2}{) \PYZhy{} 2}\PY{l+s+s2}{\PYZdq{}}\PY{p}{)} \PY{c+c1}{\PYZsh{} The shift of this axis is so that they can both be seen together.}

\PY{c+c1}{\PYZsh{} Our black lines designate the expected period boundaries}
\PY{n}{plt}\PY{o}{.}\PY{n}{axline}\PY{p}{(}\PY{p}{(}\PY{l+m+mi}{0}\PY{p}{,}\PY{o}{\PYZhy{}}\PY{l+m+mi}{2}\PY{p}{)}\PY{p}{,} \PY{p}{(}\PY{l+m+mf}{1e\PYZhy{}15}\PY{p}{,}\PY{o}{\PYZhy{}}\PY{l+m+mi}{2}\PY{p}{)}\PY{p}{,} \PY{n}{color}\PY{o}{=}\PY{l+s+s1}{\PYZsq{}}\PY{l+s+s1}{k}\PY{l+s+s1}{\PYZsq{}}\PY{p}{)}
\PY{n}{shift} \PY{o}{=} \PY{o}{\PYZhy{}} \PY{l+m+mf}{0.16875} \PY{o}{+} \PY{l+m+mi}{1}\PY{o}{/}\PY{l+m+mf}{1.2}
\PY{n}{plt}\PY{o}{.}\PY{n}{axline}\PY{p}{(}\PY{p}{(}\PY{n}{shift}\PY{p}{,}\PY{l+m+mi}{0}\PY{p}{)}\PY{p}{,} \PY{p}{(}\PY{n}{shift}\PY{p}{,}\PY{l+m+mf}{1e\PYZhy{}15}\PY{p}{)}\PY{p}{,} \PY{n}{color}\PY{o}{=}\PY{l+s+s1}{\PYZsq{}}\PY{l+s+s1}{k}\PY{l+s+s1}{\PYZsq{}}\PY{p}{,} \PY{n}{label}\PY{o}{=}\PY{l+s+s2}{\PYZdq{}}\PY{l+s+s2}{remainder of kR / 1 = 0}\PY{l+s+s2}{\PYZdq{}}\PY{p}{)}
\PY{n}{shift} \PY{o}{=} \PY{n}{shift} \PY{o}{+} \PY{l+m+mi}{1}\PY{o}{/}\PY{l+m+mf}{1.2}
\PY{n}{plt}\PY{o}{.}\PY{n}{axline}\PY{p}{(}\PY{p}{(}\PY{n}{shift}\PY{p}{,}\PY{l+m+mi}{0}\PY{p}{)}\PY{p}{,} \PY{p}{(}\PY{n}{shift}\PY{p}{,}\PY{l+m+mf}{1e\PYZhy{}15}\PY{p}{)}\PY{p}{,} \PY{n}{color}\PY{o}{=}\PY{l+s+s1}{\PYZsq{}}\PY{l+s+s1}{k}\PY{l+s+s1}{\PYZsq{}}\PY{p}{)}

\PY{n}{plt}\PY{o}{.}\PY{n}{legend}\PY{p}{(}\PY{p}{)}

\PY{n}{plt}\PY{o}{.}\PY{n}{show}\PY{p}{(}\PY{p}{)}
\end{Verbatim}
\end{tcolorbox}

    \begin{center}
    \adjustimage{max size={0.9\linewidth}{0.9\paperheight}}{figs/bvp_kM_en/output_70_0.png}
    \end{center}
    { \hspace*{\fill} \\}
    
    \begin{tcolorbox}[breakable, size=fbox, boxrule=1pt, pad at break*=1mm,colback=cellbackground, colframe=cellborder]
\prompt{In}{incolor}{ }{\boxspacing}
\begin{Verbatim}[commandchars=\\\{\}]
\PY{c+c1}{\PYZsh{} Now to show the dMn graph separately. The structure can be seen to prefer the indicated points. }
\PY{c+c1}{\PYZsh{} The actual shape of the structure is actual incorrect due to our imprecision, but this will }
\PY{c+c1}{\PYZsh{} be seen later.}

\PY{n}{jj} \PY{o}{=} \PY{p}{[}\PY{l+m+mi}{0}\PY{p}{]}

\PY{n}{ilim} \PY{o}{=} \PY{n+nb}{len}\PY{p}{(}\PY{n}{jj}\PY{p}{)}

\PY{n}{fig}\PY{p}{,} \PY{n}{ax} \PY{o}{=} \PY{n}{plt}\PY{o}{.}\PY{n}{subplots}\PY{p}{(}\PY{p}{)}
\PY{n}{plt}\PY{o}{.}\PY{n}{title}\PY{p}{(}\PY{l+s+sa}{f}\PY{l+s+s2}{\PYZdq{}}\PY{l+s+s2}{Plot of the Change in M Due to Nonlinearity}\PY{l+s+se}{\PYZbs{}n}\PY{l+s+s2}{vs R for dk = 0.45 / 2 * 10}\PY{l+s+se}{\PYZbs{}u2076}\PY{l+s+s2}{ m}\PY{l+s+se}{\PYZbs{}u207b}\PY{l+s+se}{\PYZbs{}u00b9}\PY{l+s+s2}{\PYZdq{}}\PY{p}{)}
\PY{n}{ax}\PY{o}{.}\PY{n}{set\PYZus{}ylabel}\PY{p}{(}\PY{l+s+s1}{\PYZsq{}}\PY{l+s+s1}{Change in M (dMn) due to nonlinearity (10}\PY{l+s+se}{\PYZbs{}u2074}\PY{l+s+s1}{ m}\PY{l+s+se}{\PYZbs{}u207b}\PY{l+s+se}{\PYZbs{}u00b9}\PY{l+s+s1}{)}\PY{l+s+s1}{\PYZsq{}}\PY{p}{)}
\PY{n}{ax}\PY{o}{.}\PY{n}{set\PYZus{}xlabel}\PY{p}{(}\PY{l+s+s2}{\PYZdq{}}\PY{l+s+s2}{Change dr in radius R of ring (R = 9 * 10}\PY{l+s+se}{\PYZbs{}u207b}\PY{l+s+se}{\PYZbs{}u2077}\PY{l+s+s2}{ m + dr * 10}\PY{l+s+se}{\PYZbs{}u207b}\PY{l+s+se}{\PYZbs{}u00b9}\PY{l+s+se}{\PYZbs{}u2070}\PY{l+s+s2}{ m)}\PY{l+s+s2}{\PYZdq{}}\PY{p}{)}

\PY{c+c1}{\PYZsh{} Graph}
\PY{k}{for} \PY{n}{i} \PY{o+ow}{in} \PY{n+nb}{range}\PY{p}{(}\PY{n}{ilim}\PY{p}{)}\PY{p}{:}
    \PY{n}{ax}\PY{o}{.}\PY{n}{plot}\PY{p}{(}\PY{n}{dR}\PY{p}{,} \PY{n}{dMn\PYZus{}arr}\PY{p}{[}\PY{p}{:}\PY{p}{,}\PY{n}{jj}\PY{p}{[}\PY{n}{i}\PY{p}{]}\PY{p}{]}\PY{p}{,} \PY{n}{label}\PY{o}{=}\PY{l+s+s2}{\PYZdq{}}\PY{l+s+s2}{dMn (10}\PY{l+s+se}{\PYZbs{}u2074}\PY{l+s+s2}{ m}\PY{l+s+se}{\PYZbs{}u207b}\PY{l+s+se}{\PYZbs{}u00b9}\PY{l+s+s2}{)}\PY{l+s+s2}{\PYZdq{}}\PY{p}{)}

\PY{c+c1}{\PYZsh{} Our indicator lines}
\PY{n}{plt}\PY{o}{.}\PY{n}{axline}\PY{p}{(}\PY{p}{(}\PY{l+m+mi}{0}\PY{p}{,}\PY{l+m+mi}{0}\PY{p}{)}\PY{p}{,} \PY{p}{(}\PY{l+m+mf}{1e\PYZhy{}15}\PY{p}{,}\PY{l+m+mi}{0}\PY{p}{)}\PY{p}{,} \PY{n}{color}\PY{o}{=}\PY{l+s+s1}{\PYZsq{}}\PY{l+s+s1}{k}\PY{l+s+s1}{\PYZsq{}}\PY{p}{)}
\PY{n}{shift} \PY{o}{=} \PY{o}{\PYZhy{}} \PY{l+m+mf}{0.16875} \PY{o}{+} \PY{l+m+mi}{1}\PY{o}{/}\PY{l+m+mf}{1.2}
\PY{n}{plt}\PY{o}{.}\PY{n}{axline}\PY{p}{(}\PY{p}{(}\PY{n}{shift}\PY{p}{,}\PY{l+m+mi}{0}\PY{p}{)}\PY{p}{,} \PY{p}{(}\PY{n}{shift}\PY{p}{,}\PY{l+m+mf}{1e\PYZhy{}15}\PY{p}{)}\PY{p}{,} \PY{n}{color}\PY{o}{=}\PY{l+s+s1}{\PYZsq{}}\PY{l+s+s1}{k}\PY{l+s+s1}{\PYZsq{}}\PY{p}{,} \PY{n}{label}\PY{o}{=}\PY{l+s+s2}{\PYZdq{}}\PY{l+s+s2}{remainder of kR / 1 = 0}\PY{l+s+s2}{\PYZdq{}}\PY{p}{)}
\PY{n}{shift} \PY{o}{=} \PY{n}{shift} \PY{o}{+} \PY{l+m+mi}{1}\PY{o}{/}\PY{l+m+mf}{1.2}
\PY{n}{plt}\PY{o}{.}\PY{n}{axline}\PY{p}{(}\PY{p}{(}\PY{n}{shift}\PY{p}{,}\PY{l+m+mi}{0}\PY{p}{)}\PY{p}{,} \PY{p}{(}\PY{n}{shift}\PY{p}{,}\PY{l+m+mf}{1e\PYZhy{}15}\PY{p}{)}\PY{p}{,} \PY{n}{color}\PY{o}{=}\PY{l+s+s1}{\PYZsq{}}\PY{l+s+s1}{k}\PY{l+s+s1}{\PYZsq{}}\PY{p}{)}

\PY{n}{plt}\PY{o}{.}\PY{n}{legend}\PY{p}{(}\PY{p}{)}

\PY{n}{plt}\PY{o}{.}\PY{n}{show}\PY{p}{(}\PY{p}{)}
\end{Verbatim}
\end{tcolorbox}

    \begin{center}
    \adjustimage{max size={0.9\linewidth}{0.9\paperheight}}{figs/bvp_kM_en/output_71_0.png}
    \end{center}
    { \hspace*{\fill} \\}
    
    \begin{tcolorbox}[breakable, size=fbox, boxrule=1pt, pad at break*=1mm,colback=cellbackground, colframe=cellborder]
\prompt{In}{incolor}{ }{\boxspacing}
\begin{Verbatim}[commandchars=\\\{\}]
\PY{c+c1}{\PYZsh{} I can show this for more B values to show that this period is correct, but the B oscillation will be analyzed later.}

\PY{n}{jj} \PY{o}{=} \PY{p}{[}\PY{l+m+mi}{0}\PY{p}{,} \PY{l+m+mi}{10}\PY{p}{,} \PY{l+m+mi}{20}\PY{p}{,} \PY{l+m+mi}{40}\PY{p}{]}

\PY{n}{ilim} \PY{o}{=} \PY{n+nb}{len}\PY{p}{(}\PY{n}{jj}\PY{p}{)}

\PY{n}{fig}\PY{p}{,} \PY{n}{ax} \PY{o}{=} \PY{n}{plt}\PY{o}{.}\PY{n}{subplots}\PY{p}{(}\PY{p}{)}
\PY{n}{plt}\PY{o}{.}\PY{n}{title}\PY{p}{(}\PY{l+s+sa}{f}\PY{l+s+s2}{\PYZdq{}}\PY{l+s+s2}{Plot of the Change in k Due to Nonlinearity}\PY{l+s+se}{\PYZbs{}n}\PY{l+s+s2}{vs R for dk = 0.45 / 2 * 10}\PY{l+s+se}{\PYZbs{}u2076}\PY{l+s+s2}{ m}\PY{l+s+se}{\PYZbs{}u207b}\PY{l+s+se}{\PYZbs{}u00b9}\PY{l+s+s2}{\PYZdq{}}\PY{p}{)}
\PY{n}{ax}\PY{o}{.}\PY{n}{set\PYZus{}ylabel}\PY{p}{(}\PY{l+s+s1}{\PYZsq{}}\PY{l+s+s1}{Change in k (dkn) due to nonlinearity (10}\PY{l+s+se}{\PYZbs{}u2075}\PY{l+s+s1}{ m}\PY{l+s+se}{\PYZbs{}u207b}\PY{l+s+se}{\PYZbs{}u00b9}\PY{l+s+s1}{)}\PY{l+s+s1}{\PYZsq{}}\PY{p}{)}
\PY{n}{ax}\PY{o}{.}\PY{n}{set\PYZus{}xlabel}\PY{p}{(}\PY{l+s+s2}{\PYZdq{}}\PY{l+s+s2}{Change dr in radius R of ring (R = 9 * 10}\PY{l+s+se}{\PYZbs{}u207b}\PY{l+s+se}{\PYZbs{}u2077}\PY{l+s+s2}{ m + dr * 10}\PY{l+s+se}{\PYZbs{}u207b}\PY{l+s+se}{\PYZbs{}u00b9}\PY{l+s+se}{\PYZbs{}u2070}\PY{l+s+s2}{ m)}\PY{l+s+s2}{\PYZdq{}}\PY{p}{)}

\PY{k}{for} \PY{n}{i} \PY{o+ow}{in} \PY{n+nb}{range}\PY{p}{(}\PY{n}{ilim}\PY{p}{)}\PY{p}{:}
    \PY{n}{ax}\PY{o}{.}\PY{n}{plot}\PY{p}{(}\PY{n}{dR}\PY{p}{,} \PY{n}{dkn\PYZus{}arr}\PY{p}{[}\PY{p}{:}\PY{p}{,}\PY{n}{jj}\PY{p}{[}\PY{n}{i}\PY{p}{]}\PY{p}{]}\PY{p}{,} \PY{n}{label} \PY{o}{=} \PY{l+s+sa}{f}\PY{l+s+s1}{\PYZsq{}}\PY{l+s+s1}{B = }\PY{l+s+si}{\PYZob{}}\PY{n}{B}\PY{p}{[}\PY{n}{jj}\PY{p}{[}\PY{n}{i}\PY{p}{]}\PY{p}{]}\PY{l+s+si}{:}\PY{l+s+s1}{.4f}\PY{l+s+si}{\PYZcb{}}\PY{l+s+s1}{ * B\PYZus{}max}\PY{l+s+s1}{\PYZsq{}}\PY{p}{)}

\PY{n}{shift} \PY{o}{=} \PY{o}{\PYZhy{}} \PY{l+m+mf}{0.16875} \PY{o}{+} \PY{l+m+mi}{1}\PY{o}{/}\PY{l+m+mf}{1.2}
\PY{n}{plt}\PY{o}{.}\PY{n}{axline}\PY{p}{(}\PY{p}{(}\PY{n}{shift}\PY{p}{,}\PY{l+m+mi}{0}\PY{p}{)}\PY{p}{,} \PY{p}{(}\PY{n}{shift}\PY{p}{,}\PY{l+m+mf}{1e\PYZhy{}15}\PY{p}{)}\PY{p}{,} \PY{n}{color}\PY{o}{=}\PY{l+s+s1}{\PYZsq{}}\PY{l+s+s1}{k}\PY{l+s+s1}{\PYZsq{}}\PY{p}{,} \PY{n}{label}\PY{o}{=}\PY{l+s+s2}{\PYZdq{}}\PY{l+s+s2}{remainder of kR / 1 = 0}\PY{l+s+s2}{\PYZdq{}}\PY{p}{)}
\PY{n}{shift} \PY{o}{=} \PY{n}{shift} \PY{o}{+} \PY{l+m+mi}{1}\PY{o}{/}\PY{l+m+mf}{1.2}
\PY{n}{plt}\PY{o}{.}\PY{n}{axline}\PY{p}{(}\PY{p}{(}\PY{n}{shift}\PY{p}{,}\PY{l+m+mi}{0}\PY{p}{)}\PY{p}{,} \PY{p}{(}\PY{n}{shift}\PY{p}{,}\PY{l+m+mf}{1e\PYZhy{}15}\PY{p}{)}\PY{p}{,} \PY{n}{color}\PY{o}{=}\PY{l+s+s1}{\PYZsq{}}\PY{l+s+s1}{k}\PY{l+s+s1}{\PYZsq{}}\PY{p}{)}

\PY{n}{plt}\PY{o}{.}\PY{n}{legend}\PY{p}{(}\PY{p}{)}

\PY{n}{plt}\PY{o}{.}\PY{n}{show}\PY{p}{(}\PY{p}{)}
\end{Verbatim}
\end{tcolorbox}

    \begin{center}
    \adjustimage{max size={0.9\linewidth}{0.9\paperheight}}{figs/bvp_kM_en/output_72_0.png}
    \end{center}
    { \hspace*{\fill} \\}
    
    \begin{tcolorbox}[breakable, size=fbox, boxrule=1pt, pad at break*=1mm,colback=cellbackground, colframe=cellborder]
\prompt{In}{incolor}{ }{\boxspacing}
\begin{Verbatim}[commandchars=\\\{\}]
\PY{n}{jj} \PY{o}{=} \PY{p}{[}\PY{l+m+mi}{0}\PY{p}{,} \PY{l+m+mi}{10}\PY{p}{,} \PY{l+m+mi}{20}\PY{p}{,} \PY{l+m+mi}{40}\PY{p}{]}

\PY{n}{ilim} \PY{o}{=} \PY{n+nb}{len}\PY{p}{(}\PY{n}{jj}\PY{p}{)}

\PY{n}{fig}\PY{p}{,} \PY{n}{ax} \PY{o}{=} \PY{n}{plt}\PY{o}{.}\PY{n}{subplots}\PY{p}{(}\PY{p}{)}
\PY{n}{plt}\PY{o}{.}\PY{n}{title}\PY{p}{(}\PY{l+s+sa}{f}\PY{l+s+s2}{\PYZdq{}}\PY{l+s+s2}{Plot of the Change in M Due to Nonlinearity}\PY{l+s+se}{\PYZbs{}n}\PY{l+s+s2}{vs R for dk = 0.45 / 2 * 10}\PY{l+s+se}{\PYZbs{}u2076}\PY{l+s+s2}{ m}\PY{l+s+se}{\PYZbs{}u207b}\PY{l+s+se}{\PYZbs{}u00b9}\PY{l+s+s2}{\PYZdq{}}\PY{p}{)}
\PY{n}{ax}\PY{o}{.}\PY{n}{set\PYZus{}ylabel}\PY{p}{(}\PY{l+s+s1}{\PYZsq{}}\PY{l+s+s1}{Change in M (dMn) due to nonlinearity (10}\PY{l+s+se}{\PYZbs{}u2074}\PY{l+s+s1}{ m}\PY{l+s+se}{\PYZbs{}u207b}\PY{l+s+se}{\PYZbs{}u00b9}\PY{l+s+s1}{)}\PY{l+s+s1}{\PYZsq{}}\PY{p}{)}
\PY{n}{ax}\PY{o}{.}\PY{n}{set\PYZus{}xlabel}\PY{p}{(}\PY{l+s+s2}{\PYZdq{}}\PY{l+s+s2}{Change dr in radius R of ring (R = 9 * 10}\PY{l+s+se}{\PYZbs{}u207b}\PY{l+s+se}{\PYZbs{}u2077}\PY{l+s+s2}{ m + dr * 10}\PY{l+s+se}{\PYZbs{}u207b}\PY{l+s+se}{\PYZbs{}u00b9}\PY{l+s+se}{\PYZbs{}u2070}\PY{l+s+s2}{ m)}\PY{l+s+s2}{\PYZdq{}}\PY{p}{)}

\PY{k}{for} \PY{n}{i} \PY{o+ow}{in} \PY{n+nb}{range}\PY{p}{(}\PY{n}{ilim}\PY{p}{)}\PY{p}{:}
    \PY{n}{ax}\PY{o}{.}\PY{n}{plot}\PY{p}{(}\PY{n}{dR}\PY{p}{,} \PY{n}{dMn\PYZus{}arr}\PY{p}{[}\PY{p}{:}\PY{p}{,}\PY{n}{jj}\PY{p}{[}\PY{n}{i}\PY{p}{]}\PY{p}{]}\PY{p}{,} \PY{n}{label} \PY{o}{=} \PY{l+s+sa}{f}\PY{l+s+s1}{\PYZsq{}}\PY{l+s+s1}{B = }\PY{l+s+si}{\PYZob{}}\PY{n}{B}\PY{p}{[}\PY{n}{jj}\PY{p}{[}\PY{n}{i}\PY{p}{]}\PY{p}{]}\PY{l+s+si}{:}\PY{l+s+s1}{.4f}\PY{l+s+si}{\PYZcb{}}\PY{l+s+s1}{ * B\PYZus{}max}\PY{l+s+s1}{\PYZsq{}}\PY{p}{)} 

\PY{n}{plt}\PY{o}{.}\PY{n}{axline}\PY{p}{(}\PY{p}{(}\PY{l+m+mi}{0}\PY{p}{,}\PY{l+m+mi}{0}\PY{p}{)}\PY{p}{,} \PY{p}{(}\PY{l+m+mf}{1e\PYZhy{}15}\PY{p}{,}\PY{l+m+mi}{0}\PY{p}{)}\PY{p}{,} \PY{n}{color}\PY{o}{=}\PY{l+s+s1}{\PYZsq{}}\PY{l+s+s1}{k}\PY{l+s+s1}{\PYZsq{}}\PY{p}{)}
\PY{n}{shift} \PY{o}{=} \PY{o}{\PYZhy{}} \PY{l+m+mf}{0.16875} \PY{o}{+} \PY{l+m+mi}{1}\PY{o}{/}\PY{l+m+mf}{1.2}
\PY{n}{plt}\PY{o}{.}\PY{n}{axline}\PY{p}{(}\PY{p}{(}\PY{n}{shift}\PY{p}{,}\PY{l+m+mi}{0}\PY{p}{)}\PY{p}{,} \PY{p}{(}\PY{n}{shift}\PY{p}{,}\PY{l+m+mf}{1e\PYZhy{}15}\PY{p}{)}\PY{p}{,} \PY{n}{color}\PY{o}{=}\PY{l+s+s1}{\PYZsq{}}\PY{l+s+s1}{k}\PY{l+s+s1}{\PYZsq{}}\PY{p}{,} \PY{n}{label}\PY{o}{=}\PY{l+s+s2}{\PYZdq{}}\PY{l+s+s2}{remainder of kR / 1 = 0}\PY{l+s+s2}{\PYZdq{}}\PY{p}{)}
\PY{n}{shift} \PY{o}{=} \PY{n}{shift} \PY{o}{+} \PY{l+m+mi}{1}\PY{o}{/}\PY{l+m+mf}{1.2}
\PY{n}{plt}\PY{o}{.}\PY{n}{axline}\PY{p}{(}\PY{p}{(}\PY{n}{shift}\PY{p}{,}\PY{l+m+mi}{0}\PY{p}{)}\PY{p}{,} \PY{p}{(}\PY{n}{shift}\PY{p}{,}\PY{l+m+mf}{1e\PYZhy{}15}\PY{p}{)}\PY{p}{,} \PY{n}{color}\PY{o}{=}\PY{l+s+s1}{\PYZsq{}}\PY{l+s+s1}{k}\PY{l+s+s1}{\PYZsq{}}\PY{p}{)}

\PY{n}{plt}\PY{o}{.}\PY{n}{legend}\PY{p}{(}\PY{p}{)}

\PY{n}{plt}\PY{o}{.}\PY{n}{show}\PY{p}{(}\PY{p}{)}
\end{Verbatim}
\end{tcolorbox}

    \begin{center}
    \adjustimage{max size={0.9\linewidth}{0.9\paperheight}}{figs/bvp_kM_en/output_73_0.png}
    \end{center}
    { \hspace*{\fill} \\}
    
    Neither \(\mu\) nor \(B\) change this period, although they both change
the structure of the graph. The fact that our chosen starting point of
oscillations is significant is not always obvious. Let us now extract
that skewed secant squared pattern, then move on to the \(B\) analysis,
which clearly shows all of the oscillatory behavior.

    \hypertarget{rotated-secant-squared}{%
\section{Rotated Secant Squared}\label{rotated-secant-squared}}

    \begin{tcolorbox}[breakable, size=fbox, boxrule=1pt, pad at break*=1mm,colback=cellbackground, colframe=cellborder]
\prompt{In}{incolor}{81}{\boxspacing}
\begin{Verbatim}[commandchars=\\\{\}]
\PY{n}{gind} \PY{o}{=} \PY{l+m+mi}{1}
\end{Verbatim}
\end{tcolorbox}

    \begin{tcolorbox}[breakable, size=fbox, boxrule=1pt, pad at break*=1mm,colback=cellbackground, colframe=cellborder]
\prompt{In}{incolor}{82}{\boxspacing}
\begin{Verbatim}[commandchars=\\\{\}]
\PY{n}{dk} \PY{o}{=} \PY{n}{dk\PYZus{}g}\PY{p}{[}\PY{n}{gind}\PY{p}{]}
\PY{n}{mu} \PY{o}{=} \PY{n}{mu\PYZus{}g}\PY{p}{[}\PY{n}{gind}\PY{p}{]}

\PY{n}{dkn\PYZus{}arr} \PY{o}{=} \PY{n}{np}\PY{o}{.}\PY{n}{zeros}\PY{p}{(}\PY{p}{(}\PY{n}{grid\PYZus{}size}\PY{p}{,} \PY{n}{grid\PYZus{}size}\PY{p}{)}\PY{p}{)}
\PY{n}{dMn\PYZus{}arr} \PY{o}{=} \PY{n}{np}\PY{o}{.}\PY{n}{zeros}\PY{p}{(}\PY{p}{(}\PY{n}{grid\PYZus{}size}\PY{p}{,} \PY{n}{grid\PYZus{}size}\PY{p}{)}\PY{p}{)}

\PY{k}{for} \PY{n}{r}\PY{p}{,} \PY{n}{rr} \PY{o+ow}{in} \PY{n+nb}{enumerate}\PY{p}{(}\PY{n}{dk\PYZus{}g}\PY{p}{[}\PY{n}{gind}\PY{p}{]}\PY{p}{)}\PY{p}{:}
    \PY{k}{for} \PY{n}{b}\PY{p}{,} \PY{n}{bb} \PY{o+ow}{in} \PY{n+nb}{enumerate}\PY{p}{(}\PY{n}{mu\PYZus{}g}\PY{p}{[}\PY{n}{gind}\PY{p}{]}\PY{p}{)}\PY{p}{:}
        \PY{n}{tbl\PYZus{}sel} \PY{o}{=} \PY{n}{tbl\PYZus{}red}\PY{p}{[}\PY{n}{gind}\PY{p}{]}\PY{p}{[}\PY{p}{(}\PY{n}{tbl\PYZus{}red}\PY{p}{[}\PY{n}{gind}\PY{p}{]}\PY{p}{[}\PY{l+s+s2}{\PYZdq{}}\PY{l+s+s2}{dk}\PY{l+s+s2}{\PYZdq{}}\PY{p}{]}\PY{o}{==}\PY{n}{rr}\PY{p}{)} \PY{o}{\PYZam{}}\PY{p}{(}\PY{n}{tbl\PYZus{}red}\PY{p}{[}\PY{n}{gind}\PY{p}{]}\PY{p}{[}\PY{l+s+s2}{\PYZdq{}}\PY{l+s+s2}{mu}\PY{l+s+s2}{\PYZdq{}}\PY{p}{]}\PY{o}{==}\PY{n}{bb}\PY{p}{)}\PY{p}{]}
        \PY{n}{dkn\PYZus{}arr}\PY{p}{[}\PY{n}{r}\PY{p}{,}\PY{n}{b}\PY{p}{]} \PY{o}{=} \PY{n}{tbl\PYZus{}sel}\PY{p}{[}\PY{l+s+s2}{\PYZdq{}}\PY{l+s+s2}{dkn}\PY{l+s+s2}{\PYZdq{}}\PY{p}{]}\PY{o}{.}\PY{n}{to\PYZus{}numpy}\PY{p}{(}\PY{p}{)}\PY{p}{[}\PY{l+m+mi}{0}\PY{p}{]}
        \PY{n}{dMn\PYZus{}arr}\PY{p}{[}\PY{n}{r}\PY{p}{,}\PY{n}{b}\PY{p}{]} \PY{o}{=} \PY{n}{tbl\PYZus{}sel}\PY{p}{[}\PY{l+s+s2}{\PYZdq{}}\PY{l+s+s2}{dMn}\PY{l+s+s2}{\PYZdq{}}\PY{p}{]}\PY{o}{.}\PY{n}{to\PYZus{}numpy}\PY{p}{(}\PY{p}{)}\PY{p}{[}\PY{l+m+mi}{0}\PY{p}{]}
\end{Verbatim}
\end{tcolorbox}

    \begin{tcolorbox}[breakable, size=fbox, boxrule=1pt, pad at break*=1mm,colback=cellbackground, colframe=cellborder]
\prompt{In}{incolor}{ }{\boxspacing}
\begin{Verbatim}[commandchars=\\\{\}]
\PY{c+c1}{\PYZsh{} Plot of our dkn vs dk data for a single chosen value of mu}

\PY{n}{jj} \PY{o}{=} \PY{p}{[}\PY{l+m+mi}{80}\PY{p}{]}

\PY{n}{ilim} \PY{o}{=} \PY{n+nb}{len}\PY{p}{(}\PY{n}{jj}\PY{p}{)}

\PY{n}{fig}\PY{p}{,} \PY{n}{ax} \PY{o}{=} \PY{n}{plt}\PY{o}{.}\PY{n}{subplots}\PY{p}{(}\PY{p}{)}
\PY{n}{plt}\PY{o}{.}\PY{n}{title}\PY{p}{(}\PY{l+s+sa}{f}\PY{l+s+s2}{\PYZdq{}}\PY{l+s+s2}{Plot of the Change in k due to Nonlinearity vs Linear k}\PY{l+s+se}{\PYZbs{}n}\PY{l+s+s2}{Variation For }\PY{l+s+se}{\PYZbs{}u03BC}\PY{l+s+s2}{ = }\PY{l+s+si}{\PYZob{}}\PY{n}{mu}\PY{p}{[}\PY{n}{jj}\PY{p}{[}\PY{l+m+mi}{0}\PY{p}{]}\PY{p}{]}\PY{l+s+si}{:}\PY{l+s+s2}{.5e}\PY{l+s+si}{\PYZcb{}}\PY{l+s+s2}{\PYZdq{}}\PY{p}{)}
\PY{n}{ax}\PY{o}{.}\PY{n}{set\PYZus{}ylabel}\PY{p}{(}\PY{l+s+s1}{\PYZsq{}}\PY{l+s+s1}{Change in k due to nonlinearity (\PYZhy{}dkn)  (\PYZhy{}10}\PY{l+s+se}{\PYZbs{}u2074}\PY{l+s+s1}{  m}\PY{l+s+se}{\PYZbs{}u207b}\PY{l+s+se}{\PYZbs{}u00b9}\PY{l+s+s1}{)}\PY{l+s+s1}{\PYZsq{}}\PY{p}{)}
\PY{n}{ax}\PY{o}{.}\PY{n}{set\PYZus{}xlabel}\PY{p}{(}\PY{l+s+s2}{\PYZdq{}}\PY{l+s+s2}{Small\PYZhy{}scale variation of linear wave number k (dk (}\PY{l+s+s2}{\PYZpc{}}\PY{l+s+s2}{))}\PY{l+s+s2}{\PYZdq{}}\PY{p}{)}

\PY{n}{ax}\PY{o}{.}\PY{n}{plot}\PY{p}{(}\PY{n}{dk}\PY{p}{,} \PY{o}{\PYZhy{}}\PY{n}{dkn\PYZus{}arr}\PY{p}{[}\PY{p}{:}\PY{p}{,} \PY{n}{jj}\PY{p}{[}\PY{l+m+mi}{0}\PY{p}{]}\PY{p}{]}\PY{o}{/}\PY{l+m+mi}{10000}\PY{p}{)}

\PY{n}{plt}\PY{o}{.}\PY{n}{show}\PY{p}{(}\PY{p}{)}
\end{Verbatim}
\end{tcolorbox}

    \begin{center}
    \adjustimage{max size={0.9\linewidth}{0.9\paperheight}}{figs/bvp_kM_en/output_78_0.png}
    \end{center}
    { \hspace*{\fill} \\}
    
    \begin{tcolorbox}[breakable, size=fbox, boxrule=1pt, pad at break*=1mm,colback=cellbackground, colframe=cellborder]
\prompt{In}{incolor}{85}{\boxspacing}
\begin{Verbatim}[commandchars=\\\{\}]
\PY{c+c1}{\PYZsh{} We are now going to work with the one curve shown above, and try to extract its shape.}
\PY{c+c1}{\PYZsh{} We start by finding the minimum absolute value of our curve.}
\PY{c+c1}{\PYZsh{} Note that our changes in fast oscillation period are always negative, so we switched the sign}
\PY{c+c1}{\PYZsh{} in order to work with positive numbers.}

\PY{n}{min\PYZus{}dkn} \PY{o}{=} \PY{n}{np}\PY{o}{.}\PY{n}{min}\PY{p}{(}\PY{o}{\PYZhy{}}\PY{n}{dkn\PYZus{}arr}\PY{p}{[}\PY{l+m+mi}{2}\PY{p}{:}\PY{p}{,} \PY{n}{jj}\PY{p}{[}\PY{l+m+mi}{0}\PY{p}{]}\PY{p}{]}\PY{o}{/}\PY{l+m+mi}{10000}\PY{p}{)}
\PY{n}{period} \PY{o}{=} \PY{l+m+mi}{2}\PY{o}{/}\PY{l+m+mf}{0.9}
\end{Verbatim}
\end{tcolorbox}

    \begin{tcolorbox}[breakable, size=fbox, boxrule=1pt, pad at break*=1mm,colback=cellbackground, colframe=cellborder]
\prompt{In}{incolor}{ }{\boxspacing}
\begin{Verbatim}[commandchars=\\\{\}]
\PY{c+c1}{\PYZsh{} We know that tan\PYZca{}2 + sec\PYZca{}2 = 1, so we will move the function to the x axis and}
\PY{c+c1}{\PYZsh{} match the function to tan\PYZca{}2 instead of sec\PYZca{}2. This will be convenient later, }
\PY{c+c1}{\PYZsh{} when we want to rotate this curve around our minimum.}

\PY{n}{model1} \PY{o}{=} \PY{n}{min\PYZus{}dkn}\PY{o}{*}\PY{n}{np}\PY{o}{.}\PY{n}{power}\PY{p}{(}\PY{n}{np}\PY{o}{.}\PY{n}{tan}\PY{p}{(}\PY{l+m+mf}{0.9}\PY{o}{*}\PY{n}{pi}\PY{o}{*}\PY{n}{dk}\PY{o}{/}\PY{l+m+mi}{2}\PY{p}{)}\PY{p}{,}\PY{l+m+mi}{2}\PY{p}{)}

\PY{c+c1}{\PYZsh{} Now we need to cut off values of our function which are too large, so that }
\PY{c+c1}{\PYZsh{} the plot focuses on our relevant values of dkn.}

\PY{k}{for} \PY{n}{i}\PY{p}{,}\PY{n}{j} \PY{o+ow}{in} \PY{n+nb}{enumerate}\PY{p}{(}\PY{n}{model1}\PY{p}{)}\PY{p}{:}
    \PY{k}{if} \PY{n}{j} \PY{o}{\PYZgt{}} \PY{l+m+mi}{40}\PY{p}{:} \PY{n}{model1}\PY{p}{[}\PY{n}{i}\PY{p}{]} \PY{o}{=} \PY{l+m+mi}{40}
\end{Verbatim}
\end{tcolorbox}

    \begin{tcolorbox}[breakable, size=fbox, boxrule=1pt, pad at break*=1mm,colback=cellbackground, colframe=cellborder]
\prompt{In}{incolor}{ }{\boxspacing}
\begin{Verbatim}[commandchars=\\\{\}]
\PY{c+c1}{\PYZsh{} for easy changing which two plots to use.}

\PY{n}{jj} \PY{o}{=} \PY{p}{[}\PY{l+m+mi}{80}\PY{p}{]}

\PY{n}{ilim} \PY{o}{=} \PY{n+nb}{len}\PY{p}{(}\PY{n}{jj}\PY{p}{)}

\PY{n}{fig}\PY{p}{,} \PY{n}{ax} \PY{o}{=} \PY{n}{plt}\PY{o}{.}\PY{n}{subplots}\PY{p}{(}\PY{p}{)}
\PY{n}{plt}\PY{o}{.}\PY{n}{title}\PY{p}{(}\PY{l+s+sa}{f}\PY{l+s+s2}{\PYZdq{}}\PY{l+s+s2}{Plot of the Change in k due to Nonlinearity vs Linear k}\PY{l+s+se}{\PYZbs{}n}\PY{l+s+s2}{Variation For }\PY{l+s+se}{\PYZbs{}u03BC}\PY{l+s+s2}{ = }\PY{l+s+si}{\PYZob{}}\PY{n}{mu}\PY{p}{[}\PY{n}{jj}\PY{p}{[}\PY{l+m+mi}{0}\PY{p}{]}\PY{p}{]}\PY{l+s+si}{:}\PY{l+s+s2}{.5e}\PY{l+s+si}{\PYZcb{}}\PY{l+s+s2}{\PYZdq{}}\PY{p}{)}
\PY{n}{ax}\PY{o}{.}\PY{n}{set\PYZus{}ylabel}\PY{p}{(}\PY{l+s+s1}{\PYZsq{}}\PY{l+s+s1}{Change in k due to nonlinearity (\PYZhy{}dkn)  (\PYZhy{}10}\PY{l+s+se}{\PYZbs{}u2074}\PY{l+s+s1}{  m}\PY{l+s+se}{\PYZbs{}u207b}\PY{l+s+se}{\PYZbs{}u00b9}\PY{l+s+s1}{)}\PY{l+s+s1}{\PYZsq{}}\PY{p}{)}
\PY{n}{ax}\PY{o}{.}\PY{n}{set\PYZus{}xlabel}\PY{p}{(}\PY{l+s+s2}{\PYZdq{}}\PY{l+s+s2}{Small\PYZhy{}scale variation of linear wave number k (dk (}\PY{l+s+s2}{\PYZpc{}}\PY{l+s+s2}{))}\PY{l+s+s2}{\PYZdq{}}\PY{p}{)}

\PY{n}{ax}\PY{o}{.}\PY{n}{plot}\PY{p}{(}\PY{n}{dk}\PY{p}{,} \PY{o}{\PYZhy{}}\PY{n}{dkn\PYZus{}arr}\PY{p}{[}\PY{p}{:}\PY{p}{,} \PY{n}{jj}\PY{p}{[}\PY{l+m+mi}{0}\PY{p}{]}\PY{p}{]}\PY{o}{/}\PY{l+m+mi}{10000} \PY{o}{\PYZhy{}} \PY{n}{min\PYZus{}dkn}\PY{p}{,} \PY{n}{label} \PY{o}{=} \PY{l+s+s2}{\PYZdq{}}\PY{l+s+s2}{numerical results}\PY{l+s+s2}{\PYZdq{}}\PY{p}{)} \PY{c+c1}{\PYZsh{} We are rescaling dkn and moving it to the axis.}
\PY{n}{ax}\PY{o}{.}\PY{n}{plot}\PY{p}{(}\PY{n}{dk}\PY{p}{,} \PY{n}{model1}\PY{p}{,} \PY{n}{label} \PY{o}{=} \PY{l+s+s2}{\PYZdq{}}\PY{l+s+s2}{model}\PY{l+s+s2}{\PYZdq{}}\PY{p}{)}

\PY{n}{plt}\PY{o}{.}\PY{n}{legend}\PY{p}{(}\PY{p}{)}

\PY{n}{plt}\PY{o}{.}\PY{n}{show}\PY{p}{(}\PY{p}{)}
\end{Verbatim}
\end{tcolorbox}

    \begin{center}
    \adjustimage{max size={0.9\linewidth}{0.9\paperheight}}{figs/bvp_kM_en/output_81_0.png}
    \end{center}
    { \hspace*{\fill} \\}
    
    This functional form looks somewhat similar, but it is not obvious that
it is correct yet. Now I will try to rotate this tangent squared curve
around the minimum, using a rotation matrix.

    \begin{tcolorbox}[breakable, size=fbox, boxrule=1pt, pad at break*=1mm,colback=cellbackground, colframe=cellborder]
\prompt{In}{incolor}{91}{\boxspacing}
\begin{Verbatim}[commandchars=\\\{\}]
\PY{c+c1}{\PYZsh{} tangent squared model}
\PY{n}{model1} \PY{o}{=} \PY{o}{\PYZhy{}}\PY{n}{min\PYZus{}dkn}\PY{o}{*}\PY{n}{np}\PY{o}{.}\PY{n}{power}\PY{p}{(}\PY{n}{np}\PY{o}{.}\PY{n}{tan}\PY{p}{(}\PY{l+m+mf}{0.9}\PY{o}{*}\PY{n}{pi}\PY{o}{*}\PY{n}{dk}\PY{o}{/}\PY{l+m+mi}{2}\PY{p}{)}\PY{p}{,}\PY{l+m+mi}{2}\PY{p}{)} 

\PY{n}{x2} \PY{o}{=} \PY{n}{np}\PY{o}{.}\PY{n}{zeros}\PY{p}{(}\PY{n}{grid\PYZus{}size}\PY{p}{)} \PY{c+c1}{\PYZsh{} New x values}
\PY{n}{y2} \PY{o}{=} \PY{n}{np}\PY{o}{.}\PY{n}{zeros}\PY{p}{(}\PY{n}{grid\PYZus{}size}\PY{p}{)} \PY{c+c1}{\PYZsh{} New y values}
\PY{n}{theta} \PY{o}{=} \PY{l+m+mf}{0.02025} \PY{c+c1}{\PYZsh{} This is our rotation angle, which we adjust by hand for now.}

\PY{c+c1}{\PYZsh{} Apply the rotation matrix}
\PY{k}{for} \PY{n}{i}\PY{p}{,}\PY{n}{y1} \PY{o+ow}{in} \PY{n+nb}{enumerate}\PY{p}{(}\PY{n}{model1}\PY{p}{)}\PY{p}{:}
    \PY{n}{x1} \PY{o}{=} \PY{n}{dk}\PY{p}{[}\PY{n}{i}\PY{p}{]}
    \PY{n}{x2}\PY{p}{[}\PY{n}{i}\PY{p}{]} \PY{o}{=} \PY{n}{np}\PY{o}{.}\PY{n}{cos}\PY{p}{(}\PY{n}{theta}\PY{p}{)}\PY{o}{*} \PY{n}{x1} \PY{o}{\PYZhy{}} \PY{n}{np}\PY{o}{.}\PY{n}{sin}\PY{p}{(}\PY{n}{theta}\PY{p}{)}\PY{o}{*}\PY{n}{y1}
    \PY{n}{y2}\PY{p}{[}\PY{n}{i}\PY{p}{]} \PY{o}{=} \PY{n}{np}\PY{o}{.}\PY{n}{sin}\PY{p}{(}\PY{n}{theta}\PY{p}{)}\PY{o}{*} \PY{n}{x1} \PY{o}{\PYZhy{}} \PY{n}{np}\PY{o}{.}\PY{n}{cos}\PY{p}{(}\PY{n}{theta}\PY{p}{)}\PY{o}{*}\PY{n}{y1}

\PY{c+c1}{\PYZsh{} remove values which are too large}
\PY{k}{for} \PY{n}{i}\PY{p}{,}\PY{n}{j} \PY{o+ow}{in} \PY{n+nb}{enumerate}\PY{p}{(}\PY{n}{y2}\PY{p}{)}\PY{p}{:}
    \PY{k}{if} \PY{n}{j} \PY{o}{\PYZgt{}} \PY{l+m+mi}{40}\PY{p}{:} \PY{n}{y2}\PY{p}{[}\PY{n}{i}\PY{p}{]} \PY{o}{=} \PY{k+kc}{None}
\PY{k}{for} \PY{n}{i}\PY{p}{,} \PY{n}{j} \PY{o+ow}{in} \PY{n+nb}{enumerate}\PY{p}{(}\PY{n}{x2}\PY{p}{)}\PY{p}{:}
    \PY{k}{if} \PY{n}{j} \PY{o}{\PYZgt{}} \PY{l+m+mi}{3}\PY{p}{:} \PY{n}{x2}\PY{p}{[}\PY{n}{i}\PY{p}{]}\PY{o}{=} \PY{k+kc}{None}
\end{Verbatim}
\end{tcolorbox}

    Here is my plot. Note that I have shifted the x values by a value
inserted by hand and the rotation angle has also been tuned by hand. I
will eventually need to construct a nonlinear regression script for
fitting this function, but for now I just want to know what the function
is and what parameters it has.

    \begin{tcolorbox}[breakable, size=fbox, boxrule=1pt, pad at break*=1mm,colback=cellbackground, colframe=cellborder]
\prompt{In}{incolor}{92}{\boxspacing}
\begin{Verbatim}[commandchars=\\\{\}]
\PY{n}{fig}\PY{p}{,} \PY{n}{ax} \PY{o}{=} \PY{n}{plt}\PY{o}{.}\PY{n}{subplots}\PY{p}{(}\PY{p}{)}
\PY{n}{plt}\PY{o}{.}\PY{n}{title}\PY{p}{(}\PY{l+s+sa}{f}\PY{l+s+s2}{\PYZdq{}}\PY{l+s+s2}{Plot of the Change in k due to Nonlinearity vs Linear k}\PY{l+s+se}{\PYZbs{}n}\PY{l+s+s2}{Variation For }\PY{l+s+se}{\PYZbs{}u03BC}\PY{l+s+s2}{ = }\PY{l+s+si}{\PYZob{}}\PY{n}{mu}\PY{p}{[}\PY{n}{jj}\PY{p}{[}\PY{l+m+mi}{0}\PY{p}{]}\PY{p}{]}\PY{l+s+si}{:}\PY{l+s+s2}{.5e}\PY{l+s+si}{\PYZcb{}}\PY{l+s+s2}{\PYZdq{}}\PY{p}{)}
\PY{n}{ax}\PY{o}{.}\PY{n}{set\PYZus{}ylabel}\PY{p}{(}\PY{l+s+s1}{\PYZsq{}}\PY{l+s+s1}{Change in k due to nonlinearity (\PYZhy{}dkn)  (\PYZhy{}10}\PY{l+s+se}{\PYZbs{}u2074}\PY{l+s+s1}{  m}\PY{l+s+se}{\PYZbs{}u207b}\PY{l+s+se}{\PYZbs{}u00b9}\PY{l+s+s1}{)}\PY{l+s+s1}{\PYZsq{}}\PY{p}{)}
\PY{n}{ax}\PY{o}{.}\PY{n}{set\PYZus{}xlabel}\PY{p}{(}\PY{l+s+s2}{\PYZdq{}}\PY{l+s+s2}{Small\PYZhy{}scale variation of linear wavenumber k (dk (}\PY{l+s+s2}{\PYZpc{}}\PY{l+s+s2}{))}\PY{l+s+s2}{\PYZdq{}}\PY{p}{)}

\PY{n}{ax}\PY{o}{.}\PY{n}{plot}\PY{p}{(}\PY{n}{dk}\PY{p}{,} \PY{o}{\PYZhy{}}\PY{n}{dkn\PYZus{}arr}\PY{p}{[}\PY{p}{:}\PY{p}{,} \PY{n}{jj}\PY{p}{[}\PY{l+m+mi}{0}\PY{p}{]}\PY{p}{]}\PY{o}{/}\PY{l+m+mi}{10000} \PY{o}{\PYZhy{}} \PY{n}{min\PYZus{}dkn}\PY{p}{,} \PY{n}{label} \PY{o}{=} \PY{l+s+s2}{\PYZdq{}}\PY{l+s+s2}{numerical results}\PY{l+s+s2}{\PYZdq{}}\PY{p}{)}

\PY{n}{ax}\PY{o}{.}\PY{n}{plot}\PY{p}{(}\PY{n}{x2}\PY{o}{+}\PY{l+m+mf}{0.065}\PY{p}{,} \PY{n}{y2}\PY{p}{,} \PY{n}{label} \PY{o}{=} \PY{l+s+s1}{\PYZsq{}}\PY{l+s+s1}{model}\PY{l+s+s1}{\PYZsq{}}\PY{p}{)} \PY{c+c1}{\PYZsh{} Shift the modeled skewed secant squared curve over by the set value.}

\PY{n}{plt}\PY{o}{.}\PY{n}{legend}\PY{p}{(}\PY{p}{)}

\PY{n}{plt}\PY{o}{.}\PY{n}{show}\PY{p}{(}\PY{p}{)}
\end{Verbatim}
\end{tcolorbox}

    \begin{center}
    \adjustimage{max size={0.9\linewidth}{0.9\paperheight}}{figs/bvp_kM_en/output_85_0.png}
    \end{center}
    { \hspace*{\fill} \\}
    
    This curve looks like a good fit. I just need to cut off values from the
left curve where they overlap on the right curve, and all of the values
on the left hand side of the right curve where they overlap another x
point on that curve.

    \begin{tcolorbox}[breakable, size=fbox, boxrule=1pt, pad at break*=1mm,colback=cellbackground, colframe=cellborder]
\prompt{In}{incolor}{93}{\boxspacing}
\begin{Verbatim}[commandchars=\\\{\}]
\PY{c+c1}{\PYZsh{} If one of x or y is already removed, remove the other.}
\PY{k}{for} \PY{n}{i} \PY{o+ow}{in} \PY{n+nb}{range}\PY{p}{(}\PY{n}{grid\PYZus{}size}\PY{p}{)}\PY{p}{:}
    \PY{k}{if} \PY{n}{np}\PY{o}{.}\PY{n}{isnan}\PY{p}{(}\PY{n}{y2}\PY{p}{[}\PY{n}{i}\PY{p}{]}\PY{p}{)}\PY{p}{:} \PY{n}{x2}\PY{p}{[}\PY{n}{i}\PY{p}{]} \PY{o}{=} \PY{k+kc}{None}
    \PY{k}{if} \PY{n}{np}\PY{o}{.}\PY{n}{isnan}\PY{p}{(}\PY{n}{x2}\PY{p}{[}\PY{n}{i}\PY{p}{]}\PY{p}{)}\PY{p}{:} \PY{n}{y2}\PY{p}{[}\PY{n}{i}\PY{p}{]} \PY{o}{=} \PY{k+kc}{None}
\end{Verbatim}
\end{tcolorbox}

    \begin{tcolorbox}[breakable, size=fbox, boxrule=1pt, pad at break*=1mm,colback=cellbackground, colframe=cellborder]
\prompt{In}{incolor}{ }{\boxspacing}
\begin{Verbatim}[commandchars=\\\{\}]
\PY{c+c1}{\PYZsh{} I need to find the overlaps within the x array. I know the array values are in order for x1, }
\PY{c+c1}{\PYZsh{} which means that there will be values for the first tangent squared curve, then for the second.}
\PY{c+c1}{\PYZsh{} I want to first look for the removed values, which exist for the cutoff, }
\PY{c+c1}{\PYZsh{} in order to separate these curves, since I know that removed values will appear at the two sides}
\PY{c+c1}{\PYZsh{} of each curve.}

\PY{n}{arg1} \PY{o}{=} \PY{n}{np}\PY{o}{.}\PY{n}{argwhere}\PY{p}{(}\PY{n}{np}\PY{o}{.}\PY{n}{isnan}\PY{p}{(}\PY{n}{x2}\PY{p}{)}\PY{p}{)}\PY{p}{[}\PY{l+m+mi}{0}\PY{p}{]}\PY{p}{[}\PY{l+m+mi}{0}\PY{p}{]}
\PY{n}{arg2} \PY{o}{=} \PY{n}{np}\PY{o}{.}\PY{n}{argwhere}\PY{p}{(}\PY{n}{np}\PY{o}{.}\PY{n}{isnan}\PY{p}{(}\PY{n}{y2}\PY{p}{)}\PY{p}{)}\PY{p}{[}\PY{l+m+mi}{0}\PY{p}{]}\PY{p}{[}\PY{l+m+mi}{0}\PY{p}{]}

\PY{n}{arg1\PYZus{}arr} \PY{o}{=} \PY{n}{np}\PY{o}{.}\PY{n}{argwhere}\PY{p}{(}\PY{n}{np}\PY{o}{.}\PY{n}{isnan}\PY{p}{(}\PY{n}{x2}\PY{p}{)}\PY{p}{)}\PY{p}{[}\PY{p}{:}\PY{p}{,}\PY{l+m+mi}{0}\PY{p}{]}
\PY{n}{arg2\PYZus{}arr} \PY{o}{=} \PY{n}{np}\PY{o}{.}\PY{n}{argwhere}\PY{p}{(}\PY{n}{np}\PY{o}{.}\PY{n}{isnan}\PY{p}{(}\PY{n}{y2}\PY{p}{)}\PY{p}{)}\PY{p}{[}\PY{p}{:}\PY{p}{,}\PY{l+m+mi}{0}\PY{p}{]}

\PY{c+c1}{\PYZsh{} Find the first spacing in my array of NaN indexes which is greater than one, indicating that these indexes }
\PY{c+c1}{\PYZsh{} come from different chunks of data which have been removed. }
\PY{n}{loc1} \PY{o}{=} \PY{n}{np}\PY{o}{.}\PY{n}{where}\PY{p}{(}\PY{n}{arg1\PYZus{}arr}\PY{p}{[}\PY{l+m+mi}{1}\PY{p}{:}\PY{p}{]}\PY{o}{\PYZhy{}}\PY{n}{arg1\PYZus{}arr}\PY{p}{[}\PY{p}{:}\PY{o}{\PYZhy{}}\PY{l+m+mi}{1}\PY{p}{]} \PY{o}{\PYZgt{}} \PY{l+m+mi}{1}\PY{p}{)}\PY{p}{[}\PY{l+m+mi}{0}\PY{p}{]}\PY{p}{[}\PY{l+m+mi}{0}\PY{p}{]}

\PY{c+c1}{\PYZsh{} This is then the minimum and maximum indexes of data which has not been removed, for the right curve.}
\PY{n}{loc2} \PY{o}{=} \PY{n}{arg1\PYZus{}arr}\PY{p}{[}\PY{n}{loc1}\PY{p}{]}\PY{o}{+}\PY{l+m+mi}{1}
\PY{n}{loc3} \PY{o}{=} \PY{n}{arg1\PYZus{}arr}\PY{p}{[}\PY{n}{loc1}\PY{o}{+}\PY{l+m+mi}{1}\PY{p}{]}\PY{o}{\PYZhy{}}\PY{l+m+mi}{1}
\end{Verbatim}
\end{tcolorbox}

    \begin{tcolorbox}[breakable, size=fbox, boxrule=1pt, pad at break*=1mm,colback=cellbackground, colframe=cellborder]
\prompt{In}{incolor}{95}{\boxspacing}
\begin{Verbatim}[commandchars=\\\{\}]
\PY{c+c1}{\PYZsh{} Find the minimum x value of this curve. }
\PY{n}{x2\PYZus{}min\PYZus{}2} \PY{o}{=} \PY{n}{np}\PY{o}{.}\PY{n}{min}\PY{p}{(}\PY{n}{x2}\PY{p}{[}\PY{n}{loc2}\PY{p}{:}\PY{n}{loc3}\PY{p}{]}\PY{p}{)}

\PY{c+c1}{\PYZsh{} I will now remove all data for the points which have x2 value greater than a point which appears to }
\PY{c+c1}{\PYZsh{} its right in the array, and therefore has a lower value for x1.}

\PY{c+c1}{\PYZsh{} This clears this data for the left curve.}
\PY{k}{for} \PY{n}{i} \PY{o+ow}{in} \PY{n+nb}{range}\PY{p}{(}\PY{n}{arg1}\PY{p}{)}\PY{p}{:}
    \PY{k}{if} \PY{n}{x2}\PY{p}{[}\PY{n}{i}\PY{p}{]} \PY{o}{\PYZgt{}} \PY{n}{x2\PYZus{}min\PYZus{}2}\PY{p}{:}
        \PY{n}{x2}\PY{p}{[}\PY{n}{i}\PY{p}{]} \PY{o}{=} \PY{k+kc}{None}
        \PY{n}{y2}\PY{p}{[}\PY{n}{i}\PY{p}{]} \PY{o}{=} \PY{k+kc}{None}

\PY{c+c1}{\PYZsh{} This is the index of my x2 minimum for the right curve}
\PY{n}{loc4} \PY{o}{=} \PY{n}{np}\PY{o}{.}\PY{n}{argmin}\PY{p}{(}\PY{n}{x2}\PY{p}{[}\PY{n}{loc2}\PY{p}{:}\PY{n}{loc3}\PY{p}{]}\PY{p}{)}\PY{o}{+}\PY{n}{loc2}

\PY{c+c1}{\PYZsh{} And now to clear the data for the right curve.}
\PY{k}{for} \PY{n}{i} \PY{o+ow}{in} \PY{n+nb}{range}\PY{p}{(}\PY{n}{loc2}\PY{p}{,} \PY{n}{loc4}\PY{p}{)}\PY{p}{:}
    \PY{n}{x2}\PY{p}{[}\PY{n}{i}\PY{p}{]} \PY{o}{=} \PY{k+kc}{None}
    \PY{n}{y2}\PY{p}{[}\PY{n}{i}\PY{p}{]} \PY{o}{=} \PY{k+kc}{None}
\end{Verbatim}
\end{tcolorbox}

    And now to plot the results of this data removal.

    \begin{tcolorbox}[breakable, size=fbox, boxrule=1pt, pad at break*=1mm,colback=cellbackground, colframe=cellborder]
\prompt{In}{incolor}{ }{\boxspacing}
\begin{Verbatim}[commandchars=\\\{\}]
\PY{n}{fig}\PY{p}{,} \PY{n}{ax} \PY{o}{=} \PY{n}{plt}\PY{o}{.}\PY{n}{subplots}\PY{p}{(}\PY{p}{)}
\PY{n}{plt}\PY{o}{.}\PY{n}{title}\PY{p}{(}\PY{l+s+sa}{f}\PY{l+s+s2}{\PYZdq{}}\PY{l+s+s2}{Plot of the Change in k due to Nonlinearity vs Linear k}\PY{l+s+se}{\PYZbs{}n}\PY{l+s+s2}{Variation For }\PY{l+s+se}{\PYZbs{}u03BC}\PY{l+s+s2}{ = }\PY{l+s+si}{\PYZob{}}\PY{n}{mu}\PY{p}{[}\PY{n}{jj}\PY{p}{[}\PY{l+m+mi}{0}\PY{p}{]}\PY{p}{]}\PY{l+s+si}{:}\PY{l+s+s2}{.5e}\PY{l+s+si}{\PYZcb{}}\PY{l+s+s2}{\PYZdq{}}\PY{p}{)}
\PY{n}{ax}\PY{o}{.}\PY{n}{set\PYZus{}ylabel}\PY{p}{(}\PY{l+s+s1}{\PYZsq{}}\PY{l+s+s1}{Change in k due to nonlinearity (\PYZhy{}dkn)  (\PYZhy{}10}\PY{l+s+se}{\PYZbs{}u2074}\PY{l+s+s1}{  m}\PY{l+s+se}{\PYZbs{}u207b}\PY{l+s+se}{\PYZbs{}u00b9}\PY{l+s+s1}{)}\PY{l+s+s1}{\PYZsq{}}\PY{p}{)}
\PY{n}{ax}\PY{o}{.}\PY{n}{set\PYZus{}xlabel}\PY{p}{(}\PY{l+s+s2}{\PYZdq{}}\PY{l+s+s2}{Small\PYZhy{}scale variation of linear wave number k (dk (}\PY{l+s+s2}{\PYZpc{}}\PY{l+s+s2}{))}\PY{l+s+s2}{\PYZdq{}}\PY{p}{)}

\PY{n}{ax}\PY{o}{.}\PY{n}{plot}\PY{p}{(}\PY{n}{dk}\PY{p}{,} \PY{o}{\PYZhy{}}\PY{n}{dkn\PYZus{}arr}\PY{p}{[}\PY{p}{:}\PY{p}{,} \PY{n}{jj}\PY{p}{[}\PY{l+m+mi}{0}\PY{p}{]}\PY{p}{]}\PY{o}{/}\PY{l+m+mi}{10000} \PY{o}{\PYZhy{}} \PY{n}{min\PYZus{}dkn}\PY{p}{,} \PY{n}{label} \PY{o}{=} \PY{l+s+s2}{\PYZdq{}}\PY{l+s+s2}{numerical results}\PY{l+s+s2}{\PYZdq{}}\PY{p}{)}

\PY{n}{ax}\PY{o}{.}\PY{n}{plot}\PY{p}{(}\PY{n}{x2}\PY{o}{+}\PY{l+m+mf}{0.065}\PY{p}{,} \PY{n}{y2}\PY{p}{,} \PY{n}{label} \PY{o}{=} \PY{l+s+s2}{\PYZdq{}}\PY{l+s+s2}{model}\PY{l+s+s2}{\PYZdq{}}\PY{p}{)}

\PY{n}{plt}\PY{o}{.}\PY{n}{legend}\PY{p}{(}\PY{p}{)}

\PY{n}{plt}\PY{o}{.}\PY{n}{show}\PY{p}{(}\PY{p}{)}
\end{Verbatim}
\end{tcolorbox}

    \begin{center}
    \adjustimage{max size={0.9\linewidth}{0.9\paperheight}}{figs/bvp_kM_en/output_91_0.png}
    \end{center}
    { \hspace*{\fill} \\}
    
    And the problem with multiple values for dkn for certain dk values is
gone. Now I have an algorithm to create a curve which can fit this fast
oscillation curve decently enough. Now to show it for another curve,
then to move to examining the next shape.

    \begin{tcolorbox}[breakable, size=fbox, boxrule=1pt, pad at break*=1mm,colback=cellbackground, colframe=cellborder]
\prompt{In}{incolor}{ }{\boxspacing}
\begin{Verbatim}[commandchars=\\\{\}]
\PY{n}{jj} \PY{o}{=} \PY{p}{[}\PY{l+m+mi}{40}\PY{p}{]}

\PY{n}{ilim} \PY{o}{=} \PY{n+nb}{len}\PY{p}{(}\PY{n}{jj}\PY{p}{)}

\PY{n}{fig}\PY{p}{,} \PY{n}{ax} \PY{o}{=} \PY{n}{plt}\PY{o}{.}\PY{n}{subplots}\PY{p}{(}\PY{p}{)}
\PY{n}{plt}\PY{o}{.}\PY{n}{title}\PY{p}{(}\PY{l+s+sa}{f}\PY{l+s+s2}{\PYZdq{}}\PY{l+s+s2}{Plot of the Change in k due to Nonlinearity vs Linear k}\PY{l+s+se}{\PYZbs{}n}\PY{l+s+s2}{Variation For }\PY{l+s+se}{\PYZbs{}u03BC}\PY{l+s+s2}{ = }\PY{l+s+si}{\PYZob{}}\PY{n}{mu}\PY{p}{[}\PY{n}{jj}\PY{p}{[}\PY{l+m+mi}{0}\PY{p}{]}\PY{p}{]}\PY{l+s+si}{:}\PY{l+s+s2}{.5e}\PY{l+s+si}{\PYZcb{}}\PY{l+s+s2}{\PYZdq{}}\PY{p}{)}
\PY{n}{ax}\PY{o}{.}\PY{n}{set\PYZus{}ylabel}\PY{p}{(}\PY{l+s+s1}{\PYZsq{}}\PY{l+s+s1}{Change in k due to nonlinearity (\PYZhy{}dkn)  (\PYZhy{}10}\PY{l+s+se}{\PYZbs{}u2074}\PY{l+s+s1}{  m}\PY{l+s+se}{\PYZbs{}u207b}\PY{l+s+se}{\PYZbs{}u00b9}\PY{l+s+s1}{)}\PY{l+s+s1}{\PYZsq{}}\PY{p}{)}
\PY{n}{ax}\PY{o}{.}\PY{n}{set\PYZus{}xlabel}\PY{p}{(}\PY{l+s+s2}{\PYZdq{}}\PY{l+s+s2}{Small\PYZhy{}scale variation of linear wave number k (dk (}\PY{l+s+s2}{\PYZpc{}}\PY{l+s+s2}{))}\PY{l+s+s2}{\PYZdq{}}\PY{p}{)}

\PY{n}{ax}\PY{o}{.}\PY{n}{plot}\PY{p}{(}\PY{n}{dk}\PY{p}{,} \PY{o}{\PYZhy{}}\PY{n}{dkn\PYZus{}arr}\PY{p}{[}\PY{p}{:}\PY{p}{,} \PY{n}{jj}\PY{p}{[}\PY{l+m+mi}{0}\PY{p}{]}\PY{p}{]}\PY{o}{/}\PY{l+m+mi}{10000}\PY{p}{)}

\PY{n}{plt}\PY{o}{.}\PY{n}{show}\PY{p}{(}\PY{p}{)}
\end{Verbatim}
\end{tcolorbox}

    \begin{center}
    \adjustimage{max size={0.9\linewidth}{0.9\paperheight}}{figs/bvp_kM_en/output_93_0.png}
    \end{center}
    { \hspace*{\fill} \\}
    
    \begin{tcolorbox}[breakable, size=fbox, boxrule=1pt, pad at break*=1mm,colback=cellbackground, colframe=cellborder]
\prompt{In}{incolor}{ }{\boxspacing}
\begin{Verbatim}[commandchars=\\\{\}]
\PY{n}{min\PYZus{}dkn} \PY{o}{=} \PY{n}{np}\PY{o}{.}\PY{n}{min}\PY{p}{(}\PY{o}{\PYZhy{}}\PY{n}{dkn\PYZus{}arr}\PY{p}{[}\PY{l+m+mi}{2}\PY{p}{:}\PY{p}{,} \PY{n}{jj}\PY{p}{[}\PY{l+m+mi}{0}\PY{p}{]}\PY{p}{]}\PY{o}{/}\PY{l+m+mi}{10000}\PY{p}{)}
\PY{n}{per} \PY{o}{=} \PY{l+m+mi}{2}\PY{o}{/}\PY{l+m+mf}{0.9}

\PY{n}{model1} \PY{o}{=} \PY{n}{min\PYZus{}dkn}\PY{o}{*}\PY{n}{np}\PY{o}{.}\PY{n}{power}\PY{p}{(}\PY{n}{np}\PY{o}{.}\PY{n}{tan}\PY{p}{(}\PY{l+m+mf}{0.9}\PY{o}{*}\PY{n}{pi}\PY{o}{*}\PY{n}{dk}\PY{o}{/}\PY{l+m+mi}{2}\PY{p}{)}\PY{p}{,}\PY{l+m+mi}{2}\PY{p}{)}

\PY{k}{for} \PY{n}{i}\PY{p}{,}\PY{n}{j} \PY{o+ow}{in} \PY{n+nb}{enumerate}\PY{p}{(}\PY{n}{model1}\PY{p}{)}\PY{p}{:}
    \PY{k}{if} \PY{n}{j} \PY{o}{\PYZgt{}} \PY{l+m+mi}{30}\PY{p}{:} \PY{n}{model1}\PY{p}{[}\PY{n}{i}\PY{p}{]} \PY{o}{=} \PY{l+m+mi}{30}

\PY{n}{fig}\PY{p}{,} \PY{n}{ax} \PY{o}{=} \PY{n}{plt}\PY{o}{.}\PY{n}{subplots}\PY{p}{(}\PY{p}{)}
\PY{n}{plt}\PY{o}{.}\PY{n}{title}\PY{p}{(}\PY{l+s+sa}{f}\PY{l+s+s2}{\PYZdq{}}\PY{l+s+s2}{Plot of the Change in k due to Nonlinearity vs Linear k}\PY{l+s+se}{\PYZbs{}n}\PY{l+s+s2}{Variation For }\PY{l+s+se}{\PYZbs{}u03BC}\PY{l+s+s2}{ = }\PY{l+s+si}{\PYZob{}}\PY{n}{mu}\PY{p}{[}\PY{n}{jj}\PY{p}{[}\PY{l+m+mi}{0}\PY{p}{]}\PY{p}{]}\PY{l+s+si}{:}\PY{l+s+s2}{.5e}\PY{l+s+si}{\PYZcb{}}\PY{l+s+s2}{\PYZdq{}}\PY{p}{)}
\PY{n}{ax}\PY{o}{.}\PY{n}{set\PYZus{}ylabel}\PY{p}{(}\PY{l+s+s1}{\PYZsq{}}\PY{l+s+s1}{Change in k due to nonlinearity (\PYZhy{}dkn)  (\PYZhy{}10}\PY{l+s+se}{\PYZbs{}u2074}\PY{l+s+s1}{  m}\PY{l+s+se}{\PYZbs{}u207b}\PY{l+s+se}{\PYZbs{}u00b9}\PY{l+s+s1}{)}\PY{l+s+s1}{\PYZsq{}}\PY{p}{)}
\PY{n}{ax}\PY{o}{.}\PY{n}{set\PYZus{}xlabel}\PY{p}{(}\PY{l+s+s2}{\PYZdq{}}\PY{l+s+s2}{Small\PYZhy{}scale variation of linear wave number k (dk (}\PY{l+s+s2}{\PYZpc{}}\PY{l+s+s2}{))}\PY{l+s+s2}{\PYZdq{}}\PY{p}{)}

\PY{n}{ax}\PY{o}{.}\PY{n}{plot}\PY{p}{(}\PY{n}{dk}\PY{p}{,} \PY{o}{\PYZhy{}}\PY{n}{dkn\PYZus{}arr}\PY{p}{[}\PY{p}{:}\PY{p}{,} \PY{n}{jj}\PY{p}{[}\PY{l+m+mi}{0}\PY{p}{]}\PY{p}{]}\PY{o}{/}\PY{l+m+mi}{10000} \PY{o}{\PYZhy{}} \PY{n}{min\PYZus{}dkn}\PY{p}{,} \PY{n}{label} \PY{o}{=} \PY{l+s+s2}{\PYZdq{}}\PY{l+s+s2}{numerical results}\PY{l+s+s2}{\PYZdq{}}\PY{p}{)}
\PY{n}{ax}\PY{o}{.}\PY{n}{plot}\PY{p}{(}\PY{n}{dk}\PY{p}{,} \PY{n}{model1}\PY{p}{,} \PY{n}{label} \PY{o}{=} \PY{l+s+s2}{\PYZdq{}}\PY{l+s+s2}{model}\PY{l+s+s2}{\PYZdq{}}\PY{p}{)}

\PY{n}{plt}\PY{o}{.}\PY{n}{legend}\PY{p}{(}\PY{p}{)}

\PY{n}{plt}\PY{o}{.}\PY{n}{show}\PY{p}{(}\PY{p}{)}
\end{Verbatim}
\end{tcolorbox}

    \begin{center}
    \adjustimage{max size={0.9\linewidth}{0.9\paperheight}}{figs/bvp_kM_en/output_94_0.png}
    \end{center}
    { \hspace*{\fill} \\}
    
    \begin{tcolorbox}[breakable, size=fbox, boxrule=1pt, pad at break*=1mm,colback=cellbackground, colframe=cellborder]
\prompt{In}{incolor}{100}{\boxspacing}
\begin{Verbatim}[commandchars=\\\{\}]
\PY{n}{model1} \PY{o}{=} \PY{o}{\PYZhy{}}\PY{n}{min\PYZus{}dkn}\PY{o}{*}\PY{n}{np}\PY{o}{.}\PY{n}{power}\PY{p}{(}\PY{n}{np}\PY{o}{.}\PY{n}{tan}\PY{p}{(}\PY{l+m+mf}{0.9}\PY{o}{*}\PY{n}{pi}\PY{o}{*}\PY{n}{dk}\PY{o}{/}\PY{l+m+mi}{2}\PY{p}{)}\PY{p}{,}\PY{l+m+mi}{2}\PY{p}{)}

\PY{n}{x2} \PY{o}{=} \PY{n}{np}\PY{o}{.}\PY{n}{zeros}\PY{p}{(}\PY{n}{grid\PYZus{}size}\PY{p}{)}
\PY{n}{y2} \PY{o}{=} \PY{n}{np}\PY{o}{.}\PY{n}{zeros}\PY{p}{(}\PY{n}{grid\PYZus{}size}\PY{p}{)}
\PY{n}{theta} \PY{o}{=} \PY{l+m+mf}{0.0199}

\PY{k}{for} \PY{n}{i}\PY{p}{,}\PY{n}{y1} \PY{o+ow}{in} \PY{n+nb}{enumerate}\PY{p}{(}\PY{n}{model1}\PY{p}{)}\PY{p}{:}
    \PY{n}{x1} \PY{o}{=} \PY{n}{dk}\PY{p}{[}\PY{n}{i}\PY{p}{]}
    \PY{n}{x2}\PY{p}{[}\PY{n}{i}\PY{p}{]} \PY{o}{=} \PY{n}{np}\PY{o}{.}\PY{n}{cos}\PY{p}{(}\PY{n}{theta}\PY{p}{)}\PY{o}{*} \PY{n}{x1} \PY{o}{\PYZhy{}} \PY{n}{np}\PY{o}{.}\PY{n}{sin}\PY{p}{(}\PY{n}{theta}\PY{p}{)}\PY{o}{*}\PY{n}{y1}
    \PY{n}{y2}\PY{p}{[}\PY{n}{i}\PY{p}{]} \PY{o}{=} \PY{n}{np}\PY{o}{.}\PY{n}{sin}\PY{p}{(}\PY{n}{theta}\PY{p}{)}\PY{o}{*} \PY{n}{x1} \PY{o}{\PYZhy{}} \PY{n}{np}\PY{o}{.}\PY{n}{cos}\PY{p}{(}\PY{n}{theta}\PY{p}{)}\PY{o}{*}\PY{n}{y1}


\PY{k}{for} \PY{n}{i}\PY{p}{,}\PY{n}{j} \PY{o+ow}{in} \PY{n+nb}{enumerate}\PY{p}{(}\PY{n}{y2}\PY{p}{)}\PY{p}{:}
    \PY{k}{if} \PY{n}{j} \PY{o}{\PYZgt{}} \PY{l+m+mi}{30}\PY{p}{:} \PY{n}{y2}\PY{p}{[}\PY{n}{i}\PY{p}{]} \PY{o}{=} \PY{k+kc}{None}
\PY{k}{for} \PY{n}{i}\PY{p}{,} \PY{n}{j} \PY{o+ow}{in} \PY{n+nb}{enumerate}\PY{p}{(}\PY{n}{x2}\PY{p}{)}\PY{p}{:}
    \PY{k}{if} \PY{n}{j} \PY{o}{\PYZgt{}} \PY{l+m+mi}{3}\PY{p}{:} \PY{n}{x2}\PY{p}{[}\PY{n}{i}\PY{p}{]}\PY{o}{=} \PY{k+kc}{None}

\PY{k}{for} \PY{n}{i} \PY{o+ow}{in} \PY{n+nb}{range}\PY{p}{(}\PY{n}{grid\PYZus{}size}\PY{p}{)}\PY{p}{:}
    \PY{k}{if} \PY{n}{np}\PY{o}{.}\PY{n}{isnan}\PY{p}{(}\PY{n}{y2}\PY{p}{[}\PY{n}{i}\PY{p}{]}\PY{p}{)}\PY{p}{:} \PY{n}{x2}\PY{p}{[}\PY{n}{i}\PY{p}{]} \PY{o}{=} \PY{k+kc}{None}
    \PY{k}{if} \PY{n}{np}\PY{o}{.}\PY{n}{isnan}\PY{p}{(}\PY{n}{x2}\PY{p}{[}\PY{n}{i}\PY{p}{]}\PY{p}{)}\PY{p}{:} \PY{n}{y2}\PY{p}{[}\PY{n}{i}\PY{p}{]} \PY{o}{=} \PY{k+kc}{None}

\PY{n}{arg1} \PY{o}{=} \PY{n}{np}\PY{o}{.}\PY{n}{argwhere}\PY{p}{(}\PY{n}{np}\PY{o}{.}\PY{n}{isnan}\PY{p}{(}\PY{n}{x2}\PY{p}{)}\PY{p}{)}\PY{p}{[}\PY{l+m+mi}{0}\PY{p}{]}\PY{p}{[}\PY{l+m+mi}{0}\PY{p}{]}
\PY{n}{arg2} \PY{o}{=} \PY{n}{np}\PY{o}{.}\PY{n}{argwhere}\PY{p}{(}\PY{n}{np}\PY{o}{.}\PY{n}{isnan}\PY{p}{(}\PY{n}{y2}\PY{p}{)}\PY{p}{)}\PY{p}{[}\PY{l+m+mi}{0}\PY{p}{]}\PY{p}{[}\PY{l+m+mi}{0}\PY{p}{]}
\PY{n}{arg1\PYZus{}arr} \PY{o}{=} \PY{n}{np}\PY{o}{.}\PY{n}{argwhere}\PY{p}{(}\PY{n}{np}\PY{o}{.}\PY{n}{isnan}\PY{p}{(}\PY{n}{x2}\PY{p}{)}\PY{p}{)}\PY{p}{[}\PY{p}{:}\PY{p}{,}\PY{l+m+mi}{0}\PY{p}{]}
\PY{n}{arg2\PYZus{}arr} \PY{o}{=} \PY{n}{np}\PY{o}{.}\PY{n}{argwhere}\PY{p}{(}\PY{n}{np}\PY{o}{.}\PY{n}{isnan}\PY{p}{(}\PY{n}{y2}\PY{p}{)}\PY{p}{)}\PY{p}{[}\PY{p}{:}\PY{p}{,}\PY{l+m+mi}{0}\PY{p}{]}
\PY{n}{loc1} \PY{o}{=} \PY{n}{np}\PY{o}{.}\PY{n}{where}\PY{p}{(}\PY{n}{arg1\PYZus{}arr}\PY{p}{[}\PY{l+m+mi}{1}\PY{p}{:}\PY{p}{]}\PY{o}{\PYZhy{}}\PY{n}{arg1\PYZus{}arr}\PY{p}{[}\PY{p}{:}\PY{o}{\PYZhy{}}\PY{l+m+mi}{1}\PY{p}{]} \PY{o}{\PYZgt{}} \PY{l+m+mi}{1}\PY{p}{)}\PY{p}{[}\PY{l+m+mi}{0}\PY{p}{]}\PY{p}{[}\PY{l+m+mi}{0}\PY{p}{]}
\PY{n}{loc2} \PY{o}{=} \PY{n}{arg1\PYZus{}arr}\PY{p}{[}\PY{n}{loc1}\PY{p}{]}\PY{o}{+}\PY{l+m+mi}{1}
\PY{n}{loc3} \PY{o}{=} \PY{n}{arg1\PYZus{}arr}\PY{p}{[}\PY{n}{loc1}\PY{o}{+}\PY{l+m+mi}{1}\PY{p}{]}\PY{o}{\PYZhy{}}\PY{l+m+mi}{1}
\PY{n}{loc2}\PY{p}{,}\PY{n}{loc3}
\PY{n}{x2\PYZus{}min\PYZus{}2} \PY{o}{=} \PY{n}{np}\PY{o}{.}\PY{n}{min}\PY{p}{(}\PY{n}{x2}\PY{p}{[}\PY{n}{loc2}\PY{p}{:}\PY{n}{loc3}\PY{p}{]}\PY{p}{)}
\PY{n}{x2}\PY{p}{[}\PY{n}{loc2}\PY{p}{:}\PY{n}{loc3}\PY{p}{]}
\PY{k}{for} \PY{n}{i} \PY{o+ow}{in} \PY{n+nb}{range}\PY{p}{(}\PY{n}{arg1}\PY{p}{)}\PY{p}{:}
    \PY{k}{if} \PY{n}{x2}\PY{p}{[}\PY{n}{i}\PY{p}{]} \PY{o}{\PYZgt{}} \PY{n}{x2\PYZus{}min\PYZus{}2}\PY{p}{:}
        \PY{n}{x2}\PY{p}{[}\PY{n}{i}\PY{p}{]} \PY{o}{=} \PY{k+kc}{None}
        \PY{n}{y2}\PY{p}{[}\PY{n}{i}\PY{p}{]} \PY{o}{=} \PY{k+kc}{None}

\PY{n}{loc4} \PY{o}{=} \PY{n}{np}\PY{o}{.}\PY{n}{argmin}\PY{p}{(}\PY{n}{x2}\PY{p}{[}\PY{n}{loc2}\PY{p}{:}\PY{n}{loc3}\PY{p}{]}\PY{p}{)}\PY{o}{+}\PY{n}{loc2}
\PY{k}{for} \PY{n}{i} \PY{o+ow}{in} \PY{n+nb}{range}\PY{p}{(}\PY{n}{loc2}\PY{p}{,} \PY{n}{loc4}\PY{p}{)}\PY{p}{:}
    \PY{n}{x2}\PY{p}{[}\PY{n}{i}\PY{p}{]} \PY{o}{=} \PY{k+kc}{None}
    \PY{n}{y2}\PY{p}{[}\PY{n}{i}\PY{p}{]} \PY{o}{=} \PY{k+kc}{None}

\PY{n}{arg1\PYZus{}find\PYZus{}arr} \PY{o}{=} \PY{n}{arg1\PYZus{}arr}\PY{p}{[}\PY{l+m+mi}{1}\PY{p}{:}\PY{p}{]}\PY{o}{\PYZhy{}}\PY{n}{arg1\PYZus{}arr}\PY{p}{[}\PY{p}{:}\PY{o}{\PYZhy{}}\PY{l+m+mi}{1}\PY{p}{]}
\PY{n}{ind1} \PY{o}{=} \PY{n}{np}\PY{o}{.}\PY{n}{min}\PY{p}{(}\PY{p}{[}\PY{n}{arg1}\PY{p}{,}\PY{n}{arg2}\PY{p}{]}\PY{p}{)}\PY{o}{\PYZhy{}}\PY{l+m+mi}{1}
\end{Verbatim}
\end{tcolorbox}

    \begin{tcolorbox}[breakable, size=fbox, boxrule=1pt, pad at break*=1mm,colback=cellbackground, colframe=cellborder]
\prompt{In}{incolor}{ }{\boxspacing}
\begin{Verbatim}[commandchars=\\\{\}]
\PY{n}{fig}\PY{p}{,} \PY{n}{ax} \PY{o}{=} \PY{n}{plt}\PY{o}{.}\PY{n}{subplots}\PY{p}{(}\PY{p}{)}
\PY{n}{plt}\PY{o}{.}\PY{n}{title}\PY{p}{(}\PY{l+s+sa}{f}\PY{l+s+s2}{\PYZdq{}}\PY{l+s+s2}{Plot of the Change in k due to Nonlinearity vs Linear k}\PY{l+s+se}{\PYZbs{}n}\PY{l+s+s2}{Variation For }\PY{l+s+se}{\PYZbs{}u03BC}\PY{l+s+s2}{ = }\PY{l+s+si}{\PYZob{}}\PY{n}{mu}\PY{p}{[}\PY{n}{jj}\PY{p}{[}\PY{l+m+mi}{0}\PY{p}{]}\PY{p}{]}\PY{l+s+si}{:}\PY{l+s+s2}{.5e}\PY{l+s+si}{\PYZcb{}}\PY{l+s+s2}{\PYZdq{}}\PY{p}{)}
\PY{n}{ax}\PY{o}{.}\PY{n}{set\PYZus{}ylabel}\PY{p}{(}\PY{l+s+s1}{\PYZsq{}}\PY{l+s+s1}{Change in k due to nonlinearity (\PYZhy{}dkn)  (\PYZhy{}10}\PY{l+s+se}{\PYZbs{}u2074}\PY{l+s+s1}{  m}\PY{l+s+se}{\PYZbs{}u207b}\PY{l+s+se}{\PYZbs{}u00b9}\PY{l+s+s1}{)}\PY{l+s+s1}{\PYZsq{}}\PY{p}{)}
\PY{n}{ax}\PY{o}{.}\PY{n}{set\PYZus{}xlabel}\PY{p}{(}\PY{l+s+s2}{\PYZdq{}}\PY{l+s+s2}{Small\PYZhy{}scale variation of linear wave number k (dk (}\PY{l+s+s2}{\PYZpc{}}\PY{l+s+s2}{))}\PY{l+s+s2}{\PYZdq{}}\PY{p}{)}

\PY{n}{ax}\PY{o}{.}\PY{n}{plot}\PY{p}{(}\PY{n}{dk}\PY{p}{,} \PY{o}{\PYZhy{}}\PY{n}{dkn\PYZus{}arr}\PY{p}{[}\PY{p}{:}\PY{p}{,} \PY{n}{jj}\PY{p}{[}\PY{l+m+mi}{0}\PY{p}{]}\PY{p}{]}\PY{o}{/}\PY{l+m+mi}{10000} \PY{o}{\PYZhy{}} \PY{n}{min\PYZus{}dkn}\PY{p}{,} \PY{n}{label} \PY{o}{=} \PY{l+s+s2}{\PYZdq{}}\PY{l+s+s2}{numerical results}\PY{l+s+s2}{\PYZdq{}}\PY{p}{)}

\PY{n}{ax}\PY{o}{.}\PY{n}{plot}\PY{p}{(}\PY{n}{x2}\PY{o}{+}\PY{l+m+mf}{0.033}\PY{p}{,} \PY{n}{y2}\PY{p}{,} \PY{n}{label} \PY{o}{=} \PY{l+s+s2}{\PYZdq{}}\PY{l+s+s2}{model}\PY{l+s+s2}{\PYZdq{}}\PY{p}{)}

\PY{n}{plt}\PY{o}{.}\PY{n}{legend}\PY{p}{(}\PY{p}{)}

\PY{n}{plt}\PY{o}{.}\PY{n}{show}\PY{p}{(}\PY{p}{)}
\end{Verbatim}
\end{tcolorbox}

    \begin{center}
    \adjustimage{max size={0.9\linewidth}{0.9\paperheight}}{figs/bvp_kM_en/output_96_0.png}
    \end{center}
    { \hspace*{\fill} \\}
    
    My values for the amplitude, rotation angle, and x translation are
different, but it still seems to work.

    \hypertarget{b-analysis-for-dkn-vs-r}{%
\section{\texorpdfstring{B Analysis for \(dkn\) vs
\(R\)}{B Analysis for dkn vs R}}\label{b-analysis-for-dkn-vs-r}}

    And now for the \(B\) analysis, which shows all types of oscillations in
the changes of fast and slow oscillation wave number.

    \begin{tcolorbox}[breakable, size=fbox, boxrule=1pt, pad at break*=1mm,colback=cellbackground, colframe=cellborder]
\prompt{In}{incolor}{102}{\boxspacing}
\begin{Verbatim}[commandchars=\\\{\}]
\PY{n}{gind} \PY{o}{=} \PY{l+m+mi}{0}
\end{Verbatim}
\end{tcolorbox}

    \begin{tcolorbox}[breakable, size=fbox, boxrule=1pt, pad at break*=1mm,colback=cellbackground, colframe=cellborder]
\prompt{In}{incolor}{103}{\boxspacing}
\begin{Verbatim}[commandchars=\\\{\}]
\PY{n}{g} \PY{o}{=} \PY{n}{sns}\PY{o}{.}\PY{n}{PairGrid}\PY{p}{(}\PY{n}{tbl\PYZus{}red}\PY{p}{[}\PY{l+m+mi}{0}\PY{p}{]}\PY{p}{,} \PY{n}{diag\PYZus{}sharey}\PY{o}{=}\PY{k+kc}{False}\PY{p}{,} \PY{n}{corner}\PY{o}{=}\PY{k+kc}{True}\PY{p}{)}

\PY{n}{g}\PY{o}{.}\PY{n}{map\PYZus{}diag}\PY{p}{(}\PY{n}{sns}\PY{o}{.}\PY{n}{kdeplot}\PY{p}{)}
\PY{n}{g}\PY{o}{.}\PY{n}{map\PYZus{}lower}\PY{p}{(}\PY{n}{sns}\PY{o}{.}\PY{n}{scatterplot}\PY{p}{)}
\end{Verbatim}
\end{tcolorbox}

            \begin{tcolorbox}[breakable, size=fbox, boxrule=.5pt, pad at break*=1mm, opacityfill=0]
\prompt{Out}{outcolor}{103}{\boxspacing}
\begin{Verbatim}[commandchars=\\\{\}]
<seaborn.axisgrid.PairGrid at 0x1d1a3ded6a0>
\end{Verbatim}
\end{tcolorbox}
        
    \begin{center}
    \adjustimage{max size={0.9\linewidth}{0.9\paperheight}}{figs/bvp_kM_en/output_101_1.png}
    \end{center}
    { \hspace*{\fill} \\}
    
    We start by loading our data.

    \begin{tcolorbox}[breakable, size=fbox, boxrule=1pt, pad at break*=1mm,colback=cellbackground, colframe=cellborder]
\prompt{In}{incolor}{104}{\boxspacing}
\begin{Verbatim}[commandchars=\\\{\}]
\PY{k}{del} \PY{n}{dkn\PYZus{}arr}\PY{p}{,} \PY{n}{dMn\PYZus{}arr}

\PY{n}{dkn\PYZus{}arr} \PY{o}{=} \PY{n}{np}\PY{o}{.}\PY{n}{zeros}\PY{p}{(}\PY{p}{(}\PY{n}{grid\PYZus{}size}\PY{p}{,} \PY{n}{grid\PYZus{}size}\PY{p}{)}\PY{p}{)}
\PY{n}{dMn\PYZus{}arr} \PY{o}{=} \PY{n}{np}\PY{o}{.}\PY{n}{zeros}\PY{p}{(}\PY{p}{(}\PY{n}{grid\PYZus{}size}\PY{p}{,} \PY{n}{grid\PYZus{}size}\PY{p}{)}\PY{p}{)}

\PY{k}{for} \PY{n}{r}\PY{p}{,} \PY{n}{rr} \PY{o+ow}{in} \PY{n+nb}{enumerate}\PY{p}{(}\PY{n}{R\PYZus{}g}\PY{p}{[}\PY{n}{gind}\PY{p}{]}\PY{p}{)}\PY{p}{:}
    \PY{k}{for} \PY{n}{b}\PY{p}{,} \PY{n}{bb} \PY{o+ow}{in} \PY{n+nb}{enumerate}\PY{p}{(}\PY{n}{B\PYZus{}g}\PY{p}{[}\PY{n}{gind}\PY{p}{]}\PY{p}{)}\PY{p}{:}
        \PY{n}{tbl\PYZus{}sel} \PY{o}{=} \PY{n}{tbl\PYZus{}red}\PY{p}{[}\PY{n}{gind}\PY{p}{]}\PY{p}{[}\PY{p}{(}\PY{n}{tbl\PYZus{}red}\PY{p}{[}\PY{n}{gind}\PY{p}{]}\PY{p}{[}\PY{l+s+s2}{\PYZdq{}}\PY{l+s+s2}{R}\PY{l+s+s2}{\PYZdq{}}\PY{p}{]}\PY{o}{==}\PY{n}{rr}\PY{p}{)} \PY{o}{\PYZam{}}\PY{p}{(}\PY{n}{tbl\PYZus{}red}\PY{p}{[}\PY{n}{gind}\PY{p}{]}\PY{p}{[}\PY{l+s+s2}{\PYZdq{}}\PY{l+s+s2}{B}\PY{l+s+s2}{\PYZdq{}}\PY{p}{]}\PY{o}{==}\PY{n}{bb}\PY{p}{)}\PY{p}{]}
        \PY{n}{dkn\PYZus{}arr}\PY{p}{[}\PY{n}{r}\PY{p}{,}\PY{n}{b}\PY{p}{]} \PY{o}{=} \PY{n}{tbl\PYZus{}sel}\PY{p}{[}\PY{l+s+s2}{\PYZdq{}}\PY{l+s+s2}{dkn}\PY{l+s+s2}{\PYZdq{}}\PY{p}{]}\PY{o}{.}\PY{n}{to\PYZus{}numpy}\PY{p}{(}\PY{p}{)}\PY{p}{[}\PY{l+m+mi}{0}\PY{p}{]}
        \PY{n}{dMn\PYZus{}arr}\PY{p}{[}\PY{n}{r}\PY{p}{,}\PY{n}{b}\PY{p}{]} \PY{o}{=} \PY{n}{tbl\PYZus{}sel}\PY{p}{[}\PY{l+s+s2}{\PYZdq{}}\PY{l+s+s2}{dMn}\PY{l+s+s2}{\PYZdq{}}\PY{p}{]}\PY{o}{.}\PY{n}{to\PYZus{}numpy}\PY{p}{(}\PY{p}{)}\PY{p}{[}\PY{l+m+mi}{0}\PY{p}{]}
\end{Verbatim}
\end{tcolorbox}

    \begin{tcolorbox}[breakable, size=fbox, boxrule=1pt, pad at break*=1mm,colback=cellbackground, colframe=cellborder]
\prompt{In}{incolor}{105}{\boxspacing}
\begin{Verbatim}[commandchars=\\\{\}]
\PY{c+c1}{\PYZsh{}B and B will be rescaled to show only the difference from our }
\PY{n}{scale\PYZus{}R} \PY{o}{=} \PY{l+m+mi}{1}\PY{o}{/}\PY{l+m+mi}{10000}
\PY{n}{R0} \PY{o}{=} \PY{l+m+mf}{0.9}
\PY{n}{B0} \PY{o}{=} \PY{l+m+mf}{0.5}
\PY{n}{dk} \PY{o}{=} \PY{n}{dk\PYZus{}g}\PY{p}{[}\PY{n}{gind}\PY{p}{]}\PY{p}{[}\PY{l+m+mi}{0}\PY{p}{]}

\PY{n}{R} \PY{o}{=} \PY{p}{(}\PY{n}{R\PYZus{}g}\PY{p}{[}\PY{n}{gind}\PY{p}{]} \PY{o}{\PYZhy{}} \PY{n}{R0}\PY{p}{)}\PY{o}{/}\PY{n}{scale\PYZus{}R}
\PY{n}{B} \PY{o}{=} \PY{n}{B\PYZus{}g}\PY{p}{[}\PY{n}{gind}\PY{p}{]} \PY{o}{\PYZhy{}} \PY{n}{B0}
\end{Verbatim}
\end{tcolorbox}

    \begin{tcolorbox}[breakable, size=fbox, boxrule=1pt, pad at break*=1mm,colback=cellbackground, colframe=cellborder]
\prompt{In}{incolor}{106}{\boxspacing}
\begin{Verbatim}[commandchars=\\\{\}]
\PY{c+c1}{\PYZsh{} We will learn here that our oscillations are based on the fractional part of B * R\PYZca{}2 * pi\PYZca{}2}
\PY{c+c1}{\PYZsh{} Here are a few functions for calculating our period boundaries.}

\PY{k}{def} \PY{n+nf}{find\PYZus{}count\PYZus{}B\PYZus{}osc}\PY{p}{(}\PY{n}{dR}\PY{p}{,} \PY{n}{dB}\PY{p}{)}\PY{p}{:}
    \PY{n}{R} \PY{o}{=} \PY{n}{dR} \PY{o}{*} \PY{n}{scale\PYZus{}R}  \PY{o}{+} \PY{n}{R0}
    \PY{n}{B} \PY{o}{=} \PY{n}{dB} \PY{o}{+} \PY{n}{B0}
    \PY{n}{count} \PY{o}{=}  \PY{n}{B} \PY{o}{*} \PY{n}{R}\PY{o}{*}\PY{o}{*}\PY{l+m+mi}{2} \PY{o}{*} \PY{n}{pi}\PY{o}{*}\PY{o}{*}\PY{l+m+mi}{2}
    \PY{k}{return} \PY{n}{count}

\PY{k}{def} \PY{n+nf}{find\PYZus{}count\PYZus{}k\PYZus{}osc}\PY{p}{(}\PY{n}{dR}\PY{p}{)}\PY{p}{:}
    \PY{n}{k} \PY{o}{=} \PY{p}{(}\PY{n}{k0} \PY{o}{+} \PY{n}{dk} \PY{o}{*} \PY{l+m+mf}{1e6} \PY{o}{/} \PY{l+m+mi}{2}\PY{p}{)} \PY{o}{/} \PY{l+m+mf}{1e6}
    \PY{n}{R} \PY{o}{=} \PY{n}{dR} \PY{o}{*} \PY{n}{scale\PYZus{}R} \PY{o}{+} \PY{n}{R0}
    \PY{k}{return} \PY{n}{k} \PY{o}{*} \PY{n}{R}

\PY{k}{def} \PY{n+nf}{find\PYZus{}dB\PYZus{}B\PYZus{}osc}\PY{p}{(}\PY{n}{n}\PY{p}{,} \PY{n}{dR} \PY{o}{=} \PY{l+m+mi}{0}\PY{p}{)}\PY{p}{:}
    \PY{n}{R} \PY{o}{=} \PY{n}{dR} \PY{o}{*} \PY{n}{scale\PYZus{}R}  \PY{o}{+} \PY{n}{R0}
    \PY{c+c1}{\PYZsh{} n = B * R\PYZca{}2 * pi\PYZca{}2}
    \PY{n}{dB} \PY{o}{=} \PY{n}{n} \PY{o}{/} \PY{p}{(}\PY{n}{R}\PY{o}{*}\PY{o}{*}\PY{l+m+mi}{2} \PY{o}{*} \PY{n}{pi}\PY{o}{*}\PY{o}{*}\PY{l+m+mi}{2}\PY{p}{)} \PY{o}{\PYZhy{}} \PY{n}{B0}
    \PY{k}{return} \PY{n}{dB}

\PY{k}{def} \PY{n+nf}{find\PYZus{}dR\PYZus{}k\PYZus{}osc}\PY{p}{(}\PY{n}{n}\PY{p}{)}\PY{p}{:}
    \PY{c+c1}{\PYZsh{} n = k0 * R0 + k0 dR + dk R0 + dk dR}
    \PY{n}{k0\PYZus{}diff} \PY{o}{=} \PY{n}{k0} \PY{o}{*} \PY{n}{R0} \PY{o}{\PYZpc{}} \PY{l+m+mi}{1}
    \PY{c+c1}{\PYZsh{} n \PYZhy{} k0\PYZus{}diff = k0 dR + dk R0}
    \PY{n}{dk2} \PY{o}{=} \PY{n}{dk} \PY{o}{/} \PY{l+m+mi}{2}
    \PY{n}{dR} \PY{o}{=} \PY{p}{(}\PY{n}{n} \PY{o}{\PYZhy{}} \PY{n}{k0\PYZus{}diff} \PY{o}{\PYZhy{}} \PY{n}{dk2} \PY{o}{*} \PY{n}{R0}\PY{p}{)} \PY{o}{/} \PY{n}{k0} \PY{o}{/} \PY{n}{scale\PYZus{}R} \PY{o}{*} \PY{l+m+mf}{1e6}
    \PY{k}{return} \PY{n}{dR}
\end{Verbatim}
\end{tcolorbox}

    I will start with my \(dkn\) vs \(dR\) graphs.

    \begin{tcolorbox}[breakable, size=fbox, boxrule=1pt, pad at break*=1mm,colback=cellbackground, colframe=cellborder]
\prompt{In}{incolor}{ }{\boxspacing}
\begin{Verbatim}[commandchars=\\\{\}]
\PY{c+c1}{\PYZsh{} Again, choose the indexes beforehand. This is a B value on my B node, as I will see later.}

\PY{n}{Bi} \PY{o}{=} \PY{p}{[}\PY{l+m+mi}{0}\PY{p}{]}

\PY{n}{ilim} \PY{o}{=} \PY{n+nb}{len}\PY{p}{(}\PY{n}{Bi}\PY{p}{)}

\PY{n}{fig}\PY{p}{,} \PY{n}{ax} \PY{o}{=} \PY{n}{plt}\PY{o}{.}\PY{n}{subplots}\PY{p}{(}\PY{p}{)}
\PY{n}{plt}\PY{o}{.}\PY{n}{title}\PY{p}{(}\PY{l+s+sa}{f}\PY{l+s+s2}{\PYZdq{}}\PY{l+s+s2}{Plot of the Change in k Due to Nonlinearity vs R}\PY{l+s+se}{\PYZbs{}n}\PY{l+s+s2}{for dk = 0.45 / 2 * 10}\PY{l+s+se}{\PYZbs{}u2076}\PY{l+s+s2}{ m}\PY{l+s+se}{\PYZbs{}u207b}\PY{l+s+se}{\PYZbs{}u00b9}\PY{l+s+s2}{ and B = }\PY{l+s+si}{\PYZob{}}\PY{n}{B0}\PY{o}{+}\PY{n}{Bi}\PY{p}{[}\PY{l+m+mi}{0}\PY{p}{]}\PY{l+s+si}{:}\PY{l+s+s2}{.4f}\PY{l+s+si}{\PYZcb{}}\PY{l+s+s2}{\PYZdq{}}\PY{p}{)}
\PY{n}{ax}\PY{o}{.}\PY{n}{set\PYZus{}ylabel}\PY{p}{(}\PY{l+s+s1}{\PYZsq{}}\PY{l+s+s1}{Change in k (dkn) due to nonlinearity (10}\PY{l+s+se}{\PYZbs{}u2075}\PY{l+s+s1}{ m}\PY{l+s+se}{\PYZbs{}u207b}\PY{l+s+se}{\PYZbs{}u00b9}\PY{l+s+s1}{)}\PY{l+s+s1}{\PYZsq{}}\PY{p}{)}
\PY{n}{ax}\PY{o}{.}\PY{n}{set\PYZus{}xlabel}\PY{p}{(}\PY{l+s+s2}{\PYZdq{}}\PY{l+s+s2}{Change dr in radius R of ring (R = 9 * 10}\PY{l+s+se}{\PYZbs{}u207b}\PY{l+s+se}{\PYZbs{}u2077}\PY{l+s+s2}{ m + dr * 10}\PY{l+s+se}{\PYZbs{}u207b}\PY{l+s+se}{\PYZbs{}u00b9}\PY{l+s+se}{\PYZbs{}u2070}\PY{l+s+s2}{ m)}\PY{l+s+s2}{\PYZdq{}}\PY{p}{)}

\PY{k}{for} \PY{n}{i} \PY{o+ow}{in} \PY{n+nb}{range}\PY{p}{(}\PY{n}{ilim}\PY{p}{)}\PY{p}{:}
    \PY{n}{ax}\PY{o}{.}\PY{n}{plot}\PY{p}{(}\PY{n}{R}\PY{p}{,} \PY{n}{dkn\PYZus{}arr}\PY{p}{[}\PY{p}{:}\PY{p}{,} \PY{n}{Bi}\PY{p}{[}\PY{n}{i}\PY{p}{]}\PY{p}{]}\PY{o}{/}\PY{l+m+mi}{100000}\PY{p}{)}

\PY{n}{plt}\PY{o}{.}\PY{n}{show}\PY{p}{(}\PY{p}{)}
\end{Verbatim}
\end{tcolorbox}

    \begin{center}
    \adjustimage{max size={0.9\linewidth}{0.9\paperheight}}{figs/bvp_kM_en/output_107_0.png}
    \end{center}
    { \hspace*{\fill} \\}
    
    \begin{tcolorbox}[breakable, size=fbox, boxrule=1pt, pad at break*=1mm,colback=cellbackground, colframe=cellborder]
\prompt{In}{incolor}{109}{\boxspacing}
\begin{Verbatim}[commandchars=\\\{\}]
\PY{n+nb}{print}\PY{p}{(}\PY{l+s+s2}{\PYZdq{}}\PY{l+s+s2}{Number of fast cycles:}\PY{l+s+s2}{\PYZdq{}}\PY{p}{,} \PY{n}{find\PYZus{}count\PYZus{}k\PYZus{}osc}\PY{p}{(}\PY{n}{dR}\PY{p}{[}\PY{o}{\PYZhy{}}\PY{l+m+mi}{1}\PY{p}{]}\PY{p}{)} \PY{o}{\PYZhy{}} \PY{n}{find\PYZus{}count\PYZus{}k\PYZus{}osc}\PY{p}{(}\PY{n}{dR}\PY{p}{[}\PY{l+m+mi}{0}\PY{p}{]}\PY{p}{)}\PY{p}{)}
\PY{n+nb}{print}\PY{p}{(}\PY{l+s+s2}{\PYZdq{}}\PY{l+s+s2}{Position of fast cycle starts:}\PY{l+s+s2}{\PYZdq{}}\PY{p}{,} \PY{n}{find\PYZus{}dR\PYZus{}k\PYZus{}osc}\PY{p}{(}\PY{l+m+mi}{1}\PY{p}{)}\PY{p}{,} \PY{l+s+s2}{\PYZdq{}}\PY{l+s+s2}{, }\PY{l+s+s2}{\PYZdq{}}\PY{p}{,} \PY{n}{find\PYZus{}dR\PYZus{}k\PYZus{}osc}\PY{p}{(}\PY{l+m+mi}{2}\PY{p}{)}\PY{p}{)}
\end{Verbatim}
\end{tcolorbox}

    \begin{Verbatim}[commandchars=\\\{\}]
Number of fast cycles: 2.4000449999984994
Position of fast cycle starts: 0.6645833333333333 ,  1.4979166666666663
    \end{Verbatim}

    \begin{tcolorbox}[breakable, size=fbox, boxrule=1pt, pad at break*=1mm,colback=cellbackground, colframe=cellborder]
\prompt{In}{incolor}{ }{\boxspacing}
\begin{Verbatim}[commandchars=\\\{\}]
\PY{c+c1}{\PYZsh{} Again, choose the indexes beforehand.}

\PY{n}{Bi} \PY{o}{=} \PY{p}{[}\PY{l+m+mi}{0}\PY{p}{]}

\PY{n}{ilim} \PY{o}{=} \PY{n+nb}{len}\PY{p}{(}\PY{n}{Bi}\PY{p}{)}

\PY{n}{fig}\PY{p}{,} \PY{n}{ax} \PY{o}{=} \PY{n}{plt}\PY{o}{.}\PY{n}{subplots}\PY{p}{(}\PY{p}{)}
\PY{n}{plt}\PY{o}{.}\PY{n}{title}\PY{p}{(}\PY{l+s+sa}{f}\PY{l+s+s2}{\PYZdq{}}\PY{l+s+s2}{Plot of the Change in k Due to Nonlinearity vs R}\PY{l+s+se}{\PYZbs{}n}\PY{l+s+s2}{for dk = 0.45 / 2 * 10}\PY{l+s+se}{\PYZbs{}u2076}\PY{l+s+s2}{ m}\PY{l+s+se}{\PYZbs{}u207b}\PY{l+s+se}{\PYZbs{}u00b9}\PY{l+s+s2}{ and B = }\PY{l+s+si}{\PYZob{}}\PY{n}{B0}\PY{o}{+}\PY{n}{Bi}\PY{p}{[}\PY{l+m+mi}{0}\PY{p}{]}\PY{l+s+si}{:}\PY{l+s+s2}{.4f}\PY{l+s+si}{\PYZcb{}}\PY{l+s+s2}{\PYZdq{}}\PY{p}{)}
\PY{n}{ax}\PY{o}{.}\PY{n}{set\PYZus{}ylabel}\PY{p}{(}\PY{l+s+s1}{\PYZsq{}}\PY{l+s+s1}{Change in k (dkn) due to nonlinearity (10}\PY{l+s+se}{\PYZbs{}u2075}\PY{l+s+s1}{ m}\PY{l+s+se}{\PYZbs{}u207b}\PY{l+s+se}{\PYZbs{}u00b9}\PY{l+s+s1}{)}\PY{l+s+s1}{\PYZsq{}}\PY{p}{)}
\PY{n}{ax}\PY{o}{.}\PY{n}{set\PYZus{}xlabel}\PY{p}{(}\PY{l+s+s2}{\PYZdq{}}\PY{l+s+s2}{Change dr in radius R of ring (R = 9 * 10}\PY{l+s+se}{\PYZbs{}u207b}\PY{l+s+se}{\PYZbs{}u2077}\PY{l+s+s2}{ m + dr * 10}\PY{l+s+se}{\PYZbs{}u207b}\PY{l+s+se}{\PYZbs{}u00b9}\PY{l+s+se}{\PYZbs{}u2070}\PY{l+s+s2}{ m)}\PY{l+s+s2}{\PYZdq{}}\PY{p}{)}

\PY{k}{for} \PY{n}{i} \PY{o+ow}{in} \PY{n+nb}{range}\PY{p}{(}\PY{n}{ilim}\PY{p}{)}\PY{p}{:}
    \PY{n}{ax}\PY{o}{.}\PY{n}{plot}\PY{p}{(}\PY{n}{R}\PY{p}{,} \PY{n}{dkn\PYZus{}arr}\PY{p}{[}\PY{p}{:}\PY{p}{,} \PY{n}{Bi}\PY{p}{[}\PY{n}{i}\PY{p}{]}\PY{p}{]}\PY{o}{/}\PY{l+m+mi}{100000}\PY{p}{,} \PY{n}{label} \PY{o}{=} \PY{l+s+s2}{\PYZdq{}}\PY{l+s+s2}{dkn (10}\PY{l+s+se}{\PYZbs{}u2075}\PY{l+s+s2}{ m}\PY{l+s+se}{\PYZbs{}u207b}\PY{l+s+se}{\PYZbs{}u00b9}\PY{l+s+s2}{)}\PY{l+s+s2}{\PYZdq{}}\PY{p}{)}

\PY{c+c1}{\PYZsh{} Our cycle boundaries}
\PY{n}{plt}\PY{o}{.}\PY{n}{axline}\PY{p}{(}\PY{p}{(}\PY{n}{find\PYZus{}dR\PYZus{}k\PYZus{}osc}\PY{p}{(}\PY{l+m+mi}{1}\PY{p}{)}\PY{p}{,}\PY{o}{\PYZhy{}}\PY{l+m+mf}{0.5}\PY{p}{)}\PY{p}{,} \PY{p}{(}\PY{n}{find\PYZus{}dR\PYZus{}k\PYZus{}osc}\PY{p}{(}\PY{l+m+mi}{1}\PY{p}{)}\PY{p}{,}\PY{o}{\PYZhy{}}\PY{l+m+mf}{0.6}\PY{p}{)}\PY{p}{,} \PY{n}{color}\PY{o}{=}\PY{l+s+s1}{\PYZsq{}}\PY{l+s+s1}{k}\PY{l+s+s1}{\PYZsq{}}\PY{p}{,} \PY{n}{label}\PY{o}{=}\PY{l+s+s2}{\PYZdq{}}\PY{l+s+s2}{remainder of kR / 1 = 0}\PY{l+s+s2}{\PYZdq{}}\PY{p}{)}
\PY{n}{plt}\PY{o}{.}\PY{n}{axline}\PY{p}{(}\PY{p}{(}\PY{n}{find\PYZus{}dR\PYZus{}k\PYZus{}osc}\PY{p}{(}\PY{l+m+mi}{2}\PY{p}{)}\PY{p}{,}\PY{o}{\PYZhy{}}\PY{l+m+mf}{0.5}\PY{p}{)}\PY{p}{,} \PY{p}{(}\PY{n}{find\PYZus{}dR\PYZus{}k\PYZus{}osc}\PY{p}{(}\PY{l+m+mi}{2}\PY{p}{)}\PY{p}{,}\PY{o}{\PYZhy{}}\PY{l+m+mf}{0.6}\PY{p}{)}\PY{p}{,} \PY{n}{color}\PY{o}{=}\PY{l+s+s1}{\PYZsq{}}\PY{l+s+s1}{k}\PY{l+s+s1}{\PYZsq{}}\PY{p}{)}

\PY{n}{plt}\PY{o}{.}\PY{n}{legend}\PY{p}{(}\PY{p}{)}

\PY{n}{plt}\PY{o}{.}\PY{n}{show}\PY{p}{(}\PY{p}{)}
\end{Verbatim}
\end{tcolorbox}

    \begin{center}
    \adjustimage{max size={0.9\linewidth}{0.9\paperheight}}{figs/bvp_kM_en/output_109_0.png}
    \end{center}
    { \hspace*{\fill} \\}
    
    We start here at a \(B\) node. The cycles start at the minima of our
curve.

    \begin{tcolorbox}[breakable, size=fbox, boxrule=1pt, pad at break*=1mm,colback=cellbackground, colframe=cellborder]
\prompt{In}{incolor}{ }{\boxspacing}
\begin{Verbatim}[commandchars=\\\{\}]
\PY{n}{Bi} \PY{o}{=} \PY{p}{[}\PY{l+m+mi}{0}\PY{p}{,} \PY{l+m+mi}{4}\PY{p}{,} \PY{l+m+mi}{10}\PY{p}{,} \PY{l+m+mi}{20}\PY{p}{,} \PY{l+m+mi}{30}\PY{p}{]}

\PY{n}{ilim} \PY{o}{=} \PY{n+nb}{len}\PY{p}{(}\PY{n}{Bi}\PY{p}{)}

\PY{n}{fig}\PY{p}{,} \PY{n}{ax} \PY{o}{=} \PY{n}{plt}\PY{o}{.}\PY{n}{subplots}\PY{p}{(}\PY{p}{)}
\PY{n}{plt}\PY{o}{.}\PY{n}{title}\PY{p}{(}\PY{l+s+sa}{f}\PY{l+s+s2}{\PYZdq{}}\PY{l+s+s2}{Plot of the Change in k Due to Nonlinearity vs R}\PY{l+s+se}{\PYZbs{}n}\PY{l+s+s2}{for dk = 0.45 / 2 * 10}\PY{l+s+se}{\PYZbs{}u2076}\PY{l+s+s2}{ m}\PY{l+s+se}{\PYZbs{}u207b}\PY{l+s+se}{\PYZbs{}u00b9}\PY{l+s+s2}{ for various B values}\PY{l+s+s2}{\PYZdq{}}\PY{p}{)}
\PY{n}{ax}\PY{o}{.}\PY{n}{set\PYZus{}ylabel}\PY{p}{(}\PY{l+s+s1}{\PYZsq{}}\PY{l+s+s1}{Change in k (dkn) due to nonlinearity (10}\PY{l+s+se}{\PYZbs{}u2075}\PY{l+s+s1}{ m}\PY{l+s+se}{\PYZbs{}u207b}\PY{l+s+se}{\PYZbs{}u00b9}\PY{l+s+s1}{)}\PY{l+s+s1}{\PYZsq{}}\PY{p}{)}
\PY{n}{ax}\PY{o}{.}\PY{n}{set\PYZus{}xlabel}\PY{p}{(}\PY{l+s+s2}{\PYZdq{}}\PY{l+s+s2}{Change dr in radius R of ring (R = 9 * 10}\PY{l+s+se}{\PYZbs{}u207b}\PY{l+s+se}{\PYZbs{}u2077}\PY{l+s+s2}{ m + dr * 10}\PY{l+s+se}{\PYZbs{}u207b}\PY{l+s+se}{\PYZbs{}u00b9}\PY{l+s+se}{\PYZbs{}u2070}\PY{l+s+s2}{ m)}\PY{l+s+s2}{\PYZdq{}}\PY{p}{)}

\PY{k}{for} \PY{n}{i} \PY{o+ow}{in} \PY{n+nb}{range}\PY{p}{(}\PY{n}{ilim}\PY{p}{)}\PY{p}{:}
    \PY{n}{ax}\PY{o}{.}\PY{n}{plot}\PY{p}{(}\PY{n}{R}\PY{p}{,} \PY{n}{dkn\PYZus{}arr}\PY{p}{[}\PY{p}{:}\PY{p}{,} \PY{n}{Bi}\PY{p}{[}\PY{n}{i}\PY{p}{]}\PY{p}{]}\PY{o}{/}\PY{l+m+mi}{100000}\PY{p}{,} \PY{n}{label} \PY{o}{=} \PY{l+s+sa}{f}\PY{l+s+s2}{\PYZdq{}}\PY{l+s+s2}{B = }\PY{l+s+si}{\PYZob{}}\PY{n}{B0}\PY{o}{+}\PY{n}{B}\PY{p}{[}\PY{n}{Bi}\PY{p}{[}\PY{n}{i}\PY{p}{]}\PY{p}{]}\PY{l+s+si}{:}\PY{l+s+s2}{.4f}\PY{l+s+si}{\PYZcb{}}\PY{l+s+s2}{ * B\PYZus{}max}\PY{l+s+s2}{\PYZdq{}}\PY{p}{)}
\PY{n}{plt}\PY{o}{.}\PY{n}{axline}\PY{p}{(}\PY{p}{(}\PY{n}{find\PYZus{}dR\PYZus{}k\PYZus{}osc}\PY{p}{(}\PY{l+m+mi}{1}\PY{p}{)}\PY{p}{,}\PY{o}{\PYZhy{}}\PY{l+m+mf}{0.5}\PY{p}{)}\PY{p}{,} \PY{p}{(}\PY{n}{find\PYZus{}dR\PYZus{}k\PYZus{}osc}\PY{p}{(}\PY{l+m+mi}{1}\PY{p}{)}\PY{p}{,}\PY{o}{\PYZhy{}}\PY{l+m+mf}{0.6}\PY{p}{)}\PY{p}{,} \PY{n}{color}\PY{o}{=}\PY{l+s+s1}{\PYZsq{}}\PY{l+s+s1}{k}\PY{l+s+s1}{\PYZsq{}}\PY{p}{,} \PY{n}{label}\PY{o}{=}\PY{l+s+s2}{\PYZdq{}}\PY{l+s+s2}{remainder of kR / 1 = 0}\PY{l+s+s2}{\PYZdq{}}\PY{p}{)}
\PY{n}{plt}\PY{o}{.}\PY{n}{axline}\PY{p}{(}\PY{p}{(}\PY{n}{find\PYZus{}dR\PYZus{}k\PYZus{}osc}\PY{p}{(}\PY{l+m+mi}{2}\PY{p}{)}\PY{p}{,}\PY{o}{\PYZhy{}}\PY{l+m+mf}{0.5}\PY{p}{)}\PY{p}{,} \PY{p}{(}\PY{n}{find\PYZus{}dR\PYZus{}k\PYZus{}osc}\PY{p}{(}\PY{l+m+mi}{2}\PY{p}{)}\PY{p}{,}\PY{o}{\PYZhy{}}\PY{l+m+mf}{0.6}\PY{p}{)}\PY{p}{,} \PY{n}{color}\PY{o}{=}\PY{l+s+s1}{\PYZsq{}}\PY{l+s+s1}{k}\PY{l+s+s1}{\PYZsq{}}\PY{p}{)}

\PY{n}{plt}\PY{o}{.}\PY{n}{legend}\PY{p}{(}\PY{n}{loc}\PY{o}{=}\PY{l+s+s1}{\PYZsq{}}\PY{l+s+s1}{center left}\PY{l+s+s1}{\PYZsq{}}\PY{p}{,} \PY{n}{bbox\PYZus{}to\PYZus{}anchor}\PY{o}{=}\PY{p}{(}\PY{l+m+mi}{1}\PY{p}{,} \PY{l+m+mf}{0.5}\PY{p}{)}\PY{p}{)}

\PY{n}{plt}\PY{o}{.}\PY{n}{show}\PY{p}{(}\PY{p}{)}
\end{Verbatim}
\end{tcolorbox}

    \begin{center}
    \adjustimage{max size={0.9\linewidth}{0.9\paperheight}}{figs/bvp_kM_en/output_111_0.png}
    \end{center}
    { \hspace*{\fill} \\}
    
    \begin{tcolorbox}[breakable, size=fbox, boxrule=1pt, pad at break*=1mm,colback=cellbackground, colframe=cellborder]
\prompt{In}{incolor}{113}{\boxspacing}
\begin{Verbatim}[commandchars=\\\{\}]
\PY{n}{B}\PY{p}{[}\PY{l+m+mi}{30}\PY{p}{]}
\end{Verbatim}
\end{tcolorbox}

            \begin{tcolorbox}[breakable, size=fbox, boxrule=.5pt, pad at break*=1mm, opacityfill=0]
\prompt{Out}{outcolor}{113}{\boxspacing}
\begin{Verbatim}[commandchars=\\\{\}]
0.030303030303030276
\end{Verbatim}
\end{tcolorbox}
        
    As \(B\) increases, a spike rises from this midpoint. This spike is the
same spike of our original function, but shifted over by a half cycle.
As the new pattern grows (which is our old pattern shifted by half a
cycle), the old curve will recede slightly. Around \(B = 0.53\), our two
curves have equal height.

    \begin{tcolorbox}[breakable, size=fbox, boxrule=1pt, pad at break*=1mm,colback=cellbackground, colframe=cellborder]
\prompt{In}{incolor}{ }{\boxspacing}
\begin{Verbatim}[commandchars=\\\{\}]
\PY{c+c1}{\PYZsh{} for easy changing which two plots to use.}

\PY{n}{Bi} \PY{o}{=} \PY{p}{[}\PY{l+m+mi}{0}\PY{p}{,} \PY{l+m+mi}{35}\PY{p}{,} \PY{l+m+mi}{40}\PY{p}{,} \PY{l+m+mi}{50}\PY{p}{,} \PY{l+m+mi}{55}\PY{p}{,} \PY{l+m+mi}{60}\PY{p}{]}

\PY{n}{ilim} \PY{o}{=} \PY{n+nb}{len}\PY{p}{(}\PY{n}{Bi}\PY{p}{)}

\PY{n}{fig}\PY{p}{,} \PY{n}{ax} \PY{o}{=} \PY{n}{plt}\PY{o}{.}\PY{n}{subplots}\PY{p}{(}\PY{p}{)}
\PY{n}{plt}\PY{o}{.}\PY{n}{title}\PY{p}{(}\PY{l+s+sa}{f}\PY{l+s+s2}{\PYZdq{}}\PY{l+s+s2}{Plot of the Change in k Due to Nonlinearity vs R}\PY{l+s+se}{\PYZbs{}n}\PY{l+s+s2}{for dk = 0.45 / 2 * 10}\PY{l+s+se}{\PYZbs{}u2076}\PY{l+s+s2}{ m}\PY{l+s+se}{\PYZbs{}u207b}\PY{l+s+se}{\PYZbs{}u00b9}\PY{l+s+s2}{ for various B values}\PY{l+s+s2}{\PYZdq{}}\PY{p}{)}
\PY{n}{ax}\PY{o}{.}\PY{n}{set\PYZus{}ylabel}\PY{p}{(}\PY{l+s+s1}{\PYZsq{}}\PY{l+s+s1}{Change in k (dkn) due to nonlinearity (10}\PY{l+s+se}{\PYZbs{}u2075}\PY{l+s+s1}{ m}\PY{l+s+se}{\PYZbs{}u207b}\PY{l+s+se}{\PYZbs{}u00b9}\PY{l+s+s1}{)}\PY{l+s+s1}{\PYZsq{}}\PY{p}{)}
\PY{n}{ax}\PY{o}{.}\PY{n}{set\PYZus{}xlabel}\PY{p}{(}\PY{l+s+s2}{\PYZdq{}}\PY{l+s+s2}{Change dr in radius R of ring (R = 9 * 10}\PY{l+s+se}{\PYZbs{}u207b}\PY{l+s+se}{\PYZbs{}u2077}\PY{l+s+s2}{ m + dr * 10}\PY{l+s+se}{\PYZbs{}u207b}\PY{l+s+se}{\PYZbs{}u00b9}\PY{l+s+se}{\PYZbs{}u2070}\PY{l+s+s2}{ m)}\PY{l+s+s2}{\PYZdq{}}\PY{p}{)}

\PY{k}{for} \PY{n}{i} \PY{o+ow}{in} \PY{n+nb}{range}\PY{p}{(}\PY{n}{ilim}\PY{p}{)}\PY{p}{:}
    \PY{n}{ax}\PY{o}{.}\PY{n}{plot}\PY{p}{(}\PY{n}{R}\PY{p}{,} \PY{n}{dkn\PYZus{}arr}\PY{p}{[}\PY{p}{:}\PY{p}{,} \PY{n}{Bi}\PY{p}{[}\PY{n}{i}\PY{p}{]}\PY{p}{]}\PY{o}{/}\PY{l+m+mi}{100000}\PY{p}{,} \PY{n}{label} \PY{o}{=} \PY{l+s+sa}{f}\PY{l+s+s2}{\PYZdq{}}\PY{l+s+s2}{B = }\PY{l+s+si}{\PYZob{}}\PY{n}{B0}\PY{o}{+}\PY{n}{B}\PY{p}{[}\PY{n}{Bi}\PY{p}{[}\PY{n}{i}\PY{p}{]}\PY{p}{]}\PY{l+s+si}{:}\PY{l+s+s2}{.4f}\PY{l+s+si}{\PYZcb{}}\PY{l+s+s2}{ * B\PYZus{}max}\PY{l+s+s2}{\PYZdq{}}\PY{p}{)}
\PY{n}{plt}\PY{o}{.}\PY{n}{axline}\PY{p}{(}\PY{p}{(}\PY{n}{find\PYZus{}dR\PYZus{}k\PYZus{}osc}\PY{p}{(}\PY{l+m+mi}{1}\PY{p}{)}\PY{p}{,}\PY{o}{\PYZhy{}}\PY{l+m+mf}{0.5}\PY{p}{)}\PY{p}{,} \PY{p}{(}\PY{n}{find\PYZus{}dR\PYZus{}k\PYZus{}osc}\PY{p}{(}\PY{l+m+mi}{1}\PY{p}{)}\PY{p}{,}\PY{o}{\PYZhy{}}\PY{l+m+mf}{0.6}\PY{p}{)}\PY{p}{,} \PY{n}{color}\PY{o}{=}\PY{l+s+s1}{\PYZsq{}}\PY{l+s+s1}{k}\PY{l+s+s1}{\PYZsq{}}\PY{p}{,} \PY{n}{label}\PY{o}{=}\PY{l+s+s2}{\PYZdq{}}\PY{l+s+s2}{remainder of kR / 1 = 0}\PY{l+s+s2}{\PYZdq{}}\PY{p}{)}
\PY{n}{plt}\PY{o}{.}\PY{n}{axline}\PY{p}{(}\PY{p}{(}\PY{n}{find\PYZus{}dR\PYZus{}k\PYZus{}osc}\PY{p}{(}\PY{l+m+mi}{2}\PY{p}{)}\PY{p}{,}\PY{o}{\PYZhy{}}\PY{l+m+mf}{0.5}\PY{p}{)}\PY{p}{,} \PY{p}{(}\PY{n}{find\PYZus{}dR\PYZus{}k\PYZus{}osc}\PY{p}{(}\PY{l+m+mi}{2}\PY{p}{)}\PY{p}{,}\PY{o}{\PYZhy{}}\PY{l+m+mf}{0.6}\PY{p}{)}\PY{p}{,} \PY{n}{color}\PY{o}{=}\PY{l+s+s1}{\PYZsq{}}\PY{l+s+s1}{k}\PY{l+s+s1}{\PYZsq{}}\PY{p}{)}

\PY{n}{plt}\PY{o}{.}\PY{n}{legend}\PY{p}{(}\PY{n}{loc}\PY{o}{=}\PY{l+s+s1}{\PYZsq{}}\PY{l+s+s1}{center left}\PY{l+s+s1}{\PYZsq{}}\PY{p}{,} \PY{n}{bbox\PYZus{}to\PYZus{}anchor}\PY{o}{=}\PY{p}{(}\PY{l+m+mi}{1}\PY{p}{,} \PY{l+m+mf}{0.5}\PY{p}{)}\PY{p}{)}

\PY{n}{plt}\PY{o}{.}\PY{n}{show}\PY{p}{(}\PY{p}{)}
\end{Verbatim}
\end{tcolorbox}

    \begin{center}
    \adjustimage{max size={0.9\linewidth}{0.9\paperheight}}{figs/bvp_kM_en/output_114_0.png}
    \end{center}
    { \hspace*{\fill} \\}
    
    The blue curve here is the original \(B = 0.5\) curve at our \(B\) node.
We see now that the spike of our old cycle is now receding, with the
growth of the new spike being almost complete, even at the earliest of
the depicted \(B\) values besides our reference.

    \begin{tcolorbox}[breakable, size=fbox, boxrule=1pt, pad at break*=1mm,colback=cellbackground, colframe=cellborder]
\prompt{In}{incolor}{ }{\boxspacing}
\begin{Verbatim}[commandchars=\\\{\}]
\PY{n}{Bi} \PY{o}{=} \PY{p}{[}\PY{l+m+mi}{0}\PY{p}{,} \PY{l+m+mi}{55}\PY{p}{,} \PY{l+m+mi}{60}\PY{p}{,} \PY{l+m+mi}{62}\PY{p}{]}

\PY{n}{ilim} \PY{o}{=} \PY{n+nb}{len}\PY{p}{(}\PY{n}{Bi}\PY{p}{)}

\PY{n}{fig}\PY{p}{,} \PY{n}{ax} \PY{o}{=} \PY{n}{plt}\PY{o}{.}\PY{n}{subplots}\PY{p}{(}\PY{p}{)}
\PY{n}{plt}\PY{o}{.}\PY{n}{title}\PY{p}{(}\PY{l+s+sa}{f}\PY{l+s+s2}{\PYZdq{}}\PY{l+s+s2}{Plot of the Change in k Due to Nonlinearity vs R}\PY{l+s+se}{\PYZbs{}n}\PY{l+s+s2}{for dk = 0.45 / 2 * 10}\PY{l+s+se}{\PYZbs{}u2076}\PY{l+s+s2}{ m}\PY{l+s+se}{\PYZbs{}u207b}\PY{l+s+se}{\PYZbs{}u00b9}\PY{l+s+s2}{ for various B values}\PY{l+s+s2}{\PYZdq{}}\PY{p}{)}
\PY{n}{ax}\PY{o}{.}\PY{n}{set\PYZus{}ylabel}\PY{p}{(}\PY{l+s+s1}{\PYZsq{}}\PY{l+s+s1}{Change in k (dkn) due to nonlinearity (10}\PY{l+s+se}{\PYZbs{}u2075}\PY{l+s+s1}{ m}\PY{l+s+se}{\PYZbs{}u207b}\PY{l+s+se}{\PYZbs{}u00b9}\PY{l+s+s1}{)}\PY{l+s+s1}{\PYZsq{}}\PY{p}{)}
\PY{n}{ax}\PY{o}{.}\PY{n}{set\PYZus{}xlabel}\PY{p}{(}\PY{l+s+s2}{\PYZdq{}}\PY{l+s+s2}{Change dr in radius R of ring (R = 9 * 10}\PY{l+s+se}{\PYZbs{}u207b}\PY{l+s+se}{\PYZbs{}u2077}\PY{l+s+s2}{ m + dr * 10}\PY{l+s+se}{\PYZbs{}u207b}\PY{l+s+se}{\PYZbs{}u00b9}\PY{l+s+se}{\PYZbs{}u2070}\PY{l+s+s2}{ m)}\PY{l+s+s2}{\PYZdq{}}\PY{p}{)}

\PY{k}{for} \PY{n}{i} \PY{o+ow}{in} \PY{n+nb}{range}\PY{p}{(}\PY{n}{ilim}\PY{p}{)}\PY{p}{:}
    \PY{n}{ax}\PY{o}{.}\PY{n}{plot}\PY{p}{(}\PY{n}{R}\PY{p}{,} \PY{n}{dkn\PYZus{}arr}\PY{p}{[}\PY{p}{:}\PY{p}{,} \PY{n}{Bi}\PY{p}{[}\PY{n}{i}\PY{p}{]}\PY{p}{]}\PY{o}{/}\PY{l+m+mi}{100000}\PY{p}{,} \PY{n}{label} \PY{o}{=} \PY{l+s+sa}{f}\PY{l+s+s2}{\PYZdq{}}\PY{l+s+s2}{B = }\PY{l+s+si}{\PYZob{}}\PY{n}{B0}\PY{o}{+}\PY{n}{B}\PY{p}{[}\PY{n}{Bi}\PY{p}{[}\PY{n}{i}\PY{p}{]}\PY{p}{]}\PY{l+s+si}{:}\PY{l+s+s2}{.4f}\PY{l+s+si}{\PYZcb{}}\PY{l+s+s2}{ * B\PYZus{}max}\PY{l+s+s2}{\PYZdq{}}\PY{p}{)}
    
\PY{n}{plt}\PY{o}{.}\PY{n}{axline}\PY{p}{(}\PY{p}{(}\PY{n}{find\PYZus{}dR\PYZus{}k\PYZus{}osc}\PY{p}{(}\PY{l+m+mi}{1}\PY{p}{)}\PY{p}{,}\PY{o}{\PYZhy{}}\PY{l+m+mf}{0.5}\PY{p}{)}\PY{p}{,} \PY{p}{(}\PY{n}{find\PYZus{}dR\PYZus{}k\PYZus{}osc}\PY{p}{(}\PY{l+m+mi}{1}\PY{p}{)}\PY{p}{,}\PY{o}{\PYZhy{}}\PY{l+m+mf}{0.6}\PY{p}{)}\PY{p}{,} \PY{n}{color}\PY{o}{=}\PY{l+s+s1}{\PYZsq{}}\PY{l+s+s1}{k}\PY{l+s+s1}{\PYZsq{}}\PY{p}{,} \PY{n}{label}\PY{o}{=}\PY{l+s+s2}{\PYZdq{}}\PY{l+s+s2}{remainder of kR / 1 = 0}\PY{l+s+s2}{\PYZdq{}}\PY{p}{)}
\PY{n}{plt}\PY{o}{.}\PY{n}{axline}\PY{p}{(}\PY{p}{(}\PY{n}{find\PYZus{}dR\PYZus{}k\PYZus{}osc}\PY{p}{(}\PY{l+m+mi}{2}\PY{p}{)}\PY{p}{,}\PY{o}{\PYZhy{}}\PY{l+m+mf}{0.5}\PY{p}{)}\PY{p}{,} \PY{p}{(}\PY{n}{find\PYZus{}dR\PYZus{}k\PYZus{}osc}\PY{p}{(}\PY{l+m+mi}{2}\PY{p}{)}\PY{p}{,}\PY{o}{\PYZhy{}}\PY{l+m+mf}{0.6}\PY{p}{)}\PY{p}{,} \PY{n}{color}\PY{o}{=}\PY{l+s+s1}{\PYZsq{}}\PY{l+s+s1}{k}\PY{l+s+s1}{\PYZsq{}}\PY{p}{)}

\PY{n}{plt}\PY{o}{.}\PY{n}{legend}\PY{p}{(}\PY{n}{loc}\PY{o}{=}\PY{l+s+s1}{\PYZsq{}}\PY{l+s+s1}{center left}\PY{l+s+s1}{\PYZsq{}}\PY{p}{,} \PY{n}{bbox\PYZus{}to\PYZus{}anchor}\PY{o}{=}\PY{p}{(}\PY{l+m+mi}{1}\PY{p}{,} \PY{l+m+mf}{0.5}\PY{p}{)}\PY{p}{)}

\PY{n}{plt}\PY{o}{.}\PY{n}{show}\PY{p}{(}\PY{p}{)}
\end{Verbatim}
\end{tcolorbox}

    \begin{center}
    \adjustimage{max size={0.9\linewidth}{0.9\paperheight}}{figs/bvp_kM_en/output_116_0.png}
    \end{center}
    { \hspace*{\fill} \\}
    
    \begin{tcolorbox}[breakable, size=fbox, boxrule=1pt, pad at break*=1mm,colback=cellbackground, colframe=cellborder]
\prompt{In}{incolor}{116}{\boxspacing}
\begin{Verbatim}[commandchars=\\\{\}]
\PY{n}{B}\PY{p}{[}\PY{l+m+mi}{62}\PY{p}{]}
\end{Verbatim}
\end{tcolorbox}

            \begin{tcolorbox}[breakable, size=fbox, boxrule=.5pt, pad at break*=1mm, opacityfill=0]
\prompt{Out}{outcolor}{116}{\boxspacing}
\begin{Verbatim}[commandchars=\\\{\}]
0.06262626262626259
\end{Verbatim}
\end{tcolorbox}
        
    Now we have the completed shift at around \(B = 0.563\), our next node
of \(B\), which is the midpoint in the \(B\) cycle.

    This is a graph from my full \(B\) range dataset which which shows the
two secant squared shapes at the two types of nodes in the period, with
the sum of them both at the peaks placed on top. The two peaks of this
sum are slightly less than each of the individual peaks.

    \begin{tcolorbox}[breakable, size=fbox, boxrule=1pt, pad at break*=1mm,colback=cellbackground, colframe=cellborder]
\prompt{In}{incolor}{117}{\boxspacing}
\begin{Verbatim}[commandchars=\\\{\}]
\PY{c+c1}{\PYZsh{} Load the data from the fourth dataset for full \PYZdl{}B\PYZdl{} range.}
\PY{n}{gind} \PY{o}{=} \PY{l+m+mi}{3}

\PY{k}{del} \PY{n}{dkn\PYZus{}arr}\PY{p}{,} \PY{n}{dMn\PYZus{}arr}

\PY{n}{dkn\PYZus{}arr} \PY{o}{=} \PY{n}{np}\PY{o}{.}\PY{n}{zeros}\PY{p}{(}\PY{p}{(}\PY{n}{gs\PYZus{}r2}\PY{p}{,} \PY{n}{gs\PYZus{}b2}\PY{p}{)}\PY{p}{)}
\PY{n}{dMn\PYZus{}arr} \PY{o}{=} \PY{n}{np}\PY{o}{.}\PY{n}{zeros}\PY{p}{(}\PY{p}{(}\PY{n}{gs\PYZus{}r2}\PY{p}{,} \PY{n}{gs\PYZus{}b2}\PY{p}{)}\PY{p}{)}

\PY{k}{for} \PY{n}{r}\PY{p}{,} \PY{n}{rr} \PY{o+ow}{in} \PY{n+nb}{enumerate}\PY{p}{(}\PY{n}{R\PYZus{}g}\PY{p}{[}\PY{n}{gind}\PY{p}{]}\PY{p}{)}\PY{p}{:}
    \PY{k}{for} \PY{n}{b}\PY{p}{,} \PY{n}{bb} \PY{o+ow}{in} \PY{n+nb}{enumerate}\PY{p}{(}\PY{n}{B\PYZus{}g}\PY{p}{[}\PY{n}{gind}\PY{p}{]}\PY{p}{)}\PY{p}{:}
        \PY{n}{tbl\PYZus{}sel} \PY{o}{=} \PY{n}{tbl\PYZus{}red}\PY{p}{[}\PY{n}{gind}\PY{p}{]}\PY{p}{[}\PY{p}{(}\PY{n}{tbl\PYZus{}red}\PY{p}{[}\PY{n}{gind}\PY{p}{]}\PY{p}{[}\PY{l+s+s2}{\PYZdq{}}\PY{l+s+s2}{R}\PY{l+s+s2}{\PYZdq{}}\PY{p}{]}\PY{o}{==}\PY{n}{rr}\PY{p}{)} \PY{o}{\PYZam{}}\PY{p}{(}\PY{n}{tbl\PYZus{}red}\PY{p}{[}\PY{n}{gind}\PY{p}{]}\PY{p}{[}\PY{l+s+s2}{\PYZdq{}}\PY{l+s+s2}{B}\PY{l+s+s2}{\PYZdq{}}\PY{p}{]}\PY{o}{==}\PY{n}{bb}\PY{p}{)}\PY{p}{]}
        \PY{n}{dkn\PYZus{}arr}\PY{p}{[}\PY{n}{r}\PY{p}{,}\PY{n}{b}\PY{p}{]} \PY{o}{=} \PY{n}{tbl\PYZus{}sel}\PY{p}{[}\PY{l+s+s2}{\PYZdq{}}\PY{l+s+s2}{dkn}\PY{l+s+s2}{\PYZdq{}}\PY{p}{]}\PY{o}{.}\PY{n}{to\PYZus{}numpy}\PY{p}{(}\PY{p}{)}\PY{p}{[}\PY{l+m+mi}{0}\PY{p}{]}
        \PY{n}{dMn\PYZus{}arr}\PY{p}{[}\PY{n}{r}\PY{p}{,}\PY{n}{b}\PY{p}{]} \PY{o}{=} \PY{n}{tbl\PYZus{}sel}\PY{p}{[}\PY{l+s+s2}{\PYZdq{}}\PY{l+s+s2}{dMn}\PY{l+s+s2}{\PYZdq{}}\PY{p}{]}\PY{o}{.}\PY{n}{to\PYZus{}numpy}\PY{p}{(}\PY{p}{)}\PY{p}{[}\PY{l+m+mi}{0}\PY{p}{]}

\PY{c+c1}{\PYZsh{} B and R for the fourth dataset}
\PY{n}{dk} \PY{o}{=} \PY{n}{dk\PYZus{}g}\PY{p}{[}\PY{n}{gind}\PY{p}{]}\PY{p}{[}\PY{l+m+mi}{0}\PY{p}{]}

\PY{n}{R} \PY{o}{=} \PY{p}{(}\PY{n}{R\PYZus{}g}\PY{p}{[}\PY{n}{gind}\PY{p}{]} \PY{o}{\PYZhy{}} \PY{n}{R0}\PY{p}{)}\PY{o}{/}\PY{n}{scale\PYZus{}R}
\PY{n}{B} \PY{o}{=} \PY{n}{B\PYZus{}g}\PY{p}{[}\PY{n}{gind}\PY{p}{]}
\end{Verbatim}
\end{tcolorbox}

    \begin{tcolorbox}[breakable, size=fbox, boxrule=1pt, pad at break*=1mm,colback=cellbackground, colframe=cellborder]
\prompt{In}{incolor}{ }{\boxspacing}
\begin{Verbatim}[commandchars=\\\{\}]
\PY{c+c1}{\PYZsh{} And show our plot}
\PY{n}{bmn} \PY{o}{=} \PY{l+m+mi}{85}

\PY{n}{Bi} \PY{o}{=} \PY{p}{[}\PY{l+m+mi}{4}\PY{o}{+}\PY{n}{bmn}\PY{p}{,} \PY{l+m+mi}{25}\PY{o}{+}\PY{n}{bmn}\PY{p}{,} \PY{l+m+mi}{14}\PY{o}{+}\PY{n}{bmn}\PY{p}{,} \PY{l+m+mi}{15}\PY{o}{+}\PY{n}{bmn}\PY{p}{]}

\PY{n}{ilim} \PY{o}{=} \PY{n+nb}{len}\PY{p}{(}\PY{n}{Bi}\PY{p}{)}

\PY{n}{fig}\PY{p}{,} \PY{n}{ax} \PY{o}{=} \PY{n}{plt}\PY{o}{.}\PY{n}{subplots}\PY{p}{(}\PY{p}{)}
\PY{n}{plt}\PY{o}{.}\PY{n}{title}\PY{p}{(}\PY{l+s+sa}{f}\PY{l+s+s2}{\PYZdq{}}\PY{l+s+s2}{Plot of the Change in k Due to Nonlinearity vs R}\PY{l+s+se}{\PYZbs{}n}\PY{l+s+s2}{for dk = 0.45 / 2 * 10}\PY{l+s+se}{\PYZbs{}u2076}\PY{l+s+s2}{ m}\PY{l+s+se}{\PYZbs{}u207b}\PY{l+s+se}{\PYZbs{}u00b9}\PY{l+s+s2}{ for various B values}\PY{l+s+s2}{\PYZdq{}}\PY{p}{)}
\PY{n}{ax}\PY{o}{.}\PY{n}{set\PYZus{}ylabel}\PY{p}{(}\PY{l+s+s1}{\PYZsq{}}\PY{l+s+s1}{Change in k (dkn) due to nonlinearity (10}\PY{l+s+se}{\PYZbs{}u2075}\PY{l+s+s1}{ m}\PY{l+s+se}{\PYZbs{}u207b}\PY{l+s+se}{\PYZbs{}u00b9}\PY{l+s+s1}{)}\PY{l+s+s1}{\PYZsq{}}\PY{p}{)}
\PY{n}{ax}\PY{o}{.}\PY{n}{set\PYZus{}xlabel}\PY{p}{(}\PY{l+s+s2}{\PYZdq{}}\PY{l+s+s2}{Change dr in radius R of ring (R = 9 * 10}\PY{l+s+se}{\PYZbs{}u207b}\PY{l+s+se}{\PYZbs{}u2077}\PY{l+s+s2}{ m + dr * 10}\PY{l+s+se}{\PYZbs{}u207b}\PY{l+s+se}{\PYZbs{}u00b9}\PY{l+s+se}{\PYZbs{}u2070}\PY{l+s+s2}{ m)}\PY{l+s+s2}{\PYZdq{}}\PY{p}{)}

\PY{k}{for} \PY{n}{i} \PY{o+ow}{in} \PY{n+nb}{range}\PY{p}{(}\PY{n}{ilim}\PY{p}{)}\PY{p}{:}
    \PY{n}{ax}\PY{o}{.}\PY{n}{plot}\PY{p}{(}\PY{n}{R}\PY{p}{,} \PY{n}{dkn\PYZus{}arr}\PY{p}{[}\PY{p}{:}\PY{p}{,} \PY{n}{Bi}\PY{p}{[}\PY{n}{i}\PY{p}{]}\PY{p}{]}\PY{o}{/}\PY{l+m+mi}{100000}\PY{p}{,} \PY{n}{label} \PY{o}{=} \PY{l+s+sa}{f}\PY{l+s+s2}{\PYZdq{}}\PY{l+s+s2}{B = }\PY{l+s+si}{\PYZob{}}\PY{n}{B0}\PY{o}{+}\PY{n}{B}\PY{p}{[}\PY{n}{Bi}\PY{p}{[}\PY{n}{i}\PY{p}{]}\PY{p}{]}\PY{l+s+si}{:}\PY{l+s+s2}{.4f}\PY{l+s+si}{\PYZcb{}}\PY{l+s+s2}{ * B\PYZus{}max}\PY{l+s+s2}{\PYZdq{}}\PY{p}{)}

\PY{n}{plt}\PY{o}{.}\PY{n}{legend}\PY{p}{(}\PY{n}{loc}\PY{o}{=}\PY{l+s+s1}{\PYZsq{}}\PY{l+s+s1}{center left}\PY{l+s+s1}{\PYZsq{}}\PY{p}{,} \PY{n}{bbox\PYZus{}to\PYZus{}anchor}\PY{o}{=}\PY{p}{(}\PY{l+m+mi}{1}\PY{p}{,} \PY{l+m+mf}{0.5}\PY{p}{)}\PY{p}{)}

\PY{n}{plt}\PY{o}{.}\PY{n}{show}\PY{p}{(}\PY{p}{)}
\end{Verbatim}
\end{tcolorbox}

    \begin{center}
    \adjustimage{max size={0.9\linewidth}{0.9\paperheight}}{figs/bvp_kM_en/output_121_0.png}
    \end{center}
    { \hspace*{\fill} \\}
    
    \hypertarget{dkn-vs-b}{%
\section{\texorpdfstring{\(dkn\) vs \(B\)}{dkn vs B}}\label{dkn-vs-b}}

    \begin{tcolorbox}[breakable, size=fbox, boxrule=1pt, pad at break*=1mm,colback=cellbackground, colframe=cellborder]
\prompt{In}{incolor}{119}{\boxspacing}
\begin{Verbatim}[commandchars=\\\{\}]
\PY{c+c1}{\PYZsh{} Reload the data from our first dataset, as I had temporarily used the 4th dataset previously}

\PY{n}{gind} \PY{o}{=} \PY{l+m+mi}{0}

\PY{k}{del} \PY{n}{dkn\PYZus{}arr}\PY{p}{,} \PY{n}{dMn\PYZus{}arr}

\PY{n}{dkn\PYZus{}arr} \PY{o}{=} \PY{n}{np}\PY{o}{.}\PY{n}{zeros}\PY{p}{(}\PY{p}{(}\PY{n}{grid\PYZus{}size}\PY{p}{,} \PY{n}{grid\PYZus{}size}\PY{p}{)}\PY{p}{)}
\PY{n}{dMn\PYZus{}arr} \PY{o}{=} \PY{n}{np}\PY{o}{.}\PY{n}{zeros}\PY{p}{(}\PY{p}{(}\PY{n}{grid\PYZus{}size}\PY{p}{,} \PY{n}{grid\PYZus{}size}\PY{p}{)}\PY{p}{)}

\PY{k}{for} \PY{n}{r}\PY{p}{,} \PY{n}{rr} \PY{o+ow}{in} \PY{n+nb}{enumerate}\PY{p}{(}\PY{n}{R\PYZus{}g}\PY{p}{[}\PY{n}{gind}\PY{p}{]}\PY{p}{)}\PY{p}{:}
    \PY{k}{for} \PY{n}{b}\PY{p}{,} \PY{n}{bb} \PY{o+ow}{in} \PY{n+nb}{enumerate}\PY{p}{(}\PY{n}{B\PYZus{}g}\PY{p}{[}\PY{n}{gind}\PY{p}{]}\PY{p}{)}\PY{p}{:}
        \PY{n}{tbl\PYZus{}sel} \PY{o}{=} \PY{n}{tbl\PYZus{}red}\PY{p}{[}\PY{n}{gind}\PY{p}{]}\PY{p}{[}\PY{p}{(}\PY{n}{tbl\PYZus{}red}\PY{p}{[}\PY{n}{gind}\PY{p}{]}\PY{p}{[}\PY{l+s+s2}{\PYZdq{}}\PY{l+s+s2}{R}\PY{l+s+s2}{\PYZdq{}}\PY{p}{]}\PY{o}{==}\PY{n}{rr}\PY{p}{)} \PY{o}{\PYZam{}}\PY{p}{(}\PY{n}{tbl\PYZus{}red}\PY{p}{[}\PY{n}{gind}\PY{p}{]}\PY{p}{[}\PY{l+s+s2}{\PYZdq{}}\PY{l+s+s2}{B}\PY{l+s+s2}{\PYZdq{}}\PY{p}{]}\PY{o}{==}\PY{n}{bb}\PY{p}{)}\PY{p}{]}
        \PY{n}{dkn\PYZus{}arr}\PY{p}{[}\PY{n}{r}\PY{p}{,}\PY{n}{b}\PY{p}{]} \PY{o}{=} \PY{n}{tbl\PYZus{}sel}\PY{p}{[}\PY{l+s+s2}{\PYZdq{}}\PY{l+s+s2}{dkn}\PY{l+s+s2}{\PYZdq{}}\PY{p}{]}\PY{o}{.}\PY{n}{to\PYZus{}numpy}\PY{p}{(}\PY{p}{)}\PY{p}{[}\PY{l+m+mi}{0}\PY{p}{]}
        \PY{n}{dMn\PYZus{}arr}\PY{p}{[}\PY{n}{r}\PY{p}{,}\PY{n}{b}\PY{p}{]} \PY{o}{=} \PY{n}{tbl\PYZus{}sel}\PY{p}{[}\PY{l+s+s2}{\PYZdq{}}\PY{l+s+s2}{dMn}\PY{l+s+s2}{\PYZdq{}}\PY{p}{]}\PY{o}{.}\PY{n}{to\PYZus{}numpy}\PY{p}{(}\PY{p}{)}\PY{p}{[}\PY{l+m+mi}{0}\PY{p}{]}

\PY{c+c1}{\PYZsh{}B and B will be rescaled to show only the difference from our }
\PY{n}{dk} \PY{o}{=} \PY{n}{dk\PYZus{}g}\PY{p}{[}\PY{n}{gind}\PY{p}{]}\PY{p}{[}\PY{l+m+mi}{0}\PY{p}{]}

\PY{n}{R} \PY{o}{=} \PY{p}{(}\PY{n}{R\PYZus{}g}\PY{p}{[}\PY{n}{gind}\PY{p}{]} \PY{o}{\PYZhy{}} \PY{n}{R0}\PY{p}{)}\PY{o}{/}\PY{n}{scale\PYZus{}R}
\PY{n}{B} \PY{o}{=} \PY{n}{B\PYZus{}g}\PY{p}{[}\PY{n}{gind}\PY{p}{]} \PY{o}{\PYZhy{}} \PY{n}{B0}
\end{Verbatim}
\end{tcolorbox}

    The next plots should show this same \(dkn\) analysis, but for \(dkn\)
vs \(B\) plots.

    \begin{tcolorbox}[breakable, size=fbox, boxrule=1pt, pad at break*=1mm,colback=cellbackground, colframe=cellborder]
\prompt{In}{incolor}{120}{\boxspacing}
\begin{Verbatim}[commandchars=\\\{\}]
\PY{c+c1}{\PYZsh{} For reference}
\PY{n+nb}{print}\PY{p}{(}\PY{l+s+s2}{\PYZdq{}}\PY{l+s+s2}{Position of fast cycle starts:}\PY{l+s+s2}{\PYZdq{}}\PY{p}{,} \PY{n}{find\PYZus{}dR\PYZus{}k\PYZus{}osc}\PY{p}{(}\PY{l+m+mi}{1}\PY{p}{)}\PY{p}{,} \PY{l+s+s2}{\PYZdq{}}\PY{l+s+s2}{, }\PY{l+s+s2}{\PYZdq{}}\PY{p}{,} \PY{n}{find\PYZus{}dR\PYZus{}k\PYZus{}osc}\PY{p}{(}\PY{l+m+mi}{2}\PY{p}{)}\PY{p}{)}
\end{Verbatim}
\end{tcolorbox}

    \begin{Verbatim}[commandchars=\\\{\}]
Position of fast cycle starts: 0.6645833333333333 ,  1.4979166666666663
    \end{Verbatim}

    \begin{tcolorbox}[breakable, size=fbox, boxrule=1pt, pad at break*=1mm,colback=cellbackground, colframe=cellborder]
\prompt{In}{incolor}{121}{\boxspacing}
\begin{Verbatim}[commandchars=\\\{\}]
\PY{c+c1}{\PYZsh{} To translate our dR value to an index, when we know that R[\PYZhy{}1] = 2}
\PY{c+c1}{\PYZsh{} The third number is the midway point.}
\PY{n}{find\PYZus{}dR\PYZus{}k\PYZus{}osc}\PY{p}{(}\PY{l+m+mi}{1}\PY{p}{)} \PY{o}{/} \PY{l+m+mi}{2} \PY{o}{*} \PY{l+m+mi}{100}\PY{p}{,} \PY{n}{find\PYZus{}dR\PYZus{}k\PYZus{}osc}\PY{p}{(}\PY{l+m+mi}{2}\PY{p}{)} \PY{o}{/} \PY{l+m+mi}{2} \PY{o}{*} \PY{l+m+mi}{100}\PY{p}{,} \PY{n}{find\PYZus{}dR\PYZus{}k\PYZus{}osc}\PY{p}{(}\PY{l+m+mf}{1.5}\PY{p}{)} \PY{o}{/} \PY{l+m+mi}{2} \PY{o}{*} \PY{l+m+mi}{100}\PY{p}{,}
\end{Verbatim}
\end{tcolorbox}

            \begin{tcolorbox}[breakable, size=fbox, boxrule=.5pt, pad at break*=1mm, opacityfill=0]
\prompt{Out}{outcolor}{121}{\boxspacing}
\begin{Verbatim}[commandchars=\\\{\}]
(33.229166666666664, 74.89583333333331, 54.06249999999999)
\end{Verbatim}
\end{tcolorbox}
        
    We want R indicies from 33 to 75.

    \begin{tcolorbox}[breakable, size=fbox, boxrule=1pt, pad at break*=1mm,colback=cellbackground, colframe=cellborder]
\prompt{In}{incolor}{122}{\boxspacing}
\begin{Verbatim}[commandchars=\\\{\}]
\PY{c+c1}{\PYZsh{} And for our lines marking the B oscillation boundaries.}
\PY{n}{find\PYZus{}dB\PYZus{}B\PYZus{}osc}\PY{p}{(}\PY{l+m+mi}{4}\PY{p}{)}\PY{p}{,} \PY{n}{find\PYZus{}dB\PYZus{}B\PYZus{}osc}\PY{p}{(}\PY{l+m+mi}{5}\PY{p}{)}
\end{Verbatim}
\end{tcolorbox}

            \begin{tcolorbox}[breakable, size=fbox, boxrule=.5pt, pad at break*=1mm, opacityfill=0]
\prompt{Out}{outcolor}{122}{\boxspacing}
\begin{Verbatim}[commandchars=\\\{\}]
(0.00035152415969275097, 0.12543940519961594)
\end{Verbatim}
\end{tcolorbox}
        
    \begin{tcolorbox}[breakable, size=fbox, boxrule=1pt, pad at break*=1mm,colback=cellbackground, colframe=cellborder]
\prompt{In}{incolor}{ }{\boxspacing}
\begin{Verbatim}[commandchars=\\\{\}]
\PY{c+c1}{\PYZsh{} Here are the most obvious plots to show}

\PY{n}{Ri} \PY{o}{=} \PY{p}{[}\PY{l+m+mi}{33}\PY{p}{,} \PY{l+m+mi}{54}\PY{p}{,} \PY{l+m+mi}{75}\PY{p}{]}

\PY{n}{ilim} \PY{o}{=} \PY{n+nb}{len}\PY{p}{(}\PY{n}{Ri}\PY{p}{)}

\PY{n}{fig}\PY{p}{,} \PY{n}{ax} \PY{o}{=} \PY{n}{plt}\PY{o}{.}\PY{n}{subplots}\PY{p}{(}\PY{p}{)}

\PY{n}{plt}\PY{o}{.}\PY{n}{title}\PY{p}{(}\PY{l+s+sa}{f}\PY{l+s+s2}{\PYZdq{}}\PY{l+s+s2}{Change in k Due to Nonlinearity vs B for dk = 0.45 / 2 * 10}\PY{l+s+se}{\PYZbs{}u2076}\PY{l+s+s2}{ m}\PY{l+s+se}{\PYZbs{}u207b}\PY{l+s+se}{\PYZbs{}u00b9}\PY{l+s+se}{\PYZbs{}n}\PY{l+s+s2}{for Various dR Values and R = 9 * 10}\PY{l+s+se}{\PYZbs{}u207b}\PY{l+s+se}{\PYZbs{}u2077}\PY{l+s+s2}{ m + dR}\PY{l+s+s2}{\PYZdq{}}\PY{p}{)}
\PY{n}{ax}\PY{o}{.}\PY{n}{set\PYZus{}ylabel}\PY{p}{(}\PY{l+s+s1}{\PYZsq{}}\PY{l+s+s1}{Change in k (dkn) due to nonlinearity (10}\PY{l+s+se}{\PYZbs{}u2075}\PY{l+s+s1}{ m}\PY{l+s+se}{\PYZbs{}u207b}\PY{l+s+se}{\PYZbs{}u00b9}\PY{l+s+s1}{)}\PY{l+s+s1}{\PYZsq{}}\PY{p}{)}
\PY{n}{ax}\PY{o}{.}\PY{n}{set\PYZus{}xlabel}\PY{p}{(}\PY{l+s+s2}{\PYZdq{}}\PY{l+s+s2}{Variation (dB) of magnetic field strength B (B = 0.5 + dB * B\PYZus{}max)}\PY{l+s+s2}{\PYZdq{}}\PY{p}{)}

\PY{k}{for} \PY{n}{i} \PY{o+ow}{in} \PY{n+nb}{range}\PY{p}{(}\PY{n}{ilim}\PY{p}{)}\PY{p}{:}
    \PY{n}{ax}\PY{o}{.}\PY{n}{plot}\PY{p}{(}\PY{n}{B}\PY{p}{,} \PY{n}{dkn\PYZus{}arr}\PY{p}{[}\PY{n}{Ri}\PY{p}{[}\PY{n}{i}\PY{p}{]}\PY{p}{,}\PY{p}{:}\PY{p}{]}\PY{o}{/}\PY{l+m+mi}{100000}\PY{p}{,} \PY{n}{label} \PY{o}{=} \PY{l+s+sa}{f}\PY{l+s+s2}{\PYZdq{}}\PY{l+s+s2}{dR = }\PY{l+s+si}{\PYZob{}}\PY{n}{R}\PY{p}{[}\PY{n}{Ri}\PY{p}{[}\PY{n}{i}\PY{p}{]}\PY{p}{]}\PY{l+s+si}{:}\PY{l+s+s2}{.4f}\PY{l+s+si}{\PYZcb{}}\PY{l+s+s2}{ * 10}\PY{l+s+se}{\PYZbs{}u207b}\PY{l+s+se}{\PYZbs{}u00b9}\PY{l+s+se}{\PYZbs{}u2070}\PY{l+s+s2}{ m}\PY{l+s+s2}{\PYZdq{}}\PY{p}{)}
\PY{n}{plt}\PY{o}{.}\PY{n}{axline}\PY{p}{(}\PY{p}{(}\PY{n}{find\PYZus{}dB\PYZus{}B\PYZus{}osc}\PY{p}{(}\PY{l+m+mi}{4}\PY{p}{)}\PY{p}{,}\PY{o}{\PYZhy{}}\PY{l+m+mf}{0.4}\PY{p}{)}\PY{p}{,} \PY{p}{(}\PY{n}{find\PYZus{}dB\PYZus{}B\PYZus{}osc}\PY{p}{(}\PY{l+m+mi}{4}\PY{p}{)}\PY{p}{,}\PY{o}{\PYZhy{}}\PY{l+m+mf}{0.41}\PY{p}{)}\PY{p}{,} \PY{n}{color}\PY{o}{=}\PY{l+s+s1}{\PYZsq{}}\PY{l+s+s1}{k}\PY{l+s+s1}{\PYZsq{}}\PY{p}{,} \PY{n}{label} \PY{o}{=}\PY{l+s+s2}{\PYZdq{}}\PY{l+s+s2}{period boundary}\PY{l+s+s2}{\PYZdq{}}\PY{p}{)}
\PY{n}{plt}\PY{o}{.}\PY{n}{axline}\PY{p}{(}\PY{p}{(}\PY{n}{find\PYZus{}dB\PYZus{}B\PYZus{}osc}\PY{p}{(}\PY{l+m+mf}{4.5}\PY{p}{)}\PY{p}{,}\PY{o}{\PYZhy{}}\PY{l+m+mf}{0.4}\PY{p}{)}\PY{p}{,} \PY{p}{(}\PY{n}{find\PYZus{}dB\PYZus{}B\PYZus{}osc}\PY{p}{(}\PY{l+m+mf}{4.5}\PY{p}{)}\PY{p}{,}\PY{o}{\PYZhy{}}\PY{l+m+mf}{0.41}\PY{p}{)}\PY{p}{,} \PY{n}{color}\PY{o}{=}\PY{l+s+s1}{\PYZsq{}}\PY{l+s+s1}{k}\PY{l+s+s1}{\PYZsq{}}\PY{p}{)}

\PY{n}{plt}\PY{o}{.}\PY{n}{legend}\PY{p}{(}\PY{n}{bbox\PYZus{}to\PYZus{}anchor}\PY{o}{=}\PY{p}{(}\PY{l+m+mf}{0.55}\PY{p}{,} \PY{l+m+mi}{1}\PY{p}{)}\PY{p}{)}

\PY{n}{plt}\PY{o}{.}\PY{n}{show}\PY{p}{(}\PY{p}{)}
\end{Verbatim}
\end{tcolorbox}

    \begin{center}
    \adjustimage{max size={0.9\linewidth}{0.9\paperheight}}{figs/bvp_kM_en/output_129_0.png}
    \end{center}
    { \hspace*{\fill} \\}
    
    \begin{tcolorbox}[breakable, size=fbox, boxrule=1pt, pad at break*=1mm,colback=cellbackground, colframe=cellborder]
\prompt{In}{incolor}{ }{\boxspacing}
\begin{Verbatim}[commandchars=\\\{\}]
\PY{c+c1}{\PYZsh{} To show the largest and smallest variations of dkn}

\PY{n}{Ri} \PY{o}{=} \PY{p}{[}\PY{l+m+mi}{45}\PY{p}{,} \PY{l+m+mi}{60}\PY{p}{]}

\PY{n}{ilim} \PY{o}{=} \PY{n+nb}{len}\PY{p}{(}\PY{n}{Ri}\PY{p}{)}

\PY{n}{fig}\PY{p}{,} \PY{n}{ax} \PY{o}{=} \PY{n}{plt}\PY{o}{.}\PY{n}{subplots}\PY{p}{(}\PY{p}{)}
\PY{n}{plt}\PY{o}{.}\PY{n}{title}\PY{p}{(}\PY{l+s+sa}{f}\PY{l+s+s2}{\PYZdq{}}\PY{l+s+s2}{Change in k Due to Nonlinearity vs B for dk = 0.45 / 2 * 10}\PY{l+s+se}{\PYZbs{}u2076}\PY{l+s+s2}{ m}\PY{l+s+se}{\PYZbs{}u207b}\PY{l+s+se}{\PYZbs{}u00b9}\PY{l+s+se}{\PYZbs{}n}\PY{l+s+s2}{for Various dR Values and R = 9 * 10}\PY{l+s+se}{\PYZbs{}u207b}\PY{l+s+se}{\PYZbs{}u2077}\PY{l+s+s2}{ m + dR}\PY{l+s+s2}{\PYZdq{}}\PY{p}{)}
\PY{n}{ax}\PY{o}{.}\PY{n}{set\PYZus{}ylabel}\PY{p}{(}\PY{l+s+s1}{\PYZsq{}}\PY{l+s+s1}{Change in k (dkn) due to nonlinearity (10}\PY{l+s+se}{\PYZbs{}u2075}\PY{l+s+s1}{ m}\PY{l+s+se}{\PYZbs{}u207b}\PY{l+s+se}{\PYZbs{}u00b9}\PY{l+s+s1}{)}\PY{l+s+s1}{\PYZsq{}}\PY{p}{)}
\PY{n}{ax}\PY{o}{.}\PY{n}{set\PYZus{}xlabel}\PY{p}{(}\PY{l+s+s2}{\PYZdq{}}\PY{l+s+s2}{Variation (dB) of magnetic field strength B (B = 0.5 + dB * B\PYZus{}max)}\PY{l+s+s2}{\PYZdq{}}\PY{p}{)}

\PY{k}{for} \PY{n}{i} \PY{o+ow}{in} \PY{n+nb}{range}\PY{p}{(}\PY{n}{ilim}\PY{p}{)}\PY{p}{:}
    \PY{n}{ax}\PY{o}{.}\PY{n}{plot}\PY{p}{(}\PY{n}{B}\PY{p}{,} \PY{n}{dkn\PYZus{}arr}\PY{p}{[}\PY{n}{Ri}\PY{p}{[}\PY{n}{i}\PY{p}{]}\PY{p}{,}\PY{p}{:}\PY{p}{]}\PY{o}{/}\PY{l+m+mi}{100000}\PY{p}{,} \PY{n}{label} \PY{o}{=} \PY{l+s+sa}{f}\PY{l+s+s2}{\PYZdq{}}\PY{l+s+s2}{dR = }\PY{l+s+si}{\PYZob{}}\PY{n}{R}\PY{p}{[}\PY{n}{Ri}\PY{p}{[}\PY{n}{i}\PY{p}{]}\PY{p}{]}\PY{l+s+si}{:}\PY{l+s+s2}{.4f}\PY{l+s+si}{\PYZcb{}}\PY{l+s+s2}{ * 10}\PY{l+s+se}{\PYZbs{}u207b}\PY{l+s+se}{\PYZbs{}u00b9}\PY{l+s+se}{\PYZbs{}u2070}\PY{l+s+s2}{ m}\PY{l+s+s2}{\PYZdq{}}\PY{p}{)}

\PY{n}{plt}\PY{o}{.}\PY{n}{axline}\PY{p}{(}\PY{p}{(}\PY{n}{find\PYZus{}dB\PYZus{}B\PYZus{}osc}\PY{p}{(}\PY{l+m+mi}{4}\PY{p}{)}\PY{p}{,}\PY{o}{\PYZhy{}}\PY{l+m+mf}{0.5}\PY{p}{)}\PY{p}{,} \PY{p}{(}\PY{n}{find\PYZus{}dB\PYZus{}B\PYZus{}osc}\PY{p}{(}\PY{l+m+mi}{4}\PY{p}{)}\PY{p}{,}\PY{o}{\PYZhy{}}\PY{l+m+mf}{0.6}\PY{p}{)}\PY{p}{,} \PY{n}{color}\PY{o}{=}\PY{l+s+s1}{\PYZsq{}}\PY{l+s+s1}{k}\PY{l+s+s1}{\PYZsq{}}\PY{p}{,} \PY{n}{label} \PY{o}{=} \PY{l+s+s2}{\PYZdq{}}\PY{l+s+s2}{period boundary}\PY{l+s+s2}{\PYZdq{}}\PY{p}{)}
\PY{n}{plt}\PY{o}{.}\PY{n}{axline}\PY{p}{(}\PY{p}{(}\PY{n}{find\PYZus{}dB\PYZus{}B\PYZus{}osc}\PY{p}{(}\PY{l+m+mf}{4.5}\PY{p}{)}\PY{p}{,}\PY{o}{\PYZhy{}}\PY{l+m+mf}{0.5}\PY{p}{)}\PY{p}{,} \PY{p}{(}\PY{n}{find\PYZus{}dB\PYZus{}B\PYZus{}osc}\PY{p}{(}\PY{l+m+mf}{4.5}\PY{p}{)}\PY{p}{,}\PY{o}{\PYZhy{}}\PY{l+m+mf}{0.6}\PY{p}{)}\PY{p}{,} \PY{n}{color}\PY{o}{=}\PY{l+s+s1}{\PYZsq{}}\PY{l+s+s1}{k}\PY{l+s+s1}{\PYZsq{}}\PY{p}{)}

\PY{n}{plt}\PY{o}{.}\PY{n}{legend}\PY{p}{(}\PY{n}{loc}\PY{o}{=}\PY{l+s+s1}{\PYZsq{}}\PY{l+s+s1}{center left}\PY{l+s+s1}{\PYZsq{}}\PY{p}{,} \PY{n}{bbox\PYZus{}to\PYZus{}anchor}\PY{o}{=}\PY{p}{(}\PY{o}{\PYZhy{}}\PY{l+m+mf}{0.01}\PY{p}{,} \PY{l+m+mf}{0.6}\PY{p}{)}\PY{p}{)}

\PY{n}{plt}\PY{o}{.}\PY{n}{show}\PY{p}{(}\PY{p}{)}
\end{Verbatim}
\end{tcolorbox}

    \begin{center}
    \adjustimage{max size={0.9\linewidth}{0.9\paperheight}}{figs/bvp_kM_en/output_130_0.png}
    \end{center}
    { \hspace*{\fill} \\}
    
    \begin{tcolorbox}[breakable, size=fbox, boxrule=1pt, pad at break*=1mm,colback=cellbackground, colframe=cellborder]
\prompt{In}{incolor}{ }{\boxspacing}
\begin{Verbatim}[commandchars=\\\{\}]
\PY{n}{Ri} \PY{o}{=} \PY{p}{[}\PY{l+m+mi}{33}\PY{p}{,} \PY{l+m+mi}{36}\PY{p}{,} \PY{l+m+mi}{39}\PY{p}{,} \PY{l+m+mi}{42}\PY{p}{]}

\PY{n}{ilim} \PY{o}{=} \PY{n+nb}{len}\PY{p}{(}\PY{n}{Ri}\PY{p}{)}

\PY{n}{fig}\PY{p}{,} \PY{n}{ax} \PY{o}{=} \PY{n}{plt}\PY{o}{.}\PY{n}{subplots}\PY{p}{(}\PY{p}{)}
\PY{n}{plt}\PY{o}{.}\PY{n}{title}\PY{p}{(}\PY{l+s+sa}{f}\PY{l+s+s2}{\PYZdq{}}\PY{l+s+s2}{Change in k Due to Nonlinearity vs B for dk = 0.45 / 2 * 10}\PY{l+s+se}{\PYZbs{}u2076}\PY{l+s+s2}{ m}\PY{l+s+se}{\PYZbs{}u207b}\PY{l+s+se}{\PYZbs{}u00b9}\PY{l+s+se}{\PYZbs{}n}\PY{l+s+s2}{for Various dR Values and R = 9 * 10}\PY{l+s+se}{\PYZbs{}u207b}\PY{l+s+se}{\PYZbs{}u2077}\PY{l+s+s2}{ m + dR}\PY{l+s+s2}{\PYZdq{}}\PY{p}{)}
\PY{n}{ax}\PY{o}{.}\PY{n}{set\PYZus{}ylabel}\PY{p}{(}\PY{l+s+s1}{\PYZsq{}}\PY{l+s+s1}{Change in k (dkn) due to nonlinearity (10}\PY{l+s+se}{\PYZbs{}u2075}\PY{l+s+s1}{ m}\PY{l+s+se}{\PYZbs{}u207b}\PY{l+s+se}{\PYZbs{}u00b9}\PY{l+s+s1}{)}\PY{l+s+s1}{\PYZsq{}}\PY{p}{)}
\PY{n}{ax}\PY{o}{.}\PY{n}{set\PYZus{}xlabel}\PY{p}{(}\PY{l+s+s2}{\PYZdq{}}\PY{l+s+s2}{Variation (dB) of magnetic field strength B (B = 0.5 + dB * B\PYZus{}max)}\PY{l+s+s2}{\PYZdq{}}\PY{p}{)}

\PY{k}{for} \PY{n}{i} \PY{o+ow}{in} \PY{n+nb}{range}\PY{p}{(}\PY{n}{ilim}\PY{p}{)}\PY{p}{:}
    \PY{n}{ax}\PY{o}{.}\PY{n}{plot}\PY{p}{(}\PY{n}{B}\PY{p}{,} \PY{n}{dkn\PYZus{}arr}\PY{p}{[}\PY{n}{Ri}\PY{p}{[}\PY{n}{i}\PY{p}{]}\PY{p}{,}\PY{p}{:}\PY{p}{]}\PY{o}{/}\PY{l+m+mi}{100000}\PY{p}{,} \PY{n}{label} \PY{o}{=} \PY{l+s+sa}{f}\PY{l+s+s2}{\PYZdq{}}\PY{l+s+s2}{dR = }\PY{l+s+si}{\PYZob{}}\PY{n}{R}\PY{p}{[}\PY{n}{Ri}\PY{p}{[}\PY{n}{i}\PY{p}{]}\PY{p}{]}\PY{l+s+si}{:}\PY{l+s+s2}{.4f}\PY{l+s+si}{\PYZcb{}}\PY{l+s+s2}{ * 10}\PY{l+s+se}{\PYZbs{}u207b}\PY{l+s+se}{\PYZbs{}u00b9}\PY{l+s+se}{\PYZbs{}u2070}\PY{l+s+s2}{ m}\PY{l+s+s2}{\PYZdq{}}\PY{p}{)}

\PY{n}{plt}\PY{o}{.}\PY{n}{axline}\PY{p}{(}\PY{p}{(}\PY{n}{find\PYZus{}dB\PYZus{}B\PYZus{}osc}\PY{p}{(}\PY{l+m+mi}{4}\PY{p}{)}\PY{p}{,}\PY{o}{\PYZhy{}}\PY{l+m+mf}{0.5}\PY{p}{)}\PY{p}{,} \PY{p}{(}\PY{n}{find\PYZus{}dB\PYZus{}B\PYZus{}osc}\PY{p}{(}\PY{l+m+mi}{4}\PY{p}{)}\PY{p}{,}\PY{o}{\PYZhy{}}\PY{l+m+mf}{0.6}\PY{p}{)}\PY{p}{,} \PY{n}{color}\PY{o}{=}\PY{l+s+s1}{\PYZsq{}}\PY{l+s+s1}{k}\PY{l+s+s1}{\PYZsq{}}\PY{p}{,} \PY{n}{label} \PY{o}{=} \PY{l+s+s1}{\PYZsq{}}\PY{l+s+s1}{period boundary}\PY{l+s+s1}{\PYZsq{}}\PY{p}{)}
\PY{n}{plt}\PY{o}{.}\PY{n}{axline}\PY{p}{(}\PY{p}{(}\PY{n}{find\PYZus{}dB\PYZus{}B\PYZus{}osc}\PY{p}{(}\PY{l+m+mf}{4.5}\PY{p}{)}\PY{p}{,}\PY{o}{\PYZhy{}}\PY{l+m+mf}{0.5}\PY{p}{)}\PY{p}{,} \PY{p}{(}\PY{n}{find\PYZus{}dB\PYZus{}B\PYZus{}osc}\PY{p}{(}\PY{l+m+mf}{4.5}\PY{p}{)}\PY{p}{,}\PY{o}{\PYZhy{}}\PY{l+m+mf}{0.6}\PY{p}{)}\PY{p}{,} \PY{n}{color}\PY{o}{=}\PY{l+s+s1}{\PYZsq{}}\PY{l+s+s1}{k}\PY{l+s+s1}{\PYZsq{}}\PY{p}{)}

\PY{n}{plt}\PY{o}{.}\PY{n}{legend}\PY{p}{(}\PY{n}{loc}\PY{o}{=}\PY{l+s+s1}{\PYZsq{}}\PY{l+s+s1}{center left}\PY{l+s+s1}{\PYZsq{}}\PY{p}{,} \PY{n}{bbox\PYZus{}to\PYZus{}anchor}\PY{o}{=}\PY{p}{(}\PY{l+m+mi}{1}\PY{p}{,} \PY{l+m+mf}{0.5}\PY{p}{)}\PY{p}{)}

\PY{n}{plt}\PY{o}{.}\PY{n}{show}\PY{p}{(}\PY{p}{)}
\end{Verbatim}
\end{tcolorbox}

    \begin{center}
    \adjustimage{max size={0.9\linewidth}{0.9\paperheight}}{figs/bvp_kM_en/output_131_0.png}
    \end{center}
    { \hspace*{\fill} \\}
    
    Around dB{[}45{]}, I reach the peak of my spike and drip down to a node
in \(R\), where the \(B\) oscillations disappear.

    \begin{tcolorbox}[breakable, size=fbox, boxrule=1pt, pad at break*=1mm,colback=cellbackground, colframe=cellborder]
\prompt{In}{incolor}{ }{\boxspacing}
\begin{Verbatim}[commandchars=\\\{\}]
\PY{n}{Ri} \PY{o}{=} \PY{p}{[}\PY{l+m+mi}{42}\PY{p}{,} \PY{l+m+mi}{45}\PY{p}{,} \PY{l+m+mi}{48}\PY{p}{]}

\PY{n}{ilim} \PY{o}{=} \PY{n+nb}{len}\PY{p}{(}\PY{n}{Ri}\PY{p}{)}

\PY{n}{fig}\PY{p}{,} \PY{n}{ax} \PY{o}{=} \PY{n}{plt}\PY{o}{.}\PY{n}{subplots}\PY{p}{(}\PY{p}{)}
\PY{n}{plt}\PY{o}{.}\PY{n}{title}\PY{p}{(}\PY{l+s+sa}{f}\PY{l+s+s2}{\PYZdq{}}\PY{l+s+s2}{Change in k Due to Nonlinearity vs B for dk = 0.45 / 2 * 10}\PY{l+s+se}{\PYZbs{}u2076}\PY{l+s+s2}{ m}\PY{l+s+se}{\PYZbs{}u207b}\PY{l+s+se}{\PYZbs{}u00b9}\PY{l+s+se}{\PYZbs{}n}\PY{l+s+s2}{for Various dR Values and R = 9 * 10}\PY{l+s+se}{\PYZbs{}u207b}\PY{l+s+se}{\PYZbs{}u2077}\PY{l+s+s2}{ m + dR}\PY{l+s+s2}{\PYZdq{}}\PY{p}{)}
\PY{n}{ax}\PY{o}{.}\PY{n}{set\PYZus{}ylabel}\PY{p}{(}\PY{l+s+s1}{\PYZsq{}}\PY{l+s+s1}{Change in k (dkn) due to nonlinearity (10}\PY{l+s+se}{\PYZbs{}u2075}\PY{l+s+s1}{ m}\PY{l+s+se}{\PYZbs{}u207b}\PY{l+s+se}{\PYZbs{}u00b9}\PY{l+s+s1}{)}\PY{l+s+s1}{\PYZsq{}}\PY{p}{)}
\PY{n}{ax}\PY{o}{.}\PY{n}{set\PYZus{}xlabel}\PY{p}{(}\PY{l+s+s2}{\PYZdq{}}\PY{l+s+s2}{Variation (dB) of magnetic field strength B (B = 0.5 + dB * B\PYZus{}max)}\PY{l+s+s2}{\PYZdq{}}\PY{p}{)}

\PY{k}{for} \PY{n}{i} \PY{o+ow}{in} \PY{n+nb}{range}\PY{p}{(}\PY{n}{ilim}\PY{p}{)}\PY{p}{:}
    \PY{n}{ax}\PY{o}{.}\PY{n}{plot}\PY{p}{(}\PY{n}{B}\PY{p}{,} \PY{n}{dkn\PYZus{}arr}\PY{p}{[}\PY{n}{Ri}\PY{p}{[}\PY{n}{i}\PY{p}{]}\PY{p}{,}\PY{p}{:}\PY{p}{]}\PY{o}{/}\PY{l+m+mi}{100000}\PY{p}{,} \PY{n}{label} \PY{o}{=} \PY{l+s+sa}{f}\PY{l+s+s2}{\PYZdq{}}\PY{l+s+s2}{dR = }\PY{l+s+si}{\PYZob{}}\PY{n}{R}\PY{p}{[}\PY{n}{Ri}\PY{p}{[}\PY{n}{i}\PY{p}{]}\PY{p}{]}\PY{l+s+si}{:}\PY{l+s+s2}{.4f}\PY{l+s+si}{\PYZcb{}}\PY{l+s+s2}{ * 10}\PY{l+s+se}{\PYZbs{}u207b}\PY{l+s+se}{\PYZbs{}u00b9}\PY{l+s+se}{\PYZbs{}u2070}\PY{l+s+s2}{ m}\PY{l+s+s2}{\PYZdq{}}\PY{p}{)}

\PY{n}{plt}\PY{o}{.}\PY{n}{axline}\PY{p}{(}\PY{p}{(}\PY{n}{find\PYZus{}dB\PYZus{}B\PYZus{}osc}\PY{p}{(}\PY{l+m+mi}{4}\PY{p}{)}\PY{p}{,}\PY{o}{\PYZhy{}}\PY{l+m+mf}{0.5}\PY{p}{)}\PY{p}{,} \PY{p}{(}\PY{n}{find\PYZus{}dB\PYZus{}B\PYZus{}osc}\PY{p}{(}\PY{l+m+mi}{4}\PY{p}{)}\PY{p}{,}\PY{o}{\PYZhy{}}\PY{l+m+mf}{0.6}\PY{p}{)}\PY{p}{,} \PY{n}{color}\PY{o}{=}\PY{l+s+s1}{\PYZsq{}}\PY{l+s+s1}{k}\PY{l+s+s1}{\PYZsq{}}\PY{p}{,} \PY{n}{label} \PY{o}{=} \PY{l+s+s1}{\PYZsq{}}\PY{l+s+s1}{period boundary}\PY{l+s+s1}{\PYZsq{}}\PY{p}{)}
\PY{n}{plt}\PY{o}{.}\PY{n}{axline}\PY{p}{(}\PY{p}{(}\PY{n}{find\PYZus{}dB\PYZus{}B\PYZus{}osc}\PY{p}{(}\PY{l+m+mf}{4.5}\PY{p}{)}\PY{p}{,}\PY{o}{\PYZhy{}}\PY{l+m+mf}{0.5}\PY{p}{)}\PY{p}{,} \PY{p}{(}\PY{n}{find\PYZus{}dB\PYZus{}B\PYZus{}osc}\PY{p}{(}\PY{l+m+mf}{4.5}\PY{p}{)}\PY{p}{,}\PY{o}{\PYZhy{}}\PY{l+m+mf}{0.6}\PY{p}{)}\PY{p}{,} \PY{n}{color}\PY{o}{=}\PY{l+s+s1}{\PYZsq{}}\PY{l+s+s1}{k}\PY{l+s+s1}{\PYZsq{}}\PY{p}{)}

\PY{n}{plt}\PY{o}{.}\PY{n}{legend}\PY{p}{(}\PY{n}{loc}\PY{o}{=}\PY{l+s+s1}{\PYZsq{}}\PY{l+s+s1}{center left}\PY{l+s+s1}{\PYZsq{}}\PY{p}{,} \PY{n}{bbox\PYZus{}to\PYZus{}anchor}\PY{o}{=}\PY{p}{(}\PY{o}{\PYZhy{}}\PY{l+m+mf}{0.01}\PY{p}{,} \PY{l+m+mf}{0.5}\PY{p}{)}\PY{p}{)}

\PY{n}{plt}\PY{o}{.}\PY{n}{show}\PY{p}{(}\PY{p}{)}
\end{Verbatim}
\end{tcolorbox}

    \begin{center}
    \adjustimage{max size={0.9\linewidth}{0.9\paperheight}}{figs/bvp_kM_en/output_133_0.png}
    \end{center}
    { \hspace*{\fill} \\}
    
    \begin{tcolorbox}[breakable, size=fbox, boxrule=1pt, pad at break*=1mm,colback=cellbackground, colframe=cellborder]
\prompt{In}{incolor}{ }{\boxspacing}
\begin{Verbatim}[commandchars=\\\{\}]
\PY{n}{Ri} \PY{o}{=} \PY{p}{[}\PY{l+m+mi}{43}\PY{p}{,}\PY{l+m+mi}{44}\PY{p}{,}\PY{l+m+mi}{45}\PY{p}{]}

\PY{n}{ilim} \PY{o}{=} \PY{n+nb}{len}\PY{p}{(}\PY{n}{Ri}\PY{p}{)}

\PY{n}{fig}\PY{p}{,} \PY{n}{ax} \PY{o}{=} \PY{n}{plt}\PY{o}{.}\PY{n}{subplots}\PY{p}{(}\PY{p}{)}
\PY{n}{plt}\PY{o}{.}\PY{n}{title}\PY{p}{(}\PY{l+s+sa}{f}\PY{l+s+s2}{\PYZdq{}}\PY{l+s+s2}{Change in k Due to Nonlinearity vs B for dk = 0.45 / 2 * 10}\PY{l+s+se}{\PYZbs{}u2076}\PY{l+s+s2}{ m}\PY{l+s+se}{\PYZbs{}u207b}\PY{l+s+se}{\PYZbs{}u00b9}\PY{l+s+se}{\PYZbs{}n}\PY{l+s+s2}{for Various dR Values and R = 9 * 10}\PY{l+s+se}{\PYZbs{}u207b}\PY{l+s+se}{\PYZbs{}u2077}\PY{l+s+s2}{ m + dR}\PY{l+s+s2}{\PYZdq{}}\PY{p}{)}
\PY{n}{ax}\PY{o}{.}\PY{n}{set\PYZus{}ylabel}\PY{p}{(}\PY{l+s+s1}{\PYZsq{}}\PY{l+s+s1}{Change in k (dkn) due to nonlinearity (10}\PY{l+s+se}{\PYZbs{}u2075}\PY{l+s+s1}{ m}\PY{l+s+se}{\PYZbs{}u207b}\PY{l+s+se}{\PYZbs{}u00b9}\PY{l+s+s1}{)}\PY{l+s+s1}{\PYZsq{}}\PY{p}{)}
\PY{n}{ax}\PY{o}{.}\PY{n}{set\PYZus{}xlabel}\PY{p}{(}\PY{l+s+s2}{\PYZdq{}}\PY{l+s+s2}{Variation (dB) of magnetic field strength B (B = 0.5 + dB * B\PYZus{}max)}\PY{l+s+s2}{\PYZdq{}}\PY{p}{)}

\PY{k}{for} \PY{n}{i} \PY{o+ow}{in} \PY{n+nb}{range}\PY{p}{(}\PY{n}{ilim}\PY{p}{)}\PY{p}{:}
    \PY{n}{ax}\PY{o}{.}\PY{n}{plot}\PY{p}{(}\PY{n}{B}\PY{p}{,} \PY{n}{dkn\PYZus{}arr}\PY{p}{[}\PY{n}{Ri}\PY{p}{[}\PY{n}{i}\PY{p}{]}\PY{p}{,}\PY{p}{:}\PY{p}{]}\PY{o}{/}\PY{l+m+mi}{100000}\PY{p}{,} \PY{n}{label} \PY{o}{=} \PY{l+s+sa}{f}\PY{l+s+s2}{\PYZdq{}}\PY{l+s+s2}{dR = }\PY{l+s+si}{\PYZob{}}\PY{n}{R}\PY{p}{[}\PY{n}{Ri}\PY{p}{[}\PY{n}{i}\PY{p}{]}\PY{p}{]}\PY{l+s+si}{:}\PY{l+s+s2}{.4f}\PY{l+s+si}{\PYZcb{}}\PY{l+s+s2}{ * 10}\PY{l+s+se}{\PYZbs{}u207b}\PY{l+s+se}{\PYZbs{}u00b9}\PY{l+s+se}{\PYZbs{}u2070}\PY{l+s+s2}{ m}\PY{l+s+s2}{\PYZdq{}}\PY{p}{)}
\PY{n}{plt}\PY{o}{.}\PY{n}{axline}\PY{p}{(}\PY{p}{(}\PY{n}{find\PYZus{}dB\PYZus{}B\PYZus{}osc}\PY{p}{(}\PY{l+m+mi}{4}\PY{p}{)}\PY{p}{,}\PY{o}{\PYZhy{}}\PY{l+m+mf}{0.5}\PY{p}{)}\PY{p}{,} \PY{p}{(}\PY{n}{find\PYZus{}dB\PYZus{}B\PYZus{}osc}\PY{p}{(}\PY{l+m+mi}{4}\PY{p}{)}\PY{p}{,}\PY{o}{\PYZhy{}}\PY{l+m+mf}{0.6}\PY{p}{)}\PY{p}{,} \PY{n}{color}\PY{o}{=}\PY{l+s+s1}{\PYZsq{}}\PY{l+s+s1}{k}\PY{l+s+s1}{\PYZsq{}}\PY{p}{,} \PY{n}{label} \PY{o}{=} \PY{l+s+s1}{\PYZsq{}}\PY{l+s+s1}{period boundary}\PY{l+s+s1}{\PYZsq{}}\PY{p}{)}
\PY{n}{plt}\PY{o}{.}\PY{n}{axline}\PY{p}{(}\PY{p}{(}\PY{n}{find\PYZus{}dB\PYZus{}B\PYZus{}osc}\PY{p}{(}\PY{l+m+mf}{4.5}\PY{p}{)}\PY{p}{,}\PY{o}{\PYZhy{}}\PY{l+m+mf}{0.5}\PY{p}{)}\PY{p}{,} \PY{p}{(}\PY{n}{find\PYZus{}dB\PYZus{}B\PYZus{}osc}\PY{p}{(}\PY{l+m+mf}{4.5}\PY{p}{)}\PY{p}{,}\PY{o}{\PYZhy{}}\PY{l+m+mf}{0.6}\PY{p}{)}\PY{p}{,} \PY{n}{color}\PY{o}{=}\PY{l+s+s1}{\PYZsq{}}\PY{l+s+s1}{k}\PY{l+s+s1}{\PYZsq{}}\PY{p}{)}

\PY{n}{plt}\PY{o}{.}\PY{n}{legend}\PY{p}{(}\PY{n}{loc}\PY{o}{=}\PY{l+s+s1}{\PYZsq{}}\PY{l+s+s1}{center left}\PY{l+s+s1}{\PYZsq{}}\PY{p}{,} \PY{n}{bbox\PYZus{}to\PYZus{}anchor}\PY{o}{=}\PY{p}{(}\PY{o}{\PYZhy{}}\PY{l+m+mf}{0.01}\PY{p}{,} \PY{l+m+mf}{0.5}\PY{p}{)}\PY{p}{)}

\PY{n}{plt}\PY{o}{.}\PY{n}{show}\PY{p}{(}\PY{p}{)}
\end{Verbatim}
\end{tcolorbox}

    \begin{center}
    \adjustimage{max size={0.9\linewidth}{0.9\paperheight}}{figs/bvp_kM_en/output_134_0.png}
    \end{center}
    { \hspace*{\fill} \\}
    
    The point where my curve is constant in \(B\) is after the drop in
values from the spike. Note that B{[}44{]} is already after the drop in
absolute value, but B{[}45{]} is my constant \(dkn\) value.

    \begin{tcolorbox}[breakable, size=fbox, boxrule=1pt, pad at break*=1mm,colback=cellbackground, colframe=cellborder]
\prompt{In}{incolor}{ }{\boxspacing}
\begin{Verbatim}[commandchars=\\\{\}]
\PY{n}{Ri} \PY{o}{=} \PY{p}{[}\PY{l+m+mi}{45}\PY{p}{,} \PY{l+m+mi}{48}\PY{p}{,} \PY{l+m+mi}{51}\PY{p}{,} \PY{l+m+mi}{54}\PY{p}{]}

\PY{n}{ilim} \PY{o}{=} \PY{n+nb}{len}\PY{p}{(}\PY{n}{Ri}\PY{p}{)}

\PY{n}{fig}\PY{p}{,} \PY{n}{ax} \PY{o}{=} \PY{n}{plt}\PY{o}{.}\PY{n}{subplots}\PY{p}{(}\PY{p}{)}
\PY{n}{plt}\PY{o}{.}\PY{n}{title}\PY{p}{(}\PY{l+s+sa}{f}\PY{l+s+s2}{\PYZdq{}}\PY{l+s+s2}{Change in k Due to Nonlinearity vs B for dk = 0.45 / 2 * 10}\PY{l+s+se}{\PYZbs{}u2076}\PY{l+s+s2}{ m}\PY{l+s+se}{\PYZbs{}u207b}\PY{l+s+se}{\PYZbs{}u00b9}\PY{l+s+se}{\PYZbs{}n}\PY{l+s+s2}{for Various dR Values and R = 9 * 10}\PY{l+s+se}{\PYZbs{}u207b}\PY{l+s+se}{\PYZbs{}u2077}\PY{l+s+s2}{ m + dR}\PY{l+s+s2}{\PYZdq{}}\PY{p}{)}
\PY{n}{ax}\PY{o}{.}\PY{n}{set\PYZus{}ylabel}\PY{p}{(}\PY{l+s+s1}{\PYZsq{}}\PY{l+s+s1}{Change in k (dkn) due to nonlinearity (10}\PY{l+s+se}{\PYZbs{}u2075}\PY{l+s+s1}{ m}\PY{l+s+se}{\PYZbs{}u207b}\PY{l+s+se}{\PYZbs{}u00b9}\PY{l+s+s1}{)}\PY{l+s+s1}{\PYZsq{}}\PY{p}{)}
\PY{n}{ax}\PY{o}{.}\PY{n}{set\PYZus{}xlabel}\PY{p}{(}\PY{l+s+s2}{\PYZdq{}}\PY{l+s+s2}{Variation (dB) of magnetic field strength B (B = 0.5 + dB * B\PYZus{}max)}\PY{l+s+s2}{\PYZdq{}}\PY{p}{)}

\PY{k}{for} \PY{n}{i} \PY{o+ow}{in} \PY{n+nb}{range}\PY{p}{(}\PY{n}{ilim}\PY{p}{)}\PY{p}{:}
    \PY{n}{ax}\PY{o}{.}\PY{n}{plot}\PY{p}{(}\PY{n}{B}\PY{p}{,} \PY{n}{dkn\PYZus{}arr}\PY{p}{[}\PY{n}{Ri}\PY{p}{[}\PY{n}{i}\PY{p}{]}\PY{p}{,}\PY{p}{:}\PY{p}{]}\PY{o}{/}\PY{l+m+mi}{100000}\PY{p}{,} \PY{n}{label} \PY{o}{=} \PY{l+s+sa}{f}\PY{l+s+s2}{\PYZdq{}}\PY{l+s+s2}{dR = }\PY{l+s+si}{\PYZob{}}\PY{n}{R}\PY{p}{[}\PY{n}{Ri}\PY{p}{[}\PY{n}{i}\PY{p}{]}\PY{p}{]}\PY{l+s+si}{:}\PY{l+s+s2}{.4f}\PY{l+s+si}{\PYZcb{}}\PY{l+s+s2}{ * 10}\PY{l+s+se}{\PYZbs{}u207b}\PY{l+s+se}{\PYZbs{}u00b9}\PY{l+s+se}{\PYZbs{}u2070}\PY{l+s+s2}{ m}\PY{l+s+s2}{\PYZdq{}}\PY{p}{)}

\PY{n}{plt}\PY{o}{.}\PY{n}{axline}\PY{p}{(}\PY{p}{(}\PY{n}{find\PYZus{}dB\PYZus{}B\PYZus{}osc}\PY{p}{(}\PY{l+m+mi}{4}\PY{p}{)}\PY{p}{,}\PY{o}{\PYZhy{}}\PY{l+m+mf}{0.5}\PY{p}{)}\PY{p}{,} \PY{p}{(}\PY{n}{find\PYZus{}dB\PYZus{}B\PYZus{}osc}\PY{p}{(}\PY{l+m+mi}{4}\PY{p}{)}\PY{p}{,}\PY{o}{\PYZhy{}}\PY{l+m+mf}{0.6}\PY{p}{)}\PY{p}{,} \PY{n}{color}\PY{o}{=}\PY{l+s+s1}{\PYZsq{}}\PY{l+s+s1}{k}\PY{l+s+s1}{\PYZsq{}}\PY{p}{,} \PY{n}{label} \PY{o}{=} \PY{l+s+s1}{\PYZsq{}}\PY{l+s+s1}{period boundary}\PY{l+s+s1}{\PYZsq{}}\PY{p}{)}
\PY{n}{plt}\PY{o}{.}\PY{n}{axline}\PY{p}{(}\PY{p}{(}\PY{n}{find\PYZus{}dB\PYZus{}B\PYZus{}osc}\PY{p}{(}\PY{l+m+mf}{4.5}\PY{p}{)}\PY{p}{,}\PY{o}{\PYZhy{}}\PY{l+m+mf}{0.5}\PY{p}{)}\PY{p}{,} \PY{p}{(}\PY{n}{find\PYZus{}dB\PYZus{}B\PYZus{}osc}\PY{p}{(}\PY{l+m+mf}{4.5}\PY{p}{)}\PY{p}{,}\PY{o}{\PYZhy{}}\PY{l+m+mf}{0.6}\PY{p}{)}\PY{p}{,} \PY{n}{color}\PY{o}{=}\PY{l+s+s1}{\PYZsq{}}\PY{l+s+s1}{k}\PY{l+s+s1}{\PYZsq{}}\PY{p}{)}

\PY{n}{plt}\PY{o}{.}\PY{n}{legend}\PY{p}{(}\PY{n}{loc}\PY{o}{=}\PY{l+s+s1}{\PYZsq{}}\PY{l+s+s1}{center left}\PY{l+s+s1}{\PYZsq{}}\PY{p}{,} \PY{n}{bbox\PYZus{}to\PYZus{}anchor}\PY{o}{=}\PY{p}{(}\PY{l+m+mi}{1}\PY{p}{,} \PY{l+m+mf}{0.5}\PY{p}{)}\PY{p}{)}

\PY{n}{plt}\PY{o}{.}\PY{n}{show}\PY{p}{(}\PY{p}{)}
\end{Verbatim}
\end{tcolorbox}

    \begin{center}
    \adjustimage{max size={0.9\linewidth}{0.9\paperheight}}{figs/bvp_kM_en/output_136_0.png}
    \end{center}
    { \hspace*{\fill} \\}
    
    At our midpoint of the \(R\) oscillation (the red curve above), our
pattern matches our oscillations at the period boundaries, but shifted
by half a period. This is the same behavior which we saw in the \(dkn\)
vs \(R\) plots, only with a different pattern. These cycles look like
cycloids.

    \begin{tcolorbox}[breakable, size=fbox, boxrule=1pt, pad at break*=1mm,colback=cellbackground, colframe=cellborder]
\prompt{In}{incolor}{ }{\boxspacing}
\begin{Verbatim}[commandchars=\\\{\}]
\PY{c+c1}{\PYZsh{} for easy changing which two plots to use.}

\PY{n}{Ri} \PY{o}{=} \PY{p}{[}\PY{l+m+mi}{44}\PY{p}{,} \PY{l+m+mi}{48}\PY{p}{,} \PY{l+m+mi}{52}\PY{p}{,} \PY{l+m+mi}{56}\PY{p}{,}\PY{l+m+mi}{60}\PY{p}{,} \PY{l+m+mi}{62}\PY{p}{,} \PY{l+m+mi}{63}\PY{p}{]}

\PY{n}{ilim} \PY{o}{=} \PY{n+nb}{len}\PY{p}{(}\PY{n}{Ri}\PY{p}{)}

\PY{n}{fig}\PY{p}{,} \PY{n}{ax} \PY{o}{=} \PY{n}{plt}\PY{o}{.}\PY{n}{subplots}\PY{p}{(}\PY{p}{)}
\PY{n}{plt}\PY{o}{.}\PY{n}{title}\PY{p}{(}\PY{l+s+sa}{f}\PY{l+s+s2}{\PYZdq{}}\PY{l+s+s2}{Change in k Due to Nonlinearity vs B for dk = 0.45 / 2 * 10}\PY{l+s+se}{\PYZbs{}u2076}\PY{l+s+s2}{ m}\PY{l+s+se}{\PYZbs{}u207b}\PY{l+s+se}{\PYZbs{}u00b9}\PY{l+s+se}{\PYZbs{}n}\PY{l+s+s2}{for Various dR Values and R = 9 * 10}\PY{l+s+se}{\PYZbs{}u207b}\PY{l+s+se}{\PYZbs{}u2077}\PY{l+s+s2}{ m + dR}\PY{l+s+s2}{\PYZdq{}}\PY{p}{)}
\PY{n}{ax}\PY{o}{.}\PY{n}{set\PYZus{}ylabel}\PY{p}{(}\PY{l+s+s1}{\PYZsq{}}\PY{l+s+s1}{Change in k (dkn) due to nonlinearity (10}\PY{l+s+se}{\PYZbs{}u2075}\PY{l+s+s1}{ m}\PY{l+s+se}{\PYZbs{}u207b}\PY{l+s+se}{\PYZbs{}u00b9}\PY{l+s+s1}{)}\PY{l+s+s1}{\PYZsq{}}\PY{p}{)}
\PY{n}{ax}\PY{o}{.}\PY{n}{set\PYZus{}xlabel}\PY{p}{(}\PY{l+s+s2}{\PYZdq{}}\PY{l+s+s2}{Variation (dB) of magnetic field strength B (B = 0.5 + dB * B\PYZus{}max)}\PY{l+s+s2}{\PYZdq{}}\PY{p}{)}

\PY{k}{for} \PY{n}{i} \PY{o+ow}{in} \PY{n+nb}{range}\PY{p}{(}\PY{n}{ilim}\PY{p}{)}\PY{p}{:}
    \PY{n}{ax}\PY{o}{.}\PY{n}{plot}\PY{p}{(}\PY{n}{B}\PY{p}{,} \PY{n}{dkn\PYZus{}arr}\PY{p}{[}\PY{n}{Ri}\PY{p}{[}\PY{n}{i}\PY{p}{]}\PY{p}{,}\PY{p}{:}\PY{p}{]}\PY{o}{/}\PY{l+m+mi}{100000}\PY{p}{,} \PY{n}{label} \PY{o}{=} \PY{l+s+sa}{f}\PY{l+s+s2}{\PYZdq{}}\PY{l+s+s2}{dR = }\PY{l+s+si}{\PYZob{}}\PY{n}{R}\PY{p}{[}\PY{n}{Ri}\PY{p}{[}\PY{n}{i}\PY{p}{]}\PY{p}{]}\PY{l+s+si}{:}\PY{l+s+s2}{.4f}\PY{l+s+si}{\PYZcb{}}\PY{l+s+s2}{ * 10}\PY{l+s+se}{\PYZbs{}u207b}\PY{l+s+se}{\PYZbs{}u00b9}\PY{l+s+se}{\PYZbs{}u2070}\PY{l+s+s2}{ m}\PY{l+s+s2}{\PYZdq{}}\PY{p}{)}

\PY{n}{plt}\PY{o}{.}\PY{n}{axline}\PY{p}{(}\PY{p}{(}\PY{n}{find\PYZus{}dB\PYZus{}B\PYZus{}osc}\PY{p}{(}\PY{l+m+mi}{4}\PY{p}{)}\PY{p}{,}\PY{o}{\PYZhy{}}\PY{l+m+mf}{0.5}\PY{p}{)}\PY{p}{,} \PY{p}{(}\PY{n}{find\PYZus{}dB\PYZus{}B\PYZus{}osc}\PY{p}{(}\PY{l+m+mi}{4}\PY{p}{)}\PY{p}{,}\PY{o}{\PYZhy{}}\PY{l+m+mf}{0.6}\PY{p}{)}\PY{p}{,} \PY{n}{color}\PY{o}{=}\PY{l+s+s1}{\PYZsq{}}\PY{l+s+s1}{k}\PY{l+s+s1}{\PYZsq{}}\PY{p}{,} \PY{n}{label} \PY{o}{=} \PY{l+s+s1}{\PYZsq{}}\PY{l+s+s1}{period boundary}\PY{l+s+s1}{\PYZsq{}}\PY{p}{)}
\PY{n}{plt}\PY{o}{.}\PY{n}{axline}\PY{p}{(}\PY{p}{(}\PY{n}{find\PYZus{}dB\PYZus{}B\PYZus{}osc}\PY{p}{(}\PY{l+m+mf}{4.5}\PY{p}{)}\PY{p}{,}\PY{o}{\PYZhy{}}\PY{l+m+mf}{0.5}\PY{p}{)}\PY{p}{,} \PY{p}{(}\PY{n}{find\PYZus{}dB\PYZus{}B\PYZus{}osc}\PY{p}{(}\PY{l+m+mf}{4.5}\PY{p}{)}\PY{p}{,}\PY{o}{\PYZhy{}}\PY{l+m+mf}{0.6}\PY{p}{)}\PY{p}{,} \PY{n}{color}\PY{o}{=}\PY{l+s+s1}{\PYZsq{}}\PY{l+s+s1}{k}\PY{l+s+s1}{\PYZsq{}}\PY{p}{)}

\PY{n}{plt}\PY{o}{.}\PY{n}{legend}\PY{p}{(}\PY{n}{loc}\PY{o}{=}\PY{l+s+s1}{\PYZsq{}}\PY{l+s+s1}{center left}\PY{l+s+s1}{\PYZsq{}}\PY{p}{,} \PY{n}{bbox\PYZus{}to\PYZus{}anchor}\PY{o}{=}\PY{p}{(}\PY{l+m+mi}{1}\PY{p}{,} \PY{l+m+mf}{0.5}\PY{p}{)}\PY{p}{)}

\PY{n}{plt}\PY{o}{.}\PY{n}{show}\PY{p}{(}\PY{p}{)}
\end{Verbatim}
\end{tcolorbox}

    \begin{center}
    \adjustimage{max size={0.9\linewidth}{0.9\paperheight}}{figs/bvp_kM_en/output_138_0.png}
    \end{center}
    { \hspace*{\fill} \\}
    
    The curves between pattern flipping on the side with the node loop like
the Jacobian elliptic dn functions, while on the side with the spike,
they look like two different symmetric curves, one facing up and one
facing down, with a jump connecting them. They might be cycloid arcs,
but I don't think so. I am going to try to matching the cycloid-looking
curves, then move on to the shape of the next oscillation.

    \begin{tcolorbox}[breakable, size=fbox, boxrule=1pt, pad at break*=1mm,colback=cellbackground, colframe=cellborder]
\prompt{In}{incolor}{154}{\boxspacing}
\begin{Verbatim}[commandchars=\\\{\}]
\PY{c+c1}{\PYZsh{} this is the parametric definition of a cycloid}
\PY{k}{def} \PY{n+nf}{cycl}\PY{p}{(}\PY{n}{r}\PY{p}{,} \PY{n}{t}\PY{p}{,} \PY{n}{x0}\PY{p}{)}\PY{p}{:}
    \PY{n}{x} \PY{o}{=} \PY{n}{r} \PY{o}{*} \PY{p}{(}\PY{n}{t} \PY{o}{\PYZhy{}} \PY{n}{np}\PY{o}{.}\PY{n}{sin}\PY{p}{(}\PY{n}{t}\PY{p}{)}\PY{p}{)} \PY{o}{+} \PY{n}{x0}
    \PY{n}{y} \PY{o}{=} \PY{n}{r} \PY{o}{*} \PY{p}{(}\PY{l+m+mi}{1} \PY{o}{\PYZhy{}} \PY{n}{np}\PY{o}{.}\PY{n}{cos}\PY{p}{(}\PY{n}{t}\PY{p}{)}\PY{p}{)}
    \PY{k}{return} \PY{n}{x}\PY{p}{,} \PY{n}{y}
\end{Verbatim}
\end{tcolorbox}

    \begin{tcolorbox}[breakable, size=fbox, boxrule=1pt, pad at break*=1mm,colback=cellbackground, colframe=cellborder]
\prompt{In}{incolor}{ }{\boxspacing}
\begin{Verbatim}[commandchars=\\\{\}]
\PY{n}{Ri} \PY{o}{=} \PY{l+m+mi}{33} \PY{c+c1}{\PYZsh{} This is our desired radius for analysis, when the cycloid appears}

\PY{n}{zero\PYZus{}ssh} \PY{o}{=} \PY{n}{np}\PY{o}{.}\PY{n}{nanmax}\PY{p}{(}\PY{n}{dkn\PYZus{}arr}\PY{p}{[}\PY{l+m+mi}{33}\PY{p}{]}\PY{p}{)} \PY{c+c1}{\PYZsh{} This is our smallest absolute value}
\PY{n}{arg\PYZus{}rad} \PY{o}{=} \PY{n}{np}\PY{o}{.}\PY{n}{nanargmax}\PY{p}{(}\PY{o}{\PYZhy{}}\PY{n}{dkn\PYZus{}arr}\PY{p}{[}\PY{l+m+mi}{33}\PY{p}{]}\PY{p}{)} \PY{c+c1}{\PYZsh{} This position in the array of our maximum absolute value}
\PY{c+c1}{\PYZsh{}print(zero\PYZus{}ssh, arg\PYZus{}rad)}
\end{Verbatim}
\end{tcolorbox}

    \begin{tcolorbox}[breakable, size=fbox, boxrule=1pt, pad at break*=1mm,colback=cellbackground, colframe=cellborder]
\prompt{In}{incolor}{156}{\boxspacing}
\begin{Verbatim}[commandchars=\\\{\}]
\PY{n}{Bi} \PY{o}{=} \PY{l+m+mi}{62} \PY{c+c1}{\PYZsh{} This is our maximum absolute value}
\PY{n}{rad1} \PY{o}{=} \PY{n}{B}\PY{p}{[}\PY{n}{Bi}\PY{p}{]}
\PY{n}{rad2} \PY{o}{=} \PY{n}{np}\PY{o}{.}\PY{n}{nanmax}\PY{p}{(}\PY{o}{\PYZhy{}}\PY{n}{dkn\PYZus{}arr}\PY{p}{[}\PY{n}{Ri}\PY{p}{]} \PY{o}{+} \PY{n}{zero\PYZus{}ssh}\PY{p}{)} \PY{c+c1}{\PYZsh{} We are shifting our curve so that the crests are at zero}

\PY{n}{y\PYZus{}arr} \PY{o}{=} \PY{p}{(}\PY{o}{\PYZhy{}}\PY{n}{dkn\PYZus{}arr}\PY{p}{[}\PY{n}{Ri}\PY{p}{]} \PY{o}{+} \PY{n}{zero\PYZus{}ssh}\PY{p}{)}\PY{o}{/}\PY{n}{rad2} \PY{c+c1}{\PYZsh{} rescale the y direction so that the }
\PY{n}{x\PYZus{}arr} \PY{o}{=} \PY{n}{B}\PY{o}{/}\PY{n}{rad1} \PY{o}{\PYZhy{}} \PY{l+m+mi}{1}   \PY{c+c1}{\PYZsh{} rescale the x direction, so that the oscillation goes from \PYZhy{}1 to 1}

\PY{n}{y\PYZus{}arr} \PY{o}{=} \PY{n}{y\PYZus{}arr}\PY{o}{.}\PY{n}{flatten}\PY{p}{(}\PY{p}{)}

\PY{n}{t\PYZus{}arr} \PY{o}{=} \PY{n}{np}\PY{o}{.}\PY{n}{linspace}\PY{p}{(}\PY{l+m+mi}{0}\PY{p}{,} \PY{l+m+mf}{2.7}\PY{o}{*}\PY{n}{pi}\PY{o}{/}\PY{l+m+mi}{2}\PY{p}{,} \PY{l+m+mi}{100}\PY{p}{)} \PY{c+c1}{\PYZsh{} make the t array for our cycloid}
            \PY{c+c1}{\PYZsh{} the end time here is not so important, as long as it is long enough}

\PY{n}{x\PYZus{}5}\PY{p}{,}\PY{n}{y\PYZus{}5} \PY{o}{=} \PY{n}{cycl}\PY{p}{(}\PY{l+m+mf}{0.503}\PY{p}{,}\PY{n}{t\PYZus{}arr}\PY{p}{,}\PY{o}{\PYZhy{}}\PY{n}{pi}\PY{o}{/}\PY{l+m+mi}{2}\PY{p}{)} \PY{c+c1}{\PYZsh{} x and y array for the cycloid, produced from t\PYZus{}arr}
                \PY{c+c1}{\PYZsh{} The radius was adjusted by hand, since my zero was not where I thought it was.}
                \PY{c+c1}{\PYZsh{} For a normal cycloid, the radius would be 0.5, not 0.503}

\PY{n}{x\PYZus{}5} \PY{o}{=} \PY{n}{x\PYZus{}5} \PY{o}{*}\PY{l+m+mi}{2} \PY{o}{/} \PY{n}{pi} \PY{c+c1}{\PYZsh{} rescale the x direction, so that the repeated pattern goes from \PYZhy{}1 to 1}
\end{Verbatim}
\end{tcolorbox}

    \begin{tcolorbox}[breakable, size=fbox, boxrule=1pt, pad at break*=1mm,colback=cellbackground, colframe=cellborder]
\prompt{In}{incolor}{160}{\boxspacing}
\begin{Verbatim}[commandchars=\\\{\}]
\PY{c+c1}{\PYZsh{} for easy changing which two plots to use.}

\PY{n}{Ri} \PY{o}{=} \PY{p}{[}\PY{l+m+mi}{33}\PY{p}{]}

\PY{n}{ilim} \PY{o}{=} \PY{n+nb}{len}\PY{p}{(}\PY{n}{Ri}\PY{p}{)}

\PY{n}{fig}\PY{p}{,} \PY{n}{ax} \PY{o}{=} \PY{n}{plt}\PY{o}{.}\PY{n}{subplots}\PY{p}{(}\PY{p}{)}
\PY{n}{plt}\PY{o}{.}\PY{n}{title}\PY{p}{(}\PY{l+s+sa}{f}\PY{l+s+s2}{\PYZdq{}}\PY{l+s+s2}{Change in k Due to Nonlinearity vs B,}\PY{l+s+se}{\PYZbs{}n}\PY{l+s+s2}{Normalized to Model Shape of Function}\PY{l+s+s2}{\PYZdq{}}\PY{p}{)}
\PY{n}{ax}\PY{o}{.}\PY{n}{set\PYZus{}ylabel}\PY{p}{(}\PY{l+s+s1}{\PYZsq{}}\PY{l+s+s1}{Change in k (dkn) due to nonlinearity (normalized)}\PY{l+s+s1}{\PYZsq{}}\PY{p}{)}
\PY{n}{ax}\PY{o}{.}\PY{n}{set\PYZus{}xlabel}\PY{p}{(}\PY{l+s+s2}{\PYZdq{}}\PY{l+s+s2}{Variation (dB) of magnetic field strength B (normalized)}\PY{l+s+s2}{\PYZdq{}}\PY{p}{)}

\PY{n}{ax}\PY{o}{.}\PY{n}{plot}\PY{p}{(}\PY{n}{x\PYZus{}arr}\PY{p}{,} \PY{n}{y\PYZus{}arr}\PY{p}{,} \PY{n}{label} \PY{o}{=} \PY{l+s+s2}{\PYZdq{}}\PY{l+s+s2}{target curve}\PY{l+s+s2}{\PYZdq{}}\PY{p}{)}
\PY{n}{ax}\PY{o}{.}\PY{n}{plot}\PY{p}{(}\PY{n}{x\PYZus{}5}\PY{o}{\PYZhy{}}\PY{l+m+mf}{0.01}\PY{p}{,} \PY{n}{y\PYZus{}5}\PY{o}{*}\PY{l+m+mf}{1.1} \PY{o}{\PYZhy{}}\PY{l+m+mf}{0.105}\PY{p}{,} \PY{n}{label} \PY{o}{=} \PY{l+s+s2}{\PYZdq{}}\PY{l+s+s2}{cycloid model}\PY{l+s+s2}{\PYZdq{}}\PY{p}{)} \PY{c+c1}{\PYZsh{} parameters here are adjusted, because a normal cycloid with the cusps at zero didn\PYZsq{}t work. }

\PY{n}{plt}\PY{o}{.}\PY{n}{legend}\PY{p}{(}\PY{p}{)}

\PY{n}{plt}\PY{o}{.}\PY{n}{show}\PY{p}{(}\PY{p}{)}
\end{Verbatim}
\end{tcolorbox}

    \begin{center}
    \adjustimage{max size={0.9\linewidth}{0.9\paperheight}}{figs/bvp_kM_en/output_143_0.png}
    \end{center}
    { \hspace*{\fill} \\}
    
    Well, many of the parameters were chosen by hand, but it does fit the
shape. The cusps are not at the same place as for a cycloid, instead
starting further up in the curve, and not fully going to zero. So it is
not quite a cycloid like I thought.

    \hypertarget{dmn-vs-b}{%
\section{\texorpdfstring{\(dMn\) vs \(B\)}{dMn vs B}}\label{dmn-vs-b}}

    So, we have so far found the shapes of both \(dkn\) oscillations. Now we
need to examine the same two oscillations in \(dMn\).

    \begin{tcolorbox}[breakable, size=fbox, boxrule=1pt, pad at break*=1mm,colback=cellbackground, colframe=cellborder]
\prompt{In}{incolor}{169}{\boxspacing}
\begin{Verbatim}[commandchars=\\\{\}]
\PY{c+c1}{\PYZsh{} Select which R values we want. }
\PY{c+c1}{\PYZsh{} Let us start with the peaks and nodes, like before}
\PY{n}{Ri} \PY{o}{=} \PY{p}{[}\PY{l+m+mi}{33}\PY{p}{,} \PY{l+m+mi}{54}\PY{p}{,} \PY{l+m+mi}{75}\PY{p}{]}
\PY{n}{ilim} \PY{o}{=} \PY{n+nb}{len}\PY{p}{(}\PY{n}{Ri}\PY{p}{)}

\PY{n}{fig}\PY{p}{,} \PY{n}{ax} \PY{o}{=} \PY{n}{plt}\PY{o}{.}\PY{n}{subplots}\PY{p}{(}\PY{p}{)}
\PY{n}{plt}\PY{o}{.}\PY{n}{title}\PY{p}{(}\PY{l+s+sa}{f}\PY{l+s+s2}{\PYZdq{}}\PY{l+s+s2}{Change in M Due to Nonlinearity vs B for dk = 0.45 / 2 * 10}\PY{l+s+se}{\PYZbs{}u2076}\PY{l+s+s2}{ m}\PY{l+s+se}{\PYZbs{}u207b}\PY{l+s+se}{\PYZbs{}u00b9}\PY{l+s+se}{\PYZbs{}n}\PY{l+s+s2}{for Various dR Values and R = 9 * 10}\PY{l+s+se}{\PYZbs{}u207b}\PY{l+s+se}{\PYZbs{}u2077}\PY{l+s+s2}{ m + dR}\PY{l+s+s2}{\PYZdq{}}\PY{p}{)}
\PY{n}{ax}\PY{o}{.}\PY{n}{set\PYZus{}ylabel}\PY{p}{(}\PY{l+s+s1}{\PYZsq{}}\PY{l+s+s1}{Change in M (dMn) due to nonlinearity (10}\PY{l+s+se}{\PYZbs{}u2075}\PY{l+s+s1}{ m}\PY{l+s+se}{\PYZbs{}u207b}\PY{l+s+se}{\PYZbs{}u00b9}\PY{l+s+s1}{)}\PY{l+s+s1}{\PYZsq{}}\PY{p}{)}
\PY{n}{ax}\PY{o}{.}\PY{n}{set\PYZus{}xlabel}\PY{p}{(}\PY{l+s+s2}{\PYZdq{}}\PY{l+s+s2}{Variation (dB) of magnetic field strength B (B = 0.5 + dB * B\PYZus{}max)}\PY{l+s+s2}{\PYZdq{}}\PY{p}{)}

\PY{k}{for} \PY{n}{i} \PY{o+ow}{in} \PY{n+nb}{range}\PY{p}{(}\PY{n}{ilim}\PY{p}{)}\PY{p}{:}
    \PY{n}{ax}\PY{o}{.}\PY{n}{plot}\PY{p}{(}\PY{n}{B}\PY{p}{,} \PY{n}{dMn\PYZus{}arr}\PY{p}{[}\PY{n}{Ri}\PY{p}{[}\PY{n}{i}\PY{p}{]}\PY{p}{,}\PY{p}{:}\PY{p}{]}\PY{o}{/}\PY{l+m+mi}{100000}\PY{p}{,} \PY{n}{label} \PY{o}{=} \PY{l+s+sa}{f}\PY{l+s+s2}{\PYZdq{}}\PY{l+s+s2}{dR = }\PY{l+s+si}{\PYZob{}}\PY{n}{R}\PY{p}{[}\PY{n}{Ri}\PY{p}{[}\PY{n}{i}\PY{p}{]}\PY{p}{]}\PY{l+s+si}{:}\PY{l+s+s2}{.4f}\PY{l+s+si}{\PYZcb{}}\PY{l+s+s2}{ * 10}\PY{l+s+se}{\PYZbs{}u207b}\PY{l+s+se}{\PYZbs{}u00b9}\PY{l+s+se}{\PYZbs{}u2070}\PY{l+s+s2}{ m}\PY{l+s+s2}{\PYZdq{}}\PY{p}{)}

\PY{n}{plt}\PY{o}{.}\PY{n}{legend}\PY{p}{(}\PY{n}{loc}\PY{o}{=}\PY{l+s+s1}{\PYZsq{}}\PY{l+s+s1}{center left}\PY{l+s+s1}{\PYZsq{}}\PY{p}{,} \PY{n}{bbox\PYZus{}to\PYZus{}anchor}\PY{o}{=}\PY{p}{(}\PY{l+m+mf}{0.63}\PY{p}{,} \PY{l+m+mf}{0.6}\PY{p}{)}\PY{p}{)}

\PY{n}{plt}\PY{o}{.}\PY{n}{show}\PY{p}{(}\PY{p}{)}
\end{Verbatim}
\end{tcolorbox}

    \begin{center}
    \adjustimage{max size={0.9\linewidth}{0.9\paperheight}}{figs/bvp_kM_en/output_147_0.png}
    \end{center}
    { \hspace*{\fill} \\}
    
    We have a skewed sine here, as noted before. In these graphs, there is a
discontinuity when the oscillation begins. The point which goes through
our axis in the center of our oscillation is a straight line.

    \begin{tcolorbox}[breakable, size=fbox, boxrule=1pt, pad at break*=1mm,colback=cellbackground, colframe=cellborder]
\prompt{In}{incolor}{186}{\boxspacing}
\begin{Verbatim}[commandchars=\\\{\}]
\PY{n}{Ri} \PY{o}{=} \PY{p}{[}\PY{l+m+mi}{33}\PY{p}{,} \PY{l+m+mi}{36}\PY{p}{,} \PY{l+m+mi}{39}\PY{p}{,} \PY{l+m+mi}{42}\PY{p}{]}

\PY{n}{ilim} \PY{o}{=} \PY{n+nb}{len}\PY{p}{(}\PY{n}{Ri}\PY{p}{)}

\PY{n}{fig}\PY{p}{,} \PY{n}{ax} \PY{o}{=} \PY{n}{plt}\PY{o}{.}\PY{n}{subplots}\PY{p}{(}\PY{p}{)}
\PY{n}{plt}\PY{o}{.}\PY{n}{title}\PY{p}{(}\PY{l+s+sa}{f}\PY{l+s+s2}{\PYZdq{}}\PY{l+s+s2}{Change in M Due to Nonlinearity vs B for dk = 0.45 / 2 * 10}\PY{l+s+se}{\PYZbs{}u2076}\PY{l+s+s2}{ m}\PY{l+s+se}{\PYZbs{}u207b}\PY{l+s+se}{\PYZbs{}u00b9}\PY{l+s+se}{\PYZbs{}n}\PY{l+s+s2}{for Various dR Values and R = 9 * 10}\PY{l+s+se}{\PYZbs{}u207b}\PY{l+s+se}{\PYZbs{}u2077}\PY{l+s+s2}{ m + dR}\PY{l+s+s2}{\PYZdq{}}\PY{p}{)}
\PY{n}{ax}\PY{o}{.}\PY{n}{set\PYZus{}ylabel}\PY{p}{(}\PY{l+s+s1}{\PYZsq{}}\PY{l+s+s1}{Change in M (dMn) due to nonlinearity (10}\PY{l+s+se}{\PYZbs{}u2074}\PY{l+s+s1}{ m}\PY{l+s+se}{\PYZbs{}u207b}\PY{l+s+se}{\PYZbs{}u00b9}\PY{l+s+s1}{)}\PY{l+s+s1}{\PYZsq{}}\PY{p}{)}
\PY{n}{ax}\PY{o}{.}\PY{n}{set\PYZus{}xlabel}\PY{p}{(}\PY{l+s+s2}{\PYZdq{}}\PY{l+s+s2}{Variation (dB) of magnetic field strength B (B = 0.5 + dB * B\PYZus{}max)}\PY{l+s+s2}{\PYZdq{}}\PY{p}{)}

\PY{k}{for} \PY{n}{i} \PY{o+ow}{in} \PY{n+nb}{range}\PY{p}{(}\PY{n}{ilim}\PY{p}{)}\PY{p}{:}
    \PY{n}{ax}\PY{o}{.}\PY{n}{plot}\PY{p}{(}\PY{n}{B}\PY{p}{,} \PY{n}{dMn\PYZus{}arr}\PY{p}{[}\PY{n}{Ri}\PY{p}{[}\PY{n}{i}\PY{p}{]}\PY{p}{,}\PY{p}{:}\PY{p}{]}\PY{o}{/}\PY{l+m+mi}{10000}\PY{p}{,} \PY{n}{label} \PY{o}{=} \PY{l+s+sa}{f}\PY{l+s+s2}{\PYZdq{}}\PY{l+s+s2}{dR = }\PY{l+s+si}{\PYZob{}}\PY{n}{R}\PY{p}{[}\PY{n}{Ri}\PY{p}{[}\PY{n}{i}\PY{p}{]}\PY{p}{]}\PY{l+s+si}{:}\PY{l+s+s2}{.4f}\PY{l+s+si}{\PYZcb{}}\PY{l+s+s2}{ * 10}\PY{l+s+se}{\PYZbs{}u207b}\PY{l+s+se}{\PYZbs{}u00b9}\PY{l+s+se}{\PYZbs{}u2070}\PY{l+s+s2}{ m}\PY{l+s+s2}{\PYZdq{}}\PY{p}{)}
\PY{n}{plt}\PY{o}{.}\PY{n}{axline}\PY{p}{(}\PY{p}{(}\PY{n}{find\PYZus{}dB\PYZus{}B\PYZus{}osc}\PY{p}{(}\PY{l+m+mi}{4}\PY{p}{)}\PY{p}{,}\PY{o}{\PYZhy{}}\PY{l+m+mf}{0.5}\PY{p}{)}\PY{p}{,} \PY{p}{(}\PY{n}{find\PYZus{}dB\PYZus{}B\PYZus{}osc}\PY{p}{(}\PY{l+m+mi}{4}\PY{p}{)}\PY{p}{,}\PY{o}{\PYZhy{}}\PY{l+m+mf}{0.6}\PY{p}{)}\PY{p}{,} \PY{n}{color}\PY{o}{=}\PY{l+s+s1}{\PYZsq{}}\PY{l+s+s1}{k}\PY{l+s+s1}{\PYZsq{}}\PY{p}{,} \PY{n}{label} \PY{o}{=} \PY{l+s+s2}{\PYZdq{}}\PY{l+s+s2}{period boundary}\PY{l+s+s2}{\PYZdq{}}\PY{p}{)}
\PY{n}{plt}\PY{o}{.}\PY{n}{axline}\PY{p}{(}\PY{p}{(}\PY{n}{find\PYZus{}dB\PYZus{}B\PYZus{}osc}\PY{p}{(}\PY{l+m+mf}{4.5}\PY{p}{)}\PY{p}{,}\PY{o}{\PYZhy{}}\PY{l+m+mf}{0.5}\PY{p}{)}\PY{p}{,} \PY{p}{(}\PY{n}{find\PYZus{}dB\PYZus{}B\PYZus{}osc}\PY{p}{(}\PY{l+m+mf}{4.5}\PY{p}{)}\PY{p}{,}\PY{o}{\PYZhy{}}\PY{l+m+mf}{0.6}\PY{p}{)}\PY{p}{,} \PY{n}{color}\PY{o}{=}\PY{l+s+s1}{\PYZsq{}}\PY{l+s+s1}{k}\PY{l+s+s1}{\PYZsq{}}\PY{p}{)}

\PY{n}{plt}\PY{o}{.}\PY{n}{legend}\PY{p}{(}\PY{n}{loc}\PY{o}{=}\PY{l+s+s1}{\PYZsq{}}\PY{l+s+s1}{center left}\PY{l+s+s1}{\PYZsq{}}\PY{p}{,} \PY{n}{bbox\PYZus{}to\PYZus{}anchor}\PY{o}{=}\PY{p}{(}\PY{l+m+mi}{1}\PY{p}{,} \PY{l+m+mf}{0.5}\PY{p}{)}\PY{p}{)}

\PY{n}{plt}\PY{o}{.}\PY{n}{show}\PY{p}{(}\PY{p}{)}
\end{Verbatim}
\end{tcolorbox}

    \begin{center}
    \adjustimage{max size={0.9\linewidth}{0.9\paperheight}}{figs/bvp_kM_en/output_149_0.png}
    \end{center}
    { \hspace*{\fill} \\}
    
    The discontinuity is building to the phase-shifted pattern. Both
patterns are rotating, and the cutoff is moving towards the central
point of the old oscillation, evenly from both directions.

    \begin{tcolorbox}[breakable, size=fbox, boxrule=1pt, pad at break*=1mm,colback=cellbackground, colframe=cellborder]
\prompt{In}{incolor}{177}{\boxspacing}
\begin{Verbatim}[commandchars=\\\{\}]
\PY{c+c1}{\PYZsh{} Find the index of a quarter period}
\PY{n}{find\PYZus{}dR\PYZus{}k\PYZus{}osc}\PY{p}{(}\PY{l+m+mf}{1.25}\PY{p}{)}\PY{o}{*}\PY{l+m+mi}{100}\PY{o}{/}\PY{l+m+mi}{2}
\end{Verbatim}
\end{tcolorbox}

            \begin{tcolorbox}[breakable, size=fbox, boxrule=.5pt, pad at break*=1mm, opacityfill=0]
\prompt{Out}{outcolor}{177}{\boxspacing}
\begin{Verbatim}[commandchars=\\\{\}]
43.64583333333333
\end{Verbatim}
\end{tcolorbox}
        
    \begin{tcolorbox}[breakable, size=fbox, boxrule=1pt, pad at break*=1mm,colback=cellbackground, colframe=cellborder]
\prompt{In}{incolor}{188}{\boxspacing}
\begin{Verbatim}[commandchars=\\\{\}]
\PY{n}{Ri} \PY{o}{=} \PY{p}{[}\PY{l+m+mi}{43}\PY{p}{,} \PY{l+m+mi}{44}\PY{p}{,} \PY{l+m+mi}{45}\PY{p}{]}

\PY{n}{ilim} \PY{o}{=} \PY{n+nb}{len}\PY{p}{(}\PY{n}{Ri}\PY{p}{)}

\PY{n}{y} \PY{o}{=} \PY{o}{\PYZhy{}}\PY{n}{np}\PY{o}{.}\PY{n}{sin}\PY{p}{(}\PY{l+m+mi}{2}\PY{o}{*} \PY{n}{pi} \PY{o}{*} \PY{p}{(}\PY{n}{B}\PY{o}{\PYZhy{}}\PY{n}{find\PYZus{}dB\PYZus{}B\PYZus{}osc}\PY{p}{(}\PY{l+m+mi}{4}\PY{p}{)}\PY{p}{)}\PY{o}{/} \PY{p}{(}\PY{n}{find\PYZus{}dB\PYZus{}B\PYZus{}osc}\PY{p}{(}\PY{l+m+mi}{5}\PY{p}{)}\PY{o}{\PYZhy{}}\PY{n}{find\PYZus{}dB\PYZus{}B\PYZus{}osc}\PY{p}{(}\PY{l+m+mi}{4}\PY{p}{)}\PY{p}{)}\PY{p}{)}\PY{o}{*}\PY{l+m+mf}{1.7}

\PY{n}{fig}\PY{p}{,} \PY{n}{ax} \PY{o}{=} \PY{n}{plt}\PY{o}{.}\PY{n}{subplots}\PY{p}{(}\PY{p}{)}
\PY{n}{plt}\PY{o}{.}\PY{n}{title}\PY{p}{(}\PY{l+s+sa}{f}\PY{l+s+s2}{\PYZdq{}}\PY{l+s+s2}{Change in M Due to Nonlinearity vs B for dk = 0.45 / 2 * 10}\PY{l+s+se}{\PYZbs{}u2076}\PY{l+s+s2}{ m}\PY{l+s+se}{\PYZbs{}u207b}\PY{l+s+se}{\PYZbs{}u00b9}\PY{l+s+se}{\PYZbs{}n}\PY{l+s+s2}{for Various dR Values and R = 9 * 10}\PY{l+s+se}{\PYZbs{}u207b}\PY{l+s+se}{\PYZbs{}u2077}\PY{l+s+s2}{ m + dR}\PY{l+s+s2}{\PYZdq{}}\PY{p}{)}
\PY{n}{ax}\PY{o}{.}\PY{n}{set\PYZus{}ylabel}\PY{p}{(}\PY{l+s+s1}{\PYZsq{}}\PY{l+s+s1}{Change in M (dMn) due to nonlinearity (10}\PY{l+s+se}{\PYZbs{}u2074}\PY{l+s+s1}{ m}\PY{l+s+se}{\PYZbs{}u207b}\PY{l+s+se}{\PYZbs{}u00b9}\PY{l+s+s1}{)}\PY{l+s+s1}{\PYZsq{}}\PY{p}{)}
\PY{n}{ax}\PY{o}{.}\PY{n}{set\PYZus{}xlabel}\PY{p}{(}\PY{l+s+s2}{\PYZdq{}}\PY{l+s+s2}{Variation (dB) of magnetic field strength B (B = 0.5 + dB * B\PYZus{}max)}\PY{l+s+s2}{\PYZdq{}}\PY{p}{)}

\PY{k}{for} \PY{n}{i} \PY{o+ow}{in} \PY{n+nb}{range}\PY{p}{(}\PY{n}{ilim}\PY{p}{)}\PY{p}{:}
    \PY{n}{ax}\PY{o}{.}\PY{n}{plot}\PY{p}{(}\PY{n}{B}\PY{p}{,} \PY{n}{dMn\PYZus{}arr}\PY{p}{[}\PY{n}{Ri}\PY{p}{[}\PY{n}{i}\PY{p}{]}\PY{p}{,}\PY{p}{:}\PY{p}{]}\PY{o}{/}\PY{l+m+mi}{10000}\PY{p}{,} \PY{n}{label} \PY{o}{=} \PY{l+s+sa}{f}\PY{l+s+s2}{\PYZdq{}}\PY{l+s+s2}{dR = }\PY{l+s+si}{\PYZob{}}\PY{n}{R}\PY{p}{[}\PY{n}{Ri}\PY{p}{[}\PY{n}{i}\PY{p}{]}\PY{p}{]}\PY{l+s+si}{:}\PY{l+s+s2}{.4f}\PY{l+s+si}{\PYZcb{}}\PY{l+s+s2}{ * 10}\PY{l+s+se}{\PYZbs{}u207b}\PY{l+s+se}{\PYZbs{}u00b9}\PY{l+s+se}{\PYZbs{}u2070}\PY{l+s+s2}{ m}\PY{l+s+s2}{\PYZdq{}}\PY{p}{)}
\PY{n}{ax}\PY{o}{.}\PY{n}{plot}\PY{p}{(}\PY{n}{B}\PY{p}{,} \PY{n}{y}\PY{p}{)}
\PY{n}{plt}\PY{o}{.}\PY{n}{axline}\PY{p}{(}\PY{p}{(}\PY{n}{find\PYZus{}dB\PYZus{}B\PYZus{}osc}\PY{p}{(}\PY{l+m+mi}{4}\PY{p}{)}\PY{p}{,}\PY{o}{\PYZhy{}}\PY{l+m+mf}{0.5}\PY{p}{)}\PY{p}{,} \PY{p}{(}\PY{n}{find\PYZus{}dB\PYZus{}B\PYZus{}osc}\PY{p}{(}\PY{l+m+mi}{4}\PY{p}{)}\PY{p}{,}\PY{o}{\PYZhy{}}\PY{l+m+mf}{0.6}\PY{p}{)}\PY{p}{,} \PY{n}{color}\PY{o}{=}\PY{l+s+s1}{\PYZsq{}}\PY{l+s+s1}{k}\PY{l+s+s1}{\PYZsq{}}\PY{p}{,} \PY{n}{label} \PY{o}{=}\PY{l+s+s2}{\PYZdq{}}\PY{l+s+s2}{period boundary}\PY{l+s+s2}{\PYZdq{}}\PY{p}{)}
\PY{n}{plt}\PY{o}{.}\PY{n}{axline}\PY{p}{(}\PY{p}{(}\PY{n}{find\PYZus{}dB\PYZus{}B\PYZus{}osc}\PY{p}{(}\PY{l+m+mf}{4.5}\PY{p}{)}\PY{p}{,}\PY{o}{\PYZhy{}}\PY{l+m+mf}{0.5}\PY{p}{)}\PY{p}{,} \PY{p}{(}\PY{n}{find\PYZus{}dB\PYZus{}B\PYZus{}osc}\PY{p}{(}\PY{l+m+mf}{4.5}\PY{p}{)}\PY{p}{,}\PY{o}{\PYZhy{}}\PY{l+m+mf}{0.6}\PY{p}{)}\PY{p}{,} \PY{n}{color}\PY{o}{=}\PY{l+s+s1}{\PYZsq{}}\PY{l+s+s1}{k}\PY{l+s+s1}{\PYZsq{}}\PY{p}{)}

\PY{n}{plt}\PY{o}{.}\PY{n}{legend}\PY{p}{(}\PY{p}{)}\PY{c+c1}{\PYZsh{}loc=\PYZsq{}center left\PYZsq{}, bbox\PYZus{}to\PYZus{}anchor=(1, 0.5))}

\PY{n}{plt}\PY{o}{.}\PY{n}{show}\PY{p}{(}\PY{p}{)}
\end{Verbatim}
\end{tcolorbox}

    \begin{center}
    \adjustimage{max size={0.9\linewidth}{0.9\paperheight}}{figs/bvp_kM_en/output_152_0.png}
    \end{center}
    { \hspace*{\fill} \\}
    
    At our quarter period, we loose the entirety of our old oscillations,
and settle into the new oscillations. The red sine curve here shows that
we still have a skew in our pattern.

    \begin{tcolorbox}[breakable, size=fbox, boxrule=1pt, pad at break*=1mm,colback=cellbackground, colframe=cellborder]
\prompt{In}{incolor}{190}{\boxspacing}
\begin{Verbatim}[commandchars=\\\{\}]
\PY{n}{Ri} \PY{o}{=} \PY{p}{[}\PY{l+m+mi}{45}\PY{p}{,} \PY{l+m+mi}{48}\PY{p}{,} \PY{l+m+mi}{51}\PY{p}{,} \PY{l+m+mi}{54}\PY{p}{]}

\PY{n}{ilim} \PY{o}{=} \PY{n+nb}{len}\PY{p}{(}\PY{n}{Ri}\PY{p}{)}

\PY{n}{fig}\PY{p}{,} \PY{n}{ax} \PY{o}{=} \PY{n}{plt}\PY{o}{.}\PY{n}{subplots}\PY{p}{(}\PY{p}{)}
\PY{n}{plt}\PY{o}{.}\PY{n}{title}\PY{p}{(}\PY{l+s+sa}{f}\PY{l+s+s2}{\PYZdq{}}\PY{l+s+s2}{Change in M Due to Nonlinearity vs B for dk = 0.45 / 2 * 10}\PY{l+s+se}{\PYZbs{}u2076}\PY{l+s+s2}{ m}\PY{l+s+se}{\PYZbs{}u207b}\PY{l+s+se}{\PYZbs{}u00b9}\PY{l+s+se}{\PYZbs{}n}\PY{l+s+s2}{for Various dR Values and R = 9 * 10}\PY{l+s+se}{\PYZbs{}u207b}\PY{l+s+se}{\PYZbs{}u2077}\PY{l+s+s2}{ m + dR}\PY{l+s+s2}{\PYZdq{}}\PY{p}{)}
\PY{n}{ax}\PY{o}{.}\PY{n}{set\PYZus{}ylabel}\PY{p}{(}\PY{l+s+s1}{\PYZsq{}}\PY{l+s+s1}{Change in M (dMn) due to nonlinearity (10}\PY{l+s+se}{\PYZbs{}u2074}\PY{l+s+s1}{ m}\PY{l+s+se}{\PYZbs{}u207b}\PY{l+s+se}{\PYZbs{}u00b9}\PY{l+s+s1}{)}\PY{l+s+s1}{\PYZsq{}}\PY{p}{)}
\PY{n}{ax}\PY{o}{.}\PY{n}{set\PYZus{}xlabel}\PY{p}{(}\PY{l+s+s2}{\PYZdq{}}\PY{l+s+s2}{Variation (dB) of magnetic field strength B (B = 0.5 + dB * B\PYZus{}max)}\PY{l+s+s2}{\PYZdq{}}\PY{p}{)}

\PY{k}{for} \PY{n}{i} \PY{o+ow}{in} \PY{n+nb}{range}\PY{p}{(}\PY{n}{ilim}\PY{p}{)}\PY{p}{:}
    \PY{n}{ax}\PY{o}{.}\PY{n}{plot}\PY{p}{(}\PY{n}{B}\PY{p}{,} \PY{n}{dMn\PYZus{}arr}\PY{p}{[}\PY{n}{Ri}\PY{p}{[}\PY{n}{i}\PY{p}{]}\PY{p}{,}\PY{p}{:}\PY{p}{]}\PY{o}{/}\PY{l+m+mi}{10000}\PY{p}{,} \PY{n}{label} \PY{o}{=} \PY{l+s+sa}{f}\PY{l+s+s2}{\PYZdq{}}\PY{l+s+s2}{dR = }\PY{l+s+si}{\PYZob{}}\PY{n}{R}\PY{p}{[}\PY{n}{Ri}\PY{p}{[}\PY{n}{i}\PY{p}{]}\PY{p}{]}\PY{l+s+si}{:}\PY{l+s+s2}{.4f}\PY{l+s+si}{\PYZcb{}}\PY{l+s+s2}{ * 10}\PY{l+s+se}{\PYZbs{}u207b}\PY{l+s+se}{\PYZbs{}u00b9}\PY{l+s+se}{\PYZbs{}u2070}\PY{l+s+s2}{ m}\PY{l+s+s2}{\PYZdq{}}\PY{p}{)}
    
\PY{n}{plt}\PY{o}{.}\PY{n}{axline}\PY{p}{(}\PY{p}{(}\PY{n}{find\PYZus{}dB\PYZus{}B\PYZus{}osc}\PY{p}{(}\PY{l+m+mi}{4}\PY{p}{)}\PY{p}{,}\PY{o}{\PYZhy{}}\PY{l+m+mf}{0.5}\PY{p}{)}\PY{p}{,} \PY{p}{(}\PY{n}{find\PYZus{}dB\PYZus{}B\PYZus{}osc}\PY{p}{(}\PY{l+m+mi}{4}\PY{p}{)}\PY{p}{,}\PY{o}{\PYZhy{}}\PY{l+m+mf}{0.6}\PY{p}{)}\PY{p}{,} \PY{n}{color}\PY{o}{=}\PY{l+s+s1}{\PYZsq{}}\PY{l+s+s1}{k}\PY{l+s+s1}{\PYZsq{}}\PY{p}{,} \PY{n}{label} \PY{o}{=} \PY{l+s+s2}{\PYZdq{}}\PY{l+s+s2}{period boundary}\PY{l+s+s2}{\PYZdq{}}\PY{p}{)}
\PY{n}{plt}\PY{o}{.}\PY{n}{axline}\PY{p}{(}\PY{p}{(}\PY{n}{find\PYZus{}dB\PYZus{}B\PYZus{}osc}\PY{p}{(}\PY{l+m+mf}{4.5}\PY{p}{)}\PY{p}{,}\PY{o}{\PYZhy{}}\PY{l+m+mf}{0.5}\PY{p}{)}\PY{p}{,} \PY{p}{(}\PY{n}{find\PYZus{}dB\PYZus{}B\PYZus{}osc}\PY{p}{(}\PY{l+m+mf}{4.5}\PY{p}{)}\PY{p}{,}\PY{o}{\PYZhy{}}\PY{l+m+mf}{0.6}\PY{p}{)}\PY{p}{,} \PY{n}{color}\PY{o}{=}\PY{l+s+s1}{\PYZsq{}}\PY{l+s+s1}{k}\PY{l+s+s1}{\PYZsq{}}\PY{p}{)}

\PY{n}{plt}\PY{o}{.}\PY{n}{legend}\PY{p}{(}\PY{p}{)}\PY{c+c1}{\PYZsh{}loc=\PYZsq{}center left\PYZsq{}, bbox\PYZus{}to\PYZus{}anchor=(1, 0.5))}

\PY{n}{plt}\PY{o}{.}\PY{n}{show}\PY{p}{(}\PY{p}{)}
\end{Verbatim}
\end{tcolorbox}

    \begin{center}
    \adjustimage{max size={0.9\linewidth}{0.9\paperheight}}{figs/bvp_kM_en/output_154_0.png}
    \end{center}
    { \hspace*{\fill} \\}
    
    Here there is only one skewed sine curve from the quarter period to the
half period of our oscillation, when the discontinuity appears again.

    \begin{tcolorbox}[breakable, size=fbox, boxrule=1pt, pad at break*=1mm,colback=cellbackground, colframe=cellborder]
\prompt{In}{incolor}{191}{\boxspacing}
\begin{Verbatim}[commandchars=\\\{\}]
\PY{c+c1}{\PYZsh{} Find the index of three quarters period}
\PY{n}{find\PYZus{}dR\PYZus{}k\PYZus{}osc}\PY{p}{(}\PY{l+m+mf}{1.75}\PY{p}{)}\PY{o}{*}\PY{l+m+mi}{100}\PY{o}{/}\PY{l+m+mi}{2}
\end{Verbatim}
\end{tcolorbox}

            \begin{tcolorbox}[breakable, size=fbox, boxrule=.5pt, pad at break*=1mm, opacityfill=0]
\prompt{Out}{outcolor}{191}{\boxspacing}
\begin{Verbatim}[commandchars=\\\{\}]
64.47916666666666
\end{Verbatim}
\end{tcolorbox}
        
    \begin{tcolorbox}[breakable, size=fbox, boxrule=1pt, pad at break*=1mm,colback=cellbackground, colframe=cellborder]
\prompt{In}{incolor}{192}{\boxspacing}
\begin{Verbatim}[commandchars=\\\{\}]
\PY{n}{Ri} \PY{o}{=} \PY{p}{[}\PY{l+m+mi}{56}\PY{p}{,} \PY{l+m+mi}{58}\PY{p}{,} \PY{l+m+mi}{60}\PY{p}{,} \PY{l+m+mi}{62}\PY{p}{,} \PY{l+m+mi}{63}\PY{p}{,} \PY{l+m+mi}{64}\PY{p}{]}

\PY{n}{ilim} \PY{o}{=} \PY{n+nb}{len}\PY{p}{(}\PY{n}{Ri}\PY{p}{)}

\PY{n}{fig}\PY{p}{,} \PY{n}{ax} \PY{o}{=} \PY{n}{plt}\PY{o}{.}\PY{n}{subplots}\PY{p}{(}\PY{p}{)}
\PY{n}{plt}\PY{o}{.}\PY{n}{title}\PY{p}{(}\PY{l+s+sa}{f}\PY{l+s+s2}{\PYZdq{}}\PY{l+s+s2}{Change in M Due to Nonlinearity vs B for dk = 0.45 / 2 * 10}\PY{l+s+se}{\PYZbs{}u2076}\PY{l+s+s2}{ m}\PY{l+s+se}{\PYZbs{}u207b}\PY{l+s+se}{\PYZbs{}u00b9}\PY{l+s+se}{\PYZbs{}n}\PY{l+s+s2}{for Various dR Values and R = 9 * 10}\PY{l+s+se}{\PYZbs{}u207b}\PY{l+s+se}{\PYZbs{}u2077}\PY{l+s+s2}{ m + dR}\PY{l+s+s2}{\PYZdq{}}\PY{p}{)}
\PY{n}{ax}\PY{o}{.}\PY{n}{set\PYZus{}ylabel}\PY{p}{(}\PY{l+s+s1}{\PYZsq{}}\PY{l+s+s1}{Change in M (dMn) due to nonlinearity (10}\PY{l+s+se}{\PYZbs{}u2074}\PY{l+s+s1}{ m}\PY{l+s+se}{\PYZbs{}u207b}\PY{l+s+se}{\PYZbs{}u00b9}\PY{l+s+s1}{)}\PY{l+s+s1}{\PYZsq{}}\PY{p}{)}
\PY{n}{ax}\PY{o}{.}\PY{n}{set\PYZus{}xlabel}\PY{p}{(}\PY{l+s+s2}{\PYZdq{}}\PY{l+s+s2}{Variation (dB) of magnetic field strength B (B = 0.5 + dB * B\PYZus{}max)}\PY{l+s+s2}{\PYZdq{}}\PY{p}{)}

\PY{k}{for} \PY{n}{i} \PY{o+ow}{in} \PY{n+nb}{range}\PY{p}{(}\PY{n}{ilim}\PY{p}{)}\PY{p}{:}
    \PY{n}{ax}\PY{o}{.}\PY{n}{plot}\PY{p}{(}\PY{n}{B}\PY{p}{,} \PY{n}{dMn\PYZus{}arr}\PY{p}{[}\PY{n}{Ri}\PY{p}{[}\PY{n}{i}\PY{p}{]}\PY{p}{,}\PY{p}{:}\PY{p}{]}\PY{o}{/}\PY{l+m+mi}{10000}\PY{p}{,} \PY{n}{label} \PY{o}{=} \PY{l+s+sa}{f}\PY{l+s+s2}{\PYZdq{}}\PY{l+s+s2}{dR = }\PY{l+s+si}{\PYZob{}}\PY{n}{R}\PY{p}{[}\PY{n}{Ri}\PY{p}{[}\PY{n}{i}\PY{p}{]}\PY{p}{]}\PY{l+s+si}{:}\PY{l+s+s2}{.4f}\PY{l+s+si}{\PYZcb{}}\PY{l+s+s2}{ * 10}\PY{l+s+se}{\PYZbs{}u207b}\PY{l+s+se}{\PYZbs{}u00b9}\PY{l+s+se}{\PYZbs{}u2070}\PY{l+s+s2}{ m}\PY{l+s+s2}{\PYZdq{}}\PY{p}{)}

\PY{n}{plt}\PY{o}{.}\PY{n}{axline}\PY{p}{(}\PY{p}{(}\PY{n}{find\PYZus{}dB\PYZus{}B\PYZus{}osc}\PY{p}{(}\PY{l+m+mi}{4}\PY{p}{)}\PY{p}{,}\PY{o}{\PYZhy{}}\PY{l+m+mf}{0.5}\PY{p}{)}\PY{p}{,} \PY{p}{(}\PY{n}{find\PYZus{}dB\PYZus{}B\PYZus{}osc}\PY{p}{(}\PY{l+m+mi}{4}\PY{p}{)}\PY{p}{,}\PY{o}{\PYZhy{}}\PY{l+m+mf}{0.6}\PY{p}{)}\PY{p}{,} \PY{n}{color}\PY{o}{=}\PY{l+s+s1}{\PYZsq{}}\PY{l+s+s1}{k}\PY{l+s+s1}{\PYZsq{}}\PY{p}{,} \PY{n}{label} \PY{o}{=} \PY{l+s+s2}{\PYZdq{}}\PY{l+s+s2}{period boundary}\PY{l+s+s2}{\PYZdq{}}\PY{p}{)}
\PY{n}{plt}\PY{o}{.}\PY{n}{axline}\PY{p}{(}\PY{p}{(}\PY{n}{find\PYZus{}dB\PYZus{}B\PYZus{}osc}\PY{p}{(}\PY{l+m+mf}{4.5}\PY{p}{)}\PY{p}{,}\PY{o}{\PYZhy{}}\PY{l+m+mf}{0.5}\PY{p}{)}\PY{p}{,} \PY{p}{(}\PY{n}{find\PYZus{}dB\PYZus{}B\PYZus{}osc}\PY{p}{(}\PY{l+m+mf}{4.5}\PY{p}{)}\PY{p}{,}\PY{o}{\PYZhy{}}\PY{l+m+mf}{0.6}\PY{p}{)}\PY{p}{,} \PY{n}{color}\PY{o}{=}\PY{l+s+s1}{\PYZsq{}}\PY{l+s+s1}{k}\PY{l+s+s1}{\PYZsq{}}\PY{p}{)}

\PY{n}{plt}\PY{o}{.}\PY{n}{legend}\PY{p}{(}\PY{n}{loc}\PY{o}{=}\PY{l+s+s1}{\PYZsq{}}\PY{l+s+s1}{center left}\PY{l+s+s1}{\PYZsq{}}\PY{p}{,} \PY{n}{bbox\PYZus{}to\PYZus{}anchor}\PY{o}{=}\PY{p}{(}\PY{l+m+mi}{1}\PY{p}{,} \PY{l+m+mf}{0.5}\PY{p}{)}\PY{p}{)}

\PY{n}{plt}\PY{o}{.}\PY{n}{show}\PY{p}{(}\PY{p}{)}
\end{Verbatim}
\end{tcolorbox}

    \begin{center}
    \adjustimage{max size={0.9\linewidth}{0.9\paperheight}}{figs/bvp_kM_en/output_157_0.png}
    \end{center}
    { \hspace*{\fill} \\}
    
    The error appears around our phase shift, but the pattern appears to
disappear around our 3/4 period.

    \hypertarget{dmn-vs-r}{%
\section{\texorpdfstring{\(dMn\) vs \(R\)}{dMn vs R}}\label{dmn-vs-r}}

    The final pattern is for the fast oscillations of the change of slow
oscillation wave number. We know that our \(B\) nodes are at indexes 0
and 62.

    \begin{tcolorbox}[breakable, size=fbox, boxrule=1pt, pad at break*=1mm,colback=cellbackground, colframe=cellborder]
\prompt{In}{incolor}{ }{\boxspacing}
\begin{Verbatim}[commandchars=\\\{\}]
\PY{c+c1}{\PYZsh{} Again, choose the indexes beforehand.}

\PY{n}{Bi} \PY{o}{=} \PY{p}{[}\PY{l+m+mi}{0}\PY{p}{,}\PY{l+m+mi}{62}\PY{p}{]}

\PY{n}{ilim} \PY{o}{=} \PY{n+nb}{len}\PY{p}{(}\PY{n}{Bi}\PY{p}{)}

\PY{n}{fig}\PY{p}{,} \PY{n}{ax} \PY{o}{=} \PY{n}{plt}\PY{o}{.}\PY{n}{subplots}\PY{p}{(}\PY{p}{)}
\PY{n}{plt}\PY{o}{.}\PY{n}{title}\PY{p}{(}\PY{l+s+sa}{f}\PY{l+s+s2}{\PYZdq{}}\PY{l+s+s2}{Plot of the Change in M Due to Nonlinearity vs R}\PY{l+s+se}{\PYZbs{}n}\PY{l+s+s2}{for dk = 0.45 / 2 * 10}\PY{l+s+se}{\PYZbs{}u2076}\PY{l+s+s2}{ m}\PY{l+s+se}{\PYZbs{}u207b}\PY{l+s+se}{\PYZbs{}u00b9}\PY{l+s+s2}{ for various B values}\PY{l+s+s2}{\PYZdq{}}\PY{p}{)}
\PY{n}{ax}\PY{o}{.}\PY{n}{set\PYZus{}ylabel}\PY{p}{(}\PY{l+s+s1}{\PYZsq{}}\PY{l+s+s1}{Change in M (dMn) due to nonlinearity (10}\PY{l+s+se}{\PYZbs{}u2074}\PY{l+s+s1}{ m}\PY{l+s+se}{\PYZbs{}u207b}\PY{l+s+se}{\PYZbs{}u00b9}\PY{l+s+s1}{)}\PY{l+s+s1}{\PYZsq{}}\PY{p}{)}
\PY{n}{ax}\PY{o}{.}\PY{n}{set\PYZus{}xlabel}\PY{p}{(}\PY{l+s+s2}{\PYZdq{}}\PY{l+s+s2}{Change dr in radius R of ring (R = 9 * 10}\PY{l+s+se}{\PYZbs{}u207b}\PY{l+s+se}{\PYZbs{}u2077}\PY{l+s+s2}{ m + dr * 10}\PY{l+s+se}{\PYZbs{}u207b}\PY{l+s+se}{\PYZbs{}u00b9}\PY{l+s+se}{\PYZbs{}u2070}\PY{l+s+s2}{ m)}\PY{l+s+s2}{\PYZdq{}}\PY{p}{)}

\PY{k}{for} \PY{n}{i} \PY{o+ow}{in} \PY{n+nb}{range}\PY{p}{(}\PY{n}{ilim}\PY{p}{)}\PY{p}{:}
    \PY{n}{ax}\PY{o}{.}\PY{n}{plot}\PY{p}{(}\PY{n}{R}\PY{p}{,} \PY{n}{dMn\PYZus{}arr}\PY{p}{[}\PY{p}{:}\PY{p}{,} \PY{n}{Bi}\PY{p}{[}\PY{n}{i}\PY{p}{]}\PY{p}{]}\PY{o}{/}\PY{l+m+mi}{10000}\PY{p}{,} \PY{n}{label} \PY{o}{=} \PY{l+s+sa}{f}\PY{l+s+s2}{\PYZdq{}}\PY{l+s+s2}{B = }\PY{l+s+si}{\PYZob{}}\PY{n}{B0}\PY{o}{+}\PY{n}{B}\PY{p}{[}\PY{n}{Bi}\PY{p}{[}\PY{n}{i}\PY{p}{]}\PY{p}{]}\PY{l+s+si}{:}\PY{l+s+s2}{.4f}\PY{l+s+si}{\PYZcb{}}\PY{l+s+s2}{ * B\PYZus{}max}\PY{l+s+s2}{\PYZdq{}}\PY{p}{)}

\PY{n}{plt}\PY{o}{.}\PY{n}{axline}\PY{p}{(}\PY{p}{(}\PY{n}{find\PYZus{}dR\PYZus{}k\PYZus{}osc}\PY{p}{(}\PY{l+m+mi}{1}\PY{p}{)}\PY{p}{,}\PY{l+m+mi}{0}\PY{p}{)}\PY{p}{,} \PY{p}{(}\PY{n}{find\PYZus{}dR\PYZus{}k\PYZus{}osc}\PY{p}{(}\PY{l+m+mi}{1}\PY{p}{)}\PY{p}{,}\PY{o}{\PYZhy{}}\PY{l+m+mf}{0.01}\PY{p}{)}\PY{p}{,} \PY{n}{color}\PY{o}{=}\PY{l+s+s1}{\PYZsq{}}\PY{l+s+s1}{k}\PY{l+s+s1}{\PYZsq{}}\PY{p}{,} \PY{n}{label} \PY{o}{=} \PY{l+s+s2}{\PYZdq{}}\PY{l+s+s2}{period boundary}\PY{l+s+s2}{\PYZdq{}}\PY{p}{)}
\PY{n}{plt}\PY{o}{.}\PY{n}{axline}\PY{p}{(}\PY{p}{(}\PY{n}{find\PYZus{}dR\PYZus{}k\PYZus{}osc}\PY{p}{(}\PY{l+m+mi}{2}\PY{p}{)}\PY{p}{,}\PY{l+m+mi}{0}\PY{p}{)}\PY{p}{,} \PY{p}{(}\PY{n}{find\PYZus{}dR\PYZus{}k\PYZus{}osc}\PY{p}{(}\PY{l+m+mi}{2}\PY{p}{)}\PY{p}{,}\PY{o}{\PYZhy{}}\PY{l+m+mf}{0.01}\PY{p}{)}\PY{p}{,} \PY{n}{color}\PY{o}{=}\PY{l+s+s1}{\PYZsq{}}\PY{l+s+s1}{k}\PY{l+s+s1}{\PYZsq{}}\PY{p}{)}

\PY{n}{plt}\PY{o}{.}\PY{n}{legend}\PY{p}{(}\PY{n}{loc}\PY{o}{=}\PY{l+s+s1}{\PYZsq{}}\PY{l+s+s1}{center left}\PY{l+s+s1}{\PYZsq{}}\PY{p}{,} \PY{n}{bbox\PYZus{}to\PYZus{}anchor}\PY{o}{=}\PY{p}{(}\PY{l+m+mi}{0}\PY{p}{,} \PY{l+m+mf}{0.3}\PY{p}{)}\PY{p}{)}

\PY{n}{plt}\PY{o}{.}\PY{n}{show}\PY{p}{(}\PY{p}{)}
\end{Verbatim}
\end{tcolorbox}

    \begin{center}
    \adjustimage{max size={0.9\linewidth}{0.9\paperheight}}{figs/bvp_kM_en/output_161_0.png}
    \end{center}
    { \hspace*{\fill} \\}
    
    On our actual node, \(dMn\) appears to be zero everywhere. As we move
away from this, spikes appear at the boundaries of our oscillation. Our
indexes of note are:

    \begin{tcolorbox}[breakable, size=fbox, boxrule=1pt, pad at break*=1mm,colback=cellbackground, colframe=cellborder]
\prompt{In}{incolor}{103}{\boxspacing}
\begin{Verbatim}[commandchars=\\\{\}]
\PY{n+nb}{print}\PY{p}{(} \PY{n}{find\PYZus{}dB\PYZus{}B\PYZus{}osc}\PY{p}{(}\PY{l+m+mf}{4.25}\PY{p}{)}\PY{o}{/}\PY{n}{B}\PY{p}{[}\PY{o}{\PYZhy{}}\PY{l+m+mi}{1}\PY{p}{]}\PY{o}{*}\PY{l+m+mi}{100}\PY{p}{,} \PY{n}{find\PYZus{}dB\PYZus{}B\PYZus{}osc}\PY{p}{(}\PY{l+m+mf}{4.5}\PY{p}{)}\PY{o}{/}\PY{n}{B}\PY{p}{[}\PY{o}{\PYZhy{}}\PY{l+m+mi}{1}\PY{p}{]}\PY{o}{*}\PY{l+m+mi}{100}\PY{p}{,} \PY{n}{find\PYZus{}dB\PYZus{}B\PYZus{}osc}\PY{p}{(}\PY{l+m+mf}{4.75}\PY{p}{)}\PY{o}{/}\PY{n}{B}\PY{p}{[}\PY{o}{\PYZhy{}}\PY{l+m+mi}{1}\PY{p}{]}\PY{o}{*}\PY{l+m+mi}{100}\PY{p}{)}
\end{Verbatim}
\end{tcolorbox}

    \begin{Verbatim}[commandchars=\\\{\}]
31.6234944196735 62.89546467965435 94.16743493963511
    \end{Verbatim}

    \begin{tcolorbox}[breakable, size=fbox, boxrule=1pt, pad at break*=1mm,colback=cellbackground, colframe=cellborder]
\prompt{In}{incolor}{196}{\boxspacing}
\begin{Verbatim}[commandchars=\\\{\}]
\PY{n}{Bi} \PY{o}{=} \PY{p}{[}\PY{l+m+mi}{0}\PY{p}{,} \PY{l+m+mi}{5}\PY{p}{,} \PY{l+m+mi}{10}\PY{p}{,} \PY{l+m+mi}{15}\PY{p}{,}  \PY{l+m+mi}{20}\PY{p}{,} \PY{l+m+mi}{25}\PY{p}{]}

\PY{n}{ilim} \PY{o}{=} \PY{n+nb}{len}\PY{p}{(}\PY{n}{Bi}\PY{p}{)}

\PY{n}{fig}\PY{p}{,} \PY{n}{ax} \PY{o}{=} \PY{n}{plt}\PY{o}{.}\PY{n}{subplots}\PY{p}{(}\PY{p}{)}
\PY{n}{plt}\PY{o}{.}\PY{n}{title}\PY{p}{(}\PY{l+s+sa}{f}\PY{l+s+s2}{\PYZdq{}}\PY{l+s+s2}{Plot of the Change in M Due to Nonlinearity vs R}\PY{l+s+se}{\PYZbs{}n}\PY{l+s+s2}{for dk = 0.45 / 2 * 10}\PY{l+s+se}{\PYZbs{}u2076}\PY{l+s+s2}{ m}\PY{l+s+se}{\PYZbs{}u207b}\PY{l+s+se}{\PYZbs{}u00b9}\PY{l+s+s2}{ for various B values}\PY{l+s+s2}{\PYZdq{}}\PY{p}{)}
\PY{n}{ax}\PY{o}{.}\PY{n}{set\PYZus{}ylabel}\PY{p}{(}\PY{l+s+s1}{\PYZsq{}}\PY{l+s+s1}{Change in M (dMn) due to nonlinearity (10}\PY{l+s+se}{\PYZbs{}u2074}\PY{l+s+s1}{ m}\PY{l+s+se}{\PYZbs{}u207b}\PY{l+s+se}{\PYZbs{}u00b9}\PY{l+s+s1}{)}\PY{l+s+s1}{\PYZsq{}}\PY{p}{)}
\PY{n}{ax}\PY{o}{.}\PY{n}{set\PYZus{}xlabel}\PY{p}{(}\PY{l+s+s2}{\PYZdq{}}\PY{l+s+s2}{Change dr in radius R of ring (R = 9 * 10}\PY{l+s+se}{\PYZbs{}u207b}\PY{l+s+se}{\PYZbs{}u2077}\PY{l+s+s2}{ m + dr * 10}\PY{l+s+se}{\PYZbs{}u207b}\PY{l+s+se}{\PYZbs{}u00b9}\PY{l+s+se}{\PYZbs{}u2070}\PY{l+s+s2}{ m)}\PY{l+s+s2}{\PYZdq{}}\PY{p}{)}

\PY{k}{for} \PY{n}{i} \PY{o+ow}{in} \PY{n+nb}{range}\PY{p}{(}\PY{n}{ilim}\PY{p}{)}\PY{p}{:}
    \PY{n}{ax}\PY{o}{.}\PY{n}{plot}\PY{p}{(}\PY{n}{R}\PY{p}{,} \PY{n}{dMn\PYZus{}arr}\PY{p}{[}\PY{p}{:}\PY{p}{,} \PY{n}{Bi}\PY{p}{[}\PY{n}{i}\PY{p}{]}\PY{p}{]}\PY{o}{/}\PY{l+m+mi}{10000}\PY{p}{,} \PY{n}{label} \PY{o}{=} \PY{l+s+sa}{f}\PY{l+s+s2}{\PYZdq{}}\PY{l+s+s2}{B = }\PY{l+s+si}{\PYZob{}}\PY{n}{B0}\PY{o}{+}\PY{n}{B}\PY{p}{[}\PY{n}{Bi}\PY{p}{[}\PY{n}{i}\PY{p}{]}\PY{p}{]}\PY{l+s+si}{:}\PY{l+s+s2}{.4f}\PY{l+s+si}{\PYZcb{}}\PY{l+s+s2}{ * B\PYZus{}max}\PY{l+s+s2}{\PYZdq{}}\PY{p}{)}

\PY{n}{plt}\PY{o}{.}\PY{n}{axline}\PY{p}{(}\PY{p}{(}\PY{n}{find\PYZus{}dR\PYZus{}k\PYZus{}osc}\PY{p}{(}\PY{l+m+mi}{1}\PY{p}{)}\PY{p}{,}\PY{l+m+mi}{0}\PY{p}{)}\PY{p}{,} \PY{p}{(}\PY{n}{find\PYZus{}dR\PYZus{}k\PYZus{}osc}\PY{p}{(}\PY{l+m+mi}{1}\PY{p}{)}\PY{p}{,}\PY{o}{\PYZhy{}}\PY{l+m+mf}{0.01}\PY{p}{)}\PY{p}{,} \PY{n}{color}\PY{o}{=}\PY{l+s+s1}{\PYZsq{}}\PY{l+s+s1}{k}\PY{l+s+s1}{\PYZsq{}}\PY{p}{,} \PY{n}{label} \PY{o}{=} \PY{l+s+s2}{\PYZdq{}}\PY{l+s+s2}{period boundary}\PY{l+s+s2}{\PYZdq{}}\PY{p}{)}
\PY{n}{plt}\PY{o}{.}\PY{n}{axline}\PY{p}{(}\PY{p}{(}\PY{n}{find\PYZus{}dR\PYZus{}k\PYZus{}osc}\PY{p}{(}\PY{l+m+mi}{2}\PY{p}{)}\PY{p}{,}\PY{l+m+mi}{0}\PY{p}{)}\PY{p}{,} \PY{p}{(}\PY{n}{find\PYZus{}dR\PYZus{}k\PYZus{}osc}\PY{p}{(}\PY{l+m+mi}{2}\PY{p}{)}\PY{p}{,}\PY{o}{\PYZhy{}}\PY{l+m+mf}{0.01}\PY{p}{)}\PY{p}{,} \PY{n}{color}\PY{o}{=}\PY{l+s+s1}{\PYZsq{}}\PY{l+s+s1}{k}\PY{l+s+s1}{\PYZsq{}}\PY{p}{)}

\PY{n}{plt}\PY{o}{.}\PY{n}{legend}\PY{p}{(}\PY{n}{loc}\PY{o}{=}\PY{l+s+s1}{\PYZsq{}}\PY{l+s+s1}{center left}\PY{l+s+s1}{\PYZsq{}}\PY{p}{,} \PY{n}{bbox\PYZus{}to\PYZus{}anchor}\PY{o}{=}\PY{p}{(}\PY{l+m+mi}{1}\PY{p}{,} \PY{l+m+mf}{0.5}\PY{p}{)}\PY{p}{)}

\PY{n}{plt}\PY{o}{.}\PY{n}{show}\PY{p}{(}\PY{p}{)}
\end{Verbatim}
\end{tcolorbox}

    \begin{center}
    \adjustimage{max size={0.9\linewidth}{0.9\paperheight}}{figs/bvp_kM_en/output_164_0.png}
    \end{center}
    { \hspace*{\fill} \\}
    
    Ok. That is a weird pattern. It rises up quickly and tilts towards the
axis slowly. There is a reflected pattern on the negative axis. That one
has its long curve rotating away from the axis, so that the pattern will
reverse which one is tall and which one is long.

    \begin{tcolorbox}[breakable, size=fbox, boxrule=1pt, pad at break*=1mm,colback=cellbackground, colframe=cellborder]
\prompt{In}{incolor}{197}{\boxspacing}
\begin{Verbatim}[commandchars=\\\{\}]
\PY{n}{Bi} \PY{o}{=} \PY{p}{[}\PY{l+m+mi}{30}\PY{p}{,}\PY{l+m+mi}{31}\PY{p}{,} \PY{l+m+mi}{32}\PY{p}{]}

\PY{n}{ilim} \PY{o}{=} \PY{n+nb}{len}\PY{p}{(}\PY{n}{Bi}\PY{p}{)}

\PY{n}{fig}\PY{p}{,} \PY{n}{ax} \PY{o}{=} \PY{n}{plt}\PY{o}{.}\PY{n}{subplots}\PY{p}{(}\PY{p}{)}
\PY{n}{plt}\PY{o}{.}\PY{n}{title}\PY{p}{(}\PY{l+s+sa}{f}\PY{l+s+s2}{\PYZdq{}}\PY{l+s+s2}{Plot of the Change in M Due to Nonlinearity vs R}\PY{l+s+se}{\PYZbs{}n}\PY{l+s+s2}{for dk = 0.45 / 2 * 10}\PY{l+s+se}{\PYZbs{}u2076}\PY{l+s+s2}{ m}\PY{l+s+se}{\PYZbs{}u207b}\PY{l+s+se}{\PYZbs{}u00b9}\PY{l+s+s2}{ for various B values}\PY{l+s+s2}{\PYZdq{}}\PY{p}{)}
\PY{n}{ax}\PY{o}{.}\PY{n}{set\PYZus{}ylabel}\PY{p}{(}\PY{l+s+s1}{\PYZsq{}}\PY{l+s+s1}{Change in M (dMn) due to nonlinearity (10}\PY{l+s+se}{\PYZbs{}u2074}\PY{l+s+s1}{ m}\PY{l+s+se}{\PYZbs{}u207b}\PY{l+s+se}{\PYZbs{}u00b9}\PY{l+s+s1}{)}\PY{l+s+s1}{\PYZsq{}}\PY{p}{)}
\PY{n}{ax}\PY{o}{.}\PY{n}{set\PYZus{}xlabel}\PY{p}{(}\PY{l+s+s2}{\PYZdq{}}\PY{l+s+s2}{Change dr in radius R of ring (R = 9 * 10}\PY{l+s+se}{\PYZbs{}u207b}\PY{l+s+se}{\PYZbs{}u2077}\PY{l+s+s2}{ m + dr * 10}\PY{l+s+se}{\PYZbs{}u207b}\PY{l+s+se}{\PYZbs{}u00b9}\PY{l+s+se}{\PYZbs{}u2070}\PY{l+s+s2}{ m)}\PY{l+s+s2}{\PYZdq{}}\PY{p}{)}

\PY{k}{for} \PY{n}{i} \PY{o+ow}{in} \PY{n+nb}{range}\PY{p}{(}\PY{n}{ilim}\PY{p}{)}\PY{p}{:}
    \PY{n}{ax}\PY{o}{.}\PY{n}{plot}\PY{p}{(}\PY{n}{R}\PY{p}{,} \PY{n}{dMn\PYZus{}arr}\PY{p}{[}\PY{p}{:}\PY{p}{,} \PY{n}{Bi}\PY{p}{[}\PY{n}{i}\PY{p}{]}\PY{p}{]}\PY{o}{/}\PY{l+m+mi}{10000}\PY{p}{,} \PY{n}{label} \PY{o}{=} \PY{l+s+sa}{f}\PY{l+s+s2}{\PYZdq{}}\PY{l+s+s2}{B = }\PY{l+s+si}{\PYZob{}}\PY{n}{B0}\PY{o}{+}\PY{n}{B}\PY{p}{[}\PY{n}{Bi}\PY{p}{[}\PY{n}{i}\PY{p}{]}\PY{p}{]}\PY{l+s+si}{:}\PY{l+s+s2}{.4f}\PY{l+s+si}{\PYZcb{}}\PY{l+s+s2}{ * B\PYZus{}max}\PY{l+s+s2}{\PYZdq{}}\PY{p}{)}

\PY{n}{plt}\PY{o}{.}\PY{n}{axline}\PY{p}{(}\PY{p}{(}\PY{n}{find\PYZus{}dR\PYZus{}k\PYZus{}osc}\PY{p}{(}\PY{l+m+mi}{1}\PY{p}{)}\PY{p}{,}\PY{l+m+mi}{0}\PY{p}{)}\PY{p}{,} \PY{p}{(}\PY{n}{find\PYZus{}dR\PYZus{}k\PYZus{}osc}\PY{p}{(}\PY{l+m+mi}{1}\PY{p}{)}\PY{p}{,}\PY{o}{\PYZhy{}}\PY{l+m+mf}{0.01}\PY{p}{)}\PY{p}{,} \PY{n}{color}\PY{o}{=}\PY{l+s+s1}{\PYZsq{}}\PY{l+s+s1}{k}\PY{l+s+s1}{\PYZsq{}}\PY{p}{,} \PY{n}{label} \PY{o}{=} \PY{l+s+s2}{\PYZdq{}}\PY{l+s+s2}{period boundary}\PY{l+s+s2}{\PYZdq{}}\PY{p}{)}
\PY{n}{plt}\PY{o}{.}\PY{n}{axline}\PY{p}{(}\PY{p}{(}\PY{n}{find\PYZus{}dR\PYZus{}k\PYZus{}osc}\PY{p}{(}\PY{l+m+mi}{2}\PY{p}{)}\PY{p}{,}\PY{l+m+mi}{0}\PY{p}{)}\PY{p}{,} \PY{p}{(}\PY{n}{find\PYZus{}dR\PYZus{}k\PYZus{}osc}\PY{p}{(}\PY{l+m+mi}{2}\PY{p}{)}\PY{p}{,}\PY{o}{\PYZhy{}}\PY{l+m+mf}{0.01}\PY{p}{)}\PY{p}{,} \PY{n}{color}\PY{o}{=}\PY{l+s+s1}{\PYZsq{}}\PY{l+s+s1}{k}\PY{l+s+s1}{\PYZsq{}}\PY{p}{)}

\PY{n}{plt}\PY{o}{.}\PY{n}{legend}\PY{p}{(}\PY{n}{loc}\PY{o}{=}\PY{l+s+s1}{\PYZsq{}}\PY{l+s+s1}{center left}\PY{l+s+s1}{\PYZsq{}}\PY{p}{,} \PY{n}{bbox\PYZus{}to\PYZus{}anchor}\PY{o}{=}\PY{p}{(}\PY{l+m+mi}{1}\PY{p}{,} \PY{l+m+mf}{0.5}\PY{p}{)}\PY{p}{)}

\PY{n}{plt}\PY{o}{.}\PY{n}{show}\PY{p}{(}\PY{p}{)}
\end{Verbatim}
\end{tcolorbox}

    \begin{center}
    \adjustimage{max size={0.9\linewidth}{0.9\paperheight}}{figs/bvp_kM_en/output_166_0.png}
    \end{center}
    { \hspace*{\fill} \\}
    
    I think the angle is now 45 degrees.

    \begin{tcolorbox}[breakable, size=fbox, boxrule=1pt, pad at break*=1mm,colback=cellbackground, colframe=cellborder]
\prompt{In}{incolor}{198}{\boxspacing}
\begin{Verbatim}[commandchars=\\\{\}]
\PY{c+c1}{\PYZsh{} for easy changing which two plots to use.}

\PY{n}{Bi} \PY{o}{=} \PY{p}{[}\PY{l+m+mi}{35}\PY{p}{,} \PY{l+m+mi}{40}\PY{p}{,} \PY{l+m+mi}{45}\PY{p}{,} \PY{l+m+mi}{50}\PY{p}{,} \PY{l+m+mi}{55}\PY{p}{,} \PY{l+m+mi}{60}\PY{p}{]}

\PY{n}{ilim} \PY{o}{=} \PY{n+nb}{len}\PY{p}{(}\PY{n}{Bi}\PY{p}{)}

\PY{n}{fig}\PY{p}{,} \PY{n}{ax} \PY{o}{=} \PY{n}{plt}\PY{o}{.}\PY{n}{subplots}\PY{p}{(}\PY{p}{)}
\PY{n}{plt}\PY{o}{.}\PY{n}{title}\PY{p}{(}\PY{l+s+sa}{f}\PY{l+s+s2}{\PYZdq{}}\PY{l+s+s2}{Plot of the Change in M Due to Nonlinearity vs R}\PY{l+s+se}{\PYZbs{}n}\PY{l+s+s2}{for dk = 0.45 / 2 * 10}\PY{l+s+se}{\PYZbs{}u2076}\PY{l+s+s2}{ m}\PY{l+s+se}{\PYZbs{}u207b}\PY{l+s+se}{\PYZbs{}u00b9}\PY{l+s+s2}{ for various B values}\PY{l+s+s2}{\PYZdq{}}\PY{p}{)}
\PY{n}{ax}\PY{o}{.}\PY{n}{set\PYZus{}ylabel}\PY{p}{(}\PY{l+s+s1}{\PYZsq{}}\PY{l+s+s1}{Change in M (dMn) due to nonlinearity (10}\PY{l+s+se}{\PYZbs{}u2074}\PY{l+s+s1}{ m}\PY{l+s+se}{\PYZbs{}u207b}\PY{l+s+se}{\PYZbs{}u00b9}\PY{l+s+s1}{)}\PY{l+s+s1}{\PYZsq{}}\PY{p}{)}
\PY{n}{ax}\PY{o}{.}\PY{n}{set\PYZus{}xlabel}\PY{p}{(}\PY{l+s+s2}{\PYZdq{}}\PY{l+s+s2}{Change dr in radius R of ring (R = 9 * 10}\PY{l+s+se}{\PYZbs{}u207b}\PY{l+s+se}{\PYZbs{}u2077}\PY{l+s+s2}{ m + dr * 10}\PY{l+s+se}{\PYZbs{}u207b}\PY{l+s+se}{\PYZbs{}u00b9}\PY{l+s+se}{\PYZbs{}u2070}\PY{l+s+s2}{ m)}\PY{l+s+s2}{\PYZdq{}}\PY{p}{)}

\PY{k}{for} \PY{n}{i} \PY{o+ow}{in} \PY{n+nb}{range}\PY{p}{(}\PY{n}{ilim}\PY{p}{)}\PY{p}{:}
    \PY{n}{ax}\PY{o}{.}\PY{n}{plot}\PY{p}{(}\PY{n}{R}\PY{p}{,} \PY{n}{dMn\PYZus{}arr}\PY{p}{[}\PY{p}{:}\PY{p}{,} \PY{n}{Bi}\PY{p}{[}\PY{n}{i}\PY{p}{]}\PY{p}{]}\PY{o}{/}\PY{l+m+mi}{10000}\PY{p}{,} \PY{n}{label} \PY{o}{=} \PY{l+s+sa}{f}\PY{l+s+s2}{\PYZdq{}}\PY{l+s+s2}{B = }\PY{l+s+si}{\PYZob{}}\PY{n}{B0}\PY{o}{+}\PY{n}{B}\PY{p}{[}\PY{n}{Bi}\PY{p}{[}\PY{n}{i}\PY{p}{]}\PY{p}{]}\PY{l+s+si}{:}\PY{l+s+s2}{.4f}\PY{l+s+si}{\PYZcb{}}\PY{l+s+s2}{ * B\PYZus{}max}\PY{l+s+s2}{\PYZdq{}}\PY{p}{)}

\PY{n}{plt}\PY{o}{.}\PY{n}{axline}\PY{p}{(}\PY{p}{(}\PY{n}{find\PYZus{}dR\PYZus{}k\PYZus{}osc}\PY{p}{(}\PY{l+m+mi}{1}\PY{p}{)}\PY{p}{,}\PY{l+m+mi}{0}\PY{p}{)}\PY{p}{,} \PY{p}{(}\PY{n}{find\PYZus{}dR\PYZus{}k\PYZus{}osc}\PY{p}{(}\PY{l+m+mi}{1}\PY{p}{)}\PY{p}{,}\PY{o}{\PYZhy{}}\PY{l+m+mf}{0.01}\PY{p}{)}\PY{p}{,} \PY{n}{color}\PY{o}{=}\PY{l+s+s1}{\PYZsq{}}\PY{l+s+s1}{k}\PY{l+s+s1}{\PYZsq{}}\PY{p}{,} \PY{n}{label} \PY{o}{=} \PY{l+s+s2}{\PYZdq{}}\PY{l+s+s2}{period boundary}\PY{l+s+s2}{\PYZdq{}}\PY{p}{)}
\PY{n}{plt}\PY{o}{.}\PY{n}{axline}\PY{p}{(}\PY{p}{(}\PY{n}{find\PYZus{}dR\PYZus{}k\PYZus{}osc}\PY{p}{(}\PY{l+m+mi}{2}\PY{p}{)}\PY{p}{,}\PY{l+m+mi}{0}\PY{p}{)}\PY{p}{,} \PY{p}{(}\PY{n}{find\PYZus{}dR\PYZus{}k\PYZus{}osc}\PY{p}{(}\PY{l+m+mi}{2}\PY{p}{)}\PY{p}{,}\PY{o}{\PYZhy{}}\PY{l+m+mf}{0.01}\PY{p}{)}\PY{p}{,} \PY{n}{color}\PY{o}{=}\PY{l+s+s1}{\PYZsq{}}\PY{l+s+s1}{k}\PY{l+s+s1}{\PYZsq{}}\PY{p}{)}

\PY{n}{plt}\PY{o}{.}\PY{n}{legend}\PY{p}{(}\PY{n}{loc}\PY{o}{=}\PY{l+s+s1}{\PYZsq{}}\PY{l+s+s1}{center left}\PY{l+s+s1}{\PYZsq{}}\PY{p}{,} \PY{n}{bbox\PYZus{}to\PYZus{}anchor}\PY{o}{=}\PY{p}{(}\PY{l+m+mi}{1}\PY{p}{,} \PY{l+m+mf}{0.5}\PY{p}{)}\PY{p}{)}

\PY{n}{plt}\PY{o}{.}\PY{n}{show}\PY{p}{(}\PY{p}{)}
\end{Verbatim}
\end{tcolorbox}

    \begin{center}
    \adjustimage{max size={0.9\linewidth}{0.9\paperheight}}{figs/bvp_kM_en/output_168_0.png}
    \end{center}
    { \hspace*{\fill} \\}
    
    \begin{tcolorbox}[breakable, size=fbox, boxrule=1pt, pad at break*=1mm,colback=cellbackground, colframe=cellborder]
\prompt{In}{incolor}{199}{\boxspacing}
\begin{Verbatim}[commandchars=\\\{\}]
\PY{c+c1}{\PYZsh{} for easy changing which two plots to use.}

\PY{n}{Bi} \PY{o}{=} \PY{p}{[}\PY{l+m+mi}{60}\PY{p}{,} \PY{l+m+mi}{61}\PY{p}{,} \PY{l+m+mi}{62}\PY{p}{,} \PY{l+m+mi}{63}\PY{p}{]}

\PY{n}{ilim} \PY{o}{=} \PY{n+nb}{len}\PY{p}{(}\PY{n}{Bi}\PY{p}{)}

\PY{n}{fig}\PY{p}{,} \PY{n}{ax} \PY{o}{=} \PY{n}{plt}\PY{o}{.}\PY{n}{subplots}\PY{p}{(}\PY{p}{)}
\PY{n}{plt}\PY{o}{.}\PY{n}{title}\PY{p}{(}\PY{l+s+sa}{f}\PY{l+s+s2}{\PYZdq{}}\PY{l+s+s2}{Plot of the Change in M Due to Nonlinearity vs R}\PY{l+s+se}{\PYZbs{}n}\PY{l+s+s2}{for dk = 0.45 / 2 * 10}\PY{l+s+se}{\PYZbs{}u2076}\PY{l+s+s2}{ m}\PY{l+s+se}{\PYZbs{}u207b}\PY{l+s+se}{\PYZbs{}u00b9}\PY{l+s+s2}{ for various B values}\PY{l+s+s2}{\PYZdq{}}\PY{p}{)}
\PY{n}{ax}\PY{o}{.}\PY{n}{set\PYZus{}ylabel}\PY{p}{(}\PY{l+s+s1}{\PYZsq{}}\PY{l+s+s1}{Change in M (dMn) due to nonlinearity (10}\PY{l+s+se}{\PYZbs{}u2074}\PY{l+s+s1}{ m}\PY{l+s+se}{\PYZbs{}u207b}\PY{l+s+se}{\PYZbs{}u00b9}\PY{l+s+s1}{)}\PY{l+s+s1}{\PYZsq{}}\PY{p}{)}
\PY{n}{ax}\PY{o}{.}\PY{n}{set\PYZus{}xlabel}\PY{p}{(}\PY{l+s+s2}{\PYZdq{}}\PY{l+s+s2}{Change dr in radius R of ring (R = 9 * 10}\PY{l+s+se}{\PYZbs{}u207b}\PY{l+s+se}{\PYZbs{}u2077}\PY{l+s+s2}{ m + dr * 10}\PY{l+s+se}{\PYZbs{}u207b}\PY{l+s+se}{\PYZbs{}u00b9}\PY{l+s+se}{\PYZbs{}u2070}\PY{l+s+s2}{ m)}\PY{l+s+s2}{\PYZdq{}}\PY{p}{)}

\PY{k}{for} \PY{n}{i} \PY{o+ow}{in} \PY{n+nb}{range}\PY{p}{(}\PY{n}{ilim}\PY{p}{)}\PY{p}{:}
    \PY{n}{ax}\PY{o}{.}\PY{n}{plot}\PY{p}{(}\PY{n}{R}\PY{p}{,} \PY{n}{dMn\PYZus{}arr}\PY{p}{[}\PY{p}{:}\PY{p}{,} \PY{n}{Bi}\PY{p}{[}\PY{n}{i}\PY{p}{]}\PY{p}{]}\PY{o}{/}\PY{l+m+mi}{10000}\PY{p}{,} \PY{n}{label} \PY{o}{=} \PY{l+s+sa}{f}\PY{l+s+s2}{\PYZdq{}}\PY{l+s+s2}{B = }\PY{l+s+si}{\PYZob{}}\PY{n}{B0}\PY{o}{+}\PY{n}{B}\PY{p}{[}\PY{n}{Bi}\PY{p}{[}\PY{n}{i}\PY{p}{]}\PY{p}{]}\PY{l+s+si}{:}\PY{l+s+s2}{.4f}\PY{l+s+si}{\PYZcb{}}\PY{l+s+s2}{ * B\PYZus{}max}\PY{l+s+s2}{\PYZdq{}}\PY{p}{)}

\PY{n}{plt}\PY{o}{.}\PY{n}{axline}\PY{p}{(}\PY{p}{(}\PY{n}{find\PYZus{}dR\PYZus{}k\PYZus{}osc}\PY{p}{(}\PY{l+m+mi}{1}\PY{p}{)}\PY{p}{,}\PY{l+m+mi}{0}\PY{p}{)}\PY{p}{,} \PY{p}{(}\PY{n}{find\PYZus{}dR\PYZus{}k\PYZus{}osc}\PY{p}{(}\PY{l+m+mi}{1}\PY{p}{)}\PY{p}{,}\PY{o}{\PYZhy{}}\PY{l+m+mf}{0.01}\PY{p}{)}\PY{p}{,} \PY{n}{color}\PY{o}{=}\PY{l+s+s1}{\PYZsq{}}\PY{l+s+s1}{k}\PY{l+s+s1}{\PYZsq{}}\PY{p}{,} \PY{n}{label} \PY{o}{=} \PY{l+s+s2}{\PYZdq{}}\PY{l+s+s2}{period boundary}\PY{l+s+s2}{\PYZdq{}}\PY{p}{)}
\PY{n}{plt}\PY{o}{.}\PY{n}{axline}\PY{p}{(}\PY{p}{(}\PY{n}{find\PYZus{}dR\PYZus{}k\PYZus{}osc}\PY{p}{(}\PY{l+m+mi}{2}\PY{p}{)}\PY{p}{,}\PY{l+m+mi}{0}\PY{p}{)}\PY{p}{,} \PY{p}{(}\PY{n}{find\PYZus{}dR\PYZus{}k\PYZus{}osc}\PY{p}{(}\PY{l+m+mi}{2}\PY{p}{)}\PY{p}{,}\PY{o}{\PYZhy{}}\PY{l+m+mf}{0.01}\PY{p}{)}\PY{p}{,} \PY{n}{color}\PY{o}{=}\PY{l+s+s1}{\PYZsq{}}\PY{l+s+s1}{k}\PY{l+s+s1}{\PYZsq{}}\PY{p}{)}

\PY{n}{plt}\PY{o}{.}\PY{n}{legend}\PY{p}{(}\PY{n}{loc}\PY{o}{=}\PY{l+s+s1}{\PYZsq{}}\PY{l+s+s1}{center left}\PY{l+s+s1}{\PYZsq{}}\PY{p}{,} \PY{n}{bbox\PYZus{}to\PYZus{}anchor}\PY{o}{=}\PY{p}{(}\PY{l+m+mi}{1}\PY{p}{,} \PY{l+m+mf}{0.5}\PY{p}{)}\PY{p}{)}

\PY{n}{plt}\PY{o}{.}\PY{n}{show}\PY{p}{(}\PY{p}{)}
\end{Verbatim}
\end{tcolorbox}

    \begin{center}
    \adjustimage{max size={0.9\linewidth}{0.9\paperheight}}{figs/bvp_kM_en/output_169_0.png}
    \end{center}
    { \hspace*{\fill} \\}
    
    After it crosses the \(R\) axis, it starts again, but mirrored across
the axis. The red and blue curves here are the mirrored curves.

    \hypertarget{full-b-range}{%
\section{\texorpdfstring{Full \(B\)
Range}{Full B Range}}\label{full-b-range}}

    Now lets check for \(B\) dependence on the shape of these oscillations.
This data is stored in table 3.

    \begin{tcolorbox}[breakable, size=fbox, boxrule=1pt, pad at break*=1mm,colback=cellbackground, colframe=cellborder]
\prompt{In}{incolor}{200}{\boxspacing}
\begin{Verbatim}[commandchars=\\\{\}]
\PY{n}{g} \PY{o}{=} \PY{n}{sns}\PY{o}{.}\PY{n}{PairGrid}\PY{p}{(}\PY{n}{tbl\PYZus{}red}\PY{p}{[}\PY{l+m+mi}{3}\PY{p}{]}\PY{p}{,} \PY{n}{diag\PYZus{}sharey}\PY{o}{=}\PY{k+kc}{False}\PY{p}{,} \PY{n}{corner}\PY{o}{=}\PY{k+kc}{True}\PY{p}{)}

\PY{n}{g}\PY{o}{.}\PY{n}{map\PYZus{}diag}\PY{p}{(}\PY{n}{sns}\PY{o}{.}\PY{n}{kdeplot}\PY{p}{)}
\PY{n}{g}\PY{o}{.}\PY{n}{map\PYZus{}lower}\PY{p}{(}\PY{n}{sns}\PY{o}{.}\PY{n}{scatterplot}\PY{p}{)}
\end{Verbatim}
\end{tcolorbox}

            \begin{tcolorbox}[breakable, size=fbox, boxrule=.5pt, pad at break*=1mm, opacityfill=0]
\prompt{Out}{outcolor}{200}{\boxspacing}
\begin{Verbatim}[commandchars=\\\{\}]
<seaborn.axisgrid.PairGrid at 0x1d1a0214680>
\end{Verbatim}
\end{tcolorbox}
        
    \begin{center}
    \adjustimage{max size={0.9\linewidth}{0.9\paperheight}}{figs/bvp_kM_en/output_173_1.png}
    \end{center}
    { \hspace*{\fill} \\}
    
    The periodic nature of our dataset is obvious from this plot.

    \begin{tcolorbox}[breakable, size=fbox, boxrule=1pt, pad at break*=1mm,colback=cellbackground, colframe=cellborder]
\prompt{In}{incolor}{201}{\boxspacing}
\begin{Verbatim}[commandchars=\\\{\}]
\PY{n}{gind} \PY{o}{=} \PY{l+m+mi}{3}
\end{Verbatim}
\end{tcolorbox}

    \begin{tcolorbox}[breakable, size=fbox, boxrule=1pt, pad at break*=1mm,colback=cellbackground, colframe=cellborder]
\prompt{In}{incolor}{202}{\boxspacing}
\begin{Verbatim}[commandchars=\\\{\}]
\PY{k}{del} \PY{n}{dkn\PYZus{}arr}\PY{p}{,} \PY{n}{dMn\PYZus{}arr}

\PY{n}{dkn\PYZus{}arr} \PY{o}{=} \PY{n}{np}\PY{o}{.}\PY{n}{zeros}\PY{p}{(}\PY{p}{(}\PY{n}{gs\PYZus{}r2}\PY{p}{,} \PY{n}{gs\PYZus{}b2}\PY{p}{)}\PY{p}{)}
\PY{n}{dMn\PYZus{}arr} \PY{o}{=} \PY{n}{np}\PY{o}{.}\PY{n}{zeros}\PY{p}{(}\PY{p}{(}\PY{n}{gs\PYZus{}r2}\PY{p}{,} \PY{n}{gs\PYZus{}b2}\PY{p}{)}\PY{p}{)}

\PY{k}{for} \PY{n}{r}\PY{p}{,} \PY{n}{rr} \PY{o+ow}{in} \PY{n+nb}{enumerate}\PY{p}{(}\PY{n}{R\PYZus{}g}\PY{p}{[}\PY{n}{gind}\PY{p}{]}\PY{p}{)}\PY{p}{:}
    \PY{k}{for} \PY{n}{b}\PY{p}{,} \PY{n}{bb} \PY{o+ow}{in} \PY{n+nb}{enumerate}\PY{p}{(}\PY{n}{B\PYZus{}g}\PY{p}{[}\PY{n}{gind}\PY{p}{]}\PY{p}{)}\PY{p}{:}
        \PY{n}{tbl\PYZus{}sel} \PY{o}{=} \PY{n}{tbl\PYZus{}red}\PY{p}{[}\PY{n}{gind}\PY{p}{]}\PY{p}{[}\PY{p}{(}\PY{n}{tbl\PYZus{}red}\PY{p}{[}\PY{n}{gind}\PY{p}{]}\PY{p}{[}\PY{l+s+s2}{\PYZdq{}}\PY{l+s+s2}{R}\PY{l+s+s2}{\PYZdq{}}\PY{p}{]}\PY{o}{==}\PY{n}{rr}\PY{p}{)} \PY{o}{\PYZam{}}\PY{p}{(}\PY{n}{tbl\PYZus{}red}\PY{p}{[}\PY{n}{gind}\PY{p}{]}\PY{p}{[}\PY{l+s+s2}{\PYZdq{}}\PY{l+s+s2}{B}\PY{l+s+s2}{\PYZdq{}}\PY{p}{]}\PY{o}{==}\PY{n}{bb}\PY{p}{)}\PY{p}{]}
        \PY{n}{dkn\PYZus{}arr}\PY{p}{[}\PY{n}{r}\PY{p}{,}\PY{n}{b}\PY{p}{]} \PY{o}{=} \PY{n}{tbl\PYZus{}sel}\PY{p}{[}\PY{l+s+s2}{\PYZdq{}}\PY{l+s+s2}{dkn}\PY{l+s+s2}{\PYZdq{}}\PY{p}{]}\PY{o}{.}\PY{n}{to\PYZus{}numpy}\PY{p}{(}\PY{p}{)}\PY{p}{[}\PY{l+m+mi}{0}\PY{p}{]}
        \PY{n}{dMn\PYZus{}arr}\PY{p}{[}\PY{n}{r}\PY{p}{,}\PY{n}{b}\PY{p}{]} \PY{o}{=} \PY{n}{tbl\PYZus{}sel}\PY{p}{[}\PY{l+s+s2}{\PYZdq{}}\PY{l+s+s2}{dMn}\PY{l+s+s2}{\PYZdq{}}\PY{p}{]}\PY{o}{.}\PY{n}{to\PYZus{}numpy}\PY{p}{(}\PY{p}{)}\PY{p}{[}\PY{l+m+mi}{0}\PY{p}{]}
\end{Verbatim}
\end{tcolorbox}

    \begin{tcolorbox}[breakable, size=fbox, boxrule=1pt, pad at break*=1mm,colback=cellbackground, colframe=cellborder]
\prompt{In}{incolor}{203}{\boxspacing}
\begin{Verbatim}[commandchars=\\\{\}]
\PY{c+c1}{\PYZsh{}B and B will be rescaled to show only the difference from our }
\PY{n}{scale\PYZus{}R} \PY{o}{=} \PY{l+m+mi}{1}\PY{o}{/}\PY{l+m+mi}{10000}
\PY{n}{R0} \PY{o}{=} \PY{l+m+mf}{0.9}
\PY{c+c1}{\PYZsh{}B0 = 0.5}
\PY{n}{dk} \PY{o}{=} \PY{n}{dk\PYZus{}g}\PY{p}{[}\PY{n}{gind}\PY{p}{]}\PY{p}{[}\PY{l+m+mi}{0}\PY{p}{]}

\PY{n}{R} \PY{o}{=} \PY{p}{(}\PY{n}{R\PYZus{}g}\PY{p}{[}\PY{n}{gind}\PY{p}{]} \PY{o}{\PYZhy{}} \PY{n}{R0}\PY{p}{)}\PY{o}{/}\PY{n}{scale\PYZus{}R}
\PY{n}{B} \PY{o}{=} \PY{n}{B\PYZus{}g}\PY{p}{[}\PY{n}{gind}\PY{p}{]} \PY{c+c1}{\PYZsh{}\PYZhy{} B0}
\end{Verbatim}
\end{tcolorbox}

    \begin{tcolorbox}[breakable, size=fbox, boxrule=1pt, pad at break*=1mm,colback=cellbackground, colframe=cellborder]
\prompt{In}{incolor}{204}{\boxspacing}
\begin{Verbatim}[commandchars=\\\{\}]
\PY{c+c1}{\PYZsh{} We will learn here that our oscillations are based on the fractional part of B * R\PYZca{}2 * pi\PYZca{}2}
\PY{c+c1}{\PYZsh{} Here are a few functions for calculating our period boundaries.}

\PY{k}{def} \PY{n+nf}{find\PYZus{}count\PYZus{}B\PYZus{}osc}\PY{p}{(}\PY{n}{dR}\PY{p}{,} \PY{n}{dB}\PY{p}{)}\PY{p}{:}
    \PY{n}{R} \PY{o}{=} \PY{n}{dR} \PY{o}{*} \PY{n}{scale\PYZus{}R}  \PY{o}{+} \PY{n}{R0}
    \PY{n}{B} \PY{o}{=} \PY{n}{dB}
    \PY{n}{count} \PY{o}{=}  \PY{n}{B} \PY{o}{*} \PY{n}{R}\PY{o}{*}\PY{o}{*}\PY{l+m+mi}{2} \PY{o}{*} \PY{n}{pi}\PY{o}{*}\PY{o}{*}\PY{l+m+mi}{2}
    \PY{k}{return} \PY{n}{count}

\PY{k}{def} \PY{n+nf}{find\PYZus{}dB\PYZus{}B\PYZus{}osc}\PY{p}{(}\PY{n}{n}\PY{p}{,} \PY{n}{dR} \PY{o}{=} \PY{l+m+mi}{0}\PY{p}{)}\PY{p}{:}
    \PY{n}{R} \PY{o}{=} \PY{n}{dR} \PY{o}{*} \PY{n}{scale\PYZus{}R}  \PY{o}{+} \PY{n}{R0}
    \PY{c+c1}{\PYZsh{} n = B * R\PYZca{}2 * pi\PYZca{}2}
    \PY{n}{dB} \PY{o}{=} \PY{n}{n} \PY{o}{/} \PY{p}{(}\PY{n}{R}\PY{o}{*}\PY{o}{*}\PY{l+m+mi}{2} \PY{o}{*} \PY{n}{pi}\PY{o}{*}\PY{o}{*}\PY{l+m+mi}{2}\PY{p}{)}
    \PY{k}{return} \PY{n}{dB}
\end{Verbatim}
\end{tcolorbox}

    \begin{tcolorbox}[breakable, size=fbox, boxrule=1pt, pad at break*=1mm,colback=cellbackground, colframe=cellborder]
\prompt{In}{incolor}{205}{\boxspacing}
\begin{Verbatim}[commandchars=\\\{\}]
\PY{n}{intervalB} \PY{o}{=} \PY{p}{(}\PY{n}{B}\PY{p}{[}\PY{o}{\PYZhy{}}\PY{l+m+mi}{1}\PY{p}{]}\PY{o}{\PYZhy{}}\PY{n}{B}\PY{p}{[}\PY{l+m+mi}{0}\PY{p}{]}\PY{p}{)}\PY{o}{/}\PY{l+m+mi}{300}
\PY{n}{find\PYZus{}dB\PYZus{}B\PYZus{}osc}\PY{p}{(}\PY{l+m+mi}{1}\PY{p}{)} \PY{o}{/} \PY{n}{intervalB}\PY{p}{,} \PY{n}{find\PYZus{}dB\PYZus{}B\PYZus{}osc}\PY{p}{(}\PY{l+m+mi}{2}\PY{p}{)} \PY{o}{/} \PY{n}{intervalB}\PY{p}{,} \PY{n}{find\PYZus{}dB\PYZus{}B\PYZus{}osc}\PY{p}{(}\PY{l+m+mi}{3}\PY{p}{)} \PY{o}{/} \PY{n}{intervalB}\PY{p}{,} \PY{n}{find\PYZus{}dB\PYZus{}B\PYZus{}osc}\PY{p}{(}\PY{l+m+mi}{4}\PY{p}{)} \PY{o}{/} \PY{n}{intervalB}
\end{Verbatim}
\end{tcolorbox}

            \begin{tcolorbox}[breakable, size=fbox, boxrule=.5pt, pad at break*=1mm, opacityfill=0]
\prompt{Out}{outcolor}{205}{\boxspacing}
\begin{Verbatim}[commandchars=\\\{\}]
(42.16445428312018, 84.32890856624036, 126.49336284936051, 168.65781713248072)
\end{Verbatim}
\end{tcolorbox}
        
    \begin{tcolorbox}[breakable, size=fbox, boxrule=1pt, pad at break*=1mm,colback=cellbackground, colframe=cellborder]
\prompt{In}{incolor}{207}{\boxspacing}
\begin{Verbatim}[commandchars=\\\{\}]
\PY{c+c1}{\PYZsh{} Again, choose the indicies beforehand.}

\PY{n}{Bi} \PY{o}{=} \PY{p}{[}\PY{l+m+mi}{47}\PY{p}{]}

\PY{n}{ilim} \PY{o}{=} \PY{n+nb}{len}\PY{p}{(}\PY{n}{Bi}\PY{p}{)}

\PY{n}{fig}\PY{p}{,} \PY{n}{ax} \PY{o}{=} \PY{n}{plt}\PY{o}{.}\PY{n}{subplots}\PY{p}{(}\PY{p}{)}
\PY{n}{plt}\PY{o}{.}\PY{n}{title}\PY{p}{(}\PY{l+s+sa}{f}\PY{l+s+s2}{\PYZdq{}}\PY{l+s+s2}{Plot of the Change in k Due to Nonlinearity vs R}\PY{l+s+se}{\PYZbs{}n}\PY{l+s+s2}{for dk = 0.45 / 2 * 10}\PY{l+s+se}{\PYZbs{}u2076}\PY{l+s+s2}{ m}\PY{l+s+se}{\PYZbs{}u207b}\PY{l+s+se}{\PYZbs{}u00b9}\PY{l+s+s2}{ for B = }\PY{l+s+si}{\PYZob{}}\PY{n}{B}\PY{p}{[}\PY{n}{Bi}\PY{p}{[}\PY{n}{i}\PY{p}{]}\PY{p}{]}\PY{l+s+si}{:}\PY{l+s+s2}{.4f}\PY{l+s+si}{\PYZcb{}}\PY{l+s+s2}{ * B\PYZus{}max}\PY{l+s+s2}{\PYZdq{}}\PY{p}{)}
\PY{n}{ax}\PY{o}{.}\PY{n}{set\PYZus{}ylabel}\PY{p}{(}\PY{l+s+s1}{\PYZsq{}}\PY{l+s+s1}{Change in k (dkn) due to nonlinearity (10}\PY{l+s+se}{\PYZbs{}u2075}\PY{l+s+s1}{ m}\PY{l+s+se}{\PYZbs{}u207b}\PY{l+s+se}{\PYZbs{}u00b9}\PY{l+s+s1}{)}\PY{l+s+s1}{\PYZsq{}}\PY{p}{)}
\PY{n}{ax}\PY{o}{.}\PY{n}{set\PYZus{}xlabel}\PY{p}{(}\PY{l+s+s2}{\PYZdq{}}\PY{l+s+s2}{Change dr in radius R of ring (R = 9 * 10}\PY{l+s+se}{\PYZbs{}u207b}\PY{l+s+se}{\PYZbs{}u2077}\PY{l+s+s2}{ m + dr * 10}\PY{l+s+se}{\PYZbs{}u207b}\PY{l+s+se}{\PYZbs{}u00b9}\PY{l+s+se}{\PYZbs{}u2070}\PY{l+s+s2}{ m)}\PY{l+s+s2}{\PYZdq{}}\PY{p}{)}

\PY{k}{for} \PY{n}{i} \PY{o+ow}{in} \PY{n+nb}{range}\PY{p}{(}\PY{n}{ilim}\PY{p}{)}\PY{p}{:}
    \PY{n}{ax}\PY{o}{.}\PY{n}{plot}\PY{p}{(}\PY{n}{R}\PY{p}{,} \PY{n}{dkn\PYZus{}arr}\PY{p}{[}\PY{p}{:}\PY{p}{,} \PY{n}{Bi}\PY{p}{[}\PY{n}{i}\PY{p}{]}\PY{p}{]}\PY{o}{/}\PY{l+m+mi}{100000}\PY{p}{)}
\PY{n}{plt}\PY{o}{.}\PY{n}{axline}\PY{p}{(}\PY{p}{(}\PY{n}{find\PYZus{}dR\PYZus{}k\PYZus{}osc}\PY{p}{(}\PY{l+m+mi}{1}\PY{p}{)}\PY{p}{,}\PY{o}{\PYZhy{}}\PY{l+m+mf}{0.5}\PY{p}{)}\PY{p}{,} \PY{p}{(}\PY{n}{find\PYZus{}dR\PYZus{}k\PYZus{}osc}\PY{p}{(}\PY{l+m+mi}{1}\PY{p}{)}\PY{p}{,}\PY{o}{\PYZhy{}}\PY{l+m+mf}{0.6}\PY{p}{)}\PY{p}{,} \PY{n}{color}\PY{o}{=}\PY{l+s+s1}{\PYZsq{}}\PY{l+s+s1}{k}\PY{l+s+s1}{\PYZsq{}}\PY{p}{)}
\PY{n}{plt}\PY{o}{.}\PY{n}{axline}\PY{p}{(}\PY{p}{(}\PY{n}{find\PYZus{}dR\PYZus{}k\PYZus{}osc}\PY{p}{(}\PY{l+m+mi}{2}\PY{p}{)}\PY{p}{,}\PY{o}{\PYZhy{}}\PY{l+m+mf}{0.5}\PY{p}{)}\PY{p}{,} \PY{p}{(}\PY{n}{find\PYZus{}dR\PYZus{}k\PYZus{}osc}\PY{p}{(}\PY{l+m+mi}{2}\PY{p}{)}\PY{p}{,}\PY{o}{\PYZhy{}}\PY{l+m+mf}{0.6}\PY{p}{)}\PY{p}{,} \PY{n}{color}\PY{o}{=}\PY{l+s+s1}{\PYZsq{}}\PY{l+s+s1}{k}\PY{l+s+s1}{\PYZsq{}}\PY{p}{)}

\PY{n}{plt}\PY{o}{.}\PY{n}{show}\PY{p}{(}\PY{p}{)}
\end{Verbatim}
\end{tcolorbox}

    \begin{center}
    \adjustimage{max size={0.9\linewidth}{0.9\paperheight}}{figs/bvp_kM_en/output_180_0.png}
    \end{center}
    { \hspace*{\fill} \\}
    
    \begin{tcolorbox}[breakable, size=fbox, boxrule=1pt, pad at break*=1mm,colback=cellbackground, colframe=cellborder]
\prompt{In}{incolor}{ }{\boxspacing}
\begin{Verbatim}[commandchars=\\\{\}]
\PY{c+c1}{\PYZsh{} Again, choose the indexes beforehand.}

\PY{n}{Bi} \PY{o}{=} \PY{p}{[}\PY{l+m+mi}{40}\PY{p}{,} \PY{l+m+mi}{45}\PY{p}{,} \PY{l+m+mi}{50}\PY{p}{,} \PY{l+m+mi}{55}\PY{p}{,} \PY{l+m+mi}{60}\PY{p}{,} \PY{l+m+mi}{65}\PY{p}{]}

\PY{n}{ilim} \PY{o}{=} \PY{n+nb}{len}\PY{p}{(}\PY{n}{Bi}\PY{p}{)}

\PY{n}{fig}\PY{p}{,} \PY{n}{ax} \PY{o}{=} \PY{n}{plt}\PY{o}{.}\PY{n}{subplots}\PY{p}{(}\PY{p}{)}
\PY{n}{plt}\PY{o}{.}\PY{n}{title}\PY{p}{(}\PY{l+s+sa}{f}\PY{l+s+s2}{\PYZdq{}}\PY{l+s+s2}{Plot of the Change in k Due to Nonlinearity vs R}\PY{l+s+se}{\PYZbs{}n}\PY{l+s+s2}{for dk = 0.45 / 2 * 10}\PY{l+s+se}{\PYZbs{}u2076}\PY{l+s+s2}{ m}\PY{l+s+se}{\PYZbs{}u207b}\PY{l+s+se}{\PYZbs{}u00b9}\PY{l+s+s2}{ for various B values}\PY{l+s+s2}{\PYZdq{}}\PY{p}{)}
\PY{n}{ax}\PY{o}{.}\PY{n}{set\PYZus{}ylabel}\PY{p}{(}\PY{l+s+s1}{\PYZsq{}}\PY{l+s+s1}{Change in k (dkn) due to nonlinearity (10}\PY{l+s+se}{\PYZbs{}u2075}\PY{l+s+s1}{ m}\PY{l+s+se}{\PYZbs{}u207b}\PY{l+s+se}{\PYZbs{}u00b9}\PY{l+s+s1}{)}\PY{l+s+s1}{\PYZsq{}}\PY{p}{)}
\PY{n}{ax}\PY{o}{.}\PY{n}{set\PYZus{}xlabel}\PY{p}{(}\PY{l+s+s2}{\PYZdq{}}\PY{l+s+s2}{Change dr in radius R of ring (R = 9 * 10}\PY{l+s+se}{\PYZbs{}u207b}\PY{l+s+se}{\PYZbs{}u2077}\PY{l+s+s2}{ m + dr * 10}\PY{l+s+se}{\PYZbs{}u207b}\PY{l+s+se}{\PYZbs{}u00b9}\PY{l+s+se}{\PYZbs{}u2070}\PY{l+s+s2}{ m)}\PY{l+s+s2}{\PYZdq{}}\PY{p}{)}

\PY{k}{for} \PY{n}{i} \PY{o+ow}{in} \PY{n+nb}{range}\PY{p}{(}\PY{n}{ilim}\PY{p}{)}\PY{p}{:}
    \PY{n}{ax}\PY{o}{.}\PY{n}{plot}\PY{p}{(}\PY{n}{R}\PY{p}{,} \PY{n}{dkn\PYZus{}arr}\PY{p}{[}\PY{p}{:}\PY{p}{,} \PY{n}{Bi}\PY{p}{[}\PY{n}{i}\PY{p}{]}\PY{p}{]}\PY{o}{/}\PY{l+m+mi}{100000}\PY{p}{,} \PY{n}{label} \PY{o}{=} \PY{l+s+sa}{f}\PY{l+s+s2}{\PYZdq{}}\PY{l+s+s2}{B = }\PY{l+s+si}{\PYZob{}}\PY{n}{B}\PY{p}{[}\PY{n}{Bi}\PY{p}{[}\PY{n}{i}\PY{p}{]}\PY{p}{]}\PY{l+s+si}{:}\PY{l+s+s2}{.4f}\PY{l+s+si}{\PYZcb{}}\PY{l+s+s2}{ * B\PYZus{}max}\PY{l+s+s2}{\PYZdq{}}\PY{p}{)}

\PY{n}{plt}\PY{o}{.}\PY{n}{axline}\PY{p}{(}\PY{p}{(}\PY{n}{find\PYZus{}dR\PYZus{}k\PYZus{}osc}\PY{p}{(}\PY{l+m+mi}{1}\PY{p}{)}\PY{p}{,}\PY{o}{\PYZhy{}}\PY{l+m+mf}{0.5}\PY{p}{)}\PY{p}{,} \PY{p}{(}\PY{n}{find\PYZus{}dR\PYZus{}k\PYZus{}osc}\PY{p}{(}\PY{l+m+mi}{1}\PY{p}{)}\PY{p}{,}\PY{o}{\PYZhy{}}\PY{l+m+mf}{0.6}\PY{p}{)}\PY{p}{,} \PY{n}{color}\PY{o}{=}\PY{l+s+s1}{\PYZsq{}}\PY{l+s+s1}{k}\PY{l+s+s1}{\PYZsq{}}\PY{p}{,} \PY{n}{label} \PY{o}{=} \PY{l+s+s2}{\PYZdq{}}\PY{l+s+s2}{period boundary}\PY{l+s+s2}{\PYZdq{}}\PY{p}{)}
\PY{n}{plt}\PY{o}{.}\PY{n}{axline}\PY{p}{(}\PY{p}{(}\PY{n}{find\PYZus{}dR\PYZus{}k\PYZus{}osc}\PY{p}{(}\PY{l+m+mi}{2}\PY{p}{)}\PY{p}{,}\PY{o}{\PYZhy{}}\PY{l+m+mf}{0.5}\PY{p}{)}\PY{p}{,} \PY{p}{(}\PY{n}{find\PYZus{}dR\PYZus{}k\PYZus{}osc}\PY{p}{(}\PY{l+m+mi}{2}\PY{p}{)}\PY{p}{,}\PY{o}{\PYZhy{}}\PY{l+m+mf}{0.6}\PY{p}{)}\PY{p}{,} \PY{n}{color}\PY{o}{=}\PY{l+s+s1}{\PYZsq{}}\PY{l+s+s1}{k}\PY{l+s+s1}{\PYZsq{}}\PY{p}{)}

\PY{n}{plt}\PY{o}{.}\PY{n}{legend}\PY{p}{(}\PY{n}{loc}\PY{o}{=}\PY{l+s+s1}{\PYZsq{}}\PY{l+s+s1}{center left}\PY{l+s+s1}{\PYZsq{}}\PY{p}{,} \PY{n}{bbox\PYZus{}to\PYZus{}anchor}\PY{o}{=}\PY{p}{(}\PY{l+m+mi}{1}\PY{p}{,} \PY{l+m+mf}{0.5}\PY{p}{)}\PY{p}{)}

\PY{n}{plt}\PY{o}{.}\PY{n}{show}\PY{p}{(}\PY{p}{)}
\end{Verbatim}
\end{tcolorbox}

    \begin{center}
    \adjustimage{max size={0.9\linewidth}{0.9\paperheight}}{figs/bvp_kM_en/output_181_0.png}
    \end{center}
    { \hspace*{\fill} \\}
    
    \begin{tcolorbox}[breakable, size=fbox, boxrule=1pt, pad at break*=1mm,colback=cellbackground, colframe=cellborder]
\prompt{In}{incolor}{ }{\boxspacing}
\begin{Verbatim}[commandchars=\\\{\}]
\PY{c+c1}{\PYZsh{} Again, choose the indexes beforehand.}

\PY{n}{Bi} \PY{o}{=} \PY{p}{[}\PY{l+m+mi}{0}\PY{p}{,} \PY{l+m+mi}{5}\PY{p}{,} \PY{l+m+mi}{10}\PY{p}{,} \PY{l+m+mi}{20}\PY{p}{,} \PY{l+m+mi}{280}\PY{p}{,} \PY{l+m+mi}{290}\PY{p}{,} \PY{l+m+mi}{299}\PY{p}{]}

\PY{n}{ilim} \PY{o}{=} \PY{n+nb}{len}\PY{p}{(}\PY{n}{Bi}\PY{p}{)}

\PY{n}{fig}\PY{p}{,} \PY{n}{ax} \PY{o}{=} \PY{n}{plt}\PY{o}{.}\PY{n}{subplots}\PY{p}{(}\PY{p}{)}
\PY{n}{plt}\PY{o}{.}\PY{n}{title}\PY{p}{(}\PY{l+s+sa}{f}\PY{l+s+s2}{\PYZdq{}}\PY{l+s+s2}{Plot of the Change in k Due to Nonlinearity vs R}\PY{l+s+se}{\PYZbs{}n}\PY{l+s+s2}{for dk = 0.45 / 2 * 10}\PY{l+s+se}{\PYZbs{}u2076}\PY{l+s+s2}{ m}\PY{l+s+se}{\PYZbs{}u207b}\PY{l+s+se}{\PYZbs{}u00b9}\PY{l+s+s2}{ for various B values}\PY{l+s+s2}{\PYZdq{}}\PY{p}{)}
\PY{n}{ax}\PY{o}{.}\PY{n}{set\PYZus{}ylabel}\PY{p}{(}\PY{l+s+s1}{\PYZsq{}}\PY{l+s+s1}{Change in k (dkn) due to nonlinearity (10}\PY{l+s+se}{\PYZbs{}u2075}\PY{l+s+s1}{ m}\PY{l+s+se}{\PYZbs{}u207b}\PY{l+s+se}{\PYZbs{}u00b9}\PY{l+s+s1}{)}\PY{l+s+s1}{\PYZsq{}}\PY{p}{)}
\PY{n}{ax}\PY{o}{.}\PY{n}{set\PYZus{}xlabel}\PY{p}{(}\PY{l+s+s2}{\PYZdq{}}\PY{l+s+s2}{Change dr in radius R of ring (R = 9 * 10}\PY{l+s+se}{\PYZbs{}u207b}\PY{l+s+se}{\PYZbs{}u2077}\PY{l+s+s2}{ m + dr * 10}\PY{l+s+se}{\PYZbs{}u207b}\PY{l+s+se}{\PYZbs{}u00b9}\PY{l+s+se}{\PYZbs{}u2070}\PY{l+s+s2}{ m)}\PY{l+s+s2}{\PYZdq{}}\PY{p}{)}

\PY{k}{for} \PY{n}{i} \PY{o+ow}{in} \PY{n+nb}{range}\PY{p}{(}\PY{n}{ilim}\PY{p}{)}\PY{p}{:}
    \PY{n}{ax}\PY{o}{.}\PY{n}{plot}\PY{p}{(}\PY{n}{R}\PY{p}{,} \PY{n}{dkn\PYZus{}arr}\PY{p}{[}\PY{p}{:}\PY{p}{,} \PY{n}{Bi}\PY{p}{[}\PY{n}{i}\PY{p}{]}\PY{p}{]}\PY{o}{/}\PY{l+m+mi}{100000}\PY{p}{,} \PY{n}{label} \PY{o}{=} \PY{l+s+sa}{f}\PY{l+s+s2}{\PYZdq{}}\PY{l+s+s2}{B = }\PY{l+s+si}{\PYZob{}}\PY{n}{B}\PY{p}{[}\PY{n}{Bi}\PY{p}{[}\PY{n}{i}\PY{p}{]}\PY{p}{]}\PY{l+s+si}{:}\PY{l+s+s2}{.4f}\PY{l+s+si}{\PYZcb{}}\PY{l+s+s2}{ * B\PYZus{}max}\PY{l+s+s2}{\PYZdq{}}\PY{p}{)}

\PY{n}{plt}\PY{o}{.}\PY{n}{axline}\PY{p}{(}\PY{p}{(}\PY{n}{find\PYZus{}dR\PYZus{}k\PYZus{}osc}\PY{p}{(}\PY{l+m+mi}{1}\PY{p}{)}\PY{p}{,}\PY{o}{\PYZhy{}}\PY{l+m+mf}{0.5}\PY{p}{)}\PY{p}{,} \PY{p}{(}\PY{n}{find\PYZus{}dR\PYZus{}k\PYZus{}osc}\PY{p}{(}\PY{l+m+mi}{1}\PY{p}{)}\PY{p}{,}\PY{o}{\PYZhy{}}\PY{l+m+mf}{0.6}\PY{p}{)}\PY{p}{,} \PY{n}{color}\PY{o}{=}\PY{l+s+s1}{\PYZsq{}}\PY{l+s+s1}{k}\PY{l+s+s1}{\PYZsq{}}\PY{p}{,} \PY{n}{label} \PY{o}{=} \PY{l+s+s2}{\PYZdq{}}\PY{l+s+s2}{period boundary}\PY{l+s+s2}{\PYZdq{}}\PY{p}{)}
\PY{n}{plt}\PY{o}{.}\PY{n}{axline}\PY{p}{(}\PY{p}{(}\PY{n}{find\PYZus{}dR\PYZus{}k\PYZus{}osc}\PY{p}{(}\PY{l+m+mi}{2}\PY{p}{)}\PY{p}{,}\PY{o}{\PYZhy{}}\PY{l+m+mf}{0.5}\PY{p}{)}\PY{p}{,} \PY{p}{(}\PY{n}{find\PYZus{}dR\PYZus{}k\PYZus{}osc}\PY{p}{(}\PY{l+m+mi}{2}\PY{p}{)}\PY{p}{,}\PY{o}{\PYZhy{}}\PY{l+m+mf}{0.6}\PY{p}{)}\PY{p}{,} \PY{n}{color}\PY{o}{=}\PY{l+s+s1}{\PYZsq{}}\PY{l+s+s1}{k}\PY{l+s+s1}{\PYZsq{}}\PY{p}{)}

\PY{n}{plt}\PY{o}{.}\PY{n}{legend}\PY{p}{(}\PY{n}{loc}\PY{o}{=}\PY{l+s+s1}{\PYZsq{}}\PY{l+s+s1}{center left}\PY{l+s+s1}{\PYZsq{}}\PY{p}{,} \PY{n}{bbox\PYZus{}to\PYZus{}anchor}\PY{o}{=}\PY{p}{(}\PY{l+m+mi}{1}\PY{p}{,} \PY{l+m+mf}{0.5}\PY{p}{)}\PY{p}{)}

\PY{n}{plt}\PY{o}{.}\PY{n}{show}\PY{p}{(}\PY{p}{)}
\end{Verbatim}
\end{tcolorbox}

    \begin{center}
    \adjustimage{max size={0.9\linewidth}{0.9\paperheight}}{figs/bvp_kM_en/output_182_0.png}
    \end{center}
    { \hspace*{\fill} \\}
    
    \begin{tcolorbox}[breakable, size=fbox, boxrule=1pt, pad at break*=1mm,colback=cellbackground, colframe=cellborder]
\prompt{In}{incolor}{ }{\boxspacing}
\begin{Verbatim}[commandchars=\\\{\}]
\PY{c+c1}{\PYZsh{} Again, choose the indexes beforehand.}

\PY{n}{Bi} \PY{o}{=} \PY{p}{[}\PY{l+m+mi}{0}\PY{p}{,} \PY{l+m+mi}{5}\PY{p}{,} \PY{l+m+mi}{10}\PY{p}{,} \PY{l+m+mi}{280}\PY{p}{,} \PY{l+m+mi}{290}\PY{p}{,} \PY{l+m+mi}{299}\PY{p}{]}

\PY{n}{ilim} \PY{o}{=} \PY{n+nb}{len}\PY{p}{(}\PY{n}{Bi}\PY{p}{)}

\PY{n}{fig}\PY{p}{,} \PY{n}{ax} \PY{o}{=} \PY{n}{plt}\PY{o}{.}\PY{n}{subplots}\PY{p}{(}\PY{p}{)}
\PY{n}{plt}\PY{o}{.}\PY{n}{title}\PY{p}{(}\PY{l+s+sa}{f}\PY{l+s+s2}{\PYZdq{}}\PY{l+s+s2}{Plot of the Change in M Due to Nonlinearity vs R}\PY{l+s+se}{\PYZbs{}n}\PY{l+s+s2}{for dk = 0.45 / 2 * 10}\PY{l+s+se}{\PYZbs{}u2076}\PY{l+s+s2}{ m}\PY{l+s+se}{\PYZbs{}u207b}\PY{l+s+se}{\PYZbs{}u00b9}\PY{l+s+s2}{ for various B values}\PY{l+s+s2}{\PYZdq{}}\PY{p}{)}
\PY{n}{ax}\PY{o}{.}\PY{n}{set\PYZus{}ylabel}\PY{p}{(}\PY{l+s+s1}{\PYZsq{}}\PY{l+s+s1}{Change in M (dMn) due to nonlinearity (10}\PY{l+s+se}{\PYZbs{}u2075}\PY{l+s+s1}{ m}\PY{l+s+se}{\PYZbs{}u207b}\PY{l+s+se}{\PYZbs{}u00b9}\PY{l+s+s1}{)}\PY{l+s+s1}{\PYZsq{}}\PY{p}{)}
\PY{n}{ax}\PY{o}{.}\PY{n}{set\PYZus{}xlabel}\PY{p}{(}\PY{l+s+s2}{\PYZdq{}}\PY{l+s+s2}{Change dr in radius R of ring (R = 9 * 10}\PY{l+s+se}{\PYZbs{}u207b}\PY{l+s+se}{\PYZbs{}u2077}\PY{l+s+s2}{ m + dr * 10}\PY{l+s+se}{\PYZbs{}u207b}\PY{l+s+se}{\PYZbs{}u00b9}\PY{l+s+se}{\PYZbs{}u2070}\PY{l+s+s2}{ m)}\PY{l+s+s2}{\PYZdq{}}\PY{p}{)}

\PY{k}{for} \PY{n}{i} \PY{o+ow}{in} \PY{n+nb}{range}\PY{p}{(}\PY{n}{ilim}\PY{p}{)}\PY{p}{:}
    \PY{n}{ax}\PY{o}{.}\PY{n}{plot}\PY{p}{(}\PY{n}{R}\PY{p}{,} \PY{n}{dMn\PYZus{}arr}\PY{p}{[}\PY{p}{:}\PY{p}{,} \PY{n}{Bi}\PY{p}{[}\PY{n}{i}\PY{p}{]}\PY{p}{]}\PY{o}{/}\PY{l+m+mi}{100000}\PY{p}{,} \PY{n}{label} \PY{o}{=} \PY{l+s+sa}{f}\PY{l+s+s2}{\PYZdq{}}\PY{l+s+s2}{B = }\PY{l+s+si}{\PYZob{}}\PY{n}{B}\PY{p}{[}\PY{n}{Bi}\PY{p}{[}\PY{n}{i}\PY{p}{]}\PY{p}{]}\PY{l+s+si}{:}\PY{l+s+s2}{.4f}\PY{l+s+si}{\PYZcb{}}\PY{l+s+s2}{ * B\PYZus{}max}\PY{l+s+s2}{\PYZdq{}}\PY{p}{)}

\PY{n}{plt}\PY{o}{.}\PY{n}{axline}\PY{p}{(}\PY{p}{(}\PY{n}{find\PYZus{}dR\PYZus{}k\PYZus{}osc}\PY{p}{(}\PY{l+m+mi}{1}\PY{p}{)}\PY{p}{,}\PY{l+m+mi}{0}\PY{p}{)}\PY{p}{,} \PY{p}{(}\PY{n}{find\PYZus{}dR\PYZus{}k\PYZus{}osc}\PY{p}{(}\PY{l+m+mi}{1}\PY{p}{)}\PY{p}{,}\PY{o}{\PYZhy{}}\PY{l+m+mf}{0.01}\PY{p}{)}\PY{p}{,} \PY{n}{color}\PY{o}{=}\PY{l+s+s1}{\PYZsq{}}\PY{l+s+s1}{k}\PY{l+s+s1}{\PYZsq{}}\PY{p}{,} \PY{n}{label} \PY{o}{=} \PY{l+s+s2}{\PYZdq{}}\PY{l+s+s2}{period boundary}\PY{l+s+s2}{\PYZdq{}}\PY{p}{)}
\PY{n}{plt}\PY{o}{.}\PY{n}{axline}\PY{p}{(}\PY{p}{(}\PY{n}{find\PYZus{}dR\PYZus{}k\PYZus{}osc}\PY{p}{(}\PY{l+m+mi}{2}\PY{p}{)}\PY{p}{,}\PY{l+m+mi}{0}\PY{p}{)}\PY{p}{,} \PY{p}{(}\PY{n}{find\PYZus{}dR\PYZus{}k\PYZus{}osc}\PY{p}{(}\PY{l+m+mi}{2}\PY{p}{)}\PY{p}{,}\PY{o}{\PYZhy{}}\PY{l+m+mf}{0.01}\PY{p}{)}\PY{p}{,} \PY{n}{color}\PY{o}{=}\PY{l+s+s1}{\PYZsq{}}\PY{l+s+s1}{k}\PY{l+s+s1}{\PYZsq{}}\PY{p}{)}

\PY{n}{plt}\PY{o}{.}\PY{n}{legend}\PY{p}{(}\PY{n}{loc}\PY{o}{=}\PY{l+s+s1}{\PYZsq{}}\PY{l+s+s1}{center left}\PY{l+s+s1}{\PYZsq{}}\PY{p}{,} \PY{n}{bbox\PYZus{}to\PYZus{}anchor}\PY{o}{=}\PY{p}{(}\PY{l+m+mi}{1}\PY{p}{,} \PY{l+m+mf}{0.5}\PY{p}{)}\PY{p}{)}

\PY{n}{plt}\PY{o}{.}\PY{n}{show}\PY{p}{(}\PY{p}{)}
\end{Verbatim}
\end{tcolorbox}

    \begin{center}
    \adjustimage{max size={0.9\linewidth}{0.9\paperheight}}{figs/bvp_kM_en/output_183_0.png}
    \end{center}
    { \hspace*{\fill} \\}
    
    No obvious change between indexes at the start and indexes at the end.

    \begin{tcolorbox}[breakable, size=fbox, boxrule=1pt, pad at break*=1mm,colback=cellbackground, colframe=cellborder]
\prompt{In}{incolor}{212}{\boxspacing}
\begin{Verbatim}[commandchars=\\\{\}]
\PY{n}{Ri} \PY{o}{=} \PY{p}{[} \PY{l+m+mi}{2}\PY{p}{,} \PY{l+m+mi}{6}\PY{p}{,} \PY{l+m+mi}{20}\PY{p}{]}

\PY{n}{ilim} \PY{o}{=} \PY{n+nb}{len}\PY{p}{(}\PY{n}{Ri}\PY{p}{)}

\PY{n}{fig}\PY{p}{,} \PY{n}{ax} \PY{o}{=} \PY{n}{plt}\PY{o}{.}\PY{n}{subplots}\PY{p}{(}\PY{p}{)}
\PY{n}{plt}\PY{o}{.}\PY{n}{title}\PY{p}{(}\PY{l+s+sa}{f}\PY{l+s+s2}{\PYZdq{}}\PY{l+s+s2}{Change in k Due to Nonlinearity vs B for dk = 0.45 / 2 * 10}\PY{l+s+se}{\PYZbs{}u2076}\PY{l+s+s2}{ m}\PY{l+s+se}{\PYZbs{}u207b}\PY{l+s+se}{\PYZbs{}u00b9}\PY{l+s+se}{\PYZbs{}n}\PY{l+s+s2}{for Various dR Values and R = 9 * 10}\PY{l+s+se}{\PYZbs{}u207b}\PY{l+s+se}{\PYZbs{}u2077}\PY{l+s+s2}{ m + dR}\PY{l+s+s2}{\PYZdq{}}\PY{p}{)}
\PY{n}{ax}\PY{o}{.}\PY{n}{set\PYZus{}ylabel}\PY{p}{(}\PY{l+s+s1}{\PYZsq{}}\PY{l+s+s1}{Change in k (dkn) due to nonlinearity (10}\PY{l+s+se}{\PYZbs{}u2075}\PY{l+s+s1}{ m}\PY{l+s+se}{\PYZbs{}u207b}\PY{l+s+se}{\PYZbs{}u00b9}\PY{l+s+s1}{)}\PY{l+s+s1}{\PYZsq{}}\PY{p}{)}
\PY{n}{ax}\PY{o}{.}\PY{n}{set\PYZus{}xlabel}\PY{p}{(}\PY{l+s+s2}{\PYZdq{}}\PY{l+s+s2}{Magnetic field strength B (B = B * B\PYZus{}max)}\PY{l+s+s2}{\PYZdq{}}\PY{p}{)}

\PY{k}{for} \PY{n}{i} \PY{o+ow}{in} \PY{n+nb}{range}\PY{p}{(}\PY{n}{ilim}\PY{p}{)}\PY{p}{:}
    \PY{n}{ax}\PY{o}{.}\PY{n}{plot}\PY{p}{(}\PY{n}{B}\PY{p}{,} \PY{n}{dkn\PYZus{}arr}\PY{p}{[}\PY{n}{Ri}\PY{p}{[}\PY{n}{i}\PY{p}{]}\PY{p}{,}\PY{p}{:}\PY{p}{]}\PY{o}{/}\PY{l+m+mi}{100000}\PY{p}{,} \PY{n}{label} \PY{o}{=} \PY{l+s+sa}{f}\PY{l+s+s2}{\PYZdq{}}\PY{l+s+s2}{dR = }\PY{l+s+si}{\PYZob{}}\PY{n}{R}\PY{p}{[}\PY{n}{Ri}\PY{p}{[}\PY{n}{i}\PY{p}{]}\PY{p}{]}\PY{l+s+si}{:}\PY{l+s+s2}{.4f}\PY{l+s+si}{\PYZcb{}}\PY{l+s+s2}{ * 10}\PY{l+s+se}{\PYZbs{}u207b}\PY{l+s+se}{\PYZbs{}u00b9}\PY{l+s+se}{\PYZbs{}u2070}\PY{l+s+s2}{ m}\PY{l+s+s2}{\PYZdq{}}\PY{p}{)}
    
\PY{n}{plt}\PY{o}{.}\PY{n}{axline}\PY{p}{(}\PY{p}{(}\PY{n}{find\PYZus{}dB\PYZus{}B\PYZus{}osc}\PY{p}{(}\PY{l+m+mi}{2}\PY{p}{)}\PY{p}{,}\PY{o}{\PYZhy{}}\PY{l+m+mf}{0.5}\PY{p}{)}\PY{p}{,} \PY{p}{(}\PY{n}{find\PYZus{}dB\PYZus{}B\PYZus{}osc}\PY{p}{(}\PY{l+m+mi}{2}\PY{p}{)}\PY{p}{,}\PY{o}{\PYZhy{}}\PY{l+m+mf}{0.6}\PY{p}{)}\PY{p}{,} \PY{n}{color}\PY{o}{=}\PY{l+s+s1}{\PYZsq{}}\PY{l+s+s1}{k}\PY{l+s+s1}{\PYZsq{}}\PY{p}{,} \PY{n}{label} \PY{o}{=} \PY{l+s+s2}{\PYZdq{}}\PY{l+s+s2}{period boundary}\PY{l+s+s2}{\PYZdq{}}\PY{p}{)}    
\PY{n}{plt}\PY{o}{.}\PY{n}{axline}\PY{p}{(}\PY{p}{(}\PY{n}{find\PYZus{}dB\PYZus{}B\PYZus{}osc}\PY{p}{(}\PY{l+m+mi}{4}\PY{p}{)}\PY{p}{,}\PY{o}{\PYZhy{}}\PY{l+m+mf}{0.5}\PY{p}{)}\PY{p}{,} \PY{p}{(}\PY{n}{find\PYZus{}dB\PYZus{}B\PYZus{}osc}\PY{p}{(}\PY{l+m+mi}{4}\PY{p}{)}\PY{p}{,}\PY{o}{\PYZhy{}}\PY{l+m+mf}{0.6}\PY{p}{)}\PY{p}{,} \PY{n}{color}\PY{o}{=}\PY{l+s+s1}{\PYZsq{}}\PY{l+s+s1}{k}\PY{l+s+s1}{\PYZsq{}}\PY{p}{)}
\PY{n}{plt}\PY{o}{.}\PY{n}{axline}\PY{p}{(}\PY{p}{(}\PY{n}{find\PYZus{}dB\PYZus{}B\PYZus{}osc}\PY{p}{(}\PY{l+m+mi}{6}\PY{p}{)}\PY{p}{,}\PY{o}{\PYZhy{}}\PY{l+m+mf}{0.5}\PY{p}{)}\PY{p}{,} \PY{p}{(}\PY{n}{find\PYZus{}dB\PYZus{}B\PYZus{}osc}\PY{p}{(}\PY{l+m+mi}{6}\PY{p}{)}\PY{p}{,}\PY{o}{\PYZhy{}}\PY{l+m+mf}{0.6}\PY{p}{)}\PY{p}{,} \PY{n}{color}\PY{o}{=}\PY{l+s+s1}{\PYZsq{}}\PY{l+s+s1}{k}\PY{l+s+s1}{\PYZsq{}}\PY{p}{)}

\PY{n}{plt}\PY{o}{.}\PY{n}{legend}\PY{p}{(}\PY{n}{loc}\PY{o}{=}\PY{l+s+s1}{\PYZsq{}}\PY{l+s+s1}{center left}\PY{l+s+s1}{\PYZsq{}}\PY{p}{,} \PY{n}{bbox\PYZus{}to\PYZus{}anchor}\PY{o}{=}\PY{p}{(}\PY{l+m+mi}{1}\PY{p}{,} \PY{l+m+mf}{0.5}\PY{p}{)}\PY{p}{)}

\PY{n}{plt}\PY{o}{.}\PY{n}{show}\PY{p}{(}\PY{p}{)}
\end{Verbatim}
\end{tcolorbox}

    \begin{center}
    \adjustimage{max size={0.9\linewidth}{0.9\paperheight}}{figs/bvp_kM_en/output_185_0.png}
    \end{center}
    { \hspace*{\fill} \\}
    
    \begin{tcolorbox}[breakable, size=fbox, boxrule=1pt, pad at break*=1mm,colback=cellbackground, colframe=cellborder]
\prompt{In}{incolor}{213}{\boxspacing}
\begin{Verbatim}[commandchars=\\\{\}]
\PY{n}{Ri} \PY{o}{=} \PY{p}{[}\PY{l+m+mi}{0}\PY{p}{,}\PY{l+m+mi}{1}\PY{p}{,}\PY{l+m+mi}{4}\PY{p}{,}\PY{l+m+mi}{5}\PY{p}{]}

\PY{n}{ilim} \PY{o}{=} \PY{n+nb}{len}\PY{p}{(}\PY{n}{Ri}\PY{p}{)}

\PY{n}{fig}\PY{p}{,} \PY{n}{ax} \PY{o}{=} \PY{n}{plt}\PY{o}{.}\PY{n}{subplots}\PY{p}{(}\PY{p}{)}
\PY{n}{plt}\PY{o}{.}\PY{n}{title}\PY{p}{(}\PY{l+s+sa}{f}\PY{l+s+s2}{\PYZdq{}}\PY{l+s+s2}{Change in k Due to Nonlinearity vs B for dk = 0.45 / 2 * 10}\PY{l+s+se}{\PYZbs{}u2076}\PY{l+s+s2}{ m}\PY{l+s+se}{\PYZbs{}u207b}\PY{l+s+se}{\PYZbs{}u00b9}\PY{l+s+se}{\PYZbs{}n}\PY{l+s+s2}{for Various dR Values and R = 9 * 10}\PY{l+s+se}{\PYZbs{}u207b}\PY{l+s+se}{\PYZbs{}u2077}\PY{l+s+s2}{ m + dR}\PY{l+s+s2}{\PYZdq{}}\PY{p}{)}
\PY{n}{ax}\PY{o}{.}\PY{n}{set\PYZus{}ylabel}\PY{p}{(}\PY{l+s+s1}{\PYZsq{}}\PY{l+s+s1}{Change in k (dkn) due to nonlinearity (10}\PY{l+s+se}{\PYZbs{}u2075}\PY{l+s+s1}{ m}\PY{l+s+se}{\PYZbs{}u207b}\PY{l+s+se}{\PYZbs{}u00b9}\PY{l+s+s1}{)}\PY{l+s+s1}{\PYZsq{}}\PY{p}{)}
\PY{n}{ax}\PY{o}{.}\PY{n}{set\PYZus{}xlabel}\PY{p}{(}\PY{l+s+s2}{\PYZdq{}}\PY{l+s+s2}{Magnetic field strength B (B = B * B\PYZus{}max)}\PY{l+s+s2}{\PYZdq{}}\PY{p}{)}

\PY{k}{for} \PY{n}{i} \PY{o+ow}{in} \PY{n+nb}{range}\PY{p}{(}\PY{n}{ilim}\PY{p}{)}\PY{p}{:}
    \PY{n}{ax}\PY{o}{.}\PY{n}{plot}\PY{p}{(}\PY{n}{B}\PY{p}{,} \PY{n}{dkn\PYZus{}arr}\PY{p}{[}\PY{n}{Ri}\PY{p}{[}\PY{n}{i}\PY{p}{]}\PY{p}{,}\PY{p}{:}\PY{p}{]}\PY{o}{/}\PY{l+m+mi}{100000}\PY{p}{,} \PY{n}{label} \PY{o}{=} \PY{l+s+sa}{f}\PY{l+s+s2}{\PYZdq{}}\PY{l+s+s2}{dR = }\PY{l+s+si}{\PYZob{}}\PY{n}{R}\PY{p}{[}\PY{n}{Ri}\PY{p}{[}\PY{n}{i}\PY{p}{]}\PY{p}{]}\PY{l+s+si}{:}\PY{l+s+s2}{.4f}\PY{l+s+si}{\PYZcb{}}\PY{l+s+s2}{ * 10}\PY{l+s+se}{\PYZbs{}u207b}\PY{l+s+se}{\PYZbs{}u00b9}\PY{l+s+se}{\PYZbs{}u2070}\PY{l+s+s2}{ m}\PY{l+s+s2}{\PYZdq{}}\PY{p}{)}

\PY{n}{plt}\PY{o}{.}\PY{n}{axline}\PY{p}{(}\PY{p}{(}\PY{n}{find\PYZus{}dB\PYZus{}B\PYZus{}osc}\PY{p}{(}\PY{l+m+mi}{2}\PY{p}{)}\PY{p}{,}\PY{o}{\PYZhy{}}\PY{l+m+mf}{0.5}\PY{p}{)}\PY{p}{,} \PY{p}{(}\PY{n}{find\PYZus{}dB\PYZus{}B\PYZus{}osc}\PY{p}{(}\PY{l+m+mi}{2}\PY{p}{)}\PY{p}{,}\PY{o}{\PYZhy{}}\PY{l+m+mf}{0.6}\PY{p}{)}\PY{p}{,} \PY{n}{color}\PY{o}{=}\PY{l+s+s1}{\PYZsq{}}\PY{l+s+s1}{k}\PY{l+s+s1}{\PYZsq{}}\PY{p}{,} \PY{n}{label} \PY{o}{=} \PY{l+s+s2}{\PYZdq{}}\PY{l+s+s2}{period boundary}\PY{l+s+s2}{\PYZdq{}}\PY{p}{)}    
\PY{n}{plt}\PY{o}{.}\PY{n}{axline}\PY{p}{(}\PY{p}{(}\PY{n}{find\PYZus{}dB\PYZus{}B\PYZus{}osc}\PY{p}{(}\PY{l+m+mi}{4}\PY{p}{)}\PY{p}{,}\PY{o}{\PYZhy{}}\PY{l+m+mf}{0.5}\PY{p}{)}\PY{p}{,} \PY{p}{(}\PY{n}{find\PYZus{}dB\PYZus{}B\PYZus{}osc}\PY{p}{(}\PY{l+m+mi}{4}\PY{p}{)}\PY{p}{,}\PY{o}{\PYZhy{}}\PY{l+m+mf}{0.6}\PY{p}{)}\PY{p}{,} \PY{n}{color}\PY{o}{=}\PY{l+s+s1}{\PYZsq{}}\PY{l+s+s1}{k}\PY{l+s+s1}{\PYZsq{}}\PY{p}{)}
\PY{n}{plt}\PY{o}{.}\PY{n}{axline}\PY{p}{(}\PY{p}{(}\PY{n}{find\PYZus{}dB\PYZus{}B\PYZus{}osc}\PY{p}{(}\PY{l+m+mi}{6}\PY{p}{)}\PY{p}{,}\PY{o}{\PYZhy{}}\PY{l+m+mf}{0.5}\PY{p}{)}\PY{p}{,} \PY{p}{(}\PY{n}{find\PYZus{}dB\PYZus{}B\PYZus{}osc}\PY{p}{(}\PY{l+m+mi}{6}\PY{p}{)}\PY{p}{,}\PY{o}{\PYZhy{}}\PY{l+m+mf}{0.6}\PY{p}{)}\PY{p}{,} \PY{n}{color}\PY{o}{=}\PY{l+s+s1}{\PYZsq{}}\PY{l+s+s1}{k}\PY{l+s+s1}{\PYZsq{}}\PY{p}{)}

\PY{n}{plt}\PY{o}{.}\PY{n}{legend}\PY{p}{(}\PY{n}{loc}\PY{o}{=}\PY{l+s+s1}{\PYZsq{}}\PY{l+s+s1}{center left}\PY{l+s+s1}{\PYZsq{}}\PY{p}{,} \PY{n}{bbox\PYZus{}to\PYZus{}anchor}\PY{o}{=}\PY{p}{(}\PY{l+m+mi}{1}\PY{p}{,} \PY{l+m+mf}{0.5}\PY{p}{)}\PY{p}{)}

\PY{n}{plt}\PY{o}{.}\PY{n}{show}\PY{p}{(}\PY{p}{)}
\end{Verbatim}
\end{tcolorbox}

    \begin{center}
    \adjustimage{max size={0.9\linewidth}{0.9\paperheight}}{figs/bvp_kM_en/output_186_0.png}
    \end{center}
    { \hspace*{\fill} \\}
    
    \begin{tcolorbox}[breakable, size=fbox, boxrule=1pt, pad at break*=1mm,colback=cellbackground, colframe=cellborder]
\prompt{In}{incolor}{214}{\boxspacing}
\begin{Verbatim}[commandchars=\\\{\}]
\PY{n}{Ri} \PY{o}{=} \PY{p}{[} \PY{l+m+mi}{1}\PY{p}{,} \PY{l+m+mi}{2}\PY{p}{,} \PY{l+m+mi}{3}\PY{p}{,} \PY{l+m+mi}{4}\PY{p}{,} \PY{l+m+mi}{5}\PY{p}{]}

\PY{n}{ilim} \PY{o}{=} \PY{n+nb}{len}\PY{p}{(}\PY{n}{Ri}\PY{p}{)}

\PY{n}{fig}\PY{p}{,} \PY{n}{ax} \PY{o}{=} \PY{n}{plt}\PY{o}{.}\PY{n}{subplots}\PY{p}{(}\PY{p}{)}
\PY{n}{plt}\PY{o}{.}\PY{n}{title}\PY{p}{(}\PY{l+s+sa}{f}\PY{l+s+s2}{\PYZdq{}}\PY{l+s+s2}{Change in M Due to Nonlinearity vs B for dk = 0.45 / 2 * 10}\PY{l+s+se}{\PYZbs{}u2076}\PY{l+s+s2}{ m}\PY{l+s+se}{\PYZbs{}u207b}\PY{l+s+se}{\PYZbs{}u00b9}\PY{l+s+se}{\PYZbs{}n}\PY{l+s+s2}{for Various dR Values and R = 9 * 10}\PY{l+s+se}{\PYZbs{}u207b}\PY{l+s+se}{\PYZbs{}u2077}\PY{l+s+s2}{ m + dR}\PY{l+s+s2}{\PYZdq{}}\PY{p}{)}
\PY{n}{ax}\PY{o}{.}\PY{n}{set\PYZus{}ylabel}\PY{p}{(}\PY{l+s+s1}{\PYZsq{}}\PY{l+s+s1}{Change in M (dMn) due to nonlinearity (10}\PY{l+s+se}{\PYZbs{}u2075}\PY{l+s+s1}{ m}\PY{l+s+se}{\PYZbs{}u207b}\PY{l+s+se}{\PYZbs{}u00b9}\PY{l+s+s1}{)}\PY{l+s+s1}{\PYZsq{}}\PY{p}{)}
\PY{n}{ax}\PY{o}{.}\PY{n}{set\PYZus{}xlabel}\PY{p}{(}\PY{l+s+s2}{\PYZdq{}}\PY{l+s+s2}{Magnetic field strength B (B = B * B\PYZus{}max)}\PY{l+s+s2}{\PYZdq{}}\PY{p}{)}

\PY{k}{for} \PY{n}{i} \PY{o+ow}{in} \PY{n+nb}{range}\PY{p}{(}\PY{n}{ilim}\PY{p}{)}\PY{p}{:}
    \PY{n}{ax}\PY{o}{.}\PY{n}{plot}\PY{p}{(}\PY{n}{B}\PY{p}{,} \PY{n}{dMn\PYZus{}arr}\PY{p}{[}\PY{n}{Ri}\PY{p}{[}\PY{n}{i}\PY{p}{]}\PY{p}{,}\PY{p}{:}\PY{p}{]}\PY{o}{/}\PY{l+m+mi}{100000}\PY{p}{,} \PY{n}{label} \PY{o}{=} \PY{l+s+sa}{f}\PY{l+s+s2}{\PYZdq{}}\PY{l+s+s2}{dR = }\PY{l+s+si}{\PYZob{}}\PY{n}{R}\PY{p}{[}\PY{n}{Ri}\PY{p}{[}\PY{n}{i}\PY{p}{]}\PY{p}{]}\PY{l+s+si}{:}\PY{l+s+s2}{.4f}\PY{l+s+si}{\PYZcb{}}\PY{l+s+s2}{ * 10}\PY{l+s+se}{\PYZbs{}u207b}\PY{l+s+se}{\PYZbs{}u00b9}\PY{l+s+se}{\PYZbs{}u2070}\PY{l+s+s2}{ m}\PY{l+s+s2}{\PYZdq{}}\PY{p}{)}

\PY{n}{plt}\PY{o}{.}\PY{n}{axline}\PY{p}{(}\PY{p}{(}\PY{n}{find\PYZus{}dB\PYZus{}B\PYZus{}osc}\PY{p}{(}\PY{l+m+mi}{2}\PY{p}{)}\PY{p}{,}\PY{l+m+mi}{0}\PY{p}{)}\PY{p}{,} \PY{p}{(}\PY{n}{find\PYZus{}dB\PYZus{}B\PYZus{}osc}\PY{p}{(}\PY{l+m+mi}{2}\PY{p}{)}\PY{p}{,}\PY{o}{\PYZhy{}}\PY{l+m+mf}{0.01}\PY{p}{)}\PY{p}{,} \PY{n}{color}\PY{o}{=}\PY{l+s+s1}{\PYZsq{}}\PY{l+s+s1}{k}\PY{l+s+s1}{\PYZsq{}}\PY{p}{,} \PY{n}{label} \PY{o}{=} \PY{l+s+s2}{\PYZdq{}}\PY{l+s+s2}{period boundary}\PY{l+s+s2}{\PYZdq{}}\PY{p}{)}    
\PY{n}{plt}\PY{o}{.}\PY{n}{axline}\PY{p}{(}\PY{p}{(}\PY{n}{find\PYZus{}dB\PYZus{}B\PYZus{}osc}\PY{p}{(}\PY{l+m+mi}{4}\PY{p}{)}\PY{p}{,}\PY{l+m+mi}{0}\PY{p}{)}\PY{p}{,} \PY{p}{(}\PY{n}{find\PYZus{}dB\PYZus{}B\PYZus{}osc}\PY{p}{(}\PY{l+m+mi}{4}\PY{p}{)}\PY{p}{,}\PY{o}{\PYZhy{}}\PY{l+m+mf}{0.01}\PY{p}{)}\PY{p}{,} \PY{n}{color}\PY{o}{=}\PY{l+s+s1}{\PYZsq{}}\PY{l+s+s1}{k}\PY{l+s+s1}{\PYZsq{}}\PY{p}{)}
\PY{n}{plt}\PY{o}{.}\PY{n}{axline}\PY{p}{(}\PY{p}{(}\PY{n}{find\PYZus{}dB\PYZus{}B\PYZus{}osc}\PY{p}{(}\PY{l+m+mi}{6}\PY{p}{)}\PY{p}{,}\PY{l+m+mi}{0}\PY{p}{)}\PY{p}{,} \PY{p}{(}\PY{n}{find\PYZus{}dB\PYZus{}B\PYZus{}osc}\PY{p}{(}\PY{l+m+mi}{6}\PY{p}{)}\PY{p}{,}\PY{o}{\PYZhy{}}\PY{l+m+mf}{0.01}\PY{p}{)}\PY{p}{,} \PY{n}{color}\PY{o}{=}\PY{l+s+s1}{\PYZsq{}}\PY{l+s+s1}{k}\PY{l+s+s1}{\PYZsq{}}\PY{p}{)}

\PY{n}{plt}\PY{o}{.}\PY{n}{legend}\PY{p}{(}\PY{n}{loc}\PY{o}{=}\PY{l+s+s1}{\PYZsq{}}\PY{l+s+s1}{center left}\PY{l+s+s1}{\PYZsq{}}\PY{p}{,} \PY{n}{bbox\PYZus{}to\PYZus{}anchor}\PY{o}{=}\PY{p}{(}\PY{l+m+mi}{1}\PY{p}{,} \PY{l+m+mf}{0.5}\PY{p}{)}\PY{p}{)}

\PY{n}{plt}\PY{o}{.}\PY{n}{show}\PY{p}{(}\PY{p}{)}
\end{Verbatim}
\end{tcolorbox}

    \begin{center}
    \adjustimage{max size={0.9\linewidth}{0.9\paperheight}}{figs/bvp_kM_en/output_187_0.png}
    \end{center}
    { \hspace*{\fill} \\}
    
    The pattern still appears as expected. The periodicity of the slow
oscillations matches the \(B\pi^2R^2 = B\pi^2R^2 + n\) for integers
\(n\) pattern as before.

    \begin{tcolorbox}[breakable, size=fbox, boxrule=1pt, pad at break*=1mm,colback=cellbackground, colframe=cellborder]
\prompt{In}{incolor}{215}{\boxspacing}
\begin{Verbatim}[commandchars=\\\{\}]
\PY{c+c1}{\PYZsh{} Here is one point on which to start the period.}
\PY{n}{B}\PY{p}{[}\PY{l+m+mi}{89}\PY{p}{]}\PY{p}{,} \PY{n}{find\PYZus{}dB\PYZus{}B\PYZus{}osc}\PY{p}{(}\PY{l+m+mi}{3}\PY{p}{)}
\end{Verbatim}
\end{tcolorbox}

            \begin{tcolorbox}[breakable, size=fbox, boxrule=.5pt, pad at break*=1mm, opacityfill=0]
\prompt{Out}{outcolor}{215}{\boxspacing}
\begin{Verbatim}[commandchars=\\\{\}]
(0.3749163879598662, 0.3752636431197695)
\end{Verbatim}
\end{tcolorbox}
        
    Now for a single oscillation in the \(B\) direction.

    \begin{tcolorbox}[breakable, size=fbox, boxrule=1pt, pad at break*=1mm,colback=cellbackground, colframe=cellborder]
\prompt{In}{incolor}{216}{\boxspacing}
\begin{Verbatim}[commandchars=\\\{\}]
\PY{n}{Ri} \PY{o}{=} \PY{p}{[}\PY{l+m+mi}{1}\PY{p}{,} \PY{l+m+mi}{2}\PY{p}{,} \PY{l+m+mi}{7}\PY{p}{,} \PY{l+m+mi}{13}\PY{p}{,} \PY{l+m+mi}{14}\PY{p}{,} \PY{l+m+mi}{19}\PY{p}{,} \PY{l+m+mi}{20}\PY{p}{,} \PY{l+m+mi}{25}\PY{p}{,} \PY{l+m+mi}{26}\PY{p}{]}

\PY{n}{ilim} \PY{o}{=} \PY{n+nb}{len}\PY{p}{(}\PY{n}{Ri}\PY{p}{)}
\PY{n}{bmn} \PY{o}{=} \PY{l+m+mi}{88}
\PY{n}{bmx} \PY{o}{=} \PY{n}{bmn} \PY{o}{+} \PY{l+m+mi}{45}

\PY{n}{fig}\PY{p}{,} \PY{n}{ax} \PY{o}{=} \PY{n}{plt}\PY{o}{.}\PY{n}{subplots}\PY{p}{(}\PY{p}{)}
\PY{n}{plt}\PY{o}{.}\PY{n}{title}\PY{p}{(}\PY{l+s+sa}{f}\PY{l+s+s2}{\PYZdq{}}\PY{l+s+s2}{Change in M Due to Nonlinearity vs B for dk = 0.45 / 2 * 10}\PY{l+s+se}{\PYZbs{}u2076}\PY{l+s+s2}{ m}\PY{l+s+se}{\PYZbs{}u207b}\PY{l+s+se}{\PYZbs{}u00b9}\PY{l+s+se}{\PYZbs{}n}\PY{l+s+s2}{for Various dR Values and R = 9 * 10}\PY{l+s+se}{\PYZbs{}u207b}\PY{l+s+se}{\PYZbs{}u2077}\PY{l+s+s2}{ m + dR}\PY{l+s+s2}{\PYZdq{}}\PY{p}{)}
\PY{n}{ax}\PY{o}{.}\PY{n}{set\PYZus{}ylabel}\PY{p}{(}\PY{l+s+s1}{\PYZsq{}}\PY{l+s+s1}{Change in M (dMn) due to nonlinearity (10}\PY{l+s+se}{\PYZbs{}u2075}\PY{l+s+s1}{ m}\PY{l+s+se}{\PYZbs{}u207b}\PY{l+s+se}{\PYZbs{}u00b9}\PY{l+s+s1}{)}\PY{l+s+s1}{\PYZsq{}}\PY{p}{)}
\PY{n}{ax}\PY{o}{.}\PY{n}{set\PYZus{}xlabel}\PY{p}{(}\PY{l+s+s2}{\PYZdq{}}\PY{l+s+s2}{Magnetic field strength B (B = B * B\PYZus{}max)}\PY{l+s+s2}{\PYZdq{}}\PY{p}{)}

\PY{k}{for} \PY{n}{i} \PY{o+ow}{in} \PY{n+nb}{range}\PY{p}{(}\PY{n}{ilim}\PY{p}{)}\PY{p}{:}
    \PY{n}{ax}\PY{o}{.}\PY{n}{plot}\PY{p}{(}\PY{n}{B}\PY{p}{[}\PY{n}{bmn}\PY{p}{:}\PY{n}{bmx}\PY{p}{]}\PY{p}{,} \PY{n}{dMn\PYZus{}arr}\PY{p}{[}\PY{n}{Ri}\PY{p}{[}\PY{n}{i}\PY{p}{]}\PY{p}{,}\PY{n}{bmn}\PY{p}{:}\PY{n}{bmx}\PY{p}{]}\PY{o}{/}\PY{l+m+mi}{100000}\PY{p}{,} \PY{n}{label} \PY{o}{=} \PY{l+s+sa}{f}\PY{l+s+s2}{\PYZdq{}}\PY{l+s+s2}{dR = }\PY{l+s+si}{\PYZob{}}\PY{n}{R}\PY{p}{[}\PY{n}{Ri}\PY{p}{[}\PY{n}{i}\PY{p}{]}\PY{p}{]}\PY{l+s+si}{:}\PY{l+s+s2}{.4f}\PY{l+s+si}{\PYZcb{}}\PY{l+s+s2}{ * 10}\PY{l+s+se}{\PYZbs{}u207b}\PY{l+s+se}{\PYZbs{}u00b9}\PY{l+s+se}{\PYZbs{}u2070}\PY{l+s+s2}{ m}\PY{l+s+s2}{\PYZdq{}}\PY{p}{)}

\PY{n}{plt}\PY{o}{.}\PY{n}{axline}\PY{p}{(}\PY{p}{(}\PY{n}{find\PYZus{}dB\PYZus{}B\PYZus{}osc}\PY{p}{(}\PY{l+m+mi}{4}\PY{p}{)}\PY{p}{,}\PY{l+m+mi}{0}\PY{p}{)}\PY{p}{,} \PY{p}{(}\PY{n}{find\PYZus{}dB\PYZus{}B\PYZus{}osc}\PY{p}{(}\PY{l+m+mi}{4}\PY{p}{)}\PY{p}{,}\PY{o}{\PYZhy{}}\PY{l+m+mf}{0.01}\PY{p}{)}\PY{p}{,} \PY{n}{color}\PY{o}{=}\PY{l+s+s1}{\PYZsq{}}\PY{l+s+s1}{k}\PY{l+s+s1}{\PYZsq{}}\PY{p}{,} \PY{n}{label} \PY{o}{=} \PY{l+s+s2}{\PYZdq{}}\PY{l+s+s2}{period boundary}\PY{l+s+s2}{\PYZdq{}}\PY{p}{)}
\PY{n}{plt}\PY{o}{.}\PY{n}{axline}\PY{p}{(}\PY{p}{(}\PY{n}{find\PYZus{}dB\PYZus{}B\PYZus{}osc}\PY{p}{(}\PY{l+m+mi}{3}\PY{p}{)}\PY{p}{,}\PY{l+m+mi}{0}\PY{p}{)}\PY{p}{,} \PY{p}{(}\PY{n}{find\PYZus{}dB\PYZus{}B\PYZus{}osc}\PY{p}{(}\PY{l+m+mi}{3}\PY{p}{)}\PY{p}{,}\PY{o}{\PYZhy{}}\PY{l+m+mf}{0.01}\PY{p}{)}\PY{p}{,} \PY{n}{color}\PY{o}{=}\PY{l+s+s1}{\PYZsq{}}\PY{l+s+s1}{k}\PY{l+s+s1}{\PYZsq{}}\PY{p}{)}

\PY{n}{plt}\PY{o}{.}\PY{n}{legend}\PY{p}{(}\PY{n}{loc}\PY{o}{=}\PY{l+s+s1}{\PYZsq{}}\PY{l+s+s1}{center left}\PY{l+s+s1}{\PYZsq{}}\PY{p}{,} \PY{n}{bbox\PYZus{}to\PYZus{}anchor}\PY{o}{=}\PY{p}{(}\PY{l+m+mi}{1}\PY{p}{,} \PY{l+m+mf}{0.5}\PY{p}{)}\PY{p}{)}

\PY{n}{plt}\PY{o}{.}\PY{n}{show}\PY{p}{(}\PY{p}{)}
\end{Verbatim}
\end{tcolorbox}

    \begin{center}
    \adjustimage{max size={0.9\linewidth}{0.9\paperheight}}{figs/bvp_kM_en/output_191_0.png}
    \end{center}
    { \hspace*{\fill} \\}
    
    And the same for \(dkn\).

    \begin{tcolorbox}[breakable, size=fbox, boxrule=1pt, pad at break*=1mm,colback=cellbackground, colframe=cellborder]
\prompt{In}{incolor}{217}{\boxspacing}
\begin{Verbatim}[commandchars=\\\{\}]
\PY{n}{Ri} \PY{o}{=} \PY{p}{[}\PY{l+m+mi}{1}\PY{p}{,} \PY{l+m+mi}{2}\PY{p}{,} \PY{l+m+mi}{7}\PY{p}{,} \PY{l+m+mi}{13}\PY{p}{,} \PY{l+m+mi}{14}\PY{p}{,} \PY{l+m+mi}{19}\PY{p}{,} \PY{l+m+mi}{20}\PY{p}{,} \PY{l+m+mi}{25}\PY{p}{,} \PY{l+m+mi}{26}\PY{p}{]}
\PY{c+c1}{\PYZsh{}Ri = [ 19, 26]}


\PY{n}{ilim} \PY{o}{=} \PY{n+nb}{len}\PY{p}{(}\PY{n}{Ri}\PY{p}{)}
\PY{n}{bmn} \PY{o}{=} \PY{l+m+mi}{88}
\PY{n}{bmx} \PY{o}{=} \PY{n}{bmn} \PY{o}{+} \PY{l+m+mi}{45}

\PY{n}{fig}\PY{p}{,} \PY{n}{ax} \PY{o}{=} \PY{n}{plt}\PY{o}{.}\PY{n}{subplots}\PY{p}{(}\PY{p}{)}
\PY{n}{plt}\PY{o}{.}\PY{n}{title}\PY{p}{(}\PY{l+s+sa}{f}\PY{l+s+s2}{\PYZdq{}}\PY{l+s+s2}{Change in k Due to Nonlinearity vs B for dk = 0.45 / 2 * 10}\PY{l+s+se}{\PYZbs{}u2076}\PY{l+s+s2}{ m}\PY{l+s+se}{\PYZbs{}u207b}\PY{l+s+se}{\PYZbs{}u00b9}\PY{l+s+se}{\PYZbs{}n}\PY{l+s+s2}{for Various dR Values and R = 9 * 10}\PY{l+s+se}{\PYZbs{}u207b}\PY{l+s+se}{\PYZbs{}u2077}\PY{l+s+s2}{ m + dR}\PY{l+s+s2}{\PYZdq{}}\PY{p}{)}
\PY{n}{ax}\PY{o}{.}\PY{n}{set\PYZus{}ylabel}\PY{p}{(}\PY{l+s+s1}{\PYZsq{}}\PY{l+s+s1}{Change in k (dkn) due to nonlinearity (10}\PY{l+s+se}{\PYZbs{}u2075}\PY{l+s+s1}{ m}\PY{l+s+se}{\PYZbs{}u207b}\PY{l+s+se}{\PYZbs{}u00b9}\PY{l+s+s1}{)}\PY{l+s+s1}{\PYZsq{}}\PY{p}{)}
\PY{n}{ax}\PY{o}{.}\PY{n}{set\PYZus{}xlabel}\PY{p}{(}\PY{l+s+s2}{\PYZdq{}}\PY{l+s+s2}{Magnetic field strength B (B = B * B\PYZus{}max)}\PY{l+s+s2}{\PYZdq{}}\PY{p}{)}

\PY{k}{for} \PY{n}{i} \PY{o+ow}{in} \PY{n+nb}{range}\PY{p}{(}\PY{n}{ilim}\PY{p}{)}\PY{p}{:}
    \PY{n}{ax}\PY{o}{.}\PY{n}{plot}\PY{p}{(}\PY{n}{B}\PY{p}{[}\PY{n}{bmn}\PY{p}{:}\PY{n}{bmx}\PY{p}{]}\PY{p}{,} \PY{n}{dkn\PYZus{}arr}\PY{p}{[}\PY{n}{Ri}\PY{p}{[}\PY{n}{i}\PY{p}{]}\PY{p}{,}\PY{n}{bmn}\PY{p}{:}\PY{n}{bmx}\PY{p}{]}\PY{o}{/}\PY{l+m+mi}{100000}\PY{p}{,} \PY{n}{label} \PY{o}{=} \PY{l+s+sa}{f}\PY{l+s+s2}{\PYZdq{}}\PY{l+s+s2}{dR = }\PY{l+s+si}{\PYZob{}}\PY{n}{R}\PY{p}{[}\PY{n}{Ri}\PY{p}{[}\PY{n}{i}\PY{p}{]}\PY{p}{]}\PY{l+s+si}{:}\PY{l+s+s2}{.4f}\PY{l+s+si}{\PYZcb{}}\PY{l+s+s2}{ * 10}\PY{l+s+se}{\PYZbs{}u207b}\PY{l+s+se}{\PYZbs{}u00b9}\PY{l+s+se}{\PYZbs{}u2070}\PY{l+s+s2}{ m}\PY{l+s+s2}{\PYZdq{}}\PY{p}{)}

\PY{n}{plt}\PY{o}{.}\PY{n}{axline}\PY{p}{(}\PY{p}{(}\PY{n}{find\PYZus{}dB\PYZus{}B\PYZus{}osc}\PY{p}{(}\PY{l+m+mi}{4}\PY{p}{)}\PY{p}{,}\PY{o}{\PYZhy{}}\PY{l+m+mf}{.4}\PY{p}{)}\PY{p}{,} \PY{p}{(}\PY{n}{find\PYZus{}dB\PYZus{}B\PYZus{}osc}\PY{p}{(}\PY{l+m+mi}{4}\PY{p}{)}\PY{p}{,}\PY{o}{\PYZhy{}}\PY{l+m+mf}{0.401}\PY{p}{)}\PY{p}{,} \PY{n}{color}\PY{o}{=}\PY{l+s+s1}{\PYZsq{}}\PY{l+s+s1}{k}\PY{l+s+s1}{\PYZsq{}}\PY{p}{,} \PY{n}{label} \PY{o}{=} \PY{l+s+s2}{\PYZdq{}}\PY{l+s+s2}{period boundary}\PY{l+s+s2}{\PYZdq{}}\PY{p}{)}
\PY{n}{plt}\PY{o}{.}\PY{n}{axline}\PY{p}{(}\PY{p}{(}\PY{n}{find\PYZus{}dB\PYZus{}B\PYZus{}osc}\PY{p}{(}\PY{l+m+mi}{3}\PY{p}{)}\PY{p}{,}\PY{o}{\PYZhy{}}\PY{l+m+mf}{.4}\PY{p}{)}\PY{p}{,} \PY{p}{(}\PY{n}{find\PYZus{}dB\PYZus{}B\PYZus{}osc}\PY{p}{(}\PY{l+m+mi}{3}\PY{p}{)}\PY{p}{,}\PY{o}{\PYZhy{}}\PY{l+m+mf}{0.401}\PY{p}{)}\PY{p}{,} \PY{n}{color}\PY{o}{=}\PY{l+s+s1}{\PYZsq{}}\PY{l+s+s1}{k}\PY{l+s+s1}{\PYZsq{}}\PY{p}{)}

\PY{n}{plt}\PY{o}{.}\PY{n}{legend}\PY{p}{(}\PY{n}{loc}\PY{o}{=}\PY{l+s+s1}{\PYZsq{}}\PY{l+s+s1}{center left}\PY{l+s+s1}{\PYZsq{}}\PY{p}{,} \PY{n}{bbox\PYZus{}to\PYZus{}anchor}\PY{o}{=}\PY{p}{(}\PY{l+m+mi}{1}\PY{p}{,} \PY{l+m+mf}{0.5}\PY{p}{)}\PY{p}{)}

\PY{n}{plt}\PY{o}{.}\PY{n}{show}\PY{p}{(}\PY{p}{)}
\end{Verbatim}
\end{tcolorbox}

    \begin{center}
    \adjustimage{max size={0.9\linewidth}{0.9\paperheight}}{figs/bvp_kM_en/output_193_0.png}
    \end{center}
    { \hspace*{\fill} \\}
    
    Our period guess appears correct. Now to draw the expected nodes of the
graphs.

    \begin{tcolorbox}[breakable, size=fbox, boxrule=1pt, pad at break*=1mm,colback=cellbackground, colframe=cellborder]
\prompt{In}{incolor}{224}{\boxspacing}
\begin{Verbatim}[commandchars=\\\{\}]
\PY{n}{Ri} \PY{o}{=} \PY{p}{[}\PY{l+m+mi}{1}\PY{p}{,} \PY{l+m+mi}{2}\PY{p}{,} \PY{l+m+mi}{7}\PY{p}{,} \PY{l+m+mi}{13}\PY{p}{,} \PY{l+m+mi}{14}\PY{p}{,} \PY{l+m+mi}{19}\PY{p}{,} \PY{l+m+mi}{20}\PY{p}{,} \PY{l+m+mi}{25}\PY{p}{,} \PY{l+m+mi}{26}\PY{p}{]}

\PY{n}{ilim} \PY{o}{=} \PY{n+nb}{len}\PY{p}{(}\PY{n}{Ri}\PY{p}{)}
\PY{n}{bmn} \PY{o}{=} \PY{l+m+mi}{88}
\PY{n}{bmx} \PY{o}{=} \PY{n}{bmn} \PY{o}{+} \PY{l+m+mi}{45}

\PY{n}{fig}\PY{p}{,} \PY{n}{ax} \PY{o}{=} \PY{n}{plt}\PY{o}{.}\PY{n}{subplots}\PY{p}{(}\PY{p}{)}
\PY{n}{plt}\PY{o}{.}\PY{n}{title}\PY{p}{(}\PY{l+s+sa}{f}\PY{l+s+s2}{\PYZdq{}}\PY{l+s+s2}{Change in M Due to Nonlinearity vs Number of Oscillations for}\PY{l+s+se}{\PYZbs{}n}\PY{l+s+s2}{dk = 0.45 / 2 * 10}\PY{l+s+se}{\PYZbs{}u2076}\PY{l+s+s2}{ m}\PY{l+s+se}{\PYZbs{}u207b}\PY{l+s+se}{\PYZbs{}u00b9}\PY{l+s+s2}{ for Various dR Values and R = 9 * 10}\PY{l+s+se}{\PYZbs{}u207b}\PY{l+s+se}{\PYZbs{}u2077}\PY{l+s+s2}{ m + dR}\PY{l+s+s2}{\PYZdq{}}\PY{p}{)}
\PY{n}{ax}\PY{o}{.}\PY{n}{set\PYZus{}ylabel}\PY{p}{(}\PY{l+s+s1}{\PYZsq{}}\PY{l+s+s1}{Change in M (dMn) due to nonlinearity (10}\PY{l+s+se}{\PYZbs{}u2075}\PY{l+s+s1}{ m}\PY{l+s+se}{\PYZbs{}u207b}\PY{l+s+se}{\PYZbs{}u00b9}\PY{l+s+s1}{)}\PY{l+s+s1}{\PYZsq{}}\PY{p}{)}
\PY{n}{ax}\PY{o}{.}\PY{n}{set\PYZus{}xlabel}\PY{p}{(}\PY{l+s+s2}{\PYZdq{}}\PY{l+s+s2}{B * π}\PY{l+s+se}{\PYZbs{}u00b2}\PY{l+s+s2}{ * R}\PY{l+s+se}{\PYZbs{}u00b2}\PY{l+s+s2}{ with R and B as proportions of their maximum values}\PY{l+s+s2}{\PYZdq{}}\PY{p}{)}

\PY{k}{for} \PY{n}{i} \PY{o+ow}{in} \PY{n+nb}{range}\PY{p}{(}\PY{n}{ilim}\PY{p}{)}\PY{p}{:}
    \PY{n}{ax}\PY{o}{.}\PY{n}{plot}\PY{p}{(}\PY{n}{B}\PY{o}{*}\PY{p}{(}\PY{n}{R}\PY{p}{[}\PY{n}{Ri}\PY{p}{[}\PY{n}{i}\PY{p}{]}\PY{p}{]}\PY{o}{*}\PY{n}{scale\PYZus{}R} \PY{o}{+}\PY{n}{R0}\PY{p}{)}\PY{o}{*}\PY{o}{*}\PY{l+m+mi}{2}\PY{o}{*}\PY{n}{pi}\PY{o}{*}\PY{o}{*}\PY{l+m+mi}{2}\PY{p}{,} \PY{n}{dMn\PYZus{}arr}\PY{p}{[}\PY{n}{Ri}\PY{p}{[}\PY{n}{i}\PY{p}{]}\PY{p}{,}\PY{p}{:}\PY{p}{]}\PY{o}{/}\PY{l+m+mi}{100000}\PY{p}{,} \PY{n}{label} \PY{o}{=} \PY{l+s+sa}{f}\PY{l+s+s2}{\PYZdq{}}\PY{l+s+s2}{dR = }\PY{l+s+si}{\PYZob{}}\PY{n}{R}\PY{p}{[}\PY{n}{Ri}\PY{p}{[}\PY{n}{i}\PY{p}{]}\PY{p}{]}\PY{l+s+si}{:}\PY{l+s+s2}{.4f}\PY{l+s+si}{\PYZcb{}}\PY{l+s+s2}{ * 10}\PY{l+s+se}{\PYZbs{}u207b}\PY{l+s+se}{\PYZbs{}u00b9}\PY{l+s+se}{\PYZbs{}u2070}\PY{l+s+s2}{ m}\PY{l+s+s2}{\PYZdq{}}\PY{p}{)}

\PY{k}{for} \PY{n}{i} \PY{o+ow}{in} \PY{n+nb}{range}\PY{p}{(}\PY{l+m+mi}{1}\PY{p}{,}\PY{l+m+mi}{8}\PY{p}{)}\PY{p}{:}
    \PY{n}{plt}\PY{o}{.}\PY{n}{axline}\PY{p}{(}\PY{p}{(}\PY{n}{i}\PY{p}{,}\PY{l+m+mi}{0}\PY{p}{)}\PY{p}{,} \PY{p}{(}\PY{n}{i}\PY{p}{,}\PY{o}{\PYZhy{}}\PY{l+m+mf}{0.01}\PY{p}{)}\PY{p}{,} \PY{n}{color}\PY{o}{=}\PY{l+s+s1}{\PYZsq{}}\PY{l+s+s1}{k}\PY{l+s+s1}{\PYZsq{}}\PY{p}{)}
    \PY{n}{plt}\PY{o}{.}\PY{n}{axline}\PY{p}{(}\PY{p}{(}\PY{n}{i}\PY{o}{+}\PY{l+m+mf}{.5}\PY{p}{,}\PY{l+m+mi}{0}\PY{p}{)}\PY{p}{,} \PY{p}{(}\PY{n}{i}\PY{o}{+}\PY{l+m+mf}{.5}\PY{p}{,}\PY{o}{\PYZhy{}}\PY{l+m+mf}{0.01}\PY{p}{)}\PY{p}{,} \PY{n}{color}\PY{o}{=}\PY{l+s+s1}{\PYZsq{}}\PY{l+s+s1}{blue}\PY{l+s+s1}{\PYZsq{}}\PY{p}{)}

\PY{n}{plt}\PY{o}{.}\PY{n}{legend}\PY{p}{(}\PY{n}{loc}\PY{o}{=}\PY{l+s+s1}{\PYZsq{}}\PY{l+s+s1}{center left}\PY{l+s+s1}{\PYZsq{}}\PY{p}{,} \PY{n}{bbox\PYZus{}to\PYZus{}anchor}\PY{o}{=}\PY{p}{(}\PY{l+m+mi}{1}\PY{p}{,} \PY{l+m+mf}{0.5}\PY{p}{)}\PY{p}{)}

\PY{n}{plt}\PY{o}{.}\PY{n}{show}\PY{p}{(}\PY{p}{)}
\end{Verbatim}
\end{tcolorbox}

    \begin{center}
    \adjustimage{max size={0.9\linewidth}{0.9\paperheight}}{figs/bvp_kM_en/output_195_0.png}
    \end{center}
    { \hspace*{\fill} \\}
    
    Here I am plotting the various \(dkn\) vs \(R\) graphs for B values
separated by one period on top of each other. A difference between
periods should be obvious by having the lower graphics be visible under
the upper graphic.

    \begin{tcolorbox}[breakable, size=fbox, boxrule=1pt, pad at break*=1mm,colback=cellbackground, colframe=cellborder]
\prompt{In}{incolor}{227}{\boxspacing}
\begin{Verbatim}[commandchars=\\\{\}]
\PY{n}{bmn} \PY{o}{=} \PY{l+m+mi}{85}
\PY{n}{Bi} \PY{o}{=} \PY{p}{[}\PY{l+m+mi}{4}\PY{o}{+}\PY{n}{bmn}\PY{o}{\PYZhy{}}\PY{l+m+mi}{84}\PY{p}{,}\PY{l+m+mi}{4}\PY{o}{+}\PY{n}{bmn}\PY{o}{\PYZhy{}}\PY{l+m+mi}{63}\PY{p}{,}\PY{l+m+mi}{4}\PY{o}{+}\PY{n}{bmn}\PY{o}{\PYZhy{}}\PY{l+m+mi}{42}\PY{p}{,}\PY{l+m+mi}{4}\PY{o}{+}\PY{n}{bmn}\PY{o}{\PYZhy{}}\PY{l+m+mi}{21}\PY{p}{,}\PY{l+m+mi}{4}\PY{o}{+}\PY{n}{bmn}\PY{p}{,} \PY{l+m+mi}{25}\PY{o}{+}\PY{n}{bmn}\PY{p}{,} \PY{l+m+mi}{46}\PY{o}{+}\PY{n}{bmn}\PY{p}{,} \PY{l+m+mi}{67}\PY{o}{+}\PY{n}{bmn}\PY{p}{,} \PY{l+m+mi}{67}\PY{o}{+}\PY{n}{bmn}\PY{o}{+}\PY{l+m+mi}{21}\PY{p}{,} \PY{l+m+mi}{67}\PY{o}{+}\PY{n}{bmn}\PY{o}{+}\PY{l+m+mi}{42}\PY{p}{,} \PY{l+m+mi}{67}\PY{o}{+}\PY{n}{bmn}\PY{o}{+}\PY{l+m+mi}{63}\PY{p}{]}\PY{c+c1}{\PYZsh{}, 99+bmn, 120+bmn]}

\PY{n}{ilim} \PY{o}{=} \PY{n+nb}{len}\PY{p}{(}\PY{n}{Bi}\PY{p}{)}

\PY{n}{fig}\PY{p}{,} \PY{n}{ax} \PY{o}{=} \PY{n}{plt}\PY{o}{.}\PY{n}{subplots}\PY{p}{(}\PY{p}{)}
\PY{n}{plt}\PY{o}{.}\PY{n}{title}\PY{p}{(}\PY{l+s+sa}{f}\PY{l+s+s2}{\PYZdq{}}\PY{l+s+s2}{Plot of the Change in k Due to Nonlinearity vs R for dk = 0.45 / 2 * 10}\PY{l+s+se}{\PYZbs{}u2076}\PY{l+s+s2}{ m}\PY{l+s+se}{\PYZbs{}u207b}\PY{l+s+se}{\PYZbs{}u00b9}\PY{l+s+se}{\PYZbs{}n}\PY{l+s+s2}{for B Values Predicted to be at a Half\PYZhy{}Period Boundary}\PY{l+s+s2}{\PYZdq{}}\PY{p}{)}
\PY{n}{ax}\PY{o}{.}\PY{n}{set\PYZus{}ylabel}\PY{p}{(}\PY{l+s+s1}{\PYZsq{}}\PY{l+s+s1}{Change in k (dkn) due to nonlinearity (10}\PY{l+s+se}{\PYZbs{}u2075}\PY{l+s+s1}{ m}\PY{l+s+se}{\PYZbs{}u207b}\PY{l+s+se}{\PYZbs{}u00b9}\PY{l+s+s1}{)}\PY{l+s+s1}{\PYZsq{}}\PY{p}{)}
\PY{n}{ax}\PY{o}{.}\PY{n}{set\PYZus{}xlabel}\PY{p}{(}\PY{l+s+s2}{\PYZdq{}}\PY{l+s+s2}{Change dr in radius R of ring (R = 9 * 10}\PY{l+s+se}{\PYZbs{}u207b}\PY{l+s+se}{\PYZbs{}u2077}\PY{l+s+s2}{ m + dr * 10}\PY{l+s+se}{\PYZbs{}u207b}\PY{l+s+se}{\PYZbs{}u00b9}\PY{l+s+se}{\PYZbs{}u2070}\PY{l+s+s2}{ m)}\PY{l+s+s2}{\PYZdq{}}\PY{p}{)}

\PY{k}{for} \PY{n}{i} \PY{o+ow}{in} \PY{n+nb}{range}\PY{p}{(}\PY{n}{ilim}\PY{p}{)}\PY{p}{:}
    \PY{n}{ax}\PY{o}{.}\PY{n}{plot}\PY{p}{(}\PY{n}{R}\PY{p}{,} \PY{n}{dkn\PYZus{}arr}\PY{p}{[}\PY{p}{:}\PY{p}{,} \PY{n}{Bi}\PY{p}{[}\PY{n}{i}\PY{p}{]}\PY{p}{]}\PY{o}{/}\PY{l+m+mi}{100000}\PY{p}{,} \PY{n}{label} \PY{o}{=} \PY{l+s+sa}{f}\PY{l+s+s2}{\PYZdq{}}\PY{l+s+s2}{B = }\PY{l+s+si}{\PYZob{}}\PY{n}{B}\PY{p}{[}\PY{n}{Bi}\PY{p}{[}\PY{n}{i}\PY{p}{]}\PY{p}{]}\PY{l+s+si}{:}\PY{l+s+s2}{.4f}\PY{l+s+si}{\PYZcb{}}\PY{l+s+s2}{ * B\PYZus{}max}\PY{l+s+s2}{\PYZdq{}}\PY{p}{)}

\PY{c+c1}{\PYZsh{}plt.axline((R[1],\PYZhy{}0.5), (R[1],\PYZhy{}0.51), color=\PYZsq{}k\PYZsq{})}
\PY{c+c1}{\PYZsh{}plt.axline((R[25],\PYZhy{}0.5), (R[25],\PYZhy{}0.51), color=\PYZsq{}k\PYZsq{})}

\PY{n}{plt}\PY{o}{.}\PY{n}{legend}\PY{p}{(}\PY{n}{loc}\PY{o}{=}\PY{l+s+s1}{\PYZsq{}}\PY{l+s+s1}{center left}\PY{l+s+s1}{\PYZsq{}}\PY{p}{,} \PY{n}{bbox\PYZus{}to\PYZus{}anchor}\PY{o}{=}\PY{p}{(}\PY{l+m+mi}{1}\PY{p}{,} \PY{l+m+mf}{0.5}\PY{p}{)}\PY{p}{)}

\PY{n}{plt}\PY{o}{.}\PY{n}{show}\PY{p}{(}\PY{p}{)}
\end{Verbatim}
\end{tcolorbox}

    \begin{center}
    \adjustimage{max size={0.9\linewidth}{0.9\paperheight}}{figs/bvp_kM_en/output_197_0.png}
    \end{center}
    { \hspace*{\fill} \\}
    
    \begin{tcolorbox}[breakable, size=fbox, boxrule=1pt, pad at break*=1mm,colback=cellbackground, colframe=cellborder]
\prompt{In}{incolor}{228}{\boxspacing}
\begin{Verbatim}[commandchars=\\\{\}]
\PY{n}{Bi} \PY{o}{=} \PY{p}{[}\PY{l+m+mi}{4}\PY{o}{+}\PY{n}{bmn}\PY{o}{\PYZhy{}}\PY{l+m+mi}{84}\PY{p}{,}\PY{l+m+mi}{4}\PY{o}{+}\PY{n}{bmn}\PY{o}{\PYZhy{}}\PY{l+m+mi}{63}\PY{p}{,}\PY{l+m+mi}{4}\PY{o}{+}\PY{n}{bmn}\PY{o}{\PYZhy{}}\PY{l+m+mi}{42}\PY{p}{,}\PY{l+m+mi}{4}\PY{o}{+}\PY{n}{bmn}\PY{o}{\PYZhy{}}\PY{l+m+mi}{21}\PY{p}{,}\PY{l+m+mi}{4}\PY{o}{+}\PY{n}{bmn}\PY{p}{,} \PY{l+m+mi}{25}\PY{o}{+}\PY{n}{bmn}\PY{p}{,} \PY{l+m+mi}{46}\PY{o}{+}\PY{n}{bmn}\PY{p}{,} \PY{l+m+mi}{67}\PY{o}{+}\PY{n}{bmn}\PY{p}{,} \PY{l+m+mi}{67}\PY{o}{+}\PY{n}{bmn}\PY{o}{+}\PY{l+m+mi}{21}\PY{p}{,} \PY{l+m+mi}{67}\PY{o}{+}\PY{n}{bmn}\PY{o}{+}\PY{l+m+mi}{42}\PY{p}{,} \PY{l+m+mi}{67}\PY{o}{+}\PY{n}{bmn}\PY{o}{+}\PY{l+m+mi}{63}\PY{p}{]}\PY{c+c1}{\PYZsh{}, 99+bmn, 120+bmn]}

\PY{n}{ilim} \PY{o}{=} \PY{n+nb}{len}\PY{p}{(}\PY{n}{Bi}\PY{p}{)}

\PY{n}{fig}\PY{p}{,} \PY{n}{ax} \PY{o}{=} \PY{n}{plt}\PY{o}{.}\PY{n}{subplots}\PY{p}{(}\PY{p}{)}
\PY{n}{plt}\PY{o}{.}\PY{n}{title}\PY{p}{(}\PY{l+s+sa}{f}\PY{l+s+s2}{\PYZdq{}}\PY{l+s+s2}{Plot of the Change in M Due to Nonlinearity vs R for dk = 0.45 / 2 * 10}\PY{l+s+se}{\PYZbs{}u2076}\PY{l+s+s2}{ m}\PY{l+s+se}{\PYZbs{}u207b}\PY{l+s+se}{\PYZbs{}u00b9}\PY{l+s+se}{\PYZbs{}n}\PY{l+s+s2}{for B Values Predicted to be at a Half\PYZhy{}Period Boundary}\PY{l+s+s2}{\PYZdq{}}\PY{p}{)}
\PY{n}{ax}\PY{o}{.}\PY{n}{set\PYZus{}ylabel}\PY{p}{(}\PY{l+s+s1}{\PYZsq{}}\PY{l+s+s1}{Change in M (dMn) due to nonlinearity (10}\PY{l+s+se}{\PYZbs{}u2075}\PY{l+s+s1}{ m}\PY{l+s+se}{\PYZbs{}u207b}\PY{l+s+se}{\PYZbs{}u00b9}\PY{l+s+s1}{)}\PY{l+s+s1}{\PYZsq{}}\PY{p}{)}
\PY{n}{ax}\PY{o}{.}\PY{n}{set\PYZus{}xlabel}\PY{p}{(}\PY{l+s+s2}{\PYZdq{}}\PY{l+s+s2}{Change dr in radius R of ring (R = 9 * 10}\PY{l+s+se}{\PYZbs{}u207b}\PY{l+s+se}{\PYZbs{}u2077}\PY{l+s+s2}{ m + dr * 10}\PY{l+s+se}{\PYZbs{}u207b}\PY{l+s+se}{\PYZbs{}u00b9}\PY{l+s+se}{\PYZbs{}u2070}\PY{l+s+s2}{ m)}\PY{l+s+s2}{\PYZdq{}}\PY{p}{)}

\PY{k}{for} \PY{n}{i} \PY{o+ow}{in} \PY{n+nb}{range}\PY{p}{(}\PY{n}{ilim}\PY{p}{)}\PY{p}{:}
    \PY{n}{ax}\PY{o}{.}\PY{n}{plot}\PY{p}{(}\PY{n}{R}\PY{p}{,} \PY{n}{dMn\PYZus{}arr}\PY{p}{[}\PY{p}{:}\PY{p}{,} \PY{n}{Bi}\PY{p}{[}\PY{n}{i}\PY{p}{]}\PY{p}{]}\PY{o}{/}\PY{l+m+mi}{100000}\PY{p}{,} \PY{n}{label} \PY{o}{=} \PY{l+s+sa}{f}\PY{l+s+s2}{\PYZdq{}}\PY{l+s+s2}{B = }\PY{l+s+si}{\PYZob{}}\PY{n}{B}\PY{p}{[}\PY{n}{Bi}\PY{p}{[}\PY{n}{i}\PY{p}{]}\PY{p}{]}\PY{l+s+si}{:}\PY{l+s+s2}{.4f}\PY{l+s+si}{\PYZcb{}}\PY{l+s+s2}{ * B\PYZus{}max}\PY{l+s+s2}{\PYZdq{}}\PY{p}{)}

\PY{n}{plt}\PY{o}{.}\PY{n}{legend}\PY{p}{(}\PY{n}{loc}\PY{o}{=}\PY{l+s+s1}{\PYZsq{}}\PY{l+s+s1}{center left}\PY{l+s+s1}{\PYZsq{}}\PY{p}{,} \PY{n}{bbox\PYZus{}to\PYZus{}anchor}\PY{o}{=}\PY{p}{(}\PY{l+m+mi}{1}\PY{p}{,} \PY{l+m+mf}{0.5}\PY{p}{)}\PY{p}{)}

\PY{n}{plt}\PY{o}{.}\PY{n}{show}\PY{p}{(}\PY{p}{)}
\end{Verbatim}
\end{tcolorbox}

    \begin{center}
    \adjustimage{max size={0.9\linewidth}{0.9\paperheight}}{figs/bvp_kM_en/output_198_0.png}
    \end{center}
    { \hspace*{\fill} \\}
    
    The data here is clearly periodic in \(B\), with the same pattern as
before.

    \hypertarget{mu-dependence-for-fast-oscillations}{%
\section{\texorpdfstring{\(\mu\) Dependence for Fast
Oscillations}{\textbackslash mu Dependence for Fast Oscillations}}\label{mu-dependence-for-fast-oscillations}}

    \begin{tcolorbox}[breakable, size=fbox, boxrule=1pt, pad at break*=1mm,colback=cellbackground, colframe=cellborder]
\prompt{In}{incolor}{127}{\boxspacing}
\begin{Verbatim}[commandchars=\\\{\}]
\PY{n}{tbl\PYZus{}red}\PY{p}{[}\PY{l+m+mi}{1}\PY{p}{]}
\end{Verbatim}
\end{tcolorbox}

            \begin{tcolorbox}[breakable, size=fbox, boxrule=.5pt, pad at break*=1mm, opacityfill=0]
\prompt{Out}{outcolor}{127}{\boxspacing}
\begin{Verbatim}[commandchars=\\\{\}]
       dk            mu            dkn         dMn
0     0.0  3.162278e-09   -7969.899509   13.457522
1     0.0  3.236689e-09  -21241.450718   14.759672
2     0.0  3.312852e-09  -21655.040102   15.081138
3     0.0  3.390807e-09  -22075.471048   15.409241
4     0.0  3.470596e-09  -22502.823629   15.744108
{\ldots}   {\ldots}           {\ldots}            {\ldots}         {\ldots}
9995  3.0  2.881350e-08 -102227.695108  100.520093
9996  3.0  2.949151e-08 -103715.570360  102.482218
9997  3.0  3.018547e-08 -105221.666897  104.480650
9998  3.0  3.089577e-08 -106746.223455  106.516039
9999  3.0  3.162278e-08 -108289.484694  108.589044

[10000 rows x 4 columns]
\end{Verbatim}
\end{tcolorbox}
        
    \begin{tcolorbox}[breakable, size=fbox, boxrule=1pt, pad at break*=1mm,colback=cellbackground, colframe=cellborder]
\prompt{In}{incolor}{229}{\boxspacing}
\begin{Verbatim}[commandchars=\\\{\}]
\PY{n}{g} \PY{o}{=} \PY{n}{sns}\PY{o}{.}\PY{n}{PairGrid}\PY{p}{(}\PY{n}{tbl\PYZus{}red}\PY{p}{[}\PY{l+m+mi}{1}\PY{p}{]}\PY{p}{,} \PY{n}{diag\PYZus{}sharey}\PY{o}{=}\PY{k+kc}{False}\PY{p}{,} \PY{n}{corner}\PY{o}{=}\PY{k+kc}{True}\PY{p}{)}

\PY{n}{g}\PY{o}{.}\PY{n}{map\PYZus{}diag}\PY{p}{(}\PY{n}{sns}\PY{o}{.}\PY{n}{kdeplot}\PY{p}{)}
\PY{n}{g}\PY{o}{.}\PY{n}{map\PYZus{}lower}\PY{p}{(}\PY{n}{sns}\PY{o}{.}\PY{n}{scatterplot}\PY{p}{)}
\end{Verbatim}
\end{tcolorbox}

            \begin{tcolorbox}[breakable, size=fbox, boxrule=.5pt, pad at break*=1mm, opacityfill=0]
\prompt{Out}{outcolor}{229}{\boxspacing}
\begin{Verbatim}[commandchars=\\\{\}]
<seaborn.axisgrid.PairGrid at 0x1d1a733fd70>
\end{Verbatim}
\end{tcolorbox}
        
    \begin{center}
    \adjustimage{max size={0.9\linewidth}{0.9\paperheight}}{figs/bvp_kM_en/output_202_1.png}
    \end{center}
    { \hspace*{\fill} \\}
    
    \begin{tcolorbox}[breakable, size=fbox, boxrule=1pt, pad at break*=1mm,colback=cellbackground, colframe=cellborder]
\prompt{In}{incolor}{230}{\boxspacing}
\begin{Verbatim}[commandchars=\\\{\}]
\PY{n}{gind} \PY{o}{=} \PY{l+m+mi}{1}
\end{Verbatim}
\end{tcolorbox}

    \begin{tcolorbox}[breakable, size=fbox, boxrule=1pt, pad at break*=1mm,colback=cellbackground, colframe=cellborder]
\prompt{In}{incolor}{231}{\boxspacing}
\begin{Verbatim}[commandchars=\\\{\}]
\PY{n}{dkn\PYZus{}arr} \PY{o}{=} \PY{n}{np}\PY{o}{.}\PY{n}{zeros}\PY{p}{(}\PY{p}{(}\PY{n}{grid\PYZus{}size}\PY{p}{,} \PY{n}{grid\PYZus{}size}\PY{p}{)}\PY{p}{)}
\PY{n}{dMn\PYZus{}arr} \PY{o}{=} \PY{n}{np}\PY{o}{.}\PY{n}{zeros}\PY{p}{(}\PY{p}{(}\PY{n}{grid\PYZus{}size}\PY{p}{,} \PY{n}{grid\PYZus{}size}\PY{p}{)}\PY{p}{)}

\PY{k}{for} \PY{n}{r}\PY{p}{,} \PY{n}{rr} \PY{o+ow}{in} \PY{n+nb}{enumerate}\PY{p}{(}\PY{n}{dk\PYZus{}g}\PY{p}{[}\PY{n}{gind}\PY{p}{]}\PY{p}{)}\PY{p}{:}
    \PY{k}{for} \PY{n}{b}\PY{p}{,} \PY{n}{bb} \PY{o+ow}{in} \PY{n+nb}{enumerate}\PY{p}{(}\PY{n}{mu\PYZus{}g}\PY{p}{[}\PY{n}{gind}\PY{p}{]}\PY{p}{)}\PY{p}{:}
        \PY{n}{tbl\PYZus{}sel} \PY{o}{=} \PY{n}{tbl\PYZus{}red}\PY{p}{[}\PY{n}{gind}\PY{p}{]}\PY{p}{[}\PY{p}{(}\PY{n}{tbl\PYZus{}red}\PY{p}{[}\PY{n}{gind}\PY{p}{]}\PY{p}{[}\PY{l+s+s2}{\PYZdq{}}\PY{l+s+s2}{dk}\PY{l+s+s2}{\PYZdq{}}\PY{p}{]}\PY{o}{==}\PY{n}{rr}\PY{p}{)} \PY{o}{\PYZam{}}\PY{p}{(}\PY{n}{tbl\PYZus{}red}\PY{p}{[}\PY{n}{gind}\PY{p}{]}\PY{p}{[}\PY{l+s+s2}{\PYZdq{}}\PY{l+s+s2}{mu}\PY{l+s+s2}{\PYZdq{}}\PY{p}{]}\PY{o}{==}\PY{n}{bb}\PY{p}{)}\PY{p}{]}
        \PY{n}{dkn\PYZus{}arr}\PY{p}{[}\PY{n}{r}\PY{p}{,}\PY{n}{b}\PY{p}{]} \PY{o}{=} \PY{n}{tbl\PYZus{}sel}\PY{p}{[}\PY{l+s+s2}{\PYZdq{}}\PY{l+s+s2}{dkn}\PY{l+s+s2}{\PYZdq{}}\PY{p}{]}\PY{o}{.}\PY{n}{to\PYZus{}numpy}\PY{p}{(}\PY{p}{)}\PY{p}{[}\PY{l+m+mi}{0}\PY{p}{]}
        \PY{n}{dMn\PYZus{}arr}\PY{p}{[}\PY{n}{r}\PY{p}{,}\PY{n}{b}\PY{p}{]} \PY{o}{=} \PY{n}{tbl\PYZus{}sel}\PY{p}{[}\PY{l+s+s2}{\PYZdq{}}\PY{l+s+s2}{dMn}\PY{l+s+s2}{\PYZdq{}}\PY{p}{]}\PY{o}{.}\PY{n}{to\PYZus{}numpy}\PY{p}{(}\PY{p}{)}\PY{p}{[}\PY{l+m+mi}{0}\PY{p}{]}
\end{Verbatim}
\end{tcolorbox}

    \begin{tcolorbox}[breakable, size=fbox, boxrule=1pt, pad at break*=1mm,colback=cellbackground, colframe=cellborder]
\prompt{In}{incolor}{232}{\boxspacing}
\begin{Verbatim}[commandchars=\\\{\}]
\PY{n}{dk} \PY{o}{=} \PY{n}{dk\PYZus{}g}\PY{p}{[}\PY{l+m+mi}{1}\PY{p}{]}
\PY{n}{mu} \PY{o}{=} \PY{n}{mu\PYZus{}g}\PY{p}{[}\PY{l+m+mi}{1}\PY{p}{]}
\end{Verbatim}
\end{tcolorbox}

    This dataset only contains the \(\mu\) dependence for the change in
\(k\) from the linear case with respect to the small-scale oscillations.
So lets see what this dependence looks like. We can visualize this data
either with stacked plots of \(dkn\) vs \(\mu\) or stacked plots of
\(dkn\) vs \(dk\).

    \begin{tcolorbox}[breakable, size=fbox, boxrule=1pt, pad at break*=1mm,colback=cellbackground, colframe=cellborder]
\prompt{In}{incolor}{ }{\boxspacing}
\begin{Verbatim}[commandchars=\\\{\}]
\PY{c+c1}{\PYZsh{} Here is the plot of dkn vs dk}
\PY{c+c1}{\PYZsh{} Note that dkn (difference between nonlinear and linear fast oscillation wave number) }
\PY{c+c1}{\PYZsh{} here is actually dkn / 10\PYZca{}4 and dk (difference in linear fast oscillation wave number }
\PY{c+c1}{\PYZsh{} from the wave number which produces a whole number of fast cycles in the linear case)}
\PY{c+c1}{\PYZsh{} is 2 * dk * R\PYZus{}max.}

\PY{n}{jj} \PY{o}{=} \PY{p}{[}\PY{l+m+mi}{10}\PY{p}{,} \PY{l+m+mi}{20}\PY{p}{,} \PY{l+m+mi}{30}\PY{p}{,} \PY{l+m+mi}{40}\PY{p}{,} \PY{l+m+mi}{50}\PY{p}{,} \PY{l+m+mi}{60}\PY{p}{,} \PY{l+m+mi}{70}\PY{p}{,} \PY{l+m+mi}{80}\PY{p}{,} \PY{l+m+mi}{90}\PY{p}{]}

\PY{n}{ilim} \PY{o}{=} \PY{n+nb}{len}\PY{p}{(}\PY{n}{jj}\PY{p}{)}

\PY{n}{fig}\PY{p}{,} \PY{n}{ax} \PY{o}{=} \PY{n}{plt}\PY{o}{.}\PY{n}{subplots}\PY{p}{(}\PY{p}{)}
\PY{n}{plt}\PY{o}{.}\PY{n}{title}\PY{p}{(}\PY{l+s+sa}{f}\PY{l+s+s2}{\PYZdq{}}\PY{l+s+s2}{Plot of the Change in k Due to Nonlinearity vs Linear k}\PY{l+s+se}{\PYZbs{}n}\PY{l+s+s2}{Variation For Various Values of }\PY{l+s+se}{\PYZbs{}u03BC}\PY{l+s+s2}{\PYZdq{}}\PY{p}{)}
\PY{n}{ax}\PY{o}{.}\PY{n}{set\PYZus{}ylabel}\PY{p}{(}\PY{l+s+s1}{\PYZsq{}}\PY{l+s+s1}{Change in k due to nonlinearity (dkn) (10}\PY{l+s+se}{\PYZbs{}u2074}\PY{l+s+s1}{ m}\PY{l+s+se}{\PYZbs{}u207b}\PY{l+s+se}{\PYZbs{}u00b9}\PY{l+s+s1}{)}\PY{l+s+s1}{\PYZsq{}}\PY{p}{)}
\PY{n}{ax}\PY{o}{.}\PY{n}{set\PYZus{}xlabel}\PY{p}{(}\PY{l+s+s2}{\PYZdq{}}\PY{l+s+s2}{Small\PYZhy{}scale variation of linear wave number k (dk (}\PY{l+s+s2}{\PYZpc{}}\PY{l+s+s2}{))}\PY{l+s+s2}{\PYZdq{}}\PY{p}{)}

\PY{k}{for} \PY{n}{i} \PY{o+ow}{in} \PY{n+nb}{range}\PY{p}{(}\PY{n}{ilim}\PY{p}{)}\PY{p}{:}
    \PY{n}{ax}\PY{o}{.}\PY{n}{plot}\PY{p}{(}\PY{n}{dk}\PY{p}{[}\PY{l+m+mi}{2}\PY{p}{:}\PY{p}{]}\PY{p}{,} \PY{n}{dkn\PYZus{}arr}\PY{p}{[}\PY{l+m+mi}{2}\PY{p}{:}\PY{p}{,} \PY{n}{jj}\PY{p}{[}\PY{n}{i}\PY{p}{]}\PY{p}{]}\PY{o}{/}\PY{l+m+mi}{10000}\PY{p}{,} \PY{n}{label} \PY{o}{=} \PY{l+s+sa}{f}\PY{l+s+s1}{\PYZsq{}}\PY{l+s+se}{\PYZbs{}u03BC}\PY{l+s+s1}{ = }\PY{l+s+si}{\PYZob{}}\PY{n}{mu}\PY{p}{[}\PY{n}{jj}\PY{p}{[}\PY{n}{i}\PY{p}{]}\PY{p}{]}\PY{l+s+si}{:}\PY{l+s+s1}{.2e}\PY{l+s+si}{\PYZcb{}}\PY{l+s+s1}{\PYZsq{}}\PY{p}{)}

\PY{n}{plt}\PY{o}{.}\PY{n}{legend}\PY{p}{(}\PY{n}{loc}\PY{o}{=}\PY{l+s+s1}{\PYZsq{}}\PY{l+s+s1}{center left}\PY{l+s+s1}{\PYZsq{}}\PY{p}{,} \PY{n}{bbox\PYZus{}to\PYZus{}anchor}\PY{o}{=}\PY{p}{(}\PY{l+m+mi}{0}\PY{p}{,} \PY{l+m+mf}{0.3}\PY{p}{)}\PY{p}{)}

\PY{n}{plt}\PY{o}{.}\PY{n}{show}\PY{p}{(}\PY{p}{)}
\end{Verbatim}
\end{tcolorbox}

    \begin{center}
    \adjustimage{max size={0.9\linewidth}{0.9\paperheight}}{figs/bvp_kM_en/output_207_0.png}
    \end{center}
    { \hspace*{\fill} \\}
    
    \begin{tcolorbox}[breakable, size=fbox, boxrule=1pt, pad at break*=1mm,colback=cellbackground, colframe=cellborder]
\prompt{In}{incolor}{239}{\boxspacing}
\begin{Verbatim}[commandchars=\\\{\}]
\PY{c+c1}{\PYZsh{} Now for plots of dkn / 10\PYZca{}4 vs mu * 10 \PYZca{}8 }
\PY{c+c1}{\PYZsh{} I have split this up into 4 graphs, as something interesting happens.}

\PY{n}{jj} \PY{o}{=} \PY{p}{[}\PY{l+m+mi}{46}\PY{p}{,} \PY{l+m+mi}{48}\PY{p}{,} \PY{l+m+mi}{51}\PY{p}{,} \PY{l+m+mi}{55}\PY{p}{,} \PY{l+m+mi}{61}\PY{p}{,} \PY{l+m+mi}{77}\PY{p}{,} \PY{l+m+mi}{90}\PY{p}{]}

\PY{n}{ilim} \PY{o}{=} \PY{n+nb}{len}\PY{p}{(}\PY{n}{jj}\PY{p}{)}

\PY{n}{fig}\PY{p}{,} \PY{n}{ax} \PY{o}{=} \PY{n}{plt}\PY{o}{.}\PY{n}{subplots}\PY{p}{(}\PY{p}{)}
\PY{n}{plt}\PY{o}{.}\PY{n}{title}\PY{p}{(}\PY{l+s+sa}{f}\PY{l+s+s2}{\PYZdq{}}\PY{l+s+s2}{Plot of the Change in k Due to Nonlinearity vs }\PY{l+s+se}{\PYZbs{}u03BC}\PY{l+s+se}{\PYZbs{}n}\PY{l+s+s2}{ For Various Values of dk}\PY{l+s+s2}{\PYZdq{}}\PY{p}{)}
\PY{n}{ax}\PY{o}{.}\PY{n}{set\PYZus{}ylabel}\PY{p}{(}\PY{l+s+s1}{\PYZsq{}}\PY{l+s+s1}{Change in k due to nonlinearity (dkn) (10}\PY{l+s+se}{\PYZbs{}u2074}\PY{l+s+s1}{ m}\PY{l+s+se}{\PYZbs{}u207b}\PY{l+s+se}{\PYZbs{}u00b9}\PY{l+s+s1}{)}\PY{l+s+s1}{\PYZsq{}}\PY{p}{)}
\PY{n}{ax}\PY{o}{.}\PY{n}{set\PYZus{}xlabel}\PY{p}{(}\PY{l+s+s2}{\PYZdq{}}\PY{l+s+s2}{Nonlinearity coefficient }\PY{l+s+se}{\PYZbs{}u03BC}\PY{l+s+s2}{ * 10}\PY{l+s+se}{\PYZbs{}u2078}\PY{l+s+s2}{\PYZdq{}}\PY{p}{)}

\PY{k}{for} \PY{n}{i} \PY{o+ow}{in} \PY{n+nb}{range}\PY{p}{(}\PY{n}{ilim}\PY{p}{)}\PY{p}{:}
    \PY{n}{ax}\PY{o}{.}\PY{n}{plot}\PY{p}{(}\PY{n}{mu}\PY{o}{*}\PY{l+m+mf}{1e8}\PY{p}{,} \PY{n}{dkn\PYZus{}arr}\PY{p}{[}\PY{n}{jj}\PY{p}{[}\PY{n}{i}\PY{p}{]}\PY{p}{,}\PY{p}{:}\PY{p}{]}\PY{o}{/}\PY{l+m+mi}{10000}\PY{p}{,} \PY{n}{label} \PY{o}{=} \PY{l+s+sa}{f}\PY{l+s+s2}{\PYZdq{}}\PY{l+s+s2}{dk = }\PY{l+s+si}{\PYZob{}}\PY{n}{dk}\PY{p}{[}\PY{n}{jj}\PY{p}{[}\PY{n}{i}\PY{p}{]}\PY{p}{]}\PY{l+s+si}{:}\PY{l+s+s2}{.3f}\PY{l+s+si}{\PYZcb{}}\PY{l+s+s2}{(*0.5*10}\PY{l+s+se}{\PYZbs{}u2076}\PY{l+s+s2}{ m}\PY{l+s+se}{\PYZbs{}u207b}\PY{l+s+se}{\PYZbs{}u00b9}\PY{l+s+s2}{)}\PY{l+s+s2}{\PYZdq{}}\PY{p}{)}

\PY{n}{plt}\PY{o}{.}\PY{n}{legend}\PY{p}{(}\PY{n}{loc}\PY{o}{=}\PY{l+s+s1}{\PYZsq{}}\PY{l+s+s1}{center left}\PY{l+s+s1}{\PYZsq{}}\PY{p}{,} \PY{n}{bbox\PYZus{}to\PYZus{}anchor}\PY{o}{=}\PY{p}{(}\PY{l+m+mi}{1}\PY{p}{,} \PY{l+m+mf}{0.5}\PY{p}{)}\PY{p}{)}

\PY{n}{plt}\PY{o}{.}\PY{n}{show}\PY{p}{(}\PY{p}{)}
\end{Verbatim}
\end{tcolorbox}

    \begin{center}
    \adjustimage{max size={0.9\linewidth}{0.9\paperheight}}{figs/bvp_kM_en/output_208_0.png}
    \end{center}
    { \hspace*{\fill} \\}
    
    \begin{tcolorbox}[breakable, size=fbox, boxrule=1pt, pad at break*=1mm,colback=cellbackground, colframe=cellborder]
\prompt{In}{incolor}{240}{\boxspacing}
\begin{Verbatim}[commandchars=\\\{\}]
\PY{c+c1}{\PYZsh{} More plots of dkn / 10\PYZca{}4 vs mu * 10 \PYZca{}8 }

\PY{n}{jj} \PY{o}{=} \PY{p}{[}\PY{l+m+mi}{46}\PY{p}{,} \PY{l+m+mi}{49}\PY{p}{,} \PY{l+m+mi}{52}\PY{p}{,} \PY{l+m+mi}{55}\PY{p}{,} \PY{l+m+mi}{58}\PY{p}{,} \PY{l+m+mi}{61}\PY{p}{]}

\PY{n}{ilim} \PY{o}{=} \PY{n+nb}{len}\PY{p}{(}\PY{n}{jj}\PY{p}{)}

\PY{n}{fig}\PY{p}{,} \PY{n}{ax} \PY{o}{=} \PY{n}{plt}\PY{o}{.}\PY{n}{subplots}\PY{p}{(}\PY{p}{)}
\PY{n}{plt}\PY{o}{.}\PY{n}{title}\PY{p}{(}\PY{l+s+sa}{f}\PY{l+s+s2}{\PYZdq{}}\PY{l+s+s2}{Plot of the Change in k Due to Nonlinearity vs }\PY{l+s+se}{\PYZbs{}u03BC}\PY{l+s+se}{\PYZbs{}n}\PY{l+s+s2}{ For Various Values of dk}\PY{l+s+s2}{\PYZdq{}}\PY{p}{)}
\PY{n}{ax}\PY{o}{.}\PY{n}{set\PYZus{}ylabel}\PY{p}{(}\PY{l+s+s1}{\PYZsq{}}\PY{l+s+s1}{Change in k due to nonlinearity (dkn) (10}\PY{l+s+se}{\PYZbs{}u2074}\PY{l+s+s1}{ m}\PY{l+s+se}{\PYZbs{}u207b}\PY{l+s+se}{\PYZbs{}u00b9}\PY{l+s+s1}{)}\PY{l+s+s1}{\PYZsq{}}\PY{p}{)}
\PY{n}{ax}\PY{o}{.}\PY{n}{set\PYZus{}xlabel}\PY{p}{(}\PY{l+s+s2}{\PYZdq{}}\PY{l+s+s2}{Nonlinearity coefficient }\PY{l+s+se}{\PYZbs{}u03BC}\PY{l+s+s2}{ * 10}\PY{l+s+se}{\PYZbs{}u2078}\PY{l+s+s2}{\PYZdq{}}\PY{p}{)}

\PY{k}{for} \PY{n}{i} \PY{o+ow}{in} \PY{n+nb}{range}\PY{p}{(}\PY{n}{ilim}\PY{p}{)}\PY{p}{:}
    \PY{n}{ax}\PY{o}{.}\PY{n}{plot}\PY{p}{(}\PY{n}{mu}\PY{o}{*}\PY{l+m+mf}{1e8}\PY{p}{,} \PY{n}{dkn\PYZus{}arr}\PY{p}{[}\PY{n}{jj}\PY{p}{[}\PY{n}{i}\PY{p}{]}\PY{p}{,}\PY{p}{:}\PY{p}{]}\PY{o}{/}\PY{l+m+mi}{10000}\PY{p}{,} \PY{n}{label} \PY{o}{=} \PY{l+s+sa}{f}\PY{l+s+s2}{\PYZdq{}}\PY{l+s+s2}{dk = }\PY{l+s+si}{\PYZob{}}\PY{n}{dk}\PY{p}{[}\PY{n}{jj}\PY{p}{[}\PY{n}{i}\PY{p}{]}\PY{p}{]}\PY{l+s+si}{:}\PY{l+s+s2}{.3f}\PY{l+s+si}{\PYZcb{}}\PY{l+s+s2}{(*0.5*10}\PY{l+s+se}{\PYZbs{}u2076}\PY{l+s+s2}{ m}\PY{l+s+se}{\PYZbs{}u207b}\PY{l+s+se}{\PYZbs{}u00b9}\PY{l+s+s2}{)}\PY{l+s+s2}{\PYZdq{}}\PY{p}{)}

\PY{n}{plt}\PY{o}{.}\PY{n}{legend}\PY{p}{(}\PY{n}{loc}\PY{o}{=}\PY{l+s+s1}{\PYZsq{}}\PY{l+s+s1}{center left}\PY{l+s+s1}{\PYZsq{}}\PY{p}{,} \PY{n}{bbox\PYZus{}to\PYZus{}anchor}\PY{o}{=}\PY{p}{(}\PY{l+m+mi}{1}\PY{p}{,} \PY{l+m+mf}{0.5}\PY{p}{)}\PY{p}{)}

\PY{n}{plt}\PY{o}{.}\PY{n}{show}\PY{p}{(}\PY{p}{)}
\end{Verbatim}
\end{tcolorbox}

    \begin{center}
    \adjustimage{max size={0.9\linewidth}{0.9\paperheight}}{figs/bvp_kM_en/output_209_0.png}
    \end{center}
    { \hspace*{\fill} \\}
    
    \begin{tcolorbox}[breakable, size=fbox, boxrule=1pt, pad at break*=1mm,colback=cellbackground, colframe=cellborder]
\prompt{In}{incolor}{241}{\boxspacing}
\begin{Verbatim}[commandchars=\\\{\}]
\PY{n+nb}{print}\PY{p}{(}\PY{l+s+sa}{f}\PY{l+s+s2}{\PYZdq{}}\PY{l+s+s2}{Our jump (or drop) first appears at dk = }\PY{l+s+si}{\PYZob{}}\PY{n}{dk}\PY{p}{[}\PY{l+m+mi}{48}\PY{p}{]}\PY{l+s+si}{\PYZcb{}}\PY{l+s+s2}{\PYZdq{}}\PY{p}{)}
\end{Verbatim}
\end{tcolorbox}

    \begin{Verbatim}[commandchars=\\\{\}]
Our jump (or drop) first appears at dk = 1.4545454545454546
    \end{Verbatim}

    \begin{tcolorbox}[breakable, size=fbox, boxrule=1pt, pad at break*=1mm,colback=cellbackground, colframe=cellborder]
\prompt{In}{incolor}{242}{\boxspacing}
\begin{Verbatim}[commandchars=\\\{\}]
\PY{c+c1}{\PYZsh{} These are for earlier dk values, where there is no drop.}

\PY{n}{jj} \PY{o}{=} \PY{p}{[}\PY{l+m+mi}{5}\PY{p}{,} \PY{l+m+mi}{10}\PY{p}{,} \PY{l+m+mi}{20}\PY{p}{,} \PY{l+m+mi}{32}\PY{p}{,} \PY{l+m+mi}{38}\PY{p}{]}

\PY{n}{ilim} \PY{o}{=} \PY{n+nb}{len}\PY{p}{(}\PY{n}{jj}\PY{p}{)}

\PY{n}{fig}\PY{p}{,} \PY{n}{ax} \PY{o}{=} \PY{n}{plt}\PY{o}{.}\PY{n}{subplots}\PY{p}{(}\PY{p}{)}
\PY{n}{plt}\PY{o}{.}\PY{n}{title}\PY{p}{(}\PY{l+s+sa}{f}\PY{l+s+s2}{\PYZdq{}}\PY{l+s+s2}{Plot of the Change in k Due to Nonlinearity vs }\PY{l+s+se}{\PYZbs{}u03BC}\PY{l+s+se}{\PYZbs{}n}\PY{l+s+s2}{ For Various Values of dk}\PY{l+s+s2}{\PYZdq{}}\PY{p}{)}
\PY{n}{ax}\PY{o}{.}\PY{n}{set\PYZus{}ylabel}\PY{p}{(}\PY{l+s+s1}{\PYZsq{}}\PY{l+s+s1}{Change in k due to nonlinearity (dkn) (10}\PY{l+s+se}{\PYZbs{}u2074}\PY{l+s+s1}{ m}\PY{l+s+se}{\PYZbs{}u207b}\PY{l+s+se}{\PYZbs{}u00b9}\PY{l+s+s1}{)}\PY{l+s+s1}{\PYZsq{}}\PY{p}{)}
\PY{n}{ax}\PY{o}{.}\PY{n}{set\PYZus{}xlabel}\PY{p}{(}\PY{l+s+s2}{\PYZdq{}}\PY{l+s+s2}{Nonlinearity coefficient }\PY{l+s+se}{\PYZbs{}u03BC}\PY{l+s+s2}{ * 10}\PY{l+s+se}{\PYZbs{}u2078}\PY{l+s+s2}{\PYZdq{}}\PY{p}{)}

\PY{k}{for} \PY{n}{i} \PY{o+ow}{in} \PY{n+nb}{range}\PY{p}{(}\PY{n}{ilim}\PY{p}{)}\PY{p}{:}
    \PY{n}{ax}\PY{o}{.}\PY{n}{plot}\PY{p}{(}\PY{n}{mu}\PY{o}{*}\PY{l+m+mf}{1e8}\PY{p}{,} \PY{n}{dkn\PYZus{}arr}\PY{p}{[}\PY{n}{jj}\PY{p}{[}\PY{n}{i}\PY{p}{]}\PY{p}{,}\PY{p}{:}\PY{p}{]}\PY{o}{/}\PY{l+m+mi}{10000}\PY{p}{,} \PY{n}{label} \PY{o}{=} \PY{l+s+sa}{f}\PY{l+s+s2}{\PYZdq{}}\PY{l+s+s2}{dk = }\PY{l+s+si}{\PYZob{}}\PY{n}{dk}\PY{p}{[}\PY{n}{jj}\PY{p}{[}\PY{n}{i}\PY{p}{]}\PY{p}{]}\PY{l+s+si}{:}\PY{l+s+s2}{.3f}\PY{l+s+si}{\PYZcb{}}\PY{l+s+s2}{(*0.5*10}\PY{l+s+se}{\PYZbs{}u2076}\PY{l+s+s2}{ m}\PY{l+s+se}{\PYZbs{}u207b}\PY{l+s+se}{\PYZbs{}u00b9}\PY{l+s+s2}{)}\PY{l+s+s2}{\PYZdq{}}\PY{p}{)}

\PY{n}{plt}\PY{o}{.}\PY{n}{legend}\PY{p}{(}\PY{n}{loc}\PY{o}{=}\PY{l+s+s1}{\PYZsq{}}\PY{l+s+s1}{center left}\PY{l+s+s1}{\PYZsq{}}\PY{p}{,} \PY{n}{bbox\PYZus{}to\PYZus{}anchor}\PY{o}{=}\PY{p}{(}\PY{l+m+mi}{1}\PY{p}{,} \PY{l+m+mf}{0.5}\PY{p}{)}\PY{p}{)}

\PY{n}{plt}\PY{o}{.}\PY{n}{show}\PY{p}{(}\PY{p}{)}
\end{Verbatim}
\end{tcolorbox}

    \begin{center}
    \adjustimage{max size={0.9\linewidth}{0.9\paperheight}}{figs/bvp_kM_en/output_211_0.png}
    \end{center}
    { \hspace*{\fill} \\}
    
    \begin{tcolorbox}[breakable, size=fbox, boxrule=1pt, pad at break*=1mm,colback=cellbackground, colframe=cellborder]
\prompt{In}{incolor}{243}{\boxspacing}
\begin{Verbatim}[commandchars=\\\{\}]
\PY{c+c1}{\PYZsh{} These occur for large dk numbers, where the drop in signal has already occured.}

\PY{n}{jj} \PY{o}{=} \PY{p}{[}\PY{l+m+mi}{73}\PY{p}{,} \PY{l+m+mi}{76}\PY{p}{,} \PY{l+m+mi}{79}\PY{p}{,} \PY{l+m+mi}{82}\PY{p}{,} \PY{l+m+mi}{85}\PY{p}{,} \PY{l+m+mi}{88}\PY{p}{,} \PY{l+m+mi}{91}\PY{p}{]}

\PY{n}{ilim} \PY{o}{=} \PY{n+nb}{len}\PY{p}{(}\PY{n}{jj}\PY{p}{)}

\PY{n}{fig}\PY{p}{,} \PY{n}{ax} \PY{o}{=} \PY{n}{plt}\PY{o}{.}\PY{n}{subplots}\PY{p}{(}\PY{p}{)}
\PY{n}{plt}\PY{o}{.}\PY{n}{title}\PY{p}{(}\PY{l+s+sa}{f}\PY{l+s+s2}{\PYZdq{}}\PY{l+s+s2}{Plot of the Change in k Due to Nonlinearity vs }\PY{l+s+se}{\PYZbs{}u03BC}\PY{l+s+se}{\PYZbs{}n}\PY{l+s+s2}{ For Various Values of dk}\PY{l+s+s2}{\PYZdq{}}\PY{p}{)}
\PY{n}{ax}\PY{o}{.}\PY{n}{set\PYZus{}ylabel}\PY{p}{(}\PY{l+s+s1}{\PYZsq{}}\PY{l+s+s1}{Change in k due to nonlinearity (dkn) (10}\PY{l+s+se}{\PYZbs{}u2074}\PY{l+s+s1}{ m}\PY{l+s+se}{\PYZbs{}u207b}\PY{l+s+se}{\PYZbs{}u00b9}\PY{l+s+s1}{)}\PY{l+s+s1}{\PYZsq{}}\PY{p}{)}
\PY{n}{ax}\PY{o}{.}\PY{n}{set\PYZus{}xlabel}\PY{p}{(}\PY{l+s+s2}{\PYZdq{}}\PY{l+s+s2}{Nonlinearity coefficient }\PY{l+s+se}{\PYZbs{}u03BC}\PY{l+s+s2}{ * 10}\PY{l+s+se}{\PYZbs{}u2078}\PY{l+s+s2}{\PYZdq{}}\PY{p}{)}

\PY{k}{for} \PY{n}{i} \PY{o+ow}{in} \PY{n+nb}{range}\PY{p}{(}\PY{n}{ilim}\PY{p}{)}\PY{p}{:}
    \PY{n}{ax}\PY{o}{.}\PY{n}{plot}\PY{p}{(}\PY{n}{mu}\PY{o}{*}\PY{l+m+mf}{1e8}\PY{p}{,} \PY{n}{dkn\PYZus{}arr}\PY{p}{[}\PY{n}{jj}\PY{p}{[}\PY{n}{i}\PY{p}{]}\PY{p}{,}\PY{p}{:}\PY{p}{]}\PY{o}{/}\PY{l+m+mf}{1e4}\PY{p}{,} \PY{n}{label} \PY{o}{=} \PY{l+s+sa}{f}\PY{l+s+s2}{\PYZdq{}}\PY{l+s+s2}{dk = }\PY{l+s+si}{\PYZob{}}\PY{n}{dk}\PY{p}{[}\PY{n}{jj}\PY{p}{[}\PY{n}{i}\PY{p}{]}\PY{p}{]}\PY{l+s+si}{:}\PY{l+s+s2}{.3f}\PY{l+s+si}{\PYZcb{}}\PY{l+s+s2}{(*0.5*10}\PY{l+s+se}{\PYZbs{}u2076}\PY{l+s+s2}{ m}\PY{l+s+se}{\PYZbs{}u207b}\PY{l+s+se}{\PYZbs{}u00b9}\PY{l+s+s2}{)}\PY{l+s+s2}{\PYZdq{}}\PY{p}{)}

\PY{n}{plt}\PY{o}{.}\PY{n}{legend}\PY{p}{(}\PY{n}{loc}\PY{o}{=}\PY{l+s+s1}{\PYZsq{}}\PY{l+s+s1}{center left}\PY{l+s+s1}{\PYZsq{}}\PY{p}{,} \PY{n}{bbox\PYZus{}to\PYZus{}anchor}\PY{o}{=}\PY{p}{(}\PY{l+m+mi}{1}\PY{p}{,} \PY{l+m+mf}{0.5}\PY{p}{)}\PY{p}{)}

\PY{n}{plt}\PY{o}{.}\PY{n}{show}\PY{p}{(}\PY{p}{)}
\end{Verbatim}
\end{tcolorbox}

    \begin{center}
    \adjustimage{max size={0.9\linewidth}{0.9\paperheight}}{figs/bvp_kM_en/output_212_0.png}
    \end{center}
    { \hspace*{\fill} \\}
    
    As dk increases from 0 to 1, the whole graph increases with dk. Around k
= 1.45, a drop appears for lower \(\mu\) values. This means that \(\mu\)
causes a phase shift as well as an amplification. To remove the effects
of \(dk\), I can grab the value with the large drop as my \(\mu\) value.
Lets try to get the maximum \(dkn\) and minimum \(dkn\) for every value
of \(\mu\).

    \begin{tcolorbox}[breakable, size=fbox, boxrule=1pt, pad at break*=1mm,colback=cellbackground, colframe=cellborder]
\prompt{In}{incolor}{245}{\boxspacing}
\begin{Verbatim}[commandchars=\\\{\}]
\PY{n}{max\PYZus{}for\PYZus{}mu} \PY{o}{=} \PY{o}{\PYZhy{}}\PY{n}{np}\PY{o}{.}\PY{n}{max}\PY{p}{(}\PY{n}{dkn\PYZus{}arr}\PY{p}{,} \PY{n}{axis}\PY{o}{=}\PY{l+m+mi}{0}\PY{p}{)}\PY{o}{/}\PY{l+m+mf}{1e5}
\PY{n}{min\PYZus{}for\PYZus{}mu} \PY{o}{=} \PY{o}{\PYZhy{}}\PY{n}{np}\PY{o}{.}\PY{n}{min}\PY{p}{(}\PY{n}{dkn\PYZus{}arr}\PY{p}{,} \PY{n}{axis}\PY{o}{=}\PY{l+m+mi}{0}\PY{p}{)}\PY{o}{/}\PY{l+m+mf}{1e5}
\end{Verbatim}
\end{tcolorbox}

    I have added in a rescaling of \(dkn\) of \(10^{-5}\) to remove the
zeros from the end. I need to not forget about this rescaling later.

    \begin{tcolorbox}[breakable, size=fbox, boxrule=1pt, pad at break*=1mm,colback=cellbackground, colframe=cellborder]
\prompt{In}{incolor}{252}{\boxspacing}
\begin{Verbatim}[commandchars=\\\{\}]
\PY{c+c1}{\PYZsh{} Now we plot the greatest and least values of \PYZhy{} dkn / 10\PYZca{}5 vs mu * 10\PYZca{}8 for all values of k}

\PY{n}{fig}\PY{p}{,} \PY{n}{ax} \PY{o}{=} \PY{n}{plt}\PY{o}{.}\PY{n}{subplots}\PY{p}{(}\PY{p}{)}
\PY{n}{plt}\PY{o}{.}\PY{n}{title}\PY{p}{(}\PY{l+s+sa}{f}\PY{l+s+s2}{\PYZdq{}}\PY{l+s+s2}{The Limits of the Change in k Due to Nonlinearity vs }\PY{l+s+se}{\PYZbs{}u03BC}\PY{l+s+s2}{\PYZdq{}}\PY{p}{)}
\PY{n}{ax}\PY{o}{.}\PY{n}{set\PYZus{}ylabel}\PY{p}{(}\PY{l+s+s1}{\PYZsq{}}\PY{l+s+s1}{Change in k due to nonlinearity (\PYZhy{}dkn) (\PYZhy{}10}\PY{l+s+se}{\PYZbs{}u2075}\PY{l+s+s1}{ m}\PY{l+s+se}{\PYZbs{}u207b}\PY{l+s+se}{\PYZbs{}u00b9}\PY{l+s+s1}{)}\PY{l+s+s1}{\PYZsq{}}\PY{p}{)}
\PY{n}{ax}\PY{o}{.}\PY{n}{set\PYZus{}xlabel}\PY{p}{(}\PY{l+s+s2}{\PYZdq{}}\PY{l+s+s2}{Nonlinearity coefficient }\PY{l+s+se}{\PYZbs{}u03BC}\PY{l+s+s2}{ * 10}\PY{l+s+se}{\PYZbs{}u2078}\PY{l+s+s2}{\PYZdq{}}\PY{p}{)}

\PY{n}{ax}\PY{o}{.}\PY{n}{plot}\PY{p}{(}\PY{n}{mu}\PY{o}{*}\PY{l+m+mf}{1e8}\PY{p}{,} \PY{n}{min\PYZus{}for\PYZus{}mu}\PY{p}{,} \PY{n}{label}\PY{o}{=}\PY{l+s+s2}{\PYZdq{}}\PY{l+s+s2}{maximum \PYZhy{}dkn * 10}\PY{l+s+se}{\PYZbs{}u207b}\PY{l+s+se}{\PYZbs{}u2075}\PY{l+s+s2}{\PYZdq{}}\PY{p}{)}
\PY{n}{ax}\PY{o}{.}\PY{n}{plot}\PY{p}{(}\PY{n}{mu}\PY{o}{*}\PY{l+m+mf}{1e8}\PY{p}{,} \PY{n}{max\PYZus{}for\PYZus{}mu}\PY{p}{,} \PY{n}{label}\PY{o}{=}\PY{l+s+s2}{\PYZdq{}}\PY{l+s+s2}{minimum \PYZhy{}dkn * 10}\PY{l+s+se}{\PYZbs{}u207b}\PY{l+s+se}{\PYZbs{}u2075}\PY{l+s+s2}{\PYZdq{}}\PY{p}{)}

\PY{n}{plt}\PY{o}{.}\PY{n}{legend}\PY{p}{(}\PY{n}{loc}\PY{o}{=}\PY{l+s+s1}{\PYZsq{}}\PY{l+s+s1}{center left}\PY{l+s+s1}{\PYZsq{}}\PY{p}{,} \PY{n}{bbox\PYZus{}to\PYZus{}anchor}\PY{o}{=}\PY{p}{(}\PY{l+m+mf}{0.5}\PY{p}{,} \PY{l+m+mf}{0.5}\PY{p}{)}\PY{p}{)}

\PY{n}{plt}\PY{o}{.}\PY{n}{show}\PY{p}{(}\PY{p}{)}
\end{Verbatim}
\end{tcolorbox}

    \begin{center}
    \adjustimage{max size={0.9\linewidth}{0.9\paperheight}}{figs/bvp_kM_en/output_216_0.png}
    \end{center}
    { \hspace*{\fill} \\}
    
    \begin{tcolorbox}[breakable, size=fbox, boxrule=1pt, pad at break*=1mm,colback=cellbackground, colframe=cellborder]
\prompt{In}{incolor}{259}{\boxspacing}
\begin{Verbatim}[commandchars=\\\{\}]
\PY{c+c1}{\PYZsh{} Again, but with logrithmic scaling, which makes power law scaling appear as a straight line.}
\PY{c+c1}{\PYZsh{} This is ln(\PYZhy{} dkn/10\PYZca{}5) vs ln(mu)}

\PY{n}{fig}\PY{p}{,} \PY{n}{ax} \PY{o}{=} \PY{n}{plt}\PY{o}{.}\PY{n}{subplots}\PY{p}{(}\PY{p}{)}
\PY{n}{plt}\PY{o}{.}\PY{n}{title}\PY{p}{(}\PY{l+s+sa}{f}\PY{l+s+s2}{\PYZdq{}}\PY{l+s+s2}{The Limits of the Change in k Due to Nonlinearity vs }\PY{l+s+se}{\PYZbs{}u03BC}\PY{l+s+s2}{\PYZdq{}}\PY{p}{)}
\PY{n}{ax}\PY{o}{.}\PY{n}{set\PYZus{}ylabel}\PY{p}{(}\PY{l+s+s1}{\PYZsq{}}\PY{l+s+s1}{Natural log of \PYZhy{}dkn * 10}\PY{l+s+se}{\PYZbs{}u2075}\PY{l+s+s1}{ (m}\PY{l+s+se}{\PYZbs{}u207b}\PY{l+s+se}{\PYZbs{}u00b9}\PY{l+s+s1}{)}\PY{l+s+s1}{\PYZsq{}}\PY{p}{)}
\PY{n}{ax}\PY{o}{.}\PY{n}{set\PYZus{}xlabel}\PY{p}{(}\PY{l+s+s2}{\PYZdq{}}\PY{l+s+s2}{Natural log of the nonlinearity coefficient }\PY{l+s+se}{\PYZbs{}u03BC}\PY{l+s+s2}{ * 10}\PY{l+s+se}{\PYZbs{}u2078}\PY{l+s+s2}{\PYZdq{}}\PY{p}{)}

\PY{n}{ax}\PY{o}{.}\PY{n}{plot}\PY{p}{(}\PY{n}{np}\PY{o}{.}\PY{n}{log}\PY{p}{(}\PY{n}{mu}\PY{p}{)}\PY{p}{,} \PY{n}{np}\PY{o}{.}\PY{n}{log}\PY{p}{(}\PY{n}{min\PYZus{}for\PYZus{}mu}\PY{p}{)}\PY{p}{,} \PY{n}{label}\PY{o}{=}\PY{l+s+s2}{\PYZdq{}}\PY{l+s+s2}{maximum \PYZhy{}dkn * 10}\PY{l+s+se}{\PYZbs{}u207b}\PY{l+s+se}{\PYZbs{}u2075}\PY{l+s+s2}{\PYZdq{}}\PY{p}{)}
\PY{n}{ax}\PY{o}{.}\PY{n}{plot}\PY{p}{(}\PY{n}{np}\PY{o}{.}\PY{n}{log}\PY{p}{(}\PY{n}{mu}\PY{p}{)}\PY{p}{,} \PY{n}{np}\PY{o}{.}\PY{n}{log}\PY{p}{(}\PY{n}{max\PYZus{}for\PYZus{}mu}\PY{p}{)}\PY{p}{,} \PY{n}{label}\PY{o}{=}\PY{l+s+s2}{\PYZdq{}}\PY{l+s+s2}{minimum \PYZhy{}dkn * 10}\PY{l+s+se}{\PYZbs{}u207b}\PY{l+s+se}{\PYZbs{}u2075}\PY{l+s+s2}{\PYZdq{}}\PY{p}{)}

\PY{n}{plt}\PY{o}{.}\PY{n}{legend}\PY{p}{(}\PY{n}{loc}\PY{o}{=}\PY{l+s+s1}{\PYZsq{}}\PY{l+s+s1}{center left}\PY{l+s+s1}{\PYZsq{}}\PY{p}{,} \PY{n}{bbox\PYZus{}to\PYZus{}anchor}\PY{o}{=}\PY{p}{(}\PY{l+m+mf}{0.55}\PY{p}{,} \PY{l+m+mf}{0.1}\PY{p}{)}\PY{p}{)}

\PY{n}{plt}\PY{o}{.}\PY{n}{show}\PY{p}{(}\PY{p}{)}
\end{Verbatim}
\end{tcolorbox}

    \begin{center}
    \adjustimage{max size={0.9\linewidth}{0.9\paperheight}}{figs/bvp_kM_en/output_217_0.png}
    \end{center}
    { \hspace*{\fill} \\}
    
    Now that I have straight lines, I can try linear regression to find
their shape.

    \begin{tcolorbox}[breakable, size=fbox, boxrule=1pt, pad at break*=1mm,colback=cellbackground, colframe=cellborder]
\prompt{In}{incolor}{ }{\boxspacing}
\begin{Verbatim}[commandchars=\\\{\}]
\PY{c+c1}{\PYZsh{} take the logarithm}
\PY{n}{x} \PY{o}{=} \PY{n}{np}\PY{o}{.}\PY{n}{log}\PY{p}{(}\PY{n}{mu}\PY{p}{)}
\PY{n}{X} \PY{o}{=} \PY{n}{np}\PY{o}{.}\PY{n}{array}\PY{p}{(}\PY{p}{[}\PY{n}{x}\PY{p}{]}\PY{p}{)}\PY{o}{.}\PY{n}{T}
\PY{n}{y2} \PY{o}{=} \PY{n}{np}\PY{o}{.}\PY{n}{log}\PY{p}{(}\PY{n}{max\PYZus{}for\PYZus{}mu}\PY{p}{)}
\PY{n}{y1} \PY{o}{=} \PY{n}{np}\PY{o}{.}\PY{n}{log}\PY{p}{(}\PY{n}{min\PYZus{}for\PYZus{}mu}\PY{p}{)}

\PY{c+c1}{\PYZsh{} fit to a line}
\PY{n}{reg1} \PY{o}{=} \PY{n}{LinearRegression}\PY{p}{(}\PY{p}{)}\PY{o}{.}\PY{n}{fit}\PY{p}{(}\PY{n}{X}\PY{p}{,} \PY{n}{y1}\PY{p}{)}
\PY{n}{reg2} \PY{o}{=} \PY{n}{LinearRegression}\PY{p}{(}\PY{p}{)}\PY{o}{.}\PY{n}{fit}\PY{p}{(}\PY{n}{X}\PY{p}{,} \PY{n}{y2}\PY{p}{)}

\PY{c+c1}{\PYZsh{} To save our results}
\PY{n}{slope\PYZus{}min} \PY{o}{=} \PY{n}{reg2}\PY{o}{.}\PY{n}{coef\PYZus{}}
\PY{n}{intercept\PYZus{}min} \PY{o}{=} \PY{n}{reg2}\PY{o}{.}\PY{n}{intercept\PYZus{}}
\PY{n}{slope\PYZus{}max} \PY{o}{=} \PY{n}{reg1}\PY{o}{.}\PY{n}{coef\PYZus{}}
\PY{n}{intercept\PYZus{}max} \PY{o}{=} \PY{n}{reg1}\PY{o}{.}\PY{n}{intercept\PYZus{}}

\PY{c+c1}{\PYZsh{} and show our results}
\PY{n+nb}{print}\PY{p}{(}\PY{l+s+s2}{\PYZdq{}}\PY{l+s+s2}{Slope dkn\PYZus{}max:}\PY{l+s+s2}{\PYZdq{}}\PY{p}{,} \PY{n}{reg1}\PY{o}{.}\PY{n}{coef\PYZus{}}\PY{p}{[}\PY{l+m+mi}{0}\PY{p}{]}\PY{p}{,} \PY{l+s+s2}{\PYZdq{}}\PY{l+s+s2}{, intercept dkn\PYZus{}max:}\PY{l+s+s2}{\PYZdq{}}\PY{p}{,} \PY{n}{reg1}\PY{o}{.}\PY{n}{intercept\PYZus{}}\PY{p}{,} \PY{l+s+s2}{\PYZdq{}}\PY{l+s+s2}{, convergence statistic:}\PY{l+s+s2}{\PYZdq{}}\PY{p}{,} \PY{n}{reg1}\PY{o}{.}\PY{n}{score}\PY{p}{(}\PY{n}{X}\PY{p}{,} \PY{n}{y1}\PY{p}{)}\PY{p}{)}
\PY{n+nb}{print}\PY{p}{(}\PY{l+s+s2}{\PYZdq{}}\PY{l+s+s2}{Slope dkn\PYZus{}mix:}\PY{l+s+s2}{\PYZdq{}}\PY{p}{,} \PY{n}{reg2}\PY{o}{.}\PY{n}{coef\PYZus{}}\PY{p}{[}\PY{l+m+mi}{0}\PY{p}{]}\PY{p}{,} \PY{l+s+s2}{\PYZdq{}}\PY{l+s+s2}{, intercept dkn\PYZus{}min:}\PY{l+s+s2}{\PYZdq{}}\PY{p}{,} \PY{n}{reg2}\PY{o}{.}\PY{n}{intercept\PYZus{}}\PY{p}{,} \PY{l+s+s2}{\PYZdq{}}\PY{l+s+s2}{, convergence statistic:}\PY{l+s+s2}{\PYZdq{}}\PY{p}{,} \PY{n}{reg2}\PY{o}{.}\PY{n}{score}\PY{p}{(}\PY{n}{X}\PY{p}{,} \PY{n}{y2}\PY{p}{)}\PY{p}{)}
\end{Verbatim}
\end{tcolorbox}

    \begin{Verbatim}[commandchars=\\\{\}]
Slope dkn\_max: 0.35515532338982736 , intercept dkn\_max: 7.7109215640459015 ,
convergence statistic: 0.9971399441064484
Slope dkn\_mix: 0.9987941677455325 , intercept dkn\_min: 16.619171811728208 ,
convergence statistic: 0.9999999229745502
    \end{Verbatim}

    \begin{tcolorbox}[breakable, size=fbox, boxrule=1pt, pad at break*=1mm,colback=cellbackground, colframe=cellborder]
\prompt{In}{incolor}{ }{\boxspacing}
\begin{Verbatim}[commandchars=\\\{\}]
\PY{c+c1}{\PYZsh{} Plot the results}

\PY{n}{line1} \PY{o}{=} \PY{n}{slope\PYZus{}max}\PY{o}{*} \PY{n}{np}\PY{o}{.}\PY{n}{log}\PY{p}{(}\PY{n}{mu}\PY{p}{)} \PY{o}{+} \PY{n}{intercept\PYZus{}max}
\PY{n}{line2} \PY{o}{=} \PY{n}{slope\PYZus{}min}\PY{o}{*} \PY{n}{np}\PY{o}{.}\PY{n}{log}\PY{p}{(}\PY{n}{mu}\PY{p}{)} \PY{o}{+} \PY{n}{intercept\PYZus{}min}

\PY{n}{fig}\PY{p}{,} \PY{n}{ax} \PY{o}{=} \PY{n}{plt}\PY{o}{.}\PY{n}{subplots}\PY{p}{(}\PY{p}{)}
\PY{n}{plt}\PY{o}{.}\PY{n}{title}\PY{p}{(}\PY{l+s+sa}{f}\PY{l+s+s2}{\PYZdq{}}\PY{l+s+s2}{The Limits of the Change in k Due to Nonlinearity vs }\PY{l+s+se}{\PYZbs{}u03BC}\PY{l+s+s2}{\PYZdq{}}\PY{p}{)}
\PY{n}{ax}\PY{o}{.}\PY{n}{set\PYZus{}ylabel}\PY{p}{(}\PY{l+s+s1}{\PYZsq{}}\PY{l+s+s1}{Natural log of \PYZhy{}dkn * 10}\PY{l+s+se}{\PYZbs{}u2075}\PY{l+s+s1}{ (m}\PY{l+s+se}{\PYZbs{}u207b}\PY{l+s+se}{\PYZbs{}u00b9}\PY{l+s+s1}{)}\PY{l+s+s1}{\PYZsq{}}\PY{p}{)}
\PY{n}{ax}\PY{o}{.}\PY{n}{set\PYZus{}xlabel}\PY{p}{(}\PY{l+s+s2}{\PYZdq{}}\PY{l+s+s2}{Natural log of the nonlinearity coefficient }\PY{l+s+se}{\PYZbs{}u03BC}\PY{l+s+s2}{ * 10}\PY{l+s+se}{\PYZbs{}u2078}\PY{l+s+s2}{\PYZdq{}}\PY{p}{)}

\PY{n}{ax}\PY{o}{.}\PY{n}{plot}\PY{p}{(}\PY{n}{np}\PY{o}{.}\PY{n}{log}\PY{p}{(}\PY{n}{mu}\PY{p}{)}\PY{p}{,} \PY{n}{np}\PY{o}{.}\PY{n}{log}\PY{p}{(}\PY{n}{min\PYZus{}for\PYZus{}mu}\PY{p}{)}\PY{p}{,} \PY{n}{label}\PY{o}{=}\PY{l+s+s2}{\PYZdq{}}\PY{l+s+s2}{maximum \PYZhy{}dkn * 10}\PY{l+s+se}{\PYZbs{}u207b}\PY{l+s+se}{\PYZbs{}u2075}\PY{l+s+s2}{\PYZdq{}}\PY{p}{)}
\PY{n}{ax}\PY{o}{.}\PY{n}{plot}\PY{p}{(}\PY{n}{np}\PY{o}{.}\PY{n}{log}\PY{p}{(}\PY{n}{mu}\PY{p}{)}\PY{p}{,} \PY{n}{np}\PY{o}{.}\PY{n}{log}\PY{p}{(}\PY{n}{max\PYZus{}for\PYZus{}mu}\PY{p}{)}\PY{p}{,} \PY{n}{label}\PY{o}{=}\PY{l+s+s2}{\PYZdq{}}\PY{l+s+s2}{minimum \PYZhy{}dkn * 10}\PY{l+s+se}{\PYZbs{}u207b}\PY{l+s+se}{\PYZbs{}u2075}\PY{l+s+s2}{\PYZdq{}}\PY{p}{)}
\PY{n}{ax}\PY{o}{.}\PY{n}{plot}\PY{p}{(}\PY{n}{np}\PY{o}{.}\PY{n}{log}\PY{p}{(}\PY{n}{mu}\PY{p}{)}\PY{p}{,} \PY{n}{line1}\PY{p}{,} \PY{n}{label}\PY{o}{=}\PY{l+s+s2}{\PYZdq{}}\PY{l+s+s2}{predicted maximum \PYZhy{}dkn * 10}\PY{l+s+se}{\PYZbs{}u207b}\PY{l+s+se}{\PYZbs{}u2075}\PY{l+s+s2}{\PYZdq{}}\PY{p}{)}
\PY{n}{ax}\PY{o}{.}\PY{n}{plot}\PY{p}{(}\PY{n}{np}\PY{o}{.}\PY{n}{log}\PY{p}{(}\PY{n}{mu}\PY{p}{)}\PY{p}{,} \PY{n}{line2}\PY{p}{,} \PY{n}{label}\PY{o}{=}\PY{l+s+s2}{\PYZdq{}}\PY{l+s+s2}{predicted minimum \PYZhy{}dkn * 10}\PY{l+s+se}{\PYZbs{}u207b}\PY{l+s+se}{\PYZbs{}u2075}\PY{l+s+s2}{\PYZdq{}}\PY{p}{)}

\PY{n}{plt}\PY{o}{.}\PY{n}{legend}\PY{p}{(}\PY{n}{loc}\PY{o}{=}\PY{l+s+s1}{\PYZsq{}}\PY{l+s+s1}{center left}\PY{l+s+s1}{\PYZsq{}}\PY{p}{,} \PY{n}{bbox\PYZus{}to\PYZus{}anchor}\PY{o}{=}\PY{p}{(}\PY{l+m+mf}{0.05}\PY{p}{,} \PY{l+m+mf}{0.55}\PY{p}{)}\PY{p}{)}

\PY{n}{plt}\PY{o}{.}\PY{n}{show}\PY{p}{(}\PY{p}{)}
\end{Verbatim}
\end{tcolorbox}

    \begin{center}
    \adjustimage{max size={0.9\linewidth}{0.9\paperheight}}{figs/bvp_kM_en/output_220_0.png}
    \end{center}
    { \hspace*{\fill} \\}
    
    \begin{tcolorbox}[breakable, size=fbox, boxrule=1pt, pad at break*=1mm,colback=cellbackground, colframe=cellborder]
\prompt{In}{incolor}{266}{\boxspacing}
\begin{Verbatim}[commandchars=\\\{\}]
\PY{c+c1}{\PYZsh{} Take the exponent of my line and fit the results.}

\PY{n}{line1} \PY{o}{=} \PY{n}{slope\PYZus{}max}\PY{o}{*} \PY{n}{np}\PY{o}{.}\PY{n}{log}\PY{p}{(}\PY{n}{mu}\PY{p}{)} \PY{o}{+} \PY{n}{intercept\PYZus{}max}
\PY{n}{line2} \PY{o}{=} \PY{n}{slope\PYZus{}min}\PY{o}{*}\PY{n}{np}\PY{o}{.}\PY{n}{log}\PY{p}{(}\PY{n}{mu}\PY{p}{)} \PY{o}{+} \PY{n}{intercept\PYZus{}min}

\PY{n}{fig}\PY{p}{,} \PY{n}{ax} \PY{o}{=} \PY{n}{plt}\PY{o}{.}\PY{n}{subplots}\PY{p}{(}\PY{p}{)}
\PY{n}{plt}\PY{o}{.}\PY{n}{title}\PY{p}{(}\PY{l+s+sa}{f}\PY{l+s+s2}{\PYZdq{}}\PY{l+s+s2}{The Limits of the Change in k Due to Nonlinearity vs }\PY{l+s+se}{\PYZbs{}u03BC}\PY{l+s+s2}{\PYZdq{}}\PY{p}{)}
\PY{n}{ax}\PY{o}{.}\PY{n}{set\PYZus{}ylabel}\PY{p}{(}\PY{l+s+s1}{\PYZsq{}}\PY{l+s+s1}{Change in k due to nonlinearity (\PYZhy{}dkn) (\PYZhy{}10}\PY{l+s+se}{\PYZbs{}u2075}\PY{l+s+s1}{ m}\PY{l+s+se}{\PYZbs{}u207b}\PY{l+s+se}{\PYZbs{}u00b9}\PY{l+s+s1}{)}\PY{l+s+s1}{\PYZsq{}}\PY{p}{)}
\PY{n}{ax}\PY{o}{.}\PY{n}{set\PYZus{}xlabel}\PY{p}{(}\PY{l+s+s2}{\PYZdq{}}\PY{l+s+s2}{Nonlinearity coefficient }\PY{l+s+se}{\PYZbs{}u03BC}\PY{l+s+s2}{ * 10}\PY{l+s+se}{\PYZbs{}u2078}\PY{l+s+s2}{\PYZdq{}}\PY{p}{)}

\PY{n}{ax}\PY{o}{.}\PY{n}{plot}\PY{p}{(}\PY{n}{mu}\PY{o}{*}\PY{l+m+mf}{1e8}\PY{p}{,} \PY{n}{min\PYZus{}for\PYZus{}mu}\PY{p}{,} \PY{n}{label}\PY{o}{=}\PY{l+s+s2}{\PYZdq{}}\PY{l+s+s2}{maximum \PYZhy{}dkn * 10}\PY{l+s+se}{\PYZbs{}u207b}\PY{l+s+se}{\PYZbs{}u2075}\PY{l+s+s2}{\PYZdq{}}\PY{p}{)}
\PY{n}{ax}\PY{o}{.}\PY{n}{plot}\PY{p}{(}\PY{n}{mu}\PY{o}{*}\PY{l+m+mf}{1e8}\PY{p}{,} \PY{n}{max\PYZus{}for\PYZus{}mu}\PY{p}{,} \PY{n}{label}\PY{o}{=}\PY{l+s+s2}{\PYZdq{}}\PY{l+s+s2}{minimum \PYZhy{}dkn * 10}\PY{l+s+se}{\PYZbs{}u207b}\PY{l+s+se}{\PYZbs{}u2075}\PY{l+s+s2}{\PYZdq{}}\PY{p}{)}
\PY{n}{ax}\PY{o}{.}\PY{n}{plot}\PY{p}{(}\PY{n}{mu}\PY{o}{*}\PY{l+m+mf}{1e8}\PY{p}{,} \PY{n}{np}\PY{o}{.}\PY{n}{exp}\PY{p}{(}\PY{n}{line1}\PY{p}{)}\PY{p}{,} \PY{n}{label}\PY{o}{=}\PY{l+s+s2}{\PYZdq{}}\PY{l+s+s2}{predicted maximum \PYZhy{}dkn * 10}\PY{l+s+se}{\PYZbs{}u207b}\PY{l+s+se}{\PYZbs{}u2075}\PY{l+s+s2}{\PYZdq{}}\PY{p}{)}
\PY{n}{ax}\PY{o}{.}\PY{n}{plot}\PY{p}{(}\PY{n}{mu}\PY{o}{*}\PY{l+m+mf}{1e8}\PY{p}{,} \PY{n}{np}\PY{o}{.}\PY{n}{exp}\PY{p}{(}\PY{n}{line2}\PY{p}{)}\PY{p}{,} \PY{n}{label}\PY{o}{=}\PY{l+s+s2}{\PYZdq{}}\PY{l+s+s2}{predicted minimum \PYZhy{}dkn * 10}\PY{l+s+se}{\PYZbs{}u207b}\PY{l+s+se}{\PYZbs{}u2075}\PY{l+s+s2}{\PYZdq{}}\PY{p}{)}

\PY{n}{plt}\PY{o}{.}\PY{n}{legend}\PY{p}{(}\PY{n}{loc}\PY{o}{=}\PY{l+s+s1}{\PYZsq{}}\PY{l+s+s1}{center left}\PY{l+s+s1}{\PYZsq{}}\PY{p}{,} \PY{n}{bbox\PYZus{}to\PYZus{}anchor}\PY{o}{=}\PY{p}{(}\PY{l+m+mi}{1}\PY{p}{,} \PY{l+m+mf}{0.5}\PY{p}{)}\PY{p}{)}

\PY{n}{plt}\PY{o}{.}\PY{n}{show}\PY{p}{(}\PY{p}{)}
\end{Verbatim}
\end{tcolorbox}

    \begin{center}
    \adjustimage{max size={0.9\linewidth}{0.9\paperheight}}{figs/bvp_kM_en/output_221_0.png}
    \end{center}
    { \hspace*{\fill} \\}
    
    \begin{tcolorbox}[breakable, size=fbox, boxrule=1pt, pad at break*=1mm,colback=cellbackground, colframe=cellborder]
\prompt{In}{incolor}{267}{\boxspacing}
\begin{Verbatim}[commandchars=\\\{\}]
\PY{n+nb}{print}\PY{p}{(}\PY{l+s+s2}{\PYZdq{}}\PY{l+s+s2}{Coefficient for dkn\PYZus{}min:}\PY{l+s+s2}{\PYZdq{}}\PY{p}{,} \PY{n}{np}\PY{o}{.}\PY{n}{exp}\PY{p}{(}\PY{n}{intercept\PYZus{}min}\PY{p}{)}\PY{p}{,} \PY{l+s+s2}{\PYZdq{}}\PY{l+s+s2}{, and dkn\PYZus{}max: }\PY{l+s+s2}{\PYZdq{}}\PY{p}{,} \PY{n}{np}\PY{o}{.}\PY{n}{exp}\PY{p}{(}\PY{n}{intercept\PYZus{}max}\PY{p}{)}\PY{p}{)}
\end{Verbatim}
\end{tcolorbox}

    \begin{Verbatim}[commandchars=\\\{\}]
Coefficient for dkn\_min: 16504965.149048917 , and dkn\_max:  2232.5987932024873
    \end{Verbatim}

    Our fits for the maximum and minimum values of \(-dkn*10^{-5}\) vs
\(\mu\) are:

\begin{equation}
L_{max} = 2232.6 * \mu^{0.355155}
\end{equation}

\begin{equation}
L_{min} = 16505000 * \mu^{0.99879}
\end{equation}

    Recalculated, we have:

\begin{equation}
dkn_{min}\left(\mu\right) = - 2.2326 \times 10^8 * \mu^{0.355155}
\end{equation}

\begin{equation}
dkn_{max}\left(\mu\right)  = - 1.6505 \times 10^{12} * \mu^{0.99879}
\end{equation}

    The relation marked here as \(dkn_{max}\) is expected to be more
significant, as it corresponds to the amplitude of our skewed squared
secant function, and not the shape of it. It is also linear, or at least
close to it. The slower than linear growth of \(dkn_{min}\) corresponds
to a linear growth of the amplitude and a further tilting of the squared
secant graph, which produces this fractional power growth. The \(dkn\)
vs \(dk\) graph reproduced below shows the uniform increase in growth
with a change in parameters (tilt and shift) for varying \(\mu\).

    \hypertarget{dk-periodicity-test}{%
\section{\texorpdfstring{\(dk\) Periodicity
Test}{dk Periodicity Test}}\label{dk-periodicity-test}}

    The following is a test of the periodicity of my function with respect
to \(dk\).

    \begin{tcolorbox}[breakable, size=fbox, boxrule=1pt, pad at break*=1mm,colback=cellbackground, colframe=cellborder]
\prompt{In}{incolor}{268}{\boxspacing}
\begin{Verbatim}[commandchars=\\\{\}]
\PY{c+c1}{\PYZsh{} I want to find my starting point for the dk oscillation}
\PY{n}{np}\PY{o}{.}\PY{n}{where}\PY{p}{(}\PY{n}{dk} \PY{o}{\PYZgt{}} \PY{l+m+mi}{2}\PY{o}{/}\PY{l+m+mf}{.9}\PY{p}{)}\PY{p}{[}\PY{l+m+mi}{0}\PY{p}{]}
\end{Verbatim}
\end{tcolorbox}

            \begin{tcolorbox}[breakable, size=fbox, boxrule=.5pt, pad at break*=1mm, opacityfill=0]
\prompt{Out}{outcolor}{268}{\boxspacing}
\begin{Verbatim}[commandchars=\\\{\}]
array([74, 75, 76, 77, 78, 79, 80, 81, 82, 83, 84, 85, 86, 87, 88, 89, 90,
       91, 92, 93, 94, 95, 96, 97, 98, 99], dtype=int64)
\end{Verbatim}
\end{tcolorbox}
        
    \begin{tcolorbox}[breakable, size=fbox, boxrule=1pt, pad at break*=1mm,colback=cellbackground, colframe=cellborder]
\prompt{In}{incolor}{269}{\boxspacing}
\begin{Verbatim}[commandchars=\\\{\}]
\PY{c+c1}{\PYZsh{} create a new object which is the same as my old dk array from the start of the previous cycle at 0}
\PY{c+c1}{\PYZsh{} to the start of the new cycle at dk = 2.222}
\PY{c+c1}{\PYZsh{} Then fill the rest of the array with the values from this first cycle + n * my period.}
\PY{n}{dk2}\PY{o}{=}\PY{n}{np}\PY{o}{.}\PY{n}{zeros}\PY{p}{(}\PY{l+m+mi}{962}\PY{p}{)}
\PY{n}{dk2}\PY{p}{[}\PY{l+m+mi}{0}\PY{p}{:}\PY{l+m+mi}{74}\PY{p}{]} \PY{o}{=} \PY{n}{dk}\PY{p}{[}\PY{l+m+mi}{0}\PY{p}{:}\PY{l+m+mi}{74}\PY{p}{]}
\PY{k}{for} \PY{n}{i} \PY{o+ow}{in} \PY{n+nb}{range}\PY{p}{(}\PY{l+m+mi}{13}\PY{p}{)}\PY{p}{:}
    \PY{n}{dk2}\PY{p}{[}\PY{l+m+mi}{74}\PY{o}{*}\PY{n}{i}\PY{p}{:}\PY{l+m+mi}{74}\PY{o}{*}\PY{p}{(}\PY{n}{i}\PY{o}{+}\PY{l+m+mi}{1}\PY{p}{)}\PY{p}{]} \PY{o}{=} \PY{n}{dk}\PY{p}{[}\PY{l+m+mi}{0}\PY{p}{:}\PY{l+m+mi}{74}\PY{p}{]}\PY{o}{+}\PY{n}{dk}\PY{p}{[}\PY{l+m+mi}{74}\PY{p}{]}\PY{o}{*}\PY{n}{i}
\end{Verbatim}
\end{tcolorbox}

    \begin{tcolorbox}[breakable, size=fbox, boxrule=1pt, pad at break*=1mm,colback=cellbackground, colframe=cellborder]
\prompt{In}{incolor}{270}{\boxspacing}
\begin{Verbatim}[commandchars=\\\{\}]
\PY{c+c1}{\PYZsh{} Fill the values of dkn at these dk values with with those of dkn from the previous cycle}
\PY{n}{dkn2} \PY{o}{=} \PY{n}{np}\PY{o}{.}\PY{n}{zeros}\PY{p}{(}\PY{p}{(}\PY{l+m+mi}{962}\PY{p}{,}\PY{l+m+mi}{100}\PY{p}{)}\PY{p}{)}
\PY{k}{for} \PY{n}{i} \PY{o+ow}{in} \PY{n+nb}{range}\PY{p}{(}\PY{l+m+mi}{13}\PY{p}{)}\PY{p}{:}
    \PY{n}{dkn2}\PY{p}{[}\PY{l+m+mi}{74}\PY{o}{*}\PY{n}{i}\PY{p}{:}\PY{l+m+mi}{74}\PY{o}{*}\PY{p}{(}\PY{n}{i}\PY{o}{+}\PY{l+m+mi}{1}\PY{p}{)}\PY{p}{]} \PY{o}{=} \PY{n}{dkn\PYZus{}arr}\PY{p}{[}\PY{l+m+mi}{0}\PY{p}{:}\PY{l+m+mi}{74}\PY{p}{]}
\end{Verbatim}
\end{tcolorbox}

    \begin{tcolorbox}[breakable, size=fbox, boxrule=1pt, pad at break*=1mm,colback=cellbackground, colframe=cellborder]
\prompt{In}{incolor}{ }{\boxspacing}
\begin{Verbatim}[commandchars=\\\{\}]
\PY{c+c1}{\PYZsh{} And compare to the actual results. }

\PY{n}{jj} \PY{o}{=} \PY{p}{[}\PY{l+m+mi}{10}\PY{p}{,} \PY{l+m+mi}{20}\PY{p}{,} \PY{l+m+mi}{30}\PY{p}{,} \PY{l+m+mi}{40}\PY{p}{,} \PY{l+m+mi}{50}\PY{p}{,} \PY{l+m+mi}{60}\PY{p}{,} \PY{l+m+mi}{70}\PY{p}{,} \PY{l+m+mi}{80}\PY{p}{,} \PY{l+m+mi}{90}\PY{p}{]}

\PY{n}{ilim} \PY{o}{=} \PY{n+nb}{len}\PY{p}{(}\PY{n}{jj}\PY{p}{)}

\PY{n}{fig}\PY{p}{,} \PY{n}{ax} \PY{o}{=} \PY{n}{plt}\PY{o}{.}\PY{n}{subplots}\PY{p}{(}\PY{p}{)}
\PY{n}{plt}\PY{o}{.}\PY{n}{title}\PY{p}{(}\PY{l+s+sa}{f}\PY{l+s+s2}{\PYZdq{}}\PY{l+s+s2}{Difference Between Predicted and Observed Change in k Due to}\PY{l+s+se}{\PYZbs{}n}\PY{l+s+s2}{Nonlinearity vs Linear k Variation For Various Values of }\PY{l+s+se}{\PYZbs{}u03BC}\PY{l+s+s2}{\PYZdq{}}\PY{p}{)}
\PY{n}{ax}\PY{o}{.}\PY{n}{set\PYZus{}ylabel}\PY{p}{(}\PY{l+s+s1}{\PYZsq{}}\PY{l+s+s1}{Difference in dkn (m}\PY{l+s+se}{\PYZbs{}u207b}\PY{l+s+se}{\PYZbs{}u00b9}\PY{l+s+s1}{)}\PY{l+s+s1}{\PYZsq{}}\PY{p}{)}
\PY{n}{ax}\PY{o}{.}\PY{n}{set\PYZus{}xlabel}\PY{p}{(}\PY{l+s+s2}{\PYZdq{}}\PY{l+s+s2}{Small\PYZhy{}scale variation of linear wave number k (dk (}\PY{l+s+s2}{\PYZpc{}}\PY{l+s+s2}{))}\PY{l+s+s2}{\PYZdq{}}\PY{p}{)}

\PY{k}{for} \PY{n}{i} \PY{o+ow}{in} \PY{n+nb}{range}\PY{p}{(}\PY{n}{ilim}\PY{p}{)}\PY{p}{:}
    \PY{n}{ax}\PY{o}{.}\PY{n}{plot}\PY{p}{(}\PY{n}{dk2}\PY{p}{[}\PY{l+m+mi}{78}\PY{p}{:}\PY{l+m+mi}{100}\PY{p}{]}\PY{p}{,} \PY{n}{dkn2}\PY{p}{[}\PY{l+m+mi}{78}\PY{p}{:}\PY{l+m+mi}{100}\PY{p}{,} \PY{n}{jj}\PY{p}{[}\PY{n}{i}\PY{p}{]}\PY{p}{]}\PY{o}{\PYZhy{}}\PY{n}{dkn\PYZus{}arr}\PY{p}{[}\PY{l+m+mi}{78}\PY{p}{:}\PY{p}{,} \PY{n}{jj}\PY{p}{[}\PY{n}{i}\PY{p}{]}\PY{p}{]}\PY{p}{,} \PY{n}{label} \PY{o}{=} \PY{l+s+sa}{f}\PY{l+s+s1}{\PYZsq{}}\PY{l+s+se}{\PYZbs{}u03BC}\PY{l+s+s1}{ = }\PY{l+s+si}{\PYZob{}}\PY{n}{mu}\PY{p}{[}\PY{n}{jj}\PY{p}{[}\PY{n}{i}\PY{p}{]}\PY{p}{]}\PY{l+s+si}{:}\PY{l+s+s1}{.2e}\PY{l+s+si}{\PYZcb{}}\PY{l+s+s1}{\PYZsq{}}\PY{p}{)}

\PY{n}{plt}\PY{o}{.}\PY{n}{legend}\PY{p}{(}\PY{p}{)}

\PY{n}{plt}\PY{o}{.}\PY{n}{show}\PY{p}{(}\PY{p}{)}
\end{Verbatim}
\end{tcolorbox}

    \begin{center}
    \adjustimage{max size={0.9\linewidth}{0.9\paperheight}}{figs/bvp_kM_en/output_231_0.png}
    \end{center}
    { \hspace*{\fill} \\}
    
    \begin{tcolorbox}[breakable, size=fbox, boxrule=1pt, pad at break*=1mm,colback=cellbackground, colframe=cellborder]
\prompt{In}{incolor}{ }{\boxspacing}
\begin{Verbatim}[commandchars=\\\{\}]
\PY{c+c1}{\PYZsh{} Shift the plot by one point, due to the fact that our repetition is not exactly on our grid.}

\PY{n}{jj} \PY{o}{=} \PY{p}{[}\PY{l+m+mi}{10}\PY{p}{,} \PY{l+m+mi}{20}\PY{p}{,} \PY{l+m+mi}{30}\PY{p}{,} \PY{l+m+mi}{40}\PY{p}{,} \PY{l+m+mi}{50}\PY{p}{,} \PY{l+m+mi}{60}\PY{p}{,} \PY{l+m+mi}{70}\PY{p}{,} \PY{l+m+mi}{80}\PY{p}{,} \PY{l+m+mi}{90}\PY{p}{]}

\PY{n}{ilim} \PY{o}{=} \PY{n+nb}{len}\PY{p}{(}\PY{n}{jj}\PY{p}{)}

\PY{c+c1}{\PYZsh{}dk = dk\PYZus{}g[4]}
\PY{c+c1}{\PYZsh{}mu = mu\PYZus{}g[4]}

\PY{n}{fig}\PY{p}{,} \PY{n}{ax} \PY{o}{=} \PY{n}{plt}\PY{o}{.}\PY{n}{subplots}\PY{p}{(}\PY{p}{)}
\PY{n}{plt}\PY{o}{.}\PY{n}{title}\PY{p}{(}\PY{l+s+sa}{f}\PY{l+s+s2}{\PYZdq{}}\PY{l+s+s2}{Difference Between Predicted and Observed Change in k Due to}\PY{l+s+se}{\PYZbs{}n}\PY{l+s+s2}{Nonlinearity vs Linear k Variation For Various Values of }\PY{l+s+se}{\PYZbs{}u03BC}\PY{l+s+s2}{\PYZdq{}}\PY{p}{)}
\PY{n}{ax}\PY{o}{.}\PY{n}{set\PYZus{}ylabel}\PY{p}{(}\PY{l+s+s1}{\PYZsq{}}\PY{l+s+s1}{Difference in dkn (m}\PY{l+s+se}{\PYZbs{}u207b}\PY{l+s+se}{\PYZbs{}u00b9}\PY{l+s+s1}{)}\PY{l+s+s1}{\PYZsq{}}\PY{p}{)}
\PY{n}{ax}\PY{o}{.}\PY{n}{set\PYZus{}xlabel}\PY{p}{(}\PY{l+s+s2}{\PYZdq{}}\PY{l+s+s2}{Small\PYZhy{}scale variation of linear wave number k (dk (}\PY{l+s+s2}{\PYZpc{}}\PY{l+s+s2}{))}\PY{l+s+s2}{\PYZdq{}}\PY{p}{)}

\PY{k}{for} \PY{n}{i} \PY{o+ow}{in} \PY{n+nb}{range}\PY{p}{(}\PY{n}{ilim}\PY{p}{)}\PY{p}{:}
    \PY{n}{ax}\PY{o}{.}\PY{n}{plot}\PY{p}{(}\PY{n}{dk2}\PY{p}{[}\PY{l+m+mi}{78}\PY{p}{:}\PY{l+m+mi}{100}\PY{p}{]}\PY{p}{,} \PY{n}{dkn2}\PY{p}{[}\PY{l+m+mi}{79}\PY{p}{:}\PY{l+m+mi}{101}\PY{p}{,} \PY{n}{jj}\PY{p}{[}\PY{n}{i}\PY{p}{]}\PY{p}{]}\PY{o}{\PYZhy{}}\PY{n}{dkn\PYZus{}arr}\PY{p}{[}\PY{l+m+mi}{78}\PY{p}{:}\PY{p}{,} \PY{n}{jj}\PY{p}{[}\PY{n}{i}\PY{p}{]}\PY{p}{]}\PY{p}{,} \PY{n}{label} \PY{o}{=} \PY{l+s+sa}{f}\PY{l+s+s1}{\PYZsq{}}\PY{l+s+se}{\PYZbs{}u03BC}\PY{l+s+s1}{ = }\PY{l+s+si}{\PYZob{}}\PY{n}{mu}\PY{p}{[}\PY{n}{jj}\PY{p}{[}\PY{n}{i}\PY{p}{]}\PY{p}{]}\PY{l+s+si}{:}\PY{l+s+s1}{.2e}\PY{l+s+si}{\PYZcb{}}\PY{l+s+s1}{\PYZsq{}}\PY{p}{)}

\PY{n}{plt}\PY{o}{.}\PY{n}{legend}\PY{p}{(}\PY{p}{)}

\PY{n}{plt}\PY{o}{.}\PY{n}{show}\PY{p}{(}\PY{p}{)}
\end{Verbatim}
\end{tcolorbox}

    \begin{center}
    \adjustimage{max size={0.9\linewidth}{0.9\paperheight}}{figs/bvp_kM_en/output_232_0.png}
    \end{center}
    { \hspace*{\fill} \\}
    
    Since the difference between plots is in the opposite direction at
opposite sides of the period boundary, my period is what I expect. I see
little evidence of a period change from \(\mu\). I am not sure if it is
actually shifting, or tilting with \(\mu\). Linear \(\mu\) dependence
occurs at the minimum.

    \hypertarget{visualization-mu-dependence}{%
\section{\texorpdfstring{Visualization \(\mu\)
dependence}{Visualization \textbackslash mu dependence}}\label{visualization-mu-dependence}}

    All of these plots worked with a dataset where there were 10 points for
each of my 4 parameters. My parameter ranges for \(R\), and \(B\) were
very small. All \(\mu\) dependencies show an amplification of the signal
(which is most likely linear) as well as a variation of the parameters
of the same functional forms already found. The periods of oscillation
are unchanged.

    \begin{tcolorbox}[breakable, size=fbox, boxrule=1pt, pad at break*=1mm,colback=cellbackground, colframe=cellborder]
\prompt{In}{incolor}{276}{\boxspacing}
\begin{Verbatim}[commandchars=\\\{\}]
\PY{n}{gind} \PY{o}{=} \PY{l+m+mi}{4}
\end{Verbatim}
\end{tcolorbox}

    \begin{tcolorbox}[breakable, size=fbox, boxrule=1pt, pad at break*=1mm,colback=cellbackground, colframe=cellborder]
\prompt{In}{incolor}{27}{\boxspacing}
\begin{Verbatim}[commandchars=\\\{\}]
\PY{n}{tbl\PYZus{}red}\PY{p}{[}\PY{n}{gind}\PY{p}{]}
\end{Verbatim}
\end{tcolorbox}

            \begin{tcolorbox}[breakable, size=fbox, boxrule=.5pt, pad at break*=1mm, opacityfill=0]
\prompt{Out}{outcolor}{27}{\boxspacing}
\begin{Verbatim}[commandchars=\\\{\}]
            R      B   dk            mu            dkn           dMn
0      0.9000  0.500  0.1  3.162278e-10    -543.384872      4.338550
1      0.9000  0.500  0.1  5.011872e-10    -860.990616      6.888751
2      0.9000  0.500  0.1  7.943282e-10   -1364.039793     10.949784
3      0.9000  0.500  0.1  1.258925e-09   -2160.525335     17.434920
4      0.9000  0.500  0.1  1.995262e-09   -3420.936071     27.837906
{\ldots}       {\ldots}    {\ldots}  {\ldots}           {\ldots}            {\ldots}           {\ldots}
14636  0.9002  0.506  0.9  5.011872e-09  -13583.041660   1825.339240
14637  0.9002  0.506  0.9  7.943282e-09  -22291.148302   2910.698893
14638  0.9002  0.506  0.9  1.258925e-08  -37789.155453   4676.204490
14639  0.9002  0.506  0.9  1.995262e-08  -70899.717850   7722.336209
14640  0.9002  0.506  0.9  3.162278e-08 -460820.125599 -13054.490179

[14641 rows x 6 columns]
\end{Verbatim}
\end{tcolorbox}
        
    \begin{tcolorbox}[breakable, size=fbox, boxrule=1pt, pad at break*=1mm,colback=cellbackground, colframe=cellborder]
\prompt{In}{incolor}{277}{\boxspacing}
\begin{Verbatim}[commandchars=\\\{\}]
\PY{n}{dkn} \PY{o}{=} \PY{n}{np}\PY{o}{.}\PY{n}{zeros}\PY{p}{(}\PY{p}{(}\PY{n}{gs2}\PY{p}{,} \PY{n}{gs2}\PY{p}{,} \PY{n}{gs2}\PY{p}{,} \PY{n}{gs2}\PY{p}{)}\PY{p}{)}
\PY{n}{dMn} \PY{o}{=} \PY{n}{np}\PY{o}{.}\PY{n}{zeros}\PY{p}{(}\PY{p}{(}\PY{n}{gs2}\PY{p}{,} \PY{n}{gs2}\PY{p}{,} \PY{n}{gs2}\PY{p}{,} \PY{n}{gs2}\PY{p}{)}\PY{p}{)}

\PY{k}{for} \PY{n}{r}\PY{p}{,} \PY{n}{rr} \PY{o+ow}{in} \PY{n+nb}{enumerate}\PY{p}{(}\PY{n}{R\PYZus{}g}\PY{p}{[}\PY{n}{gind}\PY{p}{]}\PY{p}{)}\PY{p}{:}
    \PY{k}{for} \PY{n}{b}\PY{p}{,} \PY{n}{bb} \PY{o+ow}{in} \PY{n+nb}{enumerate}\PY{p}{(}\PY{n}{B\PYZus{}g}\PY{p}{[}\PY{n}{gind}\PY{p}{]}\PY{p}{)}\PY{p}{:}
        \PY{k}{for} \PY{n}{kk}\PY{p}{,} \PY{n}{kkk} \PY{o+ow}{in} \PY{n+nb}{enumerate}\PY{p}{(}\PY{n}{dk\PYZus{}g}\PY{p}{[}\PY{n}{gind}\PY{p}{]}\PY{p}{)}\PY{p}{:}
            \PY{k}{for} \PY{n}{m}\PY{p}{,} \PY{n}{mm} \PY{o+ow}{in} \PY{n+nb}{enumerate}\PY{p}{(}\PY{n}{mu\PYZus{}g}\PY{p}{[}\PY{n}{gind}\PY{p}{]}\PY{p}{)}\PY{p}{:}
                \PY{n}{tbl\PYZus{}sel} \PY{o}{=} \PY{n}{tbl\PYZus{}red}\PY{p}{[}\PY{n}{gind}\PY{p}{]}\PY{p}{[}\PY{p}{(}\PY{n}{tbl\PYZus{}red}\PY{p}{[}\PY{n}{gind}\PY{p}{]}\PY{p}{[}\PY{l+s+s2}{\PYZdq{}}\PY{l+s+s2}{R}\PY{l+s+s2}{\PYZdq{}}\PY{p}{]}\PY{o}{==}\PY{n}{rr}\PY{p}{)} \PY{o}{\PYZam{}}\PY{p}{(}\PY{n}{tbl\PYZus{}red}\PY{p}{[}\PY{n}{gind}\PY{p}{]}\PY{p}{[}\PY{l+s+s2}{\PYZdq{}}\PY{l+s+s2}{B}\PY{l+s+s2}{\PYZdq{}}\PY{p}{]}\PY{o}{==}\PY{n}{bb}\PY{p}{)}\PY{o}{\PYZam{}}\PY{p}{(}\PY{n}{tbl\PYZus{}red}\PY{p}{[}\PY{n}{gind}\PY{p}{]}\PY{p}{[}\PY{l+s+s2}{\PYZdq{}}\PY{l+s+s2}{dk}\PY{l+s+s2}{\PYZdq{}}\PY{p}{]}\PY{o}{==}\PY{n}{kkk}\PY{p}{)} \PY{o}{\PYZam{}}\PY{p}{(}\PY{n}{tbl\PYZus{}red}\PY{p}{[}\PY{n}{gind}\PY{p}{]}\PY{p}{[}\PY{l+s+s2}{\PYZdq{}}\PY{l+s+s2}{mu}\PY{l+s+s2}{\PYZdq{}}\PY{p}{]}\PY{o}{==}\PY{n}{mm}\PY{p}{)}\PY{p}{]}
                \PY{n}{dkn}\PY{p}{[}\PY{n}{r}\PY{p}{,} \PY{n}{b}\PY{p}{,} \PY{n}{kk}\PY{p}{,} \PY{n}{m}\PY{p}{]} \PY{o}{=} \PY{n}{tbl\PYZus{}sel}\PY{p}{[}\PY{l+s+s2}{\PYZdq{}}\PY{l+s+s2}{dkn}\PY{l+s+s2}{\PYZdq{}}\PY{p}{]}\PY{o}{.}\PY{n}{to\PYZus{}numpy}\PY{p}{(}\PY{p}{)}\PY{p}{[}\PY{l+m+mi}{0}\PY{p}{]}
                \PY{n}{dMn}\PY{p}{[}\PY{n}{r}\PY{p}{,} \PY{n}{b}\PY{p}{,} \PY{n}{kk}\PY{p}{,} \PY{n}{m}\PY{p}{]} \PY{o}{=} \PY{n}{tbl\PYZus{}sel}\PY{p}{[}\PY{l+s+s2}{\PYZdq{}}\PY{l+s+s2}{dMn}\PY{l+s+s2}{\PYZdq{}}\PY{p}{]}\PY{o}{.}\PY{n}{to\PYZus{}numpy}\PY{p}{(}\PY{p}{)}\PY{p}{[}\PY{l+m+mi}{0}\PY{p}{]}
\end{Verbatim}
\end{tcolorbox}

    \begin{tcolorbox}[breakable, size=fbox, boxrule=1pt, pad at break*=1mm,colback=cellbackground, colframe=cellborder]
\prompt{In}{incolor}{299}{\boxspacing}
\begin{Verbatim}[commandchars=\\\{\}]
\PY{n}{kscale} \PY{o}{=} \PY{l+m+mf}{1e5}
\PY{n}{R} \PY{o}{=} \PY{p}{(}\PY{n}{R\PYZus{}g}\PY{p}{[}\PY{n}{gind}\PY{p}{]} \PY{o}{\PYZhy{}} \PY{l+m+mf}{0.9}\PY{p}{)}\PY{o}{*}\PY{l+m+mi}{10000}
\PY{n}{mu} \PY{o}{=} \PY{n}{mu\PYZus{}g}\PY{p}{[}\PY{n}{gind}\PY{p}{]}
\PY{n}{mscale} \PY{o}{=} \PY{l+m+mf}{1e4}
\end{Verbatim}
\end{tcolorbox}

    \begin{tcolorbox}[breakable, size=fbox, boxrule=1pt, pad at break*=1mm,colback=cellbackground, colframe=cellborder]
\prompt{In}{incolor}{ }{\boxspacing}
\begin{Verbatim}[commandchars=\\\{\}]
\PY{c+c1}{\PYZsh{} This is the plot for R dependence of dkn.}
\PY{c+c1}{\PYZsh{} mu amplifies the signal, but no other changes are visible here.}

\PY{n}{mui} \PY{o}{=} \PY{p}{[}\PY{l+m+mi}{0}\PY{p}{,}\PY{l+m+mi}{1}\PY{p}{,}\PY{l+m+mi}{2}\PY{p}{,}\PY{l+m+mi}{3}\PY{p}{,}\PY{l+m+mi}{4}\PY{p}{,}\PY{l+m+mi}{5}\PY{p}{,}\PY{l+m+mi}{6}\PY{p}{,}\PY{l+m+mi}{7}\PY{p}{,}\PY{l+m+mi}{8}\PY{p}{,}\PY{l+m+mi}{9}\PY{p}{,}\PY{l+m+mi}{10}\PY{p}{]}
\PY{n}{bi} \PY{o}{=} \PY{l+m+mi}{3}
\PY{n}{ki} \PY{o}{=} \PY{l+m+mi}{3}

\PY{n}{ilim} \PY{o}{=} \PY{n+nb}{len}\PY{p}{(}\PY{n}{mui}\PY{p}{)}


\PY{n}{fig}\PY{p}{,} \PY{n}{ax} \PY{o}{=} \PY{n}{plt}\PY{o}{.}\PY{n}{subplots}\PY{p}{(}\PY{p}{)}
\PY{n}{plt}\PY{o}{.}\PY{n}{title}\PY{p}{(}\PY{l+s+sa}{f}\PY{l+s+s2}{\PYZdq{}}\PY{l+s+s2}{Change in k Due to Nonlinearity vs R for dk = }\PY{l+s+si}{\PYZob{}}\PY{n}{dk\PYZus{}g}\PY{p}{[}\PY{n}{gind}\PY{p}{]}\PY{p}{[}\PY{n}{ki}\PY{p}{]}\PY{l+s+si}{:}\PY{l+s+s2}{.2f}\PY{l+s+si}{\PYZcb{}}\PY{l+s+s2}{ / 2 * 10}\PY{l+s+se}{\PYZbs{}u2076}\PY{l+s+s2}{ m}\PY{l+s+se}{\PYZbs{}u207b}\PY{l+s+se}{\PYZbs{}u00b9}\PY{l+s+se}{\PYZbs{}n}\PY{l+s+s2}{ and B = }\PY{l+s+si}{\PYZob{}}\PY{n}{B\PYZus{}g}\PY{p}{[}\PY{n}{gind}\PY{p}{]}\PY{p}{[}\PY{n}{bi}\PY{p}{]}\PY{l+s+si}{:}\PY{l+s+s2}{.4f}\PY{l+s+si}{\PYZcb{}}\PY{l+s+s2}{ for Various Values of }\PY{l+s+se}{\PYZbs{}u03BC}\PY{l+s+s2}{\PYZdq{}}\PY{p}{)}
\PY{n}{ax}\PY{o}{.}\PY{n}{set\PYZus{}ylabel}\PY{p}{(}\PY{l+s+s1}{\PYZsq{}}\PY{l+s+s1}{Change in k (dkn) due to nonlinearity (10}\PY{l+s+se}{\PYZbs{}u2075}\PY{l+s+s1}{ m}\PY{l+s+se}{\PYZbs{}u207b}\PY{l+s+se}{\PYZbs{}u00b9}\PY{l+s+s1}{)}\PY{l+s+s1}{\PYZsq{}}\PY{p}{)}
\PY{n}{ax}\PY{o}{.}\PY{n}{set\PYZus{}xlabel}\PY{p}{(}\PY{l+s+s2}{\PYZdq{}}\PY{l+s+s2}{Change dr in radius R of ring (R = 9 * 10}\PY{l+s+se}{\PYZbs{}u207b}\PY{l+s+se}{\PYZbs{}u2077}\PY{l+s+s2}{ m + dr * 10}\PY{l+s+se}{\PYZbs{}u207b}\PY{l+s+se}{\PYZbs{}u00b9}\PY{l+s+se}{\PYZbs{}u2070}\PY{l+s+s2}{ m)}\PY{l+s+s2}{\PYZdq{}}\PY{p}{)}

\PY{c+c1}{\PYZsh{}ax.set\PYZus{}xlabel(\PYZdq{}Variation (dB) of magnetic field strength B (B = 0.5 + dB * B\PYZus{}max)\PYZdq{})}

\PY{c+c1}{\PYZsh{}plt.title(f\PYZdq{}Plot of the Change in k Due to Nonlinearity vs Linear k\PYZbs{}nVariation For Various Values of \PYZbs{}u03BC\PYZdq{})}
\PY{c+c1}{\PYZsh{}ax.set\PYZus{}ylabel(\PYZsq{}Change in k due to nonlinearity (dkn) (m\PYZbs{}u207b\PYZbs{}u00b9)\PYZsq{})}
\PY{c+c1}{\PYZsh{}ax.set\PYZus{}xlabel(\PYZdq{}Small\PYZhy{}scale variation of linear wave number k (dk (\PYZpc{}))\PYZdq{})}

\PY{k}{for} \PY{n}{i} \PY{o+ow}{in} \PY{n+nb}{range}\PY{p}{(}\PY{n}{ilim}\PY{p}{)}\PY{p}{:}
    \PY{n}{ax}\PY{o}{.}\PY{n}{plot}\PY{p}{(}\PY{n}{R}\PY{p}{,} \PY{n}{dkn}\PY{p}{[}\PY{p}{:}\PY{p}{,} \PY{n}{bi}\PY{p}{,}\PY{n}{ki}\PY{p}{,}\PY{n}{mui}\PY{p}{[}\PY{n}{i}\PY{p}{]}\PY{p}{]}\PY{o}{/}\PY{n}{kscale}\PY{p}{,} \PY{n}{label} \PY{o}{=} \PY{l+s+sa}{f}\PY{l+s+s1}{\PYZsq{}}\PY{l+s+se}{\PYZbs{}u03BC}\PY{l+s+s1}{ = }\PY{l+s+si}{\PYZob{}}\PY{n}{mu}\PY{p}{[}\PY{n}{mui}\PY{p}{[}\PY{n}{i}\PY{p}{]}\PY{p}{]}\PY{l+s+si}{:}\PY{l+s+s1}{.3e}\PY{l+s+si}{\PYZcb{}}\PY{l+s+s1}{\PYZsq{}}\PY{p}{)}

\PY{n}{plt}\PY{o}{.}\PY{n}{legend}\PY{p}{(}\PY{n}{loc}\PY{o}{=}\PY{l+s+s1}{\PYZsq{}}\PY{l+s+s1}{center left}\PY{l+s+s1}{\PYZsq{}}\PY{p}{,} \PY{n}{bbox\PYZus{}to\PYZus{}anchor}\PY{o}{=}\PY{p}{(}\PY{l+m+mi}{1}\PY{p}{,} \PY{l+m+mf}{0.5}\PY{p}{)}\PY{p}{)}

\PY{n}{plt}\PY{o}{.}\PY{n}{show}\PY{p}{(}\PY{p}{)}
\end{Verbatim}
\end{tcolorbox}

    \begin{center}
    \adjustimage{max size={0.9\linewidth}{0.9\paperheight}}{figs/bvp_kM_en/output_240_0.png}
    \end{center}
    { \hspace*{\fill} \\}
    
    The periodic behavior of \(R\) is multiplied by some function of
\(\mu\). The large \(\mu\) point may either be a deviation from the
behavior or an error. This characteristic is not visible on the \(\mu\)
vs dkn graphs, but clearly visible here.

    \begin{tcolorbox}[breakable, size=fbox, boxrule=1pt, pad at break*=1mm,colback=cellbackground, colframe=cellborder]
\prompt{In}{incolor}{ }{\boxspacing}
\begin{Verbatim}[commandchars=\\\{\}]
\PY{c+c1}{\PYZsh{} This is the plot for B dependence of dkn.}
\PY{c+c1}{\PYZsh{} mu amplifies the signal, but no other changes are visible here.}

\PY{n}{ivar} \PY{o}{=} \PY{p}{[}\PY{l+m+mi}{0}\PY{p}{,}\PY{l+m+mi}{1}\PY{p}{,}\PY{l+m+mi}{2}\PY{p}{,}\PY{l+m+mi}{3}\PY{p}{,}\PY{l+m+mi}{4}\PY{p}{,}\PY{l+m+mi}{5}\PY{p}{,}\PY{l+m+mi}{6}\PY{p}{,}\PY{l+m+mi}{7}\PY{p}{,}\PY{l+m+mi}{8}\PY{p}{,}\PY{l+m+mi}{9}\PY{p}{,}\PY{l+m+mi}{10}\PY{p}{]}
\PY{n}{i1}  \PY{o}{=} \PY{l+m+mi}{3}
\PY{n}{i2}  \PY{o}{=} \PY{l+m+mi}{3}

\PY{n}{ilim} \PY{o}{=} \PY{n+nb}{len}\PY{p}{(}\PY{n}{ivar}\PY{p}{)}

\PY{n}{fig}\PY{p}{,} \PY{n}{ax} \PY{o}{=} \PY{n}{plt}\PY{o}{.}\PY{n}{subplots}\PY{p}{(}\PY{p}{)}
\PY{n}{plt}\PY{o}{.}\PY{n}{title}\PY{p}{(}\PY{l+s+sa}{f}\PY{l+s+s2}{\PYZdq{}}\PY{l+s+s2}{Change in k Due to Nonlinearity vs B for Various Values of }\PY{l+s+se}{\PYZbs{}u03BC}\PY{l+s+se}{\PYZbs{}n}\PY{l+s+s2}{for R = }\PY{l+s+si}{\PYZob{}}\PY{n}{R\PYZus{}g}\PY{p}{[}\PY{n}{gind}\PY{p}{]}\PY{p}{[}\PY{n}{i1}\PY{p}{]}\PY{l+s+si}{:}\PY{l+s+s2}{.5f}\PY{l+s+si}{\PYZcb{}}\PY{l+s+s2}{ * 10}\PY{l+s+se}{\PYZbs{}u207b}\PY{l+s+se}{\PYZbs{}u2076}\PY{l+s+s2}{ m and dk = }\PY{l+s+si}{\PYZob{}}\PY{n}{dk\PYZus{}g}\PY{p}{[}\PY{n}{gind}\PY{p}{]}\PY{p}{[}\PY{n}{i2}\PY{p}{]}\PY{l+s+si}{:}\PY{l+s+s2}{.2f}\PY{l+s+si}{\PYZcb{}}\PY{l+s+s2}{ / 2 * 10}\PY{l+s+se}{\PYZbs{}u2076}\PY{l+s+s2}{ m}\PY{l+s+se}{\PYZbs{}u207b}\PY{l+s+se}{\PYZbs{}u00b9}\PY{l+s+s2}{\PYZdq{}}\PY{p}{)}
\PY{n}{ax}\PY{o}{.}\PY{n}{set\PYZus{}ylabel}\PY{p}{(}\PY{l+s+s1}{\PYZsq{}}\PY{l+s+s1}{Change in k (dkn) due to nonlinearity (10}\PY{l+s+se}{\PYZbs{}u2075}\PY{l+s+s1}{ m}\PY{l+s+se}{\PYZbs{}u207b}\PY{l+s+se}{\PYZbs{}u00b9}\PY{l+s+s1}{)}\PY{l+s+s1}{\PYZsq{}}\PY{p}{)}
\PY{n}{ax}\PY{o}{.}\PY{n}{set\PYZus{}xlabel}\PY{p}{(}\PY{l+s+s2}{\PYZdq{}}\PY{l+s+s2}{Variation (dB) of magnetic field strength B (B = B * B\PYZus{}max)}\PY{l+s+s2}{\PYZdq{}}\PY{p}{)}

\PY{c+c1}{\PYZsh{}plt.title(f\PYZdq{}Plot of the Change in k Due to Nonlinearity vs Linear k\PYZbs{}nVariation For Various Values of \PYZbs{}u03BC\PYZdq{})}
\PY{c+c1}{\PYZsh{}ax.set\PYZus{}ylabel(\PYZsq{}Change in k due to nonlinearity (dkn) (m\PYZbs{}u207b\PYZbs{}u00b9)\PYZsq{})}
\PY{c+c1}{\PYZsh{}ax.set\PYZus{}xlabel(\PYZdq{}Small\PYZhy{}scale variation of linear wave number k (dk (\PYZpc{}))\PYZdq{})}

\PY{k}{for} \PY{n}{i} \PY{o+ow}{in} \PY{n+nb}{range}\PY{p}{(}\PY{n}{ilim}\PY{p}{)}\PY{p}{:}
    \PY{n}{ax}\PY{o}{.}\PY{n}{plot}\PY{p}{(}\PY{n}{B\PYZus{}g}\PY{p}{[}\PY{n}{gind}\PY{p}{]}\PY{p}{,} \PY{n}{dkn}\PY{p}{[}\PY{n}{i1}\PY{p}{,} \PY{p}{:}\PY{p}{,}\PY{n}{i2}\PY{p}{,}\PY{n}{ivar}\PY{p}{[}\PY{n}{i}\PY{p}{]}\PY{p}{]}\PY{o}{/}\PY{n}{kscale}\PY{p}{,} \PY{n}{label} \PY{o}{=} \PY{l+s+sa}{f}\PY{l+s+s1}{\PYZsq{}}\PY{l+s+se}{\PYZbs{}u03BC}\PY{l+s+s1}{ = }\PY{l+s+si}{\PYZob{}}\PY{n}{mu}\PY{p}{[}\PY{n}{ivar}\PY{p}{[}\PY{n}{i}\PY{p}{]}\PY{p}{]}\PY{l+s+si}{:}\PY{l+s+s1}{.3e}\PY{l+s+si}{\PYZcb{}}\PY{l+s+s1}{\PYZsq{}}\PY{p}{)}

\PY{n}{plt}\PY{o}{.}\PY{n}{legend}\PY{p}{(}\PY{n}{loc}\PY{o}{=}\PY{l+s+s1}{\PYZsq{}}\PY{l+s+s1}{center left}\PY{l+s+s1}{\PYZsq{}}\PY{p}{,} \PY{n}{bbox\PYZus{}to\PYZus{}anchor}\PY{o}{=}\PY{p}{(}\PY{l+m+mi}{1}\PY{p}{,} \PY{l+m+mf}{0.5}\PY{p}{)}\PY{p}{)}

\PY{n}{plt}\PY{o}{.}\PY{n}{show}\PY{p}{(}\PY{p}{)}
\end{Verbatim}
\end{tcolorbox}

    \begin{center}
    \adjustimage{max size={0.9\linewidth}{0.9\paperheight}}{figs/bvp_kM_en/output_242_0.png}
    \end{center}
    { \hspace*{\fill} \\}
    
    The scale for B dependence here is obviously still too small. A pattern
does exist, however. Only the effect of amplification with increased
\(\mu\) is visible at this \(B\) node.

    \begin{tcolorbox}[breakable, size=fbox, boxrule=1pt, pad at break*=1mm,colback=cellbackground, colframe=cellborder]
\prompt{In}{incolor}{ }{\boxspacing}
\begin{Verbatim}[commandchars=\\\{\}]
\PY{c+c1}{\PYZsh{} This is the plot for dk dependence of dkn.}
\PY{c+c1}{\PYZsh{} mu amplifies the signal as well as causing a shift to the right.}

\PY{n}{ivar} \PY{o}{=} \PY{p}{[}\PY{l+m+mi}{0}\PY{p}{,}\PY{l+m+mi}{1}\PY{p}{,}\PY{l+m+mi}{2}\PY{p}{,}\PY{l+m+mi}{3}\PY{p}{,}\PY{l+m+mi}{4}\PY{p}{,}\PY{l+m+mi}{5}\PY{p}{,}\PY{l+m+mi}{6}\PY{p}{,}\PY{l+m+mi}{7}\PY{p}{,}\PY{l+m+mi}{8}\PY{p}{,}\PY{l+m+mi}{9}\PY{p}{,}\PY{l+m+mi}{10}\PY{p}{]}
\PY{n}{i1}  \PY{o}{=} \PY{l+m+mi}{3}
\PY{n}{i2}  \PY{o}{=} \PY{l+m+mi}{3}

\PY{n}{ilim} \PY{o}{=} \PY{n+nb}{len}\PY{p}{(}\PY{n}{ivar}\PY{p}{)}

\PY{n}{fig}\PY{p}{,} \PY{n}{ax} \PY{o}{=} \PY{n}{plt}\PY{o}{.}\PY{n}{subplots}\PY{p}{(}\PY{p}{)}
\PY{n}{plt}\PY{o}{.}\PY{n}{title}\PY{p}{(}\PY{l+s+sa}{f}\PY{l+s+s2}{\PYZdq{}}\PY{l+s+s2}{Change in k Due to Nonlinearity vs dk for Various Values of }\PY{l+s+se}{\PYZbs{}u03BC}\PY{l+s+se}{\PYZbs{}n}\PY{l+s+s2}{for R = }\PY{l+s+si}{\PYZob{}}\PY{n}{R\PYZus{}g}\PY{p}{[}\PY{n}{gind}\PY{p}{]}\PY{p}{[}\PY{n}{i1}\PY{p}{]}\PY{l+s+si}{:}\PY{l+s+s2}{.5f}\PY{l+s+si}{\PYZcb{}}\PY{l+s+s2}{ * 10}\PY{l+s+se}{\PYZbs{}u207b}\PY{l+s+se}{\PYZbs{}u2076}\PY{l+s+s2}{ m and B = }\PY{l+s+si}{\PYZob{}}\PY{n}{B\PYZus{}g}\PY{p}{[}\PY{n}{gind}\PY{p}{]}\PY{p}{[}\PY{n}{i2}\PY{p}{]}\PY{l+s+si}{:}\PY{l+s+s2}{.4f}\PY{l+s+si}{\PYZcb{}}\PY{l+s+s2}{ * B\PYZus{}max}\PY{l+s+s2}{\PYZdq{}}\PY{p}{)}
\PY{n}{ax}\PY{o}{.}\PY{n}{set\PYZus{}ylabel}\PY{p}{(}\PY{l+s+s1}{\PYZsq{}}\PY{l+s+s1}{Change in k (dkn) due to nonlinearity (10}\PY{l+s+se}{\PYZbs{}u2075}\PY{l+s+s1}{ m}\PY{l+s+se}{\PYZbs{}u207b}\PY{l+s+se}{\PYZbs{}u00b9}\PY{l+s+s1}{)}\PY{l+s+s1}{\PYZsq{}}\PY{p}{)}
\PY{n}{ax}\PY{o}{.}\PY{n}{set\PYZus{}xlabel}\PY{p}{(}\PY{l+s+s2}{\PYZdq{}}\PY{l+s+s2}{Small\PYZhy{}scale variation of linear wave number k (dk (}\PY{l+s+s2}{\PYZpc{}}\PY{l+s+s2}{))}\PY{l+s+s2}{\PYZdq{}}\PY{p}{)}

\PY{k}{for} \PY{n}{i} \PY{o+ow}{in} \PY{n+nb}{range}\PY{p}{(}\PY{n}{ilim}\PY{p}{)}\PY{p}{:}
    \PY{n}{ax}\PY{o}{.}\PY{n}{plot}\PY{p}{(}\PY{n}{dk\PYZus{}g}\PY{p}{[}\PY{n}{gind}\PY{p}{]}\PY{p}{,} \PY{n}{dkn}\PY{p}{[}\PY{n}{i1}\PY{p}{,} \PY{n}{i2}\PY{p}{,}\PY{p}{:}\PY{p}{,}\PY{n}{ivar}\PY{p}{[}\PY{n}{i}\PY{p}{]}\PY{p}{]}\PY{o}{/}\PY{n}{kscale}\PY{p}{,} \PY{n}{label} \PY{o}{=} \PY{l+s+sa}{f}\PY{l+s+s1}{\PYZsq{}}\PY{l+s+se}{\PYZbs{}u03BC}\PY{l+s+s1}{ = }\PY{l+s+si}{\PYZob{}}\PY{n}{mu}\PY{p}{[}\PY{n}{ivar}\PY{p}{[}\PY{n}{i}\PY{p}{]}\PY{p}{]}\PY{l+s+si}{:}\PY{l+s+s1}{.3e}\PY{l+s+si}{\PYZcb{}}\PY{l+s+s1}{\PYZsq{}}\PY{p}{)}

\PY{n}{plt}\PY{o}{.}\PY{n}{legend}\PY{p}{(}\PY{p}{)}\PY{c+c1}{\PYZsh{}loc=\PYZsq{}center left\PYZsq{}, bbox\PYZus{}to\PYZus{}anchor=(1, 0.5))}

\PY{n}{plt}\PY{o}{.}\PY{n}{show}\PY{p}{(}\PY{p}{)}
\end{Verbatim}
\end{tcolorbox}

    \begin{center}
    \adjustimage{max size={0.9\linewidth}{0.9\paperheight}}{figs/bvp_kM_en/output_244_0.png}
    \end{center}
    { \hspace*{\fill} \\}
    
    dk has a phase shift with growing mu, not just a amplification. The dk
data requires a larger range to see the periodic shape.

    Now for the dependence of \(dMn\) on my parameters.

    \begin{tcolorbox}[breakable, size=fbox, boxrule=1pt, pad at break*=1mm,colback=cellbackground, colframe=cellborder]
\prompt{In}{incolor}{ }{\boxspacing}
\begin{Verbatim}[commandchars=\\\{\}]
\PY{c+c1}{\PYZsh{} This is the plot for dk dependence of dMn.}
\PY{c+c1}{\PYZsh{} mu amplifies the signal as well as causing a shift to the right.}
\PY{c+c1}{\PYZsh{} I don\PYZsq{}t know if there is something else visible here.}

\PY{n}{ivar} \PY{o}{=} \PY{p}{[}\PY{l+m+mi}{0}\PY{p}{,}\PY{l+m+mi}{1}\PY{p}{,}\PY{l+m+mi}{2}\PY{p}{,}\PY{l+m+mi}{3}\PY{p}{,}\PY{l+m+mi}{4}\PY{p}{,}\PY{l+m+mi}{5}\PY{p}{,}\PY{l+m+mi}{6}\PY{p}{,}\PY{l+m+mi}{7}\PY{p}{,}\PY{l+m+mi}{8}\PY{p}{,}\PY{l+m+mi}{9}\PY{p}{,}\PY{l+m+mi}{10}\PY{p}{]}
\PY{n}{i1}  \PY{o}{=} \PY{l+m+mi}{3}
\PY{n}{i2}  \PY{o}{=} \PY{l+m+mi}{3}

\PY{n}{ilim} \PY{o}{=} \PY{n+nb}{len}\PY{p}{(}\PY{n}{ivar}\PY{p}{)}

\PY{n}{fig}\PY{p}{,} \PY{n}{ax} \PY{o}{=} \PY{n}{plt}\PY{o}{.}\PY{n}{subplots}\PY{p}{(}\PY{p}{)}
\PY{n}{plt}\PY{o}{.}\PY{n}{title}\PY{p}{(}\PY{l+s+sa}{f}\PY{l+s+s2}{\PYZdq{}}\PY{l+s+s2}{Change in M Due to Nonlinearity vs dk for Various Values of }\PY{l+s+se}{\PYZbs{}u03BC}\PY{l+s+se}{\PYZbs{}n}\PY{l+s+s2}{for R = }\PY{l+s+si}{\PYZob{}}\PY{n}{R\PYZus{}g}\PY{p}{[}\PY{n}{gind}\PY{p}{]}\PY{p}{[}\PY{n}{i1}\PY{p}{]}\PY{l+s+si}{:}\PY{l+s+s2}{.5f}\PY{l+s+si}{\PYZcb{}}\PY{l+s+s2}{ * 10}\PY{l+s+se}{\PYZbs{}u207b}\PY{l+s+se}{\PYZbs{}u2076}\PY{l+s+s2}{ m and B = }\PY{l+s+si}{\PYZob{}}\PY{n}{B\PYZus{}g}\PY{p}{[}\PY{n}{gind}\PY{p}{]}\PY{p}{[}\PY{n}{i2}\PY{p}{]}\PY{l+s+si}{:}\PY{l+s+s2}{.4f}\PY{l+s+si}{\PYZcb{}}\PY{l+s+s2}{ * B\PYZus{}max}\PY{l+s+s2}{\PYZdq{}}\PY{p}{)}
\PY{n}{ax}\PY{o}{.}\PY{n}{set\PYZus{}ylabel}\PY{p}{(}\PY{l+s+s1}{\PYZsq{}}\PY{l+s+s1}{Change in M (dMn) due to nonlinearity (10}\PY{l+s+se}{\PYZbs{}u2074}\PY{l+s+s1}{ m}\PY{l+s+se}{\PYZbs{}u207b}\PY{l+s+se}{\PYZbs{}u00b9}\PY{l+s+s1}{)}\PY{l+s+s1}{\PYZsq{}}\PY{p}{)}
\PY{n}{ax}\PY{o}{.}\PY{n}{set\PYZus{}xlabel}\PY{p}{(}\PY{l+s+s2}{\PYZdq{}}\PY{l+s+s2}{Small\PYZhy{}scale variation of linear wave number k (dk (}\PY{l+s+s2}{\PYZpc{}}\PY{l+s+s2}{))}\PY{l+s+s2}{\PYZdq{}}\PY{p}{)}

\PY{k}{for} \PY{n}{i} \PY{o+ow}{in} \PY{n+nb}{range}\PY{p}{(}\PY{n}{ilim}\PY{p}{)}\PY{p}{:}
    \PY{n}{ax}\PY{o}{.}\PY{n}{plot}\PY{p}{(}\PY{n}{dk\PYZus{}g}\PY{p}{[}\PY{n}{gind}\PY{p}{]}\PY{p}{,} \PY{n}{dMn}\PY{p}{[}\PY{n}{i1}\PY{p}{,} \PY{n}{i2}\PY{p}{,}\PY{p}{:}\PY{p}{,}\PY{n}{ivar}\PY{p}{[}\PY{n}{i}\PY{p}{]}\PY{p}{]}\PY{o}{/}\PY{n}{mscale}\PY{p}{,} \PY{n}{label} \PY{o}{=} \PY{l+s+sa}{f}\PY{l+s+s1}{\PYZsq{}}\PY{l+s+se}{\PYZbs{}u03BC}\PY{l+s+s1}{ = }\PY{l+s+si}{\PYZob{}}\PY{n}{mu}\PY{p}{[}\PY{n}{ivar}\PY{p}{[}\PY{n}{i}\PY{p}{]}\PY{p}{]}\PY{l+s+si}{:}\PY{l+s+s1}{.3e}\PY{l+s+si}{\PYZcb{}}\PY{l+s+s1}{\PYZsq{}}\PY{p}{)}

\PY{n}{plt}\PY{o}{.}\PY{n}{legend}\PY{p}{(}\PY{n}{loc}\PY{o}{=}\PY{l+s+s1}{\PYZsq{}}\PY{l+s+s1}{center left}\PY{l+s+s1}{\PYZsq{}}\PY{p}{,} \PY{n}{bbox\PYZus{}to\PYZus{}anchor}\PY{o}{=}\PY{p}{(}\PY{l+m+mi}{1}\PY{p}{,} \PY{l+m+mf}{0.5}\PY{p}{)}\PY{p}{)}

\PY{n}{plt}\PY{o}{.}\PY{n}{show}\PY{p}{(}\PY{p}{)}
\end{Verbatim}
\end{tcolorbox}

    \begin{center}
    \adjustimage{max size={0.9\linewidth}{0.9\paperheight}}{figs/bvp_kM_en/output_247_0.png}
    \end{center}
    { \hspace*{\fill} \\}
    
    \begin{tcolorbox}[breakable, size=fbox, boxrule=1pt, pad at break*=1mm,colback=cellbackground, colframe=cellborder]
\prompt{In}{incolor}{296}{\boxspacing}
\begin{Verbatim}[commandchars=\\\{\}]
\PY{c+c1}{\PYZsh{} This is the plot for B dependence of dMn.}
\PY{c+c1}{\PYZsh{} mu amplifies the signal, but the node is fixed on zero.}

\PY{n}{ivar} \PY{o}{=} \PY{p}{[}\PY{l+m+mi}{0}\PY{p}{,}\PY{l+m+mi}{1}\PY{p}{,}\PY{l+m+mi}{2}\PY{p}{,}\PY{l+m+mi}{3}\PY{p}{,}\PY{l+m+mi}{4}\PY{p}{,}\PY{l+m+mi}{5}\PY{p}{,}\PY{l+m+mi}{6}\PY{p}{,}\PY{l+m+mi}{7}\PY{p}{,}\PY{l+m+mi}{8}\PY{p}{,}\PY{l+m+mi}{9}\PY{p}{,}\PY{l+m+mi}{10}\PY{p}{]}
\PY{n}{i1}  \PY{o}{=} \PY{l+m+mi}{3}
\PY{n}{i2}  \PY{o}{=} \PY{l+m+mi}{3}

\PY{n}{ilim} \PY{o}{=} \PY{n+nb}{len}\PY{p}{(}\PY{n}{ivar}\PY{p}{)}

\PY{n}{fig}\PY{p}{,} \PY{n}{ax} \PY{o}{=} \PY{n}{plt}\PY{o}{.}\PY{n}{subplots}\PY{p}{(}\PY{p}{)}
\PY{n}{plt}\PY{o}{.}\PY{n}{title}\PY{p}{(}\PY{l+s+sa}{f}\PY{l+s+s2}{\PYZdq{}}\PY{l+s+s2}{Change in M Due to Nonlinearity vs B for Various Values of }\PY{l+s+se}{\PYZbs{}u03BC}\PY{l+s+se}{\PYZbs{}n}\PY{l+s+s2}{for R = }\PY{l+s+si}{\PYZob{}}\PY{n}{R\PYZus{}g}\PY{p}{[}\PY{n}{gind}\PY{p}{]}\PY{p}{[}\PY{n}{i1}\PY{p}{]}\PY{l+s+si}{:}\PY{l+s+s2}{.5f}\PY{l+s+si}{\PYZcb{}}\PY{l+s+s2}{ * 10}\PY{l+s+se}{\PYZbs{}u207b}\PY{l+s+se}{\PYZbs{}u2076}\PY{l+s+s2}{ m and dk = }\PY{l+s+si}{\PYZob{}}\PY{n}{dk\PYZus{}g}\PY{p}{[}\PY{n}{gind}\PY{p}{]}\PY{p}{[}\PY{n}{i2}\PY{p}{]}\PY{l+s+si}{:}\PY{l+s+s2}{.2f}\PY{l+s+si}{\PYZcb{}}\PY{l+s+s2}{ / 2 * 10}\PY{l+s+se}{\PYZbs{}u2076}\PY{l+s+s2}{ m}\PY{l+s+se}{\PYZbs{}u207b}\PY{l+s+se}{\PYZbs{}u00b9}\PY{l+s+s2}{\PYZdq{}}\PY{p}{)}
\PY{n}{ax}\PY{o}{.}\PY{n}{set\PYZus{}ylabel}\PY{p}{(}\PY{l+s+s1}{\PYZsq{}}\PY{l+s+s1}{Change in M (dMn) due to nonlinearity (10}\PY{l+s+se}{\PYZbs{}u2074}\PY{l+s+s1}{ m}\PY{l+s+se}{\PYZbs{}u207b}\PY{l+s+se}{\PYZbs{}u00b9}\PY{l+s+s1}{)}\PY{l+s+s1}{\PYZsq{}}\PY{p}{)}
\PY{n}{ax}\PY{o}{.}\PY{n}{set\PYZus{}xlabel}\PY{p}{(}\PY{l+s+s2}{\PYZdq{}}\PY{l+s+s2}{Variation (dB) of magnetic field strength B (B = B * B\PYZus{}max)}\PY{l+s+s2}{\PYZdq{}}\PY{p}{)}

\PY{k}{for} \PY{n}{i} \PY{o+ow}{in} \PY{n+nb}{range}\PY{p}{(}\PY{n}{ilim}\PY{p}{)}\PY{p}{:}
    \PY{n}{ax}\PY{o}{.}\PY{n}{plot}\PY{p}{(}\PY{n}{B\PYZus{}g}\PY{p}{[}\PY{n}{gind}\PY{p}{]}\PY{p}{,} \PY{n}{dMn}\PY{p}{[}\PY{n}{i1}\PY{p}{,} \PY{p}{:}\PY{p}{,} \PY{n}{i2}\PY{p}{,}\PY{n}{ivar}\PY{p}{[}\PY{n}{i}\PY{p}{]}\PY{p}{]}\PY{o}{/}\PY{n}{mscale}\PY{p}{,} \PY{n}{label} \PY{o}{=} \PY{l+s+sa}{f}\PY{l+s+s1}{\PYZsq{}}\PY{l+s+se}{\PYZbs{}u03BC}\PY{l+s+s1}{ = }\PY{l+s+si}{\PYZob{}}\PY{n}{mu}\PY{p}{[}\PY{n}{ivar}\PY{p}{[}\PY{n}{i}\PY{p}{]}\PY{p}{]}\PY{l+s+si}{:}\PY{l+s+s1}{.3e}\PY{l+s+si}{\PYZcb{}}\PY{l+s+s1}{\PYZsq{}}\PY{p}{)}

\PY{n}{plt}\PY{o}{.}\PY{n}{legend}\PY{p}{(}\PY{n}{loc}\PY{o}{=}\PY{l+s+s1}{\PYZsq{}}\PY{l+s+s1}{center left}\PY{l+s+s1}{\PYZsq{}}\PY{p}{,} \PY{n}{bbox\PYZus{}to\PYZus{}anchor}\PY{o}{=}\PY{p}{(}\PY{l+m+mi}{1}\PY{p}{,} \PY{l+m+mf}{0.5}\PY{p}{)}\PY{p}{)}

\PY{n}{plt}\PY{o}{.}\PY{n}{show}\PY{p}{(}\PY{p}{)}
\end{Verbatim}
\end{tcolorbox}

    \begin{center}
    \adjustimage{max size={0.9\linewidth}{0.9\paperheight}}{figs/bvp_kM_en/output_248_0.png}
    \end{center}
    { \hspace*{\fill} \\}
    
    Linear or linearized. \(B\) is still too small in scale to see the
oscillations. \(\mu\) definitely amplifies the signal, but leaves the
position of the \(0\) value unchanged.

    \begin{tcolorbox}[breakable, size=fbox, boxrule=1pt, pad at break*=1mm,colback=cellbackground, colframe=cellborder]
\prompt{In}{incolor}{301}{\boxspacing}
\begin{Verbatim}[commandchars=\\\{\}]
\PY{c+c1}{\PYZsh{} This is the plot for R dependence of dMn.}
\PY{c+c1}{\PYZsh{} mu amplifies the signal, but the nodes are unchanged.}

\PY{n}{ivar} \PY{o}{=} \PY{p}{[}\PY{l+m+mi}{0}\PY{p}{,}\PY{l+m+mi}{1}\PY{p}{,}\PY{l+m+mi}{2}\PY{p}{,}\PY{l+m+mi}{3}\PY{p}{,}\PY{l+m+mi}{4}\PY{p}{,}\PY{l+m+mi}{5}\PY{p}{,}\PY{l+m+mi}{6}\PY{p}{,}\PY{l+m+mi}{7}\PY{p}{,}\PY{l+m+mi}{8}\PY{p}{,}\PY{l+m+mi}{9}\PY{p}{,}\PY{l+m+mi}{10}\PY{p}{]}
\PY{n}{i1}  \PY{o}{=} \PY{l+m+mi}{3}
\PY{n}{i2}  \PY{o}{=} \PY{l+m+mi}{3}

\PY{n}{ilim} \PY{o}{=} \PY{n+nb}{len}\PY{p}{(}\PY{n}{ivar}\PY{p}{)}

\PY{n}{fig}\PY{p}{,} \PY{n}{ax} \PY{o}{=} \PY{n}{plt}\PY{o}{.}\PY{n}{subplots}\PY{p}{(}\PY{p}{)}
\PY{n}{plt}\PY{o}{.}\PY{n}{title}\PY{p}{(}\PY{l+s+sa}{f}\PY{l+s+s2}{\PYZdq{}}\PY{l+s+s2}{Change in M Due to Nonlinearity vs R for dk = }\PY{l+s+si}{\PYZob{}}\PY{n}{dk\PYZus{}g}\PY{p}{[}\PY{n}{gind}\PY{p}{]}\PY{p}{[}\PY{n}{ki}\PY{p}{]}\PY{l+s+si}{:}\PY{l+s+s2}{.2f}\PY{l+s+si}{\PYZcb{}}\PY{l+s+s2}{ / 2 * 10}\PY{l+s+se}{\PYZbs{}u2076}\PY{l+s+s2}{ m}\PY{l+s+se}{\PYZbs{}u207b}\PY{l+s+se}{\PYZbs{}u00b9}\PY{l+s+se}{\PYZbs{}n}\PY{l+s+s2}{ and B = }\PY{l+s+si}{\PYZob{}}\PY{n}{B\PYZus{}g}\PY{p}{[}\PY{n}{gind}\PY{p}{]}\PY{p}{[}\PY{n}{bi}\PY{p}{]}\PY{l+s+si}{:}\PY{l+s+s2}{.4f}\PY{l+s+si}{\PYZcb{}}\PY{l+s+s2}{ for Various Values of }\PY{l+s+se}{\PYZbs{}u03BC}\PY{l+s+s2}{\PYZdq{}}\PY{p}{)}
\PY{n}{ax}\PY{o}{.}\PY{n}{set\PYZus{}ylabel}\PY{p}{(}\PY{l+s+s1}{\PYZsq{}}\PY{l+s+s1}{Change in M (dMn) due to nonlinearity (10}\PY{l+s+se}{\PYZbs{}u2074}\PY{l+s+s1}{ m}\PY{l+s+se}{\PYZbs{}u207b}\PY{l+s+se}{\PYZbs{}u00b9}\PY{l+s+s1}{)}\PY{l+s+s1}{\PYZsq{}}\PY{p}{)}
\PY{n}{ax}\PY{o}{.}\PY{n}{set\PYZus{}xlabel}\PY{p}{(}\PY{l+s+s2}{\PYZdq{}}\PY{l+s+s2}{Change dr in radius R of ring (R = 9 * 10}\PY{l+s+se}{\PYZbs{}u207b}\PY{l+s+se}{\PYZbs{}u2077}\PY{l+s+s2}{ m + dr * 10}\PY{l+s+se}{\PYZbs{}u207b}\PY{l+s+se}{\PYZbs{}u00b9}\PY{l+s+se}{\PYZbs{}u2070}\PY{l+s+s2}{ m)}\PY{l+s+s2}{\PYZdq{}}\PY{p}{)}

\PY{k}{for} \PY{n}{i} \PY{o+ow}{in} \PY{n+nb}{range}\PY{p}{(}\PY{n}{ilim}\PY{p}{)}\PY{p}{:}
    \PY{n}{ax}\PY{o}{.}\PY{n}{plot}\PY{p}{(}\PY{n}{R}\PY{p}{,} \PY{n}{dMn}\PY{p}{[}\PY{p}{:}\PY{p}{,} \PY{n}{i1}\PY{p}{,} \PY{n}{i2}\PY{p}{,}\PY{n}{ivar}\PY{p}{[}\PY{n}{i}\PY{p}{]}\PY{p}{]}\PY{o}{/}\PY{n}{mscale}\PY{p}{,} \PY{n}{label} \PY{o}{=} \PY{l+s+sa}{f}\PY{l+s+s1}{\PYZsq{}}\PY{l+s+se}{\PYZbs{}u03BC}\PY{l+s+s1}{ = }\PY{l+s+si}{\PYZob{}}\PY{n}{mu}\PY{p}{[}\PY{n}{ivar}\PY{p}{[}\PY{n}{i}\PY{p}{]}\PY{p}{]}\PY{l+s+si}{:}\PY{l+s+s1}{.3e}\PY{l+s+si}{\PYZcb{}}\PY{l+s+s1}{\PYZsq{}}\PY{p}{)}

\PY{n}{plt}\PY{o}{.}\PY{n}{legend}\PY{p}{(}\PY{n}{loc}\PY{o}{=}\PY{l+s+s1}{\PYZsq{}}\PY{l+s+s1}{center left}\PY{l+s+s1}{\PYZsq{}}\PY{p}{,} \PY{n}{bbox\PYZus{}to\PYZus{}anchor}\PY{o}{=}\PY{p}{(}\PY{l+m+mi}{1}\PY{p}{,} \PY{l+m+mf}{0.5}\PY{p}{)}\PY{p}{)}

\PY{n}{plt}\PY{o}{.}\PY{n}{show}\PY{p}{(}\PY{p}{)}
\end{Verbatim}
\end{tcolorbox}

    \begin{center}
    \adjustimage{max size={0.9\linewidth}{0.9\paperheight}}{figs/bvp_kM_en/output_250_0.png}
    \end{center}
    { \hspace*{\fill} \\}
    
    They are all meeting at zero at the same points. There is an erroneous
point (or possibly a phase change) for the largest \(\mu\), but the rest
appear to show a periodic structure with the same frequency and offset
for every values of \(\mu\).

    \hypertarget{large-r-analysis}{%
\section{Large R analysis}\label{large-r-analysis}}

    For my final analysis, I have fixed \(R\) to always be at the same point
in the small-scale oscillations with a periodicity based on \(kR\), and
sampled the full range of \(R\) values. The large-scale \(B\pi^2R^2\)
oscillations should be visible in this graph, as should any unaccounted
for \(R\) dependent behavior.

    \begin{tcolorbox}[breakable, size=fbox, boxrule=1pt, pad at break*=1mm,colback=cellbackground, colframe=cellborder]
\prompt{In}{incolor}{302}{\boxspacing}
\begin{Verbatim}[commandchars=\\\{\}]
\PY{n}{g} \PY{o}{=} \PY{n}{sns}\PY{o}{.}\PY{n}{PairGrid}\PY{p}{(}\PY{n}{tbl\PYZus{}red}\PY{p}{[}\PY{l+m+mi}{2}\PY{p}{]}\PY{p}{,} \PY{n}{diag\PYZus{}sharey}\PY{o}{=}\PY{k+kc}{False}\PY{p}{,} \PY{n}{corner}\PY{o}{=}\PY{k+kc}{True}\PY{p}{)}

\PY{n}{g}\PY{o}{.}\PY{n}{map\PYZus{}diag}\PY{p}{(}\PY{n}{sns}\PY{o}{.}\PY{n}{kdeplot}\PY{p}{)}
\PY{n}{g}\PY{o}{.}\PY{n}{map\PYZus{}lower}\PY{p}{(}\PY{n}{sns}\PY{o}{.}\PY{n}{scatterplot}\PY{p}{)}
\end{Verbatim}
\end{tcolorbox}

            \begin{tcolorbox}[breakable, size=fbox, boxrule=.5pt, pad at break*=1mm, opacityfill=0]
\prompt{Out}{outcolor}{302}{\boxspacing}
\begin{Verbatim}[commandchars=\\\{\}]
<seaborn.axisgrid.PairGrid at 0x1d1a84e19d0>
\end{Verbatim}
\end{tcolorbox}
        
    \begin{center}
    \adjustimage{max size={0.9\linewidth}{0.9\paperheight}}{figs/bvp_kM_en/output_254_1.png}
    \end{center}
    { \hspace*{\fill} \\}
    
    Lets start by loading our data.

    \begin{tcolorbox}[breakable, size=fbox, boxrule=1pt, pad at break*=1mm,colback=cellbackground, colframe=cellborder]
\prompt{In}{incolor}{303}{\boxspacing}
\begin{Verbatim}[commandchars=\\\{\}]
\PY{n}{gind} \PY{o}{=} \PY{l+m+mi}{2}
\end{Verbatim}
\end{tcolorbox}

    \begin{tcolorbox}[breakable, size=fbox, boxrule=1pt, pad at break*=1mm,colback=cellbackground, colframe=cellborder]
\prompt{In}{incolor}{304}{\boxspacing}
\begin{Verbatim}[commandchars=\\\{\}]
\PY{k}{del} \PY{n}{dkn\PYZus{}arr}\PY{p}{,} \PY{n}{dMn\PYZus{}arr}

\PY{n}{dkn\PYZus{}arr} \PY{o}{=} \PY{n}{np}\PY{o}{.}\PY{n}{zeros}\PY{p}{(}\PY{p}{(}\PY{n}{gs\PYZus{}r1}\PY{p}{,} \PY{n}{gs\PYZus{}b1}\PY{p}{)}\PY{p}{)}
\PY{n}{dMn\PYZus{}arr} \PY{o}{=} \PY{n}{np}\PY{o}{.}\PY{n}{zeros}\PY{p}{(}\PY{p}{(}\PY{n}{gs\PYZus{}r1}\PY{p}{,} \PY{n}{gs\PYZus{}b1}\PY{p}{)}\PY{p}{)}

\PY{k}{for} \PY{n}{r}\PY{p}{,} \PY{n}{rr} \PY{o+ow}{in} \PY{n+nb}{enumerate}\PY{p}{(}\PY{n}{R\PYZus{}g}\PY{p}{[}\PY{n}{gind}\PY{p}{]}\PY{p}{)}\PY{p}{:}
    \PY{k}{for} \PY{n}{b}\PY{p}{,} \PY{n}{bb} \PY{o+ow}{in} \PY{n+nb}{enumerate}\PY{p}{(}\PY{n}{B\PYZus{}g}\PY{p}{[}\PY{n}{gind}\PY{p}{]}\PY{p}{)}\PY{p}{:}
        \PY{n}{tbl\PYZus{}sel} \PY{o}{=} \PY{n}{tbl\PYZus{}red}\PY{p}{[}\PY{n}{gind}\PY{p}{]}\PY{p}{[}\PY{p}{(}\PY{n}{tbl\PYZus{}red}\PY{p}{[}\PY{n}{gind}\PY{p}{]}\PY{p}{[}\PY{l+s+s2}{\PYZdq{}}\PY{l+s+s2}{R}\PY{l+s+s2}{\PYZdq{}}\PY{p}{]}\PY{o}{==}\PY{n}{rr}\PY{p}{)} \PY{o}{\PYZam{}}\PY{p}{(}\PY{n}{tbl\PYZus{}red}\PY{p}{[}\PY{n}{gind}\PY{p}{]}\PY{p}{[}\PY{l+s+s2}{\PYZdq{}}\PY{l+s+s2}{B}\PY{l+s+s2}{\PYZdq{}}\PY{p}{]}\PY{o}{==}\PY{n}{bb}\PY{p}{)}\PY{p}{]}
        \PY{n}{dkn\PYZus{}arr}\PY{p}{[}\PY{n}{r}\PY{p}{,}\PY{n}{b}\PY{p}{]} \PY{o}{=} \PY{n}{tbl\PYZus{}sel}\PY{p}{[}\PY{l+s+s2}{\PYZdq{}}\PY{l+s+s2}{dkn}\PY{l+s+s2}{\PYZdq{}}\PY{p}{]}\PY{o}{.}\PY{n}{to\PYZus{}numpy}\PY{p}{(}\PY{p}{)}\PY{p}{[}\PY{l+m+mi}{0}\PY{p}{]}
        \PY{n}{dMn\PYZus{}arr}\PY{p}{[}\PY{n}{r}\PY{p}{,}\PY{n}{b}\PY{p}{]} \PY{o}{=} \PY{n}{tbl\PYZus{}sel}\PY{p}{[}\PY{l+s+s2}{\PYZdq{}}\PY{l+s+s2}{dMn}\PY{l+s+s2}{\PYZdq{}}\PY{p}{]}\PY{o}{.}\PY{n}{to\PYZus{}numpy}\PY{p}{(}\PY{p}{)}\PY{p}{[}\PY{l+m+mi}{0}\PY{p}{]}
\end{Verbatim}
\end{tcolorbox}

    \begin{tcolorbox}[breakable, size=fbox, boxrule=1pt, pad at break*=1mm,colback=cellbackground, colframe=cellborder]
\prompt{In}{incolor}{305}{\boxspacing}
\begin{Verbatim}[commandchars=\\\{\}]
\PY{n}{R} \PY{o}{=} \PY{n}{R\PYZus{}g}\PY{p}{[}\PY{n}{gind}\PY{p}{]}
\PY{n}{B} \PY{o}{=} \PY{n}{B\PYZus{}g}\PY{p}{[}\PY{n}{gind}\PY{p}{]}
\end{Verbatim}
\end{tcolorbox}

    \begin{tcolorbox}[breakable, size=fbox, boxrule=1pt, pad at break*=1mm,colback=cellbackground, colframe=cellborder]
\prompt{In}{incolor}{310}{\boxspacing}
\begin{Verbatim}[commandchars=\\\{\}]
\PY{c+c1}{\PYZsh{} Plots of dkn vs R}

\PY{n}{Bi} \PY{o}{=} \PY{p}{[}\PY{l+m+mi}{5}\PY{p}{,} \PY{l+m+mi}{14}\PY{p}{]}

\PY{n}{ilim} \PY{o}{=} \PY{n+nb}{len}\PY{p}{(}\PY{n}{Bi}\PY{p}{)}

\PY{n}{fig}\PY{p}{,} \PY{n}{ax} \PY{o}{=} \PY{n}{plt}\PY{o}{.}\PY{n}{subplots}\PY{p}{(}\PY{p}{)}
\PY{n}{plt}\PY{o}{.}\PY{n}{title}\PY{p}{(}\PY{l+s+sa}{f}\PY{l+s+s2}{\PYZdq{}}\PY{l+s+s2}{Change in k Due to Nonlinearity vs R for Various B Values}\PY{l+s+s2}{\PYZdq{}}\PY{p}{)}
\PY{n}{ax}\PY{o}{.}\PY{n}{set\PYZus{}ylabel}\PY{p}{(}\PY{l+s+s1}{\PYZsq{}}\PY{l+s+s1}{Change in k (dkn) due to nonlinearity (10}\PY{l+s+se}{\PYZbs{}u2075}\PY{l+s+s1}{ m}\PY{l+s+se}{\PYZbs{}u207b}\PY{l+s+se}{\PYZbs{}u00b9}\PY{l+s+s1}{)}\PY{l+s+s1}{\PYZsq{}}\PY{p}{)}
\PY{n}{ax}\PY{o}{.}\PY{n}{set\PYZus{}xlabel}\PY{p}{(}\PY{l+s+s2}{\PYZdq{}}\PY{l+s+s2}{Radius R of the ring (10}\PY{l+s+se}{\PYZbs{}u207b}\PY{l+s+se}{\PYZbs{}u2076}\PY{l+s+s2}{ m)}\PY{l+s+s2}{\PYZdq{}}\PY{p}{)}

\PY{k}{for} \PY{n}{bb} \PY{o+ow}{in} \PY{n}{Bi}\PY{p}{:}
    \PY{n}{ax}\PY{o}{.}\PY{n}{plot}\PY{p}{(}\PY{n}{R}\PY{p}{,} \PY{n}{dkn\PYZus{}arr}\PY{p}{[}\PY{p}{:}\PY{p}{,}\PY{n}{bb}\PY{p}{]}\PY{o}{/}\PY{l+m+mi}{100000}\PY{p}{,} \PY{n}{label} \PY{o}{=} \PY{l+s+sa}{f}\PY{l+s+s2}{\PYZdq{}}\PY{l+s+s2}{B = }\PY{l+s+si}{\PYZob{}}\PY{n}{B}\PY{p}{[}\PY{n}{bb}\PY{p}{]}\PY{l+s+si}{:}\PY{l+s+s2}{.4f}\PY{l+s+si}{\PYZcb{}}\PY{l+s+s2}{ * B\PYZus{}max}\PY{l+s+s2}{\PYZdq{}}\PY{p}{)}

\PY{n}{plt}\PY{o}{.}\PY{n}{legend}\PY{p}{(}\PY{p}{)}\PY{c+c1}{\PYZsh{}loc=\PYZsq{}center left\PYZsq{}, bbox\PYZus{}to\PYZus{}anchor=(1, 0.5))}

\PY{n}{plt}\PY{o}{.}\PY{n}{show}\PY{p}{(}\PY{p}{)}
\end{Verbatim}
\end{tcolorbox}

    \begin{center}
    \adjustimage{max size={0.9\linewidth}{0.9\paperheight}}{figs/bvp_kM_en/output_259_0.png}
    \end{center}
    { \hspace*{\fill} \\}
    
    \begin{tcolorbox}[breakable, size=fbox, boxrule=1pt, pad at break*=1mm,colback=cellbackground, colframe=cellborder]
\prompt{In}{incolor}{311}{\boxspacing}
\begin{Verbatim}[commandchars=\\\{\}]
\PY{c+c1}{\PYZsh{} Plots of dkn vs R}

\PY{n}{Bi} \PY{o}{=} \PY{p}{[}\PY{l+m+mi}{3}\PY{p}{,}\PY{l+m+mi}{6}\PY{p}{,}\PY{l+m+mi}{9}\PY{p}{,}\PY{l+m+mi}{12}\PY{p}{,}\PY{l+m+mi}{15}\PY{p}{,}\PY{l+m+mi}{18}\PY{p}{]}

\PY{n}{ilim} \PY{o}{=} \PY{n+nb}{len}\PY{p}{(}\PY{n}{Bi}\PY{p}{)}

\PY{n}{fig}\PY{p}{,} \PY{n}{ax} \PY{o}{=} \PY{n}{plt}\PY{o}{.}\PY{n}{subplots}\PY{p}{(}\PY{p}{)}
\PY{n}{plt}\PY{o}{.}\PY{n}{title}\PY{p}{(}\PY{l+s+sa}{f}\PY{l+s+s2}{\PYZdq{}}\PY{l+s+s2}{Change in k Due to Nonlinearity vs R for Various B Values}\PY{l+s+s2}{\PYZdq{}}\PY{p}{)}
\PY{n}{ax}\PY{o}{.}\PY{n}{set\PYZus{}ylabel}\PY{p}{(}\PY{l+s+s1}{\PYZsq{}}\PY{l+s+s1}{Change in k (dkn) due to nonlinearity (10}\PY{l+s+se}{\PYZbs{}u2075}\PY{l+s+s1}{ m}\PY{l+s+se}{\PYZbs{}u207b}\PY{l+s+se}{\PYZbs{}u00b9}\PY{l+s+s1}{)}\PY{l+s+s1}{\PYZsq{}}\PY{p}{)}
\PY{n}{ax}\PY{o}{.}\PY{n}{set\PYZus{}xlabel}\PY{p}{(}\PY{l+s+s2}{\PYZdq{}}\PY{l+s+s2}{Radius R of the ring (10}\PY{l+s+se}{\PYZbs{}u207b}\PY{l+s+se}{\PYZbs{}u2076}\PY{l+s+s2}{ m)}\PY{l+s+s2}{\PYZdq{}}\PY{p}{)}

\PY{k}{for} \PY{n}{bb} \PY{o+ow}{in} \PY{n}{Bi}\PY{p}{:}
    \PY{n}{ax}\PY{o}{.}\PY{n}{plot}\PY{p}{(}\PY{n}{R}\PY{p}{,} \PY{n}{dkn\PYZus{}arr}\PY{p}{[}\PY{p}{:}\PY{p}{,}\PY{n}{bb}\PY{p}{]}\PY{o}{/}\PY{l+m+mi}{100000}\PY{p}{,} \PY{n}{label} \PY{o}{=} \PY{l+s+sa}{f}\PY{l+s+s2}{\PYZdq{}}\PY{l+s+s2}{B = }\PY{l+s+si}{\PYZob{}}\PY{n}{B}\PY{p}{[}\PY{n}{bb}\PY{p}{]}\PY{l+s+si}{:}\PY{l+s+s2}{.4f}\PY{l+s+si}{\PYZcb{}}\PY{l+s+s2}{ * B\PYZus{}max}\PY{l+s+s2}{\PYZdq{}}\PY{p}{)}

\PY{n}{plt}\PY{o}{.}\PY{n}{legend}\PY{p}{(}\PY{p}{)}

\PY{n}{plt}\PY{o}{.}\PY{n}{show}\PY{p}{(}\PY{p}{)}
\end{Verbatim}
\end{tcolorbox}

    \begin{center}
    \adjustimage{max size={0.9\linewidth}{0.9\paperheight}}{figs/bvp_kM_en/output_260_0.png}
    \end{center}
    { \hspace*{\fill} \\}
    
    The oscillations here appear to have a parabolic maximum absolute value
and a linear minimum absolute value. There is a point of intersection of
all the lines at a value close to \(R=1\), but the exact location of
that point may have to do with the choise of point on the fast
oscillations. I avoided the minima, since despite the fact that it would
be the most insightful, there is a problem relating to the spikes which
appear only at the minima. They are either error, or a relic of the fact
that we are not exactly on our \(B\) node.

I now want to check that the oscillations are indeed our known
\(B\pi^2R^2\) oscillations.

    \begin{tcolorbox}[breakable, size=fbox, boxrule=1pt, pad at break*=1mm,colback=cellbackground, colframe=cellborder]
\prompt{In}{incolor}{314}{\boxspacing}
\begin{Verbatim}[commandchars=\\\{\}]
\PY{n}{Bi} \PY{o}{=} \PY{p}{[}\PY{l+m+mi}{3}\PY{p}{,}\PY{l+m+mi}{6}\PY{p}{,}\PY{l+m+mi}{9}\PY{p}{,}\PY{l+m+mi}{12}\PY{p}{,}\PY{l+m+mi}{15}\PY{p}{,}\PY{l+m+mi}{18}\PY{p}{]}

\PY{n}{ilim} \PY{o}{=} \PY{n+nb}{len}\PY{p}{(}\PY{n}{Bi}\PY{p}{)}


\PY{n}{fig}\PY{p}{,} \PY{n}{ax} \PY{o}{=} \PY{n}{plt}\PY{o}{.}\PY{n}{subplots}\PY{p}{(}\PY{p}{)}
\PY{n}{plt}\PY{o}{.}\PY{n}{title}\PY{p}{(}\PY{l+s+sa}{f}\PY{l+s+s2}{\PYZdq{}}\PY{l+s+s2}{Change in k Due to Nonlinearity vs Number of Oscillations}\PY{l+s+se}{\PYZbs{}n}\PY{l+s+s2}{for Various B Values}\PY{l+s+s2}{\PYZdq{}}\PY{p}{)}
\PY{n}{ax}\PY{o}{.}\PY{n}{set\PYZus{}ylabel}\PY{p}{(}\PY{l+s+s1}{\PYZsq{}}\PY{l+s+s1}{Change in k (dkn) due to nonlinearity (10}\PY{l+s+se}{\PYZbs{}u2075}\PY{l+s+s1}{ m}\PY{l+s+se}{\PYZbs{}u207b}\PY{l+s+se}{\PYZbs{}u00b9}\PY{l+s+s1}{)}\PY{l+s+s1}{\PYZsq{}}\PY{p}{)}
\PY{n}{ax}\PY{o}{.}\PY{n}{set\PYZus{}xlabel}\PY{p}{(}\PY{l+s+s2}{\PYZdq{}}\PY{l+s+s2}{B * π}\PY{l+s+se}{\PYZbs{}u00b2}\PY{l+s+s2}{ * R}\PY{l+s+se}{\PYZbs{}u00b2}\PY{l+s+s2}{ with R and B as proportions of their maximum values}\PY{l+s+s2}{\PYZdq{}}\PY{p}{)}

\PY{k}{for} \PY{n}{bb} \PY{o+ow}{in} \PY{n}{Bi}\PY{p}{:}
    \PY{c+c1}{\PYZsh{} x = B * pi\PYZca{}2 * R\PYZca{}2}
    \PY{n}{cycle\PYZus{}number} \PY{o}{=} \PY{n}{np}\PY{o}{.}\PY{n}{power}\PY{p}{(}\PY{n}{np}\PY{o}{.}\PY{n}{array}\PY{p}{(}\PY{n}{R}\PY{p}{)}\PY{p}{,} \PY{l+m+mi}{2}\PY{p}{)}\PY{o}{*}\PY{n}{pi}\PY{o}{*}\PY{o}{*}\PY{l+m+mi}{2}\PY{o}{*}\PY{n}{B}\PY{p}{[}\PY{n}{bb}\PY{p}{]}
    \PY{c+c1}{\PYZsh{} and plot}
    \PY{n}{ax}\PY{o}{.}\PY{n}{plot}\PY{p}{(}\PY{n}{cycle\PYZus{}number}\PY{p}{,} \PY{n}{dkn\PYZus{}arr}\PY{p}{[}\PY{p}{:}\PY{p}{,}\PY{n}{bb}\PY{p}{]}\PY{o}{/}\PY{l+m+mi}{100000}\PY{p}{,} \PY{n}{label} \PY{o}{=} \PY{l+s+sa}{f}\PY{l+s+s2}{\PYZdq{}}\PY{l+s+s2}{B = }\PY{l+s+si}{\PYZob{}}\PY{n}{B}\PY{p}{[}\PY{n}{bb}\PY{p}{]}\PY{l+s+si}{:}\PY{l+s+s2}{.4f}\PY{l+s+si}{\PYZcb{}}\PY{l+s+s2}{ * B\PYZus{}max}\PY{l+s+s2}{\PYZdq{}}\PY{p}{)}
\PY{c+c1}{\PYZsh{} Count the number of cycles.}
\PY{n}{plt}\PY{o}{.}\PY{n}{axline}\PY{p}{(}\PY{p}{(}\PY{l+m+mi}{2}\PY{p}{,}\PY{o}{\PYZhy{}}\PY{l+m+mf}{0.6}\PY{p}{)}\PY{p}{,} \PY{p}{(}\PY{l+m+mi}{2}\PY{p}{,}\PY{o}{\PYZhy{}}\PY{l+m+mf}{0.66}\PY{p}{)}\PY{p}{,} \PY{n}{color}\PY{o}{=}\PY{l+s+s1}{\PYZsq{}}\PY{l+s+s1}{k}\PY{l+s+s1}{\PYZsq{}}\PY{p}{,} \PY{n}{label}\PY{o}{=}\PY{l+s+s2}{\PYZdq{}}\PY{l+s+s2}{period boundary}\PY{l+s+s2}{\PYZdq{}}\PY{p}{)}
\PY{n}{plt}\PY{o}{.}\PY{n}{axline}\PY{p}{(}\PY{p}{(}\PY{l+m+mf}{2.5}\PY{p}{,}\PY{o}{\PYZhy{}}\PY{l+m+mf}{0.6}\PY{p}{)}\PY{p}{,} \PY{p}{(}\PY{l+m+mf}{2.5}\PY{p}{,}\PY{o}{\PYZhy{}}\PY{l+m+mf}{0.66}\PY{p}{)}\PY{p}{,} \PY{n}{color}\PY{o}{=}\PY{l+s+s1}{\PYZsq{}}\PY{l+s+s1}{b}\PY{l+s+s1}{\PYZsq{}}\PY{p}{,} \PY{n}{label} \PY{o}{=} \PY{l+s+s2}{\PYZdq{}}\PY{l+s+s2}{period midpoint}\PY{l+s+s2}{\PYZdq{}}\PY{p}{)}
\PY{n}{plt}\PY{o}{.}\PY{n}{axline}\PY{p}{(}\PY{p}{(}\PY{l+m+mi}{3}\PY{p}{,}\PY{o}{\PYZhy{}}\PY{l+m+mf}{0.6}\PY{p}{)}\PY{p}{,} \PY{p}{(}\PY{l+m+mi}{3}\PY{p}{,}\PY{o}{\PYZhy{}}\PY{l+m+mf}{0.66}\PY{p}{)}\PY{p}{,} \PY{n}{color}\PY{o}{=}\PY{l+s+s1}{\PYZsq{}}\PY{l+s+s1}{k}\PY{l+s+s1}{\PYZsq{}}\PY{p}{)}
\PY{n}{plt}\PY{o}{.}\PY{n}{axline}\PY{p}{(}\PY{p}{(}\PY{l+m+mf}{3.5}\PY{p}{,}\PY{o}{\PYZhy{}}\PY{l+m+mf}{0.6}\PY{p}{)}\PY{p}{,} \PY{p}{(}\PY{l+m+mf}{3.5}\PY{p}{,}\PY{o}{\PYZhy{}}\PY{l+m+mf}{0.66}\PY{p}{)}\PY{p}{,} \PY{n}{color}\PY{o}{=}\PY{l+s+s1}{\PYZsq{}}\PY{l+s+s1}{b}\PY{l+s+s1}{\PYZsq{}}\PY{p}{)}
\PY{n}{plt}\PY{o}{.}\PY{n}{axline}\PY{p}{(}\PY{p}{(}\PY{l+m+mi}{4}\PY{p}{,}\PY{o}{\PYZhy{}}\PY{l+m+mf}{0.6}\PY{p}{)}\PY{p}{,} \PY{p}{(}\PY{l+m+mi}{4}\PY{p}{,}\PY{o}{\PYZhy{}}\PY{l+m+mf}{0.66}\PY{p}{)}\PY{p}{,} \PY{n}{color}\PY{o}{=}\PY{l+s+s1}{\PYZsq{}}\PY{l+s+s1}{k}\PY{l+s+s1}{\PYZsq{}}\PY{p}{)}
\PY{n}{plt}\PY{o}{.}\PY{n}{axline}\PY{p}{(}\PY{p}{(}\PY{l+m+mf}{4.5}\PY{p}{,}\PY{o}{\PYZhy{}}\PY{l+m+mf}{0.6}\PY{p}{)}\PY{p}{,} \PY{p}{(}\PY{l+m+mf}{4.5}\PY{p}{,}\PY{o}{\PYZhy{}}\PY{l+m+mf}{0.66}\PY{p}{)}\PY{p}{,} \PY{n}{color}\PY{o}{=}\PY{l+s+s1}{\PYZsq{}}\PY{l+s+s1}{b}\PY{l+s+s1}{\PYZsq{}}\PY{p}{)}
\PY{n}{plt}\PY{o}{.}\PY{n}{axline}\PY{p}{(}\PY{p}{(}\PY{l+m+mi}{5}\PY{p}{,}\PY{o}{\PYZhy{}}\PY{l+m+mf}{0.6}\PY{p}{)}\PY{p}{,} \PY{p}{(}\PY{l+m+mi}{5}\PY{p}{,}\PY{o}{\PYZhy{}}\PY{l+m+mf}{0.66}\PY{p}{)}\PY{p}{,} \PY{n}{color}\PY{o}{=}\PY{l+s+s1}{\PYZsq{}}\PY{l+s+s1}{k}\PY{l+s+s1}{\PYZsq{}}\PY{p}{)}

\PY{n}{plt}\PY{o}{.}\PY{n}{legend}\PY{p}{(}\PY{p}{)}

\PY{n}{plt}\PY{o}{.}\PY{n}{show}\PY{p}{(}\PY{p}{)}
\end{Verbatim}
\end{tcolorbox}

    \begin{center}
    \adjustimage{max size={0.9\linewidth}{0.9\paperheight}}{figs/bvp_kM_en/output_262_0.png}
    \end{center}
    { \hspace*{\fill} \\}
    
    The minima are close to or on the cycle boundary, while the maxima are
close to or on the half period. The period is definitely correct. There
may be a slight translation, but that could actually be expected from
our skewed secant squared function fit.

    The next series of plots are for the \(dkn\) vs \(B\) data, with large
\(R\) variations being represented by different curves.

    \begin{tcolorbox}[breakable, size=fbox, boxrule=1pt, pad at break*=1mm,colback=cellbackground, colframe=cellborder]
\prompt{In}{incolor}{321}{\boxspacing}
\begin{Verbatim}[commandchars=\\\{\}]
\PY{c+c1}{\PYZsh{} Low R values}
\PY{n}{Ri} \PY{o}{=} \PY{p}{[}\PY{l+m+mi}{5}\PY{p}{,} \PY{l+m+mi}{10}\PY{p}{,} \PY{l+m+mi}{15}\PY{p}{,} \PY{l+m+mi}{20}\PY{p}{,} \PY{l+m+mi}{25}\PY{p}{]}
\PY{n}{Bi} \PY{o}{=} \PY{p}{[}\PY{l+m+mi}{3}\PY{p}{,}\PY{l+m+mi}{6}\PY{p}{,}\PY{l+m+mi}{9}\PY{p}{,}\PY{l+m+mi}{12}\PY{p}{,}\PY{l+m+mi}{15}\PY{p}{,}\PY{l+m+mi}{18}\PY{p}{]}

\PY{n}{ilim} \PY{o}{=} \PY{n+nb}{len}\PY{p}{(}\PY{n}{Bi}\PY{p}{)}

\PY{n}{fig}\PY{p}{,} \PY{n}{ax} \PY{o}{=} \PY{n}{plt}\PY{o}{.}\PY{n}{subplots}\PY{p}{(}\PY{p}{)}
\PY{n}{plt}\PY{o}{.}\PY{n}{title}\PY{p}{(}\PY{l+s+sa}{f}\PY{l+s+s2}{\PYZdq{}}\PY{l+s+s2}{Change in k Due to Nonlinearity vs B for Various R Values}\PY{l+s+se}{\PYZbs{}n}\PY{l+s+s2}{With their Expected Cycle Boundary Marked}\PY{l+s+s2}{\PYZdq{}}\PY{p}{)}
\PY{n}{ax}\PY{o}{.}\PY{n}{set\PYZus{}ylabel}\PY{p}{(}\PY{l+s+s1}{\PYZsq{}}\PY{l+s+s1}{Change in k (dkn) due to nonlinearity (10}\PY{l+s+se}{\PYZbs{}u2075}\PY{l+s+s1}{ m}\PY{l+s+se}{\PYZbs{}u207b}\PY{l+s+se}{\PYZbs{}u00b9}\PY{l+s+s1}{)}\PY{l+s+s1}{\PYZsq{}}\PY{p}{)}
\PY{n}{ax}\PY{o}{.}\PY{n}{set\PYZus{}xlabel}\PY{p}{(}\PY{l+s+s2}{\PYZdq{}}\PY{l+s+s2}{Magnetic field strength B (B = B * B\PYZus{}max)}\PY{l+s+s2}{\PYZdq{}}\PY{p}{)}

\PY{k}{for} \PY{n}{rr} \PY{o+ow}{in} \PY{n}{Ri}\PY{p}{:}
    \PY{n}{ax}\PY{o}{.}\PY{n}{plot}\PY{p}{(}\PY{n}{B}\PY{p}{,} \PY{n}{dkn\PYZus{}arr}\PY{p}{[}\PY{n}{rr}\PY{p}{,}\PY{p}{:}\PY{p}{]}\PY{o}{/}\PY{l+m+mi}{100000}\PY{p}{,} \PY{n}{label} \PY{o}{=} \PY{l+s+sa}{f}\PY{l+s+s2}{\PYZdq{}}\PY{l+s+s2}{R = }\PY{l+s+si}{\PYZob{}}\PY{n}{R}\PY{p}{[}\PY{n}{rr}\PY{p}{]}\PY{l+s+si}{:}\PY{l+s+s2}{.5f}\PY{l+s+si}{\PYZcb{}}\PY{l+s+s2}{ * 10}\PY{l+s+se}{\PYZbs{}u207b}\PY{l+s+se}{\PYZbs{}u2076}\PY{l+s+s2}{ m}\PY{l+s+s2}{\PYZdq{}}\PY{p}{)}
    \PY{n+nb}{print}\PY{p}{(}\PY{l+s+sa}{f}\PY{l+s+s2}{\PYZdq{}}\PY{l+s+s2}{Minimum absolute value of dkn for R = }\PY{l+s+si}{\PYZob{}}\PY{n}{R}\PY{p}{[}\PY{n}{rr}\PY{p}{]}\PY{l+s+si}{:}\PY{l+s+s2}{.5f}\PY{l+s+si}{\PYZcb{}}\PY{l+s+s2}{: }\PY{l+s+si}{\PYZob{}}\PY{o}{\PYZhy{}}\PY{n}{np}\PY{o}{.}\PY{n}{max}\PY{p}{(}\PY{+w}{ }\PY{n}{dkn\PYZus{}arr}\PY{p}{[}\PY{n}{rr}\PY{p}{,}\PY{p}{:}\PY{p}{]}\PY{p}{)}\PY{l+s+si}{:}\PY{l+s+s2}{.5f}\PY{l+s+si}{\PYZcb{}}\PY{l+s+s2}{\PYZdq{}}\PY{p}{)}
    \PY{c+c1}{\PYZsh{} The black lines printed will designate the expected cycle boundary}
    \PY{n}{plt}\PY{o}{.}\PY{n}{axline}\PY{p}{(}\PY{p}{(}\PY{l+m+mi}{2}\PY{o}{/}\PY{p}{(}\PY{n}{R}\PY{p}{[}\PY{n}{rr}\PY{p}{]}\PY{o}{*}\PY{o}{*}\PY{l+m+mi}{2} \PY{o}{*} \PY{n}{pi}\PY{o}{*}\PY{o}{*}\PY{l+m+mi}{2}\PY{p}{)}\PY{p}{,}\PY{o}{\PYZhy{}}\PY{l+m+mf}{0.45}\PY{p}{)}\PY{p}{,} \PY{p}{(}\PY{l+m+mi}{2}\PY{o}{/}\PY{p}{(}\PY{n}{R}\PY{p}{[}\PY{n}{rr}\PY{p}{]}\PY{o}{*}\PY{o}{*}\PY{l+m+mi}{2} \PY{o}{*} \PY{n}{pi}\PY{o}{*}\PY{o}{*}\PY{l+m+mi}{2}\PY{p}{)}\PY{p}{,}\PY{o}{\PYZhy{}}\PY{l+m+mf}{0.5}\PY{p}{)}\PY{p}{,} \PY{n}{color}\PY{o}{=}\PY{l+s+s1}{\PYZsq{}}\PY{l+s+s1}{k}\PY{l+s+s1}{\PYZsq{}}\PY{p}{)}

\PY{n}{plt}\PY{o}{.}\PY{n}{legend}\PY{p}{(}\PY{p}{)}

\PY{n}{plt}\PY{o}{.}\PY{n}{show}\PY{p}{(}\PY{p}{)}
\end{Verbatim}
\end{tcolorbox}

    \begin{Verbatim}[commandchars=\\\{\}]
Minimum absolute value of dkn for R = 0.57509: 37257.13817
Minimum absolute value of dkn for R = 0.57926: 36836.36141
Minimum absolute value of dkn for R = 0.58342: 36899.82919
Minimum absolute value of dkn for R = 0.58759: 37034.96513
Minimum absolute value of dkn for R = 0.59176: 37128.13734
    \end{Verbatim}

    \begin{center}
    \adjustimage{max size={0.9\linewidth}{0.9\paperheight}}{figs/bvp_kM_en/output_265_1.png}
    \end{center}
    { \hspace*{\fill} \\}
    
    \begin{tcolorbox}[breakable, size=fbox, boxrule=1pt, pad at break*=1mm,colback=cellbackground, colframe=cellborder]
\prompt{In}{incolor}{320}{\boxspacing}
\begin{Verbatim}[commandchars=\\\{\}]
\PY{c+c1}{\PYZsh{} R is now increasing }

\PY{n}{Ri} \PY{o}{=} \PY{p}{[}\PY{l+m+mi}{30}\PY{p}{,} \PY{l+m+mi}{35}\PY{p}{,} \PY{l+m+mi}{40}\PY{p}{,} \PY{l+m+mi}{45}\PY{p}{,} \PY{l+m+mi}{50}\PY{p}{]}
\PY{n}{Bi} \PY{o}{=} \PY{p}{[}\PY{l+m+mi}{3}\PY{p}{,}\PY{l+m+mi}{6}\PY{p}{,}\PY{l+m+mi}{9}\PY{p}{,}\PY{l+m+mi}{12}\PY{p}{,}\PY{l+m+mi}{15}\PY{p}{,}\PY{l+m+mi}{18}\PY{p}{]}

\PY{n}{ilim} \PY{o}{=} \PY{n+nb}{len}\PY{p}{(}\PY{n}{Bi}\PY{p}{)}

\PY{n}{fig}\PY{p}{,} \PY{n}{ax} \PY{o}{=} \PY{n}{plt}\PY{o}{.}\PY{n}{subplots}\PY{p}{(}\PY{p}{)}
\PY{n}{plt}\PY{o}{.}\PY{n}{title}\PY{p}{(}\PY{l+s+sa}{f}\PY{l+s+s2}{\PYZdq{}}\PY{l+s+s2}{Change in k Due to Nonlinearity vs B for Various R Values}\PY{l+s+se}{\PYZbs{}n}\PY{l+s+s2}{With their Expected Cycle Boundary Marked}\PY{l+s+s2}{\PYZdq{}}\PY{p}{)}
\PY{n}{ax}\PY{o}{.}\PY{n}{set\PYZus{}ylabel}\PY{p}{(}\PY{l+s+s1}{\PYZsq{}}\PY{l+s+s1}{Change in k (dkn) due to nonlinearity (10}\PY{l+s+se}{\PYZbs{}u2075}\PY{l+s+s1}{ m}\PY{l+s+se}{\PYZbs{}u207b}\PY{l+s+se}{\PYZbs{}u00b9}\PY{l+s+s1}{)}\PY{l+s+s1}{\PYZsq{}}\PY{p}{)}
\PY{n}{ax}\PY{o}{.}\PY{n}{set\PYZus{}xlabel}\PY{p}{(}\PY{l+s+s2}{\PYZdq{}}\PY{l+s+s2}{Magnetic field strength B (B = B * B\PYZus{}max)}\PY{l+s+s2}{\PYZdq{}}\PY{p}{)}

\PY{k}{for} \PY{n}{rr} \PY{o+ow}{in} \PY{n}{Ri}\PY{p}{:}
    \PY{n}{ax}\PY{o}{.}\PY{n}{plot}\PY{p}{(}\PY{n}{B}\PY{p}{,} \PY{n}{dkn\PYZus{}arr}\PY{p}{[}\PY{n}{rr}\PY{p}{,}\PY{p}{:}\PY{p}{]}\PY{o}{/}\PY{l+m+mi}{100000}\PY{p}{,} \PY{n}{label} \PY{o}{=} \PY{l+s+sa}{f}\PY{l+s+s2}{\PYZdq{}}\PY{l+s+s2}{R = }\PY{l+s+si}{\PYZob{}}\PY{n}{R}\PY{p}{[}\PY{n}{rr}\PY{p}{]}\PY{l+s+si}{:}\PY{l+s+s2}{.5f}\PY{l+s+si}{\PYZcb{}}\PY{l+s+s2}{ * 10}\PY{l+s+se}{\PYZbs{}u207b}\PY{l+s+se}{\PYZbs{}u2076}\PY{l+s+s2}{ m}\PY{l+s+s2}{\PYZdq{}}\PY{p}{)}
    \PY{n+nb}{print}\PY{p}{(}\PY{l+s+sa}{f}\PY{l+s+s2}{\PYZdq{}}\PY{l+s+s2}{Minimum absolute value of dkn for R = }\PY{l+s+si}{\PYZob{}}\PY{n}{R}\PY{p}{[}\PY{n}{rr}\PY{p}{]}\PY{l+s+si}{:}\PY{l+s+s2}{.5f}\PY{l+s+si}{\PYZcb{}}\PY{l+s+s2}{: }\PY{l+s+si}{\PYZob{}}\PY{o}{\PYZhy{}}\PY{n}{np}\PY{o}{.}\PY{n}{max}\PY{p}{(}\PY{+w}{ }\PY{n}{dkn\PYZus{}arr}\PY{p}{[}\PY{n}{rr}\PY{p}{,}\PY{p}{:}\PY{p}{]}\PY{p}{)}\PY{l+s+si}{:}\PY{l+s+s2}{.5f}\PY{l+s+si}{\PYZcb{}}\PY{l+s+s2}{\PYZdq{}}\PY{p}{)}
    \PY{c+c1}{\PYZsh{} The black lines printed here designate the expected cycle boundary}
    \PY{n}{plt}\PY{o}{.}\PY{n}{axline}\PY{p}{(}\PY{p}{(}\PY{l+m+mi}{2}\PY{o}{/}\PY{p}{(}\PY{n}{R}\PY{p}{[}\PY{n}{rr}\PY{p}{]}\PY{o}{*}\PY{o}{*}\PY{l+m+mi}{2} \PY{o}{*} \PY{n}{pi}\PY{o}{*}\PY{o}{*}\PY{l+m+mi}{2}\PY{p}{)}\PY{p}{,}\PY{o}{\PYZhy{}}\PY{l+m+mf}{0.45}\PY{p}{)}\PY{p}{,} \PY{p}{(}\PY{l+m+mi}{2}\PY{o}{/}\PY{p}{(}\PY{n}{R}\PY{p}{[}\PY{n}{rr}\PY{p}{]}\PY{o}{*}\PY{o}{*}\PY{l+m+mi}{2} \PY{o}{*} \PY{n}{pi}\PY{o}{*}\PY{o}{*}\PY{l+m+mi}{2}\PY{p}{)}\PY{p}{,}\PY{o}{\PYZhy{}}\PY{l+m+mf}{0.5}\PY{p}{)}\PY{p}{,} \PY{n}{color}\PY{o}{=}\PY{l+s+s1}{\PYZsq{}}\PY{l+s+s1}{k}\PY{l+s+s1}{\PYZsq{}}\PY{p}{)}

\PY{n}{plt}\PY{o}{.}\PY{n}{legend}\PY{p}{(}\PY{p}{)}

\PY{n}{plt}\PY{o}{.}\PY{n}{show}\PY{p}{(}\PY{p}{)}
\end{Verbatim}
\end{tcolorbox}

    \begin{Verbatim}[commandchars=\\\{\}]
Minimum absolute value of dkn for R = 0.59592: 37255.42202
Minimum absolute value of dkn for R = 0.60009: 37358.95283
Minimum absolute value of dkn for R = 0.60426: 37489.68861
Minimum absolute value of dkn for R = 0.60842: 37589.56458
Minimum absolute value of dkn for R = 0.61259: 37731.22027
    \end{Verbatim}

    \begin{center}
    \adjustimage{max size={0.9\linewidth}{0.9\paperheight}}{figs/bvp_kM_en/output_266_1.png}
    \end{center}
    { \hspace*{\fill} \\}
    
    \begin{tcolorbox}[breakable, size=fbox, boxrule=1pt, pad at break*=1mm,colback=cellbackground, colframe=cellborder]
\prompt{In}{incolor}{319}{\boxspacing}
\begin{Verbatim}[commandchars=\\\{\}]
\PY{n}{Ri} \PY{o}{=} \PY{p}{[}\PY{l+m+mi}{55}\PY{p}{,} \PY{l+m+mi}{60}\PY{p}{,} \PY{l+m+mi}{65}\PY{p}{,} \PY{l+m+mi}{70}\PY{p}{,} \PY{l+m+mi}{75}\PY{p}{,} \PY{l+m+mi}{80}\PY{p}{]}
\PY{n}{Bi} \PY{o}{=} \PY{p}{[}\PY{l+m+mi}{3}\PY{p}{,}\PY{l+m+mi}{6}\PY{p}{,}\PY{l+m+mi}{9}\PY{p}{,}\PY{l+m+mi}{12}\PY{p}{,}\PY{l+m+mi}{15}\PY{p}{,}\PY{l+m+mi}{18}\PY{p}{]}

\PY{n}{ilim} \PY{o}{=} \PY{n+nb}{len}\PY{p}{(}\PY{n}{Bi}\PY{p}{)}

\PY{n}{fig}\PY{p}{,} \PY{n}{ax} \PY{o}{=} \PY{n}{plt}\PY{o}{.}\PY{n}{subplots}\PY{p}{(}\PY{p}{)}
\PY{n}{plt}\PY{o}{.}\PY{n}{title}\PY{p}{(}\PY{l+s+sa}{f}\PY{l+s+s2}{\PYZdq{}}\PY{l+s+s2}{Change in k Due to Nonlinearity vs B for Various R Values}\PY{l+s+se}{\PYZbs{}n}\PY{l+s+s2}{With their Expected Cycle Boundary Marked}\PY{l+s+s2}{\PYZdq{}}\PY{p}{)}
\PY{n}{ax}\PY{o}{.}\PY{n}{set\PYZus{}ylabel}\PY{p}{(}\PY{l+s+s1}{\PYZsq{}}\PY{l+s+s1}{Change in k (dkn) due to nonlinearity (10}\PY{l+s+se}{\PYZbs{}u2075}\PY{l+s+s1}{ m}\PY{l+s+se}{\PYZbs{}u207b}\PY{l+s+se}{\PYZbs{}u00b9}\PY{l+s+s1}{)}\PY{l+s+s1}{\PYZsq{}}\PY{p}{)}
\PY{n}{ax}\PY{o}{.}\PY{n}{set\PYZus{}xlabel}\PY{p}{(}\PY{l+s+s2}{\PYZdq{}}\PY{l+s+s2}{Magnetic field strength B (B = B * B\PYZus{}max)}\PY{l+s+s2}{\PYZdq{}}\PY{p}{)}

\PY{k}{for} \PY{n}{rr} \PY{o+ow}{in} \PY{n}{Ri}\PY{p}{:}
    \PY{n}{ax}\PY{o}{.}\PY{n}{plot}\PY{p}{(}\PY{n}{B}\PY{p}{,} \PY{n}{dkn\PYZus{}arr}\PY{p}{[}\PY{n}{rr}\PY{p}{,}\PY{p}{:}\PY{p}{]}\PY{o}{/}\PY{l+m+mi}{100000}\PY{p}{,} \PY{n}{label} \PY{o}{=} \PY{l+s+sa}{f}\PY{l+s+s2}{\PYZdq{}}\PY{l+s+s2}{R = }\PY{l+s+si}{\PYZob{}}\PY{n}{R}\PY{p}{[}\PY{n}{rr}\PY{p}{]}\PY{l+s+si}{:}\PY{l+s+s2}{.5f}\PY{l+s+si}{\PYZcb{}}\PY{l+s+s2}{ * 10}\PY{l+s+se}{\PYZbs{}u207b}\PY{l+s+se}{\PYZbs{}u2076}\PY{l+s+s2}{ m}\PY{l+s+s2}{\PYZdq{}}\PY{p}{)}
    \PY{n+nb}{print}\PY{p}{(}\PY{l+s+sa}{f}\PY{l+s+s2}{\PYZdq{}}\PY{l+s+s2}{Minimum absolute value of dkn for R = }\PY{l+s+si}{\PYZob{}}\PY{n}{R}\PY{p}{[}\PY{n}{rr}\PY{p}{]}\PY{l+s+si}{:}\PY{l+s+s2}{.5f}\PY{l+s+si}{\PYZcb{}}\PY{l+s+s2}{: }\PY{l+s+si}{\PYZob{}}\PY{o}{\PYZhy{}}\PY{n}{np}\PY{o}{.}\PY{n}{max}\PY{p}{(}\PY{+w}{ }\PY{n}{dkn\PYZus{}arr}\PY{p}{[}\PY{n}{rr}\PY{p}{,}\PY{p}{:}\PY{p}{]}\PY{p}{)}\PY{l+s+si}{:}\PY{l+s+s2}{.5f}\PY{l+s+si}{\PYZcb{}}\PY{l+s+s2}{\PYZdq{}}\PY{p}{)}
    \PY{c+c1}{\PYZsh{} The black lines printed here designate the expected cycle boundary}
    \PY{n}{plt}\PY{o}{.}\PY{n}{axline}\PY{p}{(}\PY{p}{(}\PY{l+m+mi}{2}\PY{o}{/}\PY{p}{(}\PY{n}{R}\PY{p}{[}\PY{n}{rr}\PY{p}{]}\PY{o}{*}\PY{o}{*}\PY{l+m+mi}{2} \PY{o}{*} \PY{n}{pi}\PY{o}{*}\PY{o}{*}\PY{l+m+mi}{2}\PY{p}{)}\PY{p}{,}\PY{o}{\PYZhy{}}\PY{l+m+mf}{0.45}\PY{p}{)}\PY{p}{,} \PY{p}{(}\PY{l+m+mi}{2}\PY{o}{/}\PY{p}{(}\PY{n}{R}\PY{p}{[}\PY{n}{rr}\PY{p}{]}\PY{o}{*}\PY{o}{*}\PY{l+m+mi}{2} \PY{o}{*} \PY{n}{pi}\PY{o}{*}\PY{o}{*}\PY{l+m+mi}{2}\PY{p}{)}\PY{p}{,}\PY{o}{\PYZhy{}}\PY{l+m+mf}{0.5}\PY{p}{)}\PY{p}{,} \PY{n}{color}\PY{o}{=}\PY{l+s+s1}{\PYZsq{}}\PY{l+s+s1}{k}\PY{l+s+s1}{\PYZsq{}}\PY{p}{)}

\PY{n}{plt}\PY{o}{.}\PY{n}{legend}\PY{p}{(}\PY{p}{)}

\PY{n}{plt}\PY{o}{.}\PY{n}{show}\PY{p}{(}\PY{p}{)}
\end{Verbatim}
\end{tcolorbox}

    \begin{Verbatim}[commandchars=\\\{\}]
Minimum absolute value of dkn for R = 0.61676: 37827.62983
Minimum absolute value of dkn for R = 0.62092: 37947.87134
Minimum absolute value of dkn for R = 0.62509: 38090.30873
Minimum absolute value of dkn for R = 0.62926: 38186.86325
Minimum absolute value of dkn for R = 0.63342: 38304.87916
Minimum absolute value of dkn for R = 0.63759: 38439.78907
    \end{Verbatim}

    \begin{center}
    \adjustimage{max size={0.9\linewidth}{0.9\paperheight}}{figs/bvp_kM_en/output_267_1.png}
    \end{center}
    { \hspace*{\fill} \\}
    
    The black line appears to be hitting the minima of its curve. The
minimum absolute value appears to be increasing, although this increase
is not visible. It appears that the curve is moving to the right with a
constant speed. Changes to the shape are minimal. No new behaviors are
visible.

    Now for the large-scale oscillation dependence of the slow oscillation
wave number.

    \begin{tcolorbox}[breakable, size=fbox, boxrule=1pt, pad at break*=1mm,colback=cellbackground, colframe=cellborder]
\prompt{In}{incolor}{324}{\boxspacing}
\begin{Verbatim}[commandchars=\\\{\}]
\PY{c+c1}{\PYZsh{} Plot of dMn as a function of R\PYZca{}2}

\PY{n}{Bi} \PY{o}{=} \PY{p}{[}\PY{l+m+mi}{3}\PY{p}{,}\PY{l+m+mi}{6}\PY{p}{,}\PY{l+m+mi}{9}\PY{p}{,}\PY{l+m+mi}{12}\PY{p}{,}\PY{l+m+mi}{15}\PY{p}{,}\PY{l+m+mi}{18}\PY{p}{]}

\PY{n}{x} \PY{o}{=} \PY{p}{[}\PY{l+m+mi}{2}\PY{p}{,} \PY{l+m+mi}{3}\PY{p}{,} \PY{l+m+mi}{4}\PY{p}{,} \PY{l+m+mi}{5}\PY{p}{]}

\PY{n}{fig}\PY{p}{,} \PY{n}{ax} \PY{o}{=} \PY{n}{plt}\PY{o}{.}\PY{n}{subplots}\PY{p}{(}\PY{p}{)}
\PY{n}{plt}\PY{o}{.}\PY{n}{title}\PY{p}{(}\PY{l+s+sa}{f}\PY{l+s+s2}{\PYZdq{}}\PY{l+s+s2}{Change in M Due to Nonlinearity vs R}\PY{l+s+se}{\PYZbs{}u00b2}\PY{l+s+s2}{ for Various B Values}\PY{l+s+s2}{\PYZdq{}}\PY{p}{)}
\PY{n}{ax}\PY{o}{.}\PY{n}{set\PYZus{}ylabel}\PY{p}{(}\PY{l+s+s1}{\PYZsq{}}\PY{l+s+s1}{Change in M (dMn) due to nonlinearity (10}\PY{l+s+se}{\PYZbs{}u2075}\PY{l+s+s1}{ m}\PY{l+s+se}{\PYZbs{}u207b}\PY{l+s+se}{\PYZbs{}u00b9}\PY{l+s+s1}{)}\PY{l+s+s1}{\PYZsq{}}\PY{p}{)}
\PY{n}{ax}\PY{o}{.}\PY{n}{set\PYZus{}xlabel}\PY{p}{(}\PY{l+s+s2}{\PYZdq{}}\PY{l+s+s2}{Radius Squared R}\PY{l+s+se}{\PYZbs{}u00b2}\PY{l+s+s2}{ of the ring (10}\PY{l+s+se}{\PYZbs{}u207b}\PY{l+s+se}{\PYZbs{}u00b9}\PY{l+s+se}{\PYZbs{}u00b2}\PY{l+s+s2}{ m}\PY{l+s+se}{\PYZbs{}u00b2}\PY{l+s+s2}{)}\PY{l+s+s2}{\PYZdq{}}\PY{p}{)}

\PY{k}{for} \PY{n}{bb} \PY{o+ow}{in} \PY{n}{Bi}\PY{p}{:}
    \PY{c+c1}{\PYZsh{}cycle\PYZus{}number = np.power(np.array(R), 2)*pi**2*B[bb]}
    \PY{n}{ax}\PY{o}{.}\PY{n}{plot}\PY{p}{(}\PY{n}{np}\PY{o}{.}\PY{n}{power}\PY{p}{(}\PY{n}{np}\PY{o}{.}\PY{n}{array}\PY{p}{(}\PY{n}{R}\PY{p}{)}\PY{p}{,} \PY{l+m+mi}{2}\PY{p}{)}\PY{p}{,} \PY{n}{dMn\PYZus{}arr}\PY{p}{[}\PY{p}{:}\PY{p}{,}\PY{n}{bb}\PY{p}{]}\PY{o}{/}\PY{l+m+mi}{100000}\PY{p}{,} \PY{n}{label} \PY{o}{=} \PY{l+s+sa}{f}\PY{l+s+s2}{\PYZdq{}}\PY{l+s+s2}{B = }\PY{l+s+si}{\PYZob{}}\PY{n}{B}\PY{p}{[}\PY{n}{bb}\PY{p}{]}\PY{l+s+si}{:}\PY{l+s+s2}{.4f}\PY{l+s+si}{\PYZcb{}}\PY{l+s+s2}{ * B\PYZus{}max}\PY{l+s+s2}{\PYZdq{}}\PY{p}{)}
\PY{c+c1}{\PYZsh{}for ii in x:}
\PY{c+c1}{\PYZsh{}    plt.axline((ii,\PYZhy{}0.1), (ii,\PYZhy{}0.12), color=\PYZsq{}k\PYZsq{})}
\PY{n}{plt}\PY{o}{.}\PY{n}{axline}\PY{p}{(}\PY{p}{(}\PY{l+m+mf}{0.5}\PY{p}{,}\PY{l+m+mi}{0}\PY{p}{)}\PY{p}{,} \PY{p}{(}\PY{l+m+mf}{0.6}\PY{p}{,}\PY{l+m+mi}{0}\PY{p}{)}\PY{p}{,} \PY{n}{color}\PY{o}{=}\PY{l+s+s1}{\PYZsq{}}\PY{l+s+s1}{k}\PY{l+s+s1}{\PYZsq{}}\PY{p}{)}

\PY{n}{plt}\PY{o}{.}\PY{n}{legend}\PY{p}{(}\PY{n}{loc}\PY{o}{=}\PY{l+s+s1}{\PYZsq{}}\PY{l+s+s1}{center left}\PY{l+s+s1}{\PYZsq{}}\PY{p}{,} \PY{n}{bbox\PYZus{}to\PYZus{}anchor}\PY{o}{=}\PY{p}{(}\PY{l+m+mi}{1}\PY{p}{,} \PY{l+m+mf}{0.5}\PY{p}{)}\PY{p}{)}

\PY{n}{plt}\PY{o}{.}\PY{n}{show}\PY{p}{(}\PY{p}{)}
\end{Verbatim}
\end{tcolorbox}

    \begin{center}
    \adjustimage{max size={0.9\linewidth}{0.9\paperheight}}{figs/bvp_kM_en/output_270_0.png}
    \end{center}
    { \hspace*{\fill} \\}
    
    \begin{tcolorbox}[breakable, size=fbox, boxrule=1pt, pad at break*=1mm,colback=cellbackground, colframe=cellborder]
\prompt{In}{incolor}{325}{\boxspacing}
\begin{Verbatim}[commandchars=\\\{\}]
\PY{c+c1}{\PYZsh{} Plot of dMn as a function of the number of large\PYZhy{}scale periods}

\PY{n}{Bi} \PY{o}{=} \PY{p}{[}\PY{l+m+mi}{3}\PY{p}{,}\PY{l+m+mi}{6}\PY{p}{,}\PY{l+m+mi}{9}\PY{p}{,}\PY{l+m+mi}{12}\PY{p}{,}\PY{l+m+mi}{15}\PY{p}{,}\PY{l+m+mi}{18}\PY{p}{]}

\PY{n}{x} \PY{o}{=} \PY{p}{[}\PY{l+m+mi}{2}\PY{p}{,} \PY{l+m+mi}{3}\PY{p}{,} \PY{l+m+mi}{4}\PY{p}{,} \PY{l+m+mi}{5}\PY{p}{]}

\PY{n}{fig}\PY{p}{,} \PY{n}{ax} \PY{o}{=} \PY{n}{plt}\PY{o}{.}\PY{n}{subplots}\PY{p}{(}\PY{p}{)}
\PY{n}{plt}\PY{o}{.}\PY{n}{title}\PY{p}{(}\PY{l+s+sa}{f}\PY{l+s+s2}{\PYZdq{}}\PY{l+s+s2}{Change in M Due to Nonlinearity vs Number of Oscillations}\PY{l+s+se}{\PYZbs{}n}\PY{l+s+s2}{for Various B Values}\PY{l+s+s2}{\PYZdq{}}\PY{p}{)}
\PY{n}{ax}\PY{o}{.}\PY{n}{set\PYZus{}ylabel}\PY{p}{(}\PY{l+s+s1}{\PYZsq{}}\PY{l+s+s1}{Change in M (dMn) due to nonlinearity (10}\PY{l+s+se}{\PYZbs{}u2075}\PY{l+s+s1}{ m}\PY{l+s+se}{\PYZbs{}u207b}\PY{l+s+se}{\PYZbs{}u00b9}\PY{l+s+s1}{)}\PY{l+s+s1}{\PYZsq{}}\PY{p}{)}
\PY{n}{ax}\PY{o}{.}\PY{n}{set\PYZus{}xlabel}\PY{p}{(}\PY{l+s+s2}{\PYZdq{}}\PY{l+s+s2}{B * π}\PY{l+s+se}{\PYZbs{}u00b2}\PY{l+s+s2}{ * R}\PY{l+s+se}{\PYZbs{}u00b2}\PY{l+s+s2}{ with R and B as proportions of their maximum values}\PY{l+s+s2}{\PYZdq{}}\PY{p}{)}

\PY{k}{for} \PY{n}{bb} \PY{o+ow}{in} \PY{n}{Bi}\PY{p}{:}
    \PY{n}{cycle\PYZus{}number} \PY{o}{=} \PY{n}{np}\PY{o}{.}\PY{n}{power}\PY{p}{(}\PY{n}{np}\PY{o}{.}\PY{n}{array}\PY{p}{(}\PY{n}{R}\PY{p}{)}\PY{p}{,} \PY{l+m+mi}{2}\PY{p}{)}\PY{o}{*}\PY{n}{pi}\PY{o}{*}\PY{o}{*}\PY{l+m+mi}{2}\PY{o}{*}\PY{n}{B}\PY{p}{[}\PY{n}{bb}\PY{p}{]}
    \PY{n}{ax}\PY{o}{.}\PY{n}{plot}\PY{p}{(}\PY{n}{cycle\PYZus{}number}\PY{p}{,} \PY{n}{dMn\PYZus{}arr}\PY{p}{[}\PY{p}{:}\PY{p}{,}\PY{n}{bb}\PY{p}{]}\PY{o}{/}\PY{l+m+mi}{100000}\PY{p}{,} \PY{n}{label} \PY{o}{=} \PY{l+s+sa}{f}\PY{l+s+s2}{\PYZdq{}}\PY{l+s+s2}{B = }\PY{l+s+si}{\PYZob{}}\PY{n}{B}\PY{p}{[}\PY{n}{bb}\PY{p}{]}\PY{l+s+si}{:}\PY{l+s+s2}{.4f}\PY{l+s+si}{\PYZcb{}}\PY{l+s+s2}{ * B\PYZus{}max}\PY{l+s+s2}{\PYZdq{}}\PY{p}{)}
\PY{k}{for} \PY{n}{ii} \PY{o+ow}{in} \PY{n}{x}\PY{p}{:}
    \PY{n}{plt}\PY{o}{.}\PY{n}{axline}\PY{p}{(}\PY{p}{(}\PY{n}{ii}\PY{p}{,}\PY{o}{\PYZhy{}}\PY{l+m+mf}{0.1}\PY{p}{)}\PY{p}{,} \PY{p}{(}\PY{n}{ii}\PY{p}{,}\PY{o}{\PYZhy{}}\PY{l+m+mf}{0.12}\PY{p}{)}\PY{p}{,} \PY{n}{color}\PY{o}{=}\PY{l+s+s1}{\PYZsq{}}\PY{l+s+s1}{k}\PY{l+s+s1}{\PYZsq{}}\PY{p}{)}

\PY{n}{plt}\PY{o}{.}\PY{n}{axline}\PY{p}{(}\PY{p}{(}\PY{l+m+mi}{5}\PY{p}{,}\PY{l+m+mi}{0}\PY{p}{)}\PY{p}{,} \PY{p}{(}\PY{l+m+mi}{4}\PY{p}{,}\PY{l+m+mi}{0}\PY{p}{)}\PY{p}{,} \PY{n}{color}\PY{o}{=}\PY{l+s+s1}{\PYZsq{}}\PY{l+s+s1}{k}\PY{l+s+s1}{\PYZsq{}}\PY{p}{,} \PY{n}{label} \PY{o}{=} \PY{l+s+s2}{\PYZdq{}}\PY{l+s+s2}{period boundary}\PY{l+s+s2}{\PYZdq{}}\PY{p}{)}

\PY{n}{plt}\PY{o}{.}\PY{n}{legend}\PY{p}{(}\PY{n}{loc}\PY{o}{=}\PY{l+s+s1}{\PYZsq{}}\PY{l+s+s1}{center left}\PY{l+s+s1}{\PYZsq{}}\PY{p}{,} \PY{n}{bbox\PYZus{}to\PYZus{}anchor}\PY{o}{=}\PY{p}{(}\PY{l+m+mi}{1}\PY{p}{,} \PY{l+m+mf}{0.5}\PY{p}{)}\PY{p}{)}

\PY{n}{plt}\PY{o}{.}\PY{n}{show}\PY{p}{(}\PY{p}{)}
\end{Verbatim}
\end{tcolorbox}

    \begin{center}
    \adjustimage{max size={0.9\linewidth}{0.9\paperheight}}{figs/bvp_kM_en/output_271_0.png}
    \end{center}
    { \hspace*{\fill} \\}
    
    \begin{tcolorbox}[breakable, size=fbox, boxrule=1pt, pad at break*=1mm,colback=cellbackground, colframe=cellborder]
\prompt{In}{incolor}{326}{\boxspacing}
\begin{Verbatim}[commandchars=\\\{\}]
\PY{c+c1}{\PYZsh{} Plot of dMn as a function of the number of large\PYZhy{}scale periods}
\PY{c+c1}{\PYZsh{} for one value of B, and compared to a sine curve of equal period}

\PY{n}{Bi}  \PY{o}{=} \PY{l+m+mi}{4}

\PY{c+c1}{\PYZsh{} a sine function}
\PY{n}{x} \PY{o}{=} \PY{n}{np}\PY{o}{.}\PY{n}{linspace}\PY{p}{(}\PY{l+m+mi}{2}\PY{p}{,}\PY{l+m+mi}{5}\PY{p}{,}\PY{l+m+mi}{1000}\PY{p}{)}
\PY{n}{y} \PY{o}{=} \PY{o}{\PYZhy{}}\PY{n}{np}\PY{o}{.}\PY{n}{sin}\PY{p}{(}\PY{l+m+mi}{2}\PY{o}{*}\PY{n}{pi}\PY{o}{*}\PY{n}{x}\PY{p}{)}\PY{o}{*}\PY{l+m+mf}{0.17}

\PY{n}{fig}\PY{p}{,} \PY{n}{ax} \PY{o}{=} \PY{n}{plt}\PY{o}{.}\PY{n}{subplots}\PY{p}{(}\PY{p}{)}
\PY{n}{plt}\PY{o}{.}\PY{n}{title}\PY{p}{(}\PY{l+s+sa}{f}\PY{l+s+s2}{\PYZdq{}}\PY{l+s+s2}{Change in M Due to Nonlinearity vs Number of Oscillations}\PY{l+s+se}{\PYZbs{}n}\PY{l+s+s2}{Compared to a Sine Function}\PY{l+s+s2}{\PYZdq{}}\PY{p}{)}
\PY{n}{ax}\PY{o}{.}\PY{n}{set\PYZus{}ylabel}\PY{p}{(}\PY{l+s+s1}{\PYZsq{}}\PY{l+s+s1}{Change in M (dMn) due to nonlinearity (10}\PY{l+s+se}{\PYZbs{}u2075}\PY{l+s+s1}{ m}\PY{l+s+se}{\PYZbs{}u207b}\PY{l+s+se}{\PYZbs{}u00b9}\PY{l+s+s1}{)}\PY{l+s+s1}{\PYZsq{}}\PY{p}{)}
\PY{n}{ax}\PY{o}{.}\PY{n}{set\PYZus{}xlabel}\PY{p}{(}\PY{l+s+s2}{\PYZdq{}}\PY{l+s+s2}{B * π}\PY{l+s+se}{\PYZbs{}u00b2}\PY{l+s+s2}{ * R}\PY{l+s+se}{\PYZbs{}u00b2}\PY{l+s+s2}{ with R and B as proportions of their maximum values}\PY{l+s+s2}{\PYZdq{}}\PY{p}{)}

\PY{c+c1}{\PYZsh{} our curve}
\PY{n}{ax}\PY{o}{.}\PY{n}{plot}\PY{p}{(}\PY{n}{np}\PY{o}{.}\PY{n}{power}\PY{p}{(}\PY{n}{np}\PY{o}{.}\PY{n}{array}\PY{p}{(}\PY{n}{R}\PY{p}{)}\PY{p}{,} \PY{l+m+mi}{2}\PY{p}{)}\PY{o}{*}\PY{n}{pi}\PY{o}{*}\PY{o}{*}\PY{l+m+mi}{2}\PY{o}{*}\PY{n}{B}\PY{p}{[}\PY{n}{Bi}\PY{p}{]}\PY{p}{,} \PY{n}{dMn\PYZus{}arr}\PY{p}{[}\PY{p}{:}\PY{p}{,}\PY{n}{Bi}\PY{p}{]}\PY{o}{/}\PY{l+m+mi}{100000}\PY{p}{,} \PY{n}{label} \PY{o}{=} \PY{l+s+s2}{\PYZdq{}}\PY{l+s+s2}{dMn (10}\PY{l+s+se}{\PYZbs{}u2075}\PY{l+s+s2}{ m}\PY{l+s+se}{\PYZbs{}u207b}\PY{l+s+se}{\PYZbs{}u00b9}\PY{l+s+s2}{)}\PY{l+s+s2}{\PYZdq{}}\PY{p}{)}
\PY{c+c1}{\PYZsh{} a sine function with the same period and similar amplitude}
\PY{n}{ax}\PY{o}{.}\PY{n}{plot}\PY{p}{(}\PY{n}{x}\PY{p}{,}\PY{n}{y}\PY{p}{,} \PY{n}{label} \PY{o}{=} \PY{l+s+s2}{\PYZdq{}}\PY{l+s+s2}{\PYZhy{} 0.17 * sin(2*π*x) model}\PY{l+s+s2}{\PYZdq{}}\PY{p}{)}
\PY{c+c1}{\PYZsh{} boundaries of oscillations}
\PY{n}{plt}\PY{o}{.}\PY{n}{axline}\PY{p}{(}\PY{p}{(}\PY{l+m+mi}{2}\PY{p}{,}\PY{o}{\PYZhy{}}\PY{l+m+mf}{0.1}\PY{p}{)}\PY{p}{,} \PY{p}{(}\PY{l+m+mi}{2}\PY{p}{,}\PY{o}{\PYZhy{}}\PY{l+m+mf}{0.12}\PY{p}{)}\PY{p}{,} \PY{n}{color}\PY{o}{=}\PY{l+s+s1}{\PYZsq{}}\PY{l+s+s1}{k}\PY{l+s+s1}{\PYZsq{}}\PY{p}{,} \PY{n}{label} \PY{o}{=} \PY{l+s+s2}{\PYZdq{}}\PY{l+s+s2}{period boundary}\PY{l+s+s2}{\PYZdq{}}\PY{p}{)}
\PY{n}{plt}\PY{o}{.}\PY{n}{axline}\PY{p}{(}\PY{p}{(}\PY{l+m+mi}{3}\PY{p}{,}\PY{o}{\PYZhy{}}\PY{l+m+mf}{0.1}\PY{p}{)}\PY{p}{,} \PY{p}{(}\PY{l+m+mi}{3}\PY{p}{,}\PY{o}{\PYZhy{}}\PY{l+m+mf}{0.12}\PY{p}{)}\PY{p}{,} \PY{n}{color}\PY{o}{=}\PY{l+s+s1}{\PYZsq{}}\PY{l+s+s1}{k}\PY{l+s+s1}{\PYZsq{}}\PY{p}{)}
\PY{n}{plt}\PY{o}{.}\PY{n}{axline}\PY{p}{(}\PY{p}{(}\PY{l+m+mi}{4}\PY{p}{,}\PY{o}{\PYZhy{}}\PY{l+m+mf}{0.1}\PY{p}{)}\PY{p}{,} \PY{p}{(}\PY{l+m+mi}{4}\PY{p}{,}\PY{o}{\PYZhy{}}\PY{l+m+mf}{0.12}\PY{p}{)}\PY{p}{,} \PY{n}{color}\PY{o}{=}\PY{l+s+s1}{\PYZsq{}}\PY{l+s+s1}{k}\PY{l+s+s1}{\PYZsq{}}\PY{p}{)}
\PY{n}{plt}\PY{o}{.}\PY{n}{axline}\PY{p}{(}\PY{p}{(}\PY{l+m+mi}{5}\PY{p}{,}\PY{o}{\PYZhy{}}\PY{l+m+mf}{0.1}\PY{p}{)}\PY{p}{,} \PY{p}{(}\PY{l+m+mi}{5}\PY{p}{,}\PY{o}{\PYZhy{}}\PY{l+m+mf}{0.12}\PY{p}{)}\PY{p}{,} \PY{n}{color}\PY{o}{=}\PY{l+s+s1}{\PYZsq{}}\PY{l+s+s1}{k}\PY{l+s+s1}{\PYZsq{}}\PY{p}{)}
\PY{c+c1}{\PYZsh{} the x axis}
\PY{n}{plt}\PY{o}{.}\PY{n}{axline}\PY{p}{(}\PY{p}{(}\PY{l+m+mi}{5}\PY{p}{,}\PY{l+m+mi}{0}\PY{p}{)}\PY{p}{,} \PY{p}{(}\PY{l+m+mi}{4}\PY{p}{,}\PY{l+m+mi}{0}\PY{p}{)}\PY{p}{,} \PY{n}{color}\PY{o}{=}\PY{l+s+s1}{\PYZsq{}}\PY{l+s+s1}{k}\PY{l+s+s1}{\PYZsq{}}\PY{p}{)}

\PY{n}{plt}\PY{o}{.}\PY{n}{legend}\PY{p}{(}\PY{n}{loc}\PY{o}{=}\PY{l+s+s1}{\PYZsq{}}\PY{l+s+s1}{center left}\PY{l+s+s1}{\PYZsq{}}\PY{p}{,} \PY{n}{bbox\PYZus{}to\PYZus{}anchor}\PY{o}{=}\PY{p}{(}\PY{l+m+mi}{1}\PY{p}{,} \PY{l+m+mf}{0.5}\PY{p}{)}\PY{p}{)}

\PY{n}{plt}\PY{o}{.}\PY{n}{show}\PY{p}{(}\PY{p}{)}
\end{Verbatim}
\end{tcolorbox}

    \begin{center}
    \adjustimage{max size={0.9\linewidth}{0.9\paperheight}}{figs/bvp_kM_en/output_272_0.png}
    \end{center}
    { \hspace*{\fill} \\}
    
    The period is as expected. The structure appears have a slight tilt with
respect to a sine. The amplitude decreases slightly with increased
\(R\). This effect is very small, unlike the rapid decrease in amplitude
of the \(dkn\) oscillations.

    The final graph I wish to show is the plot of \(dMn\) vs \(B\) for
varying \(R\).

    \begin{tcolorbox}[breakable, size=fbox, boxrule=1pt, pad at break*=1mm,colback=cellbackground, colframe=cellborder]
\prompt{In}{incolor}{327}{\boxspacing}
\begin{Verbatim}[commandchars=\\\{\}]
\PY{n}{Ri} \PY{o}{=} \PY{p}{[}\PY{l+m+mi}{10}\PY{p}{,} \PY{l+m+mi}{30}\PY{p}{,} \PY{l+m+mi}{60}\PY{p}{,}  \PY{l+m+mi}{50}\PY{p}{,}\PY{l+m+mi}{100}\PY{p}{]}
\PY{n}{Bi} \PY{o}{=} \PY{p}{[}\PY{l+m+mi}{3}\PY{p}{,}\PY{l+m+mi}{6}\PY{p}{,}\PY{l+m+mi}{9}\PY{p}{,}\PY{l+m+mi}{12}\PY{p}{,}\PY{l+m+mi}{15}\PY{p}{,}\PY{l+m+mi}{18}\PY{p}{]}

\PY{n}{ilim} \PY{o}{=} \PY{n+nb}{len}\PY{p}{(}\PY{n}{Bi}\PY{p}{)}

\PY{n}{fig}\PY{p}{,} \PY{n}{ax} \PY{o}{=} \PY{n}{plt}\PY{o}{.}\PY{n}{subplots}\PY{p}{(}\PY{p}{)}
\PY{n}{plt}\PY{o}{.}\PY{n}{title}\PY{p}{(}\PY{l+s+sa}{f}\PY{l+s+s2}{\PYZdq{}}\PY{l+s+s2}{Change in M Due to Nonlinearity vs B for Various R Values}\PY{l+s+se}{\PYZbs{}n}\PY{l+s+s2}{With their Expected Cycle Boundary Marked}\PY{l+s+s2}{\PYZdq{}}\PY{p}{)}
\PY{n}{ax}\PY{o}{.}\PY{n}{set\PYZus{}ylabel}\PY{p}{(}\PY{l+s+s1}{\PYZsq{}}\PY{l+s+s1}{Change in M (dMn) due to nonlinearity (10}\PY{l+s+se}{\PYZbs{}u2075}\PY{l+s+s1}{ m}\PY{l+s+se}{\PYZbs{}u207b}\PY{l+s+se}{\PYZbs{}u00b9}\PY{l+s+s1}{)}\PY{l+s+s1}{\PYZsq{}}\PY{p}{)}
\PY{n}{ax}\PY{o}{.}\PY{n}{set\PYZus{}xlabel}\PY{p}{(}\PY{l+s+s2}{\PYZdq{}}\PY{l+s+s2}{Magnetic field strength B (B = B * B\PYZus{}max)}\PY{l+s+s2}{\PYZdq{}}\PY{p}{)}

\PY{k}{for} \PY{n}{rr} \PY{o+ow}{in} \PY{n}{Ri}\PY{p}{:}
    \PY{n}{ax}\PY{o}{.}\PY{n}{plot}\PY{p}{(}\PY{n}{B}\PY{p}{,} \PY{n}{dMn\PYZus{}arr}\PY{p}{[}\PY{n}{rr}\PY{p}{,}\PY{p}{:}\PY{p}{]}\PY{o}{/}\PY{l+m+mi}{100000}\PY{p}{,} \PY{n}{label} \PY{o}{=} \PY{l+s+sa}{f}\PY{l+s+s2}{\PYZdq{}}\PY{l+s+s2}{R = }\PY{l+s+si}{\PYZob{}}\PY{n}{R}\PY{p}{[}\PY{n}{rr}\PY{p}{]}\PY{l+s+si}{:}\PY{l+s+s2}{.5f}\PY{l+s+si}{\PYZcb{}}\PY{l+s+s2}{ * 10}\PY{l+s+se}{\PYZbs{}u207b}\PY{l+s+se}{\PYZbs{}u2076}\PY{l+s+s2}{ m}\PY{l+s+s2}{\PYZdq{}}\PY{p}{)}
    \PY{n}{plt}\PY{o}{.}\PY{n}{axline}\PY{p}{(}\PY{p}{(}\PY{l+m+mi}{2}\PY{o}{/}\PY{p}{(}\PY{n}{R}\PY{p}{[}\PY{n}{rr}\PY{p}{]}\PY{o}{*}\PY{o}{*}\PY{l+m+mi}{2} \PY{o}{*} \PY{n}{pi}\PY{o}{*}\PY{o}{*}\PY{l+m+mi}{2}\PY{p}{)}\PY{p}{,}\PY{o}{\PYZhy{}}\PY{l+m+mf}{0.1}\PY{p}{)}\PY{p}{,} \PY{p}{(}\PY{l+m+mi}{2}\PY{o}{/}\PY{p}{(}\PY{n}{R}\PY{p}{[}\PY{n}{rr}\PY{p}{]}\PY{o}{*}\PY{o}{*}\PY{l+m+mi}{2} \PY{o}{*} \PY{n}{pi}\PY{o}{*}\PY{o}{*}\PY{l+m+mi}{2}\PY{p}{)}\PY{p}{,}\PY{o}{\PYZhy{}}\PY{l+m+mf}{0.11}\PY{p}{)}\PY{p}{,} \PY{n}{color}\PY{o}{=}\PY{l+s+s1}{\PYZsq{}}\PY{l+s+s1}{k}\PY{l+s+s1}{\PYZsq{}}\PY{p}{)}
\PY{n}{plt}\PY{o}{.}\PY{n}{axline}\PY{p}{(}\PY{p}{(}\PY{l+m+mf}{0.54}\PY{p}{,}\PY{l+m+mi}{0}\PY{p}{)}\PY{p}{,} \PY{p}{(}\PY{l+m+mf}{0.55}\PY{p}{,}\PY{l+m+mi}{0}\PY{p}{)}\PY{p}{,} \PY{n}{color}\PY{o}{=}\PY{l+s+s1}{\PYZsq{}}\PY{l+s+s1}{k}\PY{l+s+s1}{\PYZsq{}}\PY{p}{,} \PY{n}{label} \PY{o}{=} \PY{l+s+s2}{\PYZdq{}}\PY{l+s+s2}{period boundary}\PY{l+s+s2}{\PYZdq{}}\PY{p}{)}

\PY{n}{plt}\PY{o}{.}\PY{n}{legend}\PY{p}{(}\PY{n}{loc}\PY{o}{=}\PY{l+s+s1}{\PYZsq{}}\PY{l+s+s1}{center left}\PY{l+s+s1}{\PYZsq{}}\PY{p}{,} \PY{n}{bbox\PYZus{}to\PYZus{}anchor}\PY{o}{=}\PY{p}{(}\PY{l+m+mi}{1}\PY{p}{,} \PY{l+m+mf}{0.5}\PY{p}{)}\PY{p}{)}

\PY{n}{plt}\PY{o}{.}\PY{n}{show}\PY{p}{(}\PY{p}{)}
\end{Verbatim}
\end{tcolorbox}

    \begin{center}
    \adjustimage{max size={0.9\linewidth}{0.9\paperheight}}{figs/bvp_kM_en/output_275_0.png}
    \end{center}
    { \hspace*{\fill} \\}
    
    \begin{tcolorbox}[breakable, size=fbox, boxrule=1pt, pad at break*=1mm,colback=cellbackground, colframe=cellborder]
\prompt{In}{incolor}{328}{\boxspacing}
\begin{Verbatim}[commandchars=\\\{\}]
\PY{n}{Ri} \PY{o}{=} \PY{p}{[}\PY{l+m+mi}{55}\PY{p}{,} \PY{l+m+mi}{60}\PY{p}{,} \PY{l+m+mi}{65}\PY{p}{,} \PY{l+m+mi}{70}\PY{p}{,} \PY{l+m+mi}{75}\PY{p}{,} \PY{l+m+mi}{80}\PY{p}{]}
\PY{n}{Ri} \PY{o}{=} \PY{p}{[}\PY{l+m+mi}{60}\PY{p}{,} \PY{l+m+mi}{70}\PY{p}{,} \PY{l+m+mi}{80}\PY{p}{,} \PY{l+m+mi}{90}\PY{p}{,} \PY{l+m+mi}{100}\PY{p}{]}
\PY{n}{Bi} \PY{o}{=} \PY{p}{[}\PY{l+m+mi}{3}\PY{p}{,}\PY{l+m+mi}{6}\PY{p}{,}\PY{l+m+mi}{9}\PY{p}{,}\PY{l+m+mi}{12}\PY{p}{,}\PY{l+m+mi}{15}\PY{p}{,}\PY{l+m+mi}{18}\PY{p}{]}

\PY{n}{ilim} \PY{o}{=} \PY{n+nb}{len}\PY{p}{(}\PY{n}{Bi}\PY{p}{)}

\PY{n}{fig}\PY{p}{,} \PY{n}{ax} \PY{o}{=} \PY{n}{plt}\PY{o}{.}\PY{n}{subplots}\PY{p}{(}\PY{p}{)}
\PY{n}{plt}\PY{o}{.}\PY{n}{title}\PY{p}{(}\PY{l+s+sa}{f}\PY{l+s+s2}{\PYZdq{}}\PY{l+s+s2}{Change in M Due to Nonlinearity vs B for Various R Values}\PY{l+s+se}{\PYZbs{}n}\PY{l+s+s2}{With their Expected Cycle Boundary Marked}\PY{l+s+s2}{\PYZdq{}}\PY{p}{)}
\PY{n}{ax}\PY{o}{.}\PY{n}{set\PYZus{}ylabel}\PY{p}{(}\PY{l+s+s1}{\PYZsq{}}\PY{l+s+s1}{Change in M (dMn) due to nonlinearity (10}\PY{l+s+se}{\PYZbs{}u2075}\PY{l+s+s1}{ m}\PY{l+s+se}{\PYZbs{}u207b}\PY{l+s+se}{\PYZbs{}u00b9}\PY{l+s+s1}{)}\PY{l+s+s1}{\PYZsq{}}\PY{p}{)}
\PY{n}{ax}\PY{o}{.}\PY{n}{set\PYZus{}xlabel}\PY{p}{(}\PY{l+s+s2}{\PYZdq{}}\PY{l+s+s2}{Magnetic field strength B (B = B * B\PYZus{}max)}\PY{l+s+s2}{\PYZdq{}}\PY{p}{)}

\PY{k}{for} \PY{n}{rr} \PY{o+ow}{in} \PY{n}{Ri}\PY{p}{:}
    \PY{n}{ax}\PY{o}{.}\PY{n}{plot}\PY{p}{(}\PY{n}{B}\PY{p}{,} \PY{n}{dMn\PYZus{}arr}\PY{p}{[}\PY{n}{rr}\PY{p}{,}\PY{p}{:}\PY{p}{]}\PY{o}{/}\PY{l+m+mi}{100000}\PY{p}{,} \PY{n}{label} \PY{o}{=} \PY{l+s+sa}{f}\PY{l+s+s2}{\PYZdq{}}\PY{l+s+s2}{R = }\PY{l+s+si}{\PYZob{}}\PY{n}{R}\PY{p}{[}\PY{n}{rr}\PY{p}{]}\PY{l+s+si}{:}\PY{l+s+s2}{.5f}\PY{l+s+si}{\PYZcb{}}\PY{l+s+s2}{ * 10}\PY{l+s+se}{\PYZbs{}u207b}\PY{l+s+se}{\PYZbs{}u2076}\PY{l+s+s2}{ m}\PY{l+s+s2}{\PYZdq{}}\PY{p}{)}
    \PY{n}{plt}\PY{o}{.}\PY{n}{axline}\PY{p}{(}\PY{p}{(}\PY{l+m+mi}{2}\PY{o}{/}\PY{p}{(}\PY{n}{R}\PY{p}{[}\PY{n}{rr}\PY{p}{]}\PY{o}{*}\PY{o}{*}\PY{l+m+mi}{2} \PY{o}{*} \PY{n}{pi}\PY{o}{*}\PY{o}{*}\PY{l+m+mi}{2}\PY{p}{)}\PY{p}{,}\PY{o}{\PYZhy{}}\PY{l+m+mf}{0.1}\PY{p}{)}\PY{p}{,} \PY{p}{(}\PY{l+m+mi}{2}\PY{o}{/}\PY{p}{(}\PY{n}{R}\PY{p}{[}\PY{n}{rr}\PY{p}{]}\PY{o}{*}\PY{o}{*}\PY{l+m+mi}{2} \PY{o}{*} \PY{n}{pi}\PY{o}{*}\PY{o}{*}\PY{l+m+mi}{2}\PY{p}{)}\PY{p}{,}\PY{o}{\PYZhy{}}\PY{l+m+mf}{0.11}\PY{p}{)}\PY{p}{,} \PY{n}{color}\PY{o}{=}\PY{l+s+s1}{\PYZsq{}}\PY{l+s+s1}{k}\PY{l+s+s1}{\PYZsq{}}\PY{p}{)}
\PY{n}{plt}\PY{o}{.}\PY{n}{axline}\PY{p}{(}\PY{p}{(}\PY{l+m+mf}{0.54}\PY{p}{,}\PY{l+m+mi}{0}\PY{p}{)}\PY{p}{,} \PY{p}{(}\PY{l+m+mf}{0.55}\PY{p}{,}\PY{l+m+mi}{0}\PY{p}{)}\PY{p}{,} \PY{n}{color}\PY{o}{=}\PY{l+s+s1}{\PYZsq{}}\PY{l+s+s1}{k}\PY{l+s+s1}{\PYZsq{}}\PY{p}{,} \PY{n}{label} \PY{o}{=} \PY{l+s+s2}{\PYZdq{}}\PY{l+s+s2}{period boundary}\PY{l+s+s2}{\PYZdq{}}\PY{p}{)}

\PY{n}{plt}\PY{o}{.}\PY{n}{legend}\PY{p}{(}\PY{n}{loc}\PY{o}{=}\PY{l+s+s1}{\PYZsq{}}\PY{l+s+s1}{center left}\PY{l+s+s1}{\PYZsq{}}\PY{p}{,} \PY{n}{bbox\PYZus{}to\PYZus{}anchor}\PY{o}{=}\PY{p}{(}\PY{l+m+mi}{1}\PY{p}{,} \PY{l+m+mf}{0.5}\PY{p}{)}\PY{p}{)}

\PY{n}{plt}\PY{o}{.}\PY{n}{show}\PY{p}{(}\PY{p}{)}
\end{Verbatim}
\end{tcolorbox}

    \begin{center}
    \adjustimage{max size={0.9\linewidth}{0.9\paperheight}}{figs/bvp_kM_en/output_276_0.png}
    \end{center}
    { \hspace*{\fill} \\}
    
    The lines here are when my pattern reaches zero, not the maxima. If we
ignore the spike, the pattern looks like it travels left with increasing
\(R\), roughly uniformly. This matches my expectations for the shift in
period. No new behaviors are visible.

    \hypertarget{conclusions}{%
\section{Conclusions}\label{conclusions}}

    There are two oscillations which repeat on the scale of \(Rk\) and
\(B\pi^2R^2\), respectively. \(R\) and \(\mu\) change the parameters
which designate the shape and amplitude of these oscillations, while
\(B\) doesn't impact this shape at all. \(dk\) variations for this model
are too small to impact this shape.

We can separate these oscillations into 4 separate oscillation
structures: 
\begin{enumerate}

\item \(Rk\) oscillations of the change in our fast oscillation
wave number when adding nonlinearity (\(dkn\))\\
The shape of these oscillations resemble a secant squared function,
rotated around the minimum of each curve. The minimum is at the period
boundary when the \(B\pi^2R^2\) oscillation is at its period boundary,
and at the period midpoint when the \(B\pi^2R^2\) oscillation is also at
its midpoint. Quarter \(B\pi^2R^2\) oscillations produce a function
which is slightly less than the maximum of both curves. Oscillation
causes the new spike to grow at the old minimum before the old spike
decays into the new minimum. 
\item \(B\pi^2R^2\) oscillations of the change
in our fast oscillation wave number when adding nonlinearity (\(dkn\))\\
The shape of these oscillations at the half-period boundaries of the
\(Rk\) oscillation appear similar to cycloids which start their next
cycle before hitting the 0 axis. This shape has the same phase-shifting
behavior as the \(Rk\) oscillations. The pattern at half-period
boundaries is shifted over by half a period from the one at whole period
boundaries. Between these boundaries, the shape can have one of two
forms. The first is a piecewise graph with two separate arc types, one
facing towards the axis and the other away from the axis, one centered
on the half-period boundaries, and the other centered on the whole
period boundaries. They are centered in terms of curvature as well as
what percentage of the graph they take up. The new arc grows out of the
troughs of the cycloid, rounding it with decreasing curvature, while the
other grows in amplitude and shrinks in curvature as the percentage of
the graph it takes up shrinks. Once it shrinks entirely, the \(kR\)
oscillation reaches a node, with \(B\pi^2R^2\) graph for this value of
\(kR\) being constant. After that, the oscillations grow in a
non-symmetric curve with symmetry similar to the Jacobian dn elliptic,
but with the sharper peaks facing the origin. The sharper peaks become
sharper with time, until we return to the cycloid-like shape, but
shifted over by a half period.
\item \(B\pi^2R^2\) oscillations of the
change in our slow oscillation wave number when adding nonlinearity
(\(dMn\))\\
The oscillation starts at the period boundary as a tilted sine. At this
point, we switch to a piecewise function. Our old function continues to
tilt and look more and more like a sawtooth, but it is cut off by a
growing pattern with a quarter phase shift. The space the old function
takes up is symmetric around the midpoints of the oscillations, while
the new function, which take up the boundaries of the oscillation, start
as a sawtooth with the \(kR\) period boundaries as their midpoints. As
the portion of the graph this new pattern takes grows, the function
curves and the end falls. At the quarter \(kR\) period point, the
previous sawtooth piecewise function disappears and we are very close to
a negative sin function. Growth in \(kR\) causes the sine to grow and
tilt towards the axis. At the half \(kR\) period point, the two sides of
the tilted sine function split.\\
Each time when the new function grows, it doesn't look like a tilted
sine which has been cutoff, but rather has an additional curve near the
cutoff which tilts away from the x-axis. The old function which
disappears doesn't have this feature, and just looks like a cutoff
tilted sine. 
\item \(Rk\) oscillations of the change in our slow
oscillation wave number when adding nonlinearity (\(dMn\))\\
At \(B\pi^2R^2\) period boundaries, the \(dMn\) vs \(Rk\) graph is
almost everywhere zero, with peaks at the period boundaries. At the
half-period boundaries of \(B\pi^2R^2\) oscillations, the spikes are at
the half-period boundaries of \(Rk\) oscillations. As we increase \(B\)
from the \(B\) period boundaries, we see a curve appear, with a shorter,
curved side, and a long straight side, sort of like half of our secant
squared curve. When it rises out of our period boundaries, it has the
long end pointed away from the axis and short end along the axis, above
the axis, but not as far away as the long end. The short end is to the
left of the period boundary, with only the greatest points of the long
end to the right of the period boundary. This piece is connected
piecewise to another copy of the curve with negative values. This curve
has the long end along the x axis and the short end pointing away from
the axis, to the left of the long end. As \(B\) increases, both curves
move away from the axis and rotate, the vertical J with its long end
falling, and the horizontal curve with the long end rotating away from
the axis. The curve changes its shape during the rotation, so that the
short end starts to curve towards the direction of the long end. By the
quarter period, the two shapes appear as mirror opposite of each other.
The rotation continues to the point where the tall curve is not long,
and the long curve is now tall, mirrored and moved over by a half period
from the old pattern. The pattern flattens everywhere except for the
half-period lines, before flipping over and starting again.
\end{enumerate}

\textbf{\(\mu\) amplitude dependence.} The amplitude of all the oscillations
appear to have a roughly linear dependence on the value of \(\mu\).
\(\mu\) also impacts the parameters of the functional forms of the
oscillations. The increase in rotation angle of the secant squared
pattern with increasing \(\mu\) causes the maximum \(dkn\) value to
increase at a noticeably less than linear rate (with our best fit being
\(\mu^{0.355155}\)), while the minimum increases at a rate consistent
with linear dependence (with our best fit being \(\mu^{0.99879}\)). It
also causes a noticeable shift in the point where the left curve is cut
off in favor of the right curve.

\textbf{\(R\) shape dependence.} Our point in the fast oscillations for analysis
of the large \(R\) dependence is not the most convenient for analysis,
but the results are most consistent with a change in parameters of our
functional form with \(R\), most notably the tilt of the secant squared
curve. Our choice of point is between the quarter period and half period
of the \(kR\) oscillations, where the tilted negative sine appears for
the \(B\pi^2R^2\) oscillations. This point is to the right of a whole
period rotated secant squared minimum, but to the left of the
half-period minimum. This produces the dn-like pattern. As \(R\)
decreases, the rotation angle decreases, and the point moves up the
curve on the left and down the curve on the right. Increasing \(R\)
produces the opposite effect, until reaching the point of intersection
of the two switching patterns. \(R\) dependence on amplitude appears
minimal, if any, as any apparent amplitude change could just be a
parameter change.

    


    % Add a bibliography block to the postdoc
    
    
    
\end{document}
